%DIF 1c1
%DIF LATEXDIFF DIFFERENCE FILE
%DIF DEL /Users/ericg/Documents/original-paper/paper.tex   Thu Oct 24 11:37:13 2024
%DIF ADD ../../paper.tex                                   Thu Oct 24 22:28:17 2024
%DIF < \documentclass[acmsmall,anonymous,review,screen]{acmart}
%DIF -------
\documentclass[acmsmall,review,screen]{acmart} %DIF > 
%DIF -------
\let\Bbbk\relax

\usepackage{quiver}
\usepackage{mathpartir}
% \usepackage{tikz-cd}
\usepackage{enumitem}
\usepackage{wrapfig}
\usepackage{fancyvrb}
\usepackage{comment}
\usepackage{array}

%DIF 13a13-14
\theoremstyle{remark} %DIF > 
\newtheorem{remark}{Remark} %DIF > 
%DIF -------
\newenvironment{DIFnomarkup}{}{}
%DIF 14a16-18
\pagestyle{empty} %DIF > 
\settopmatter{printfolios=false} %DIF > 
 %DIF > 
%DIF -------

%% Rights management information.  This information is sent to you
%% when you complete the rights form.  These commands have SAMPLE
%% values in them; it is your responsibility as an author to replace
%% the commands and values with those provided to you when you
%% complete the rights form.
\setcopyright{acmcopyright}
\copyrightyear{2018}
\acmYear{2018}
\acmDOI{10.1145/1122445.1122456}

%% These commands are for a PROCEEDINGS abstract or paper.
\acmConference[Woodstock '18]{Woodstock '18: ACM Symposium on Neural
  Gaze Detection}{June 03--05, 2018}{Woodstock, NY}
\acmBooktitle{Woodstock '18: ACM Symposium on Neural Gaze Detection,
  June 03--05, 2018, Woodstock, NY}
\acmPrice{15.00}
\acmISBN{978-1-4503-XXXX-X/18/06}

\newcommand{\To}{\Rightarrow}
\newcommand{\inl}{\mathsf{inl}}
\newcommand{\inr}{\mathsf{inr}}
\newcommand{\alt}{\mathrel{\bf \,\mid\,}}

\newcommand{\ob}{\text{Ob}}


\newcommand{\extlc}{\text{Ext-}\lambda}
\newcommand{\extlcm}{\text{Ext-}\lambda^{-\text{trans}}}
\newcommand{\extlcmm}{\text{Ext-}\lambda^{-\text{trans}-\text{cast}}}
\newcommand{\extlcprime}{\text{Ext-}\lambda'}
\newcommand{\intlc}{\text{Int-}\lambda}
\newcommand{\intlcbisim}{\text{Int$_\approx$-}\lambda}
\newcommand{\erase}[1]{\lfloor {#1} \rfloor}


\newcommand{\uarrowl}{\mathrel{\rotatebox[origin=c]{-30}{$\leftarrowtail$}}}
\newcommand{\uarrowr}{\mathrel{\rotatebox[origin=c]{60}{$\leftarrowtail$}}}
\newcommand{\darrowl}{\mathrel{\rotatebox[origin=c]{30}{$\twoheadleftarrow$}}}
\newcommand{\darrowr}{\mathrel{\rotatebox[origin=c]{120}{$\twoheadleftarrow$}}}
\newcommand{\vuarrow}{\mathrel{\rotatebox[origin=c]{-90}{$\leftarrowtail$}}}
\newcommand{\vdarrow}{\mathrel{\rotatebox[origin=c]{90}{$\twoheadleftarrow$}}}

% Types, terms, and precision 

\newcommand{\dyn}{{?}}
\newcommand{\nat}{\text{Nat}}
\newcommand{\bool}{\text{Bool}}
\newcommand{\ra}{\rightharpoonup}
\newcommand{\Ret}[1]{\mathsf{Ret\,}{#1}}
\newcommand{\hole}[1]{\bullet \colon {#1}}
\newcommand{\dyntodyn}{\dyn \ra\, \dyn}


\newcommand{\up}[2]{\langle{#2}\uarrowl{#1}\rangle}
\newcommand{\dn}[2]{\langle{#1}\darrowl{#2}\rangle}

\newcommand{\upc}[1]{\text{up}\,{#1}\,}
\newcommand{\dnc}[1]{\text{dn}\,{#1}\,}


% \newcommand{\ret}{\mathsf{ret}}
\newcommand{\err}{\mho}
\newcommand{\zro}{\textsf{zro}}
\newcommand{\suc}{\textsf{suc}}
\newcommand{\lda}[2]{\lambda {#1} . {#2}}

\newcommand{\injarr}[1]{\textsf{Inj}_\ra ({#1})}
\newcommand{\injnat}[1]{\textsf{Inj}_\text{nat} ({#1})}
\newcommand{\casenat}[4]{\text{Case}_\text{nat} ({#1}) \{ \text{no} \to {#2} \alt \text{nat}({#3}) \to {#4} \}}
\newcommand{\casearr}[4]{\text{Case}_\ra ({#1}) \{ \text{no} \to {#2} \alt \text{fun}({#3}) \to {#4} \}}
\newcommand{\casedyn}[5]{\text{Case}_\Dyn ({#1}) \{ \text{nat}({#2}) \to {#3} \alt \text{fun}({#4}) \to {#5} \}}
%\newcommand{\bind}[3]{\text{var } {#1} = {#2} \text{ in } {#3}}
\newcommand{\matchnat}[4]{\text{match } {#1} \text{ with } \{ \textsf {zero} \Rightarrow {#2} \alt \textsf{ suc } {#3} \Rightarrow {#4} \}}

\newcommand{\refl}{\text{refl}}

\newcommand{\Lift}{\text{Lift}}

%\newcommand{\rel}{\circ\hspace{-4px}-\hspace{-4px}\bullet}
%% \newcommand{\rel}{\mathrel{\circ\mkern-6mu-\mkern-6mu\bullet}}
%% \def\rel{\mathrel{\joinrel=}}
\newcommand{\rel}{\mathrel{{\relbar}\hspace{-2px}{\relbar}\hspace{-2px}{\shortmid}\hspace{-2px}{\relbar}\hspace{-2px}{\relbar}}}
% \newcommand{\wand}{\mathrel{-\mkern-6mu*}}
\newcommand{\wand}{\mathrel{\multimap}}
\newcommand{\ltdyn}{\sqsubseteq}
\newcommand{\gtdyn}{\sqsupseteq}
\newcommand{\equidyn}{\mathrel{\gtdyn\ltdyn}}
\newcommand{\gamlt}{\Gamma^\ltdyn}
\newcommand{\deltalt}{\Delta^\ltdyn}
\newcommand{\relcomp}{\odot}

\newcommand{\hasty}[3]{{#1} \vdash {#2} \colon {#3}}
\newcommand{\vhasty}[3]{{#1} \vdash^v {#2} \colon {#3}}
\newcommand{\phasty}[3]{{#1} \vdash^c {#2} \colon {#3}}
\newcommand{\etmprec}[4]{{#1} \vdash {#2} \ltdyn_e {#3} \colon {#4}}
\newcommand{\itmprec}[4]{{#1} \vdash {#2} \ltdyn_i {#3} \colon {#4}}
\newcommand{\etmequidyn}[4]{{#1} \vdash {#2} \equidyn_e {#3} \colon {#4}}
\newcommand{\itmequidyn}[4]{{#1} \vdash {#2} \equidyn_i {#3} \colon {#4}}
\newcommand{\synbisim}{\approx_{\text{syn}}}

\newcommand{\Dwn}{\Downarrow}
\newcommand{\qte}[1]{\text{quote}({#1})}

\newcommand{\elab}[1]{\text{Elab}({#1})}

% Perturbations
\newcommand{\pertp}{\text{Pert}^\text{P}}
\newcommand{\perte}{\text{Pert}^\text{E}}
\newcommand{\pertdyn}[2]{\text{pert-dyn}({#1}, {#2})}
\newcommand{\delaypert}[1]{\text{delay-pert}({#1})}

\newcommand{\pertc}{\text{Pert}_{\text{C}}}
\newcommand{\pertv}{\text{Pert}_{\text{V}}}


% SGDT and Intensional Stuff

\newcommand{\later}{{\vartriangleright}}
\newcommand{\laterhs}{{\later}}
\newcommand{\type}{\texttt{Type}}
\newcommand{\lob}{\text{L\"{o}b}}
\newcommand{\tick}{\mathsf{tick}}
\newcommand{\nxt}{\mathsf{next}}
\newcommand{\fix}{\mathsf{fix}}
\newcommand{\kpa}{\kappa}

% Model-related stuff
\newcommand{\calV}{\mathcal{V}}
\newcommand{\calE}{\mathcal{E}}
\newcommand{\calVk}{\mathcal{V}_k}
\newcommand{\calEk}{\mathcal{E}_k}
\newcommand{\tok}{\mathrel{\mathop{\to}\limits^{\textrm{\footnotesize k}}}}
\newcommand{\timesk}{\mathrel{\mathop{\times}\limits^{\textrm{k}}}}
\newcommand{\Fk}{\mathrel{\mathop{F}\limits^{\textrm{k}}}}
\newcommand{\Uk}{\mathrel{\mathop{U}\limits^{\textrm{k}}}}

\newcommand{\op}[1]{{#1}^{\textrm{op}}}
\newcommand{\calC}{\mathcal{C}}
\newcommand{\Set}{\mathsf{Set}}
\newcommand{\ErrDom}{\mathsf{ErrDom}}
\newcommand{\PreDom}{\mathsf{PreDom}}
\newcommand{\Yo}{\mathsf{Yo}}
\newcommand{\Hom}{\mathsf{Hom}}
\newcommand{\calS}{\mathcal{S}}
\newcommand{\gfix}{\texttt{gfix}}
\newcommand{\qfix}{\texttt{qfix}}
\newcommand{\calU}{\mathcal{U}}
\newcommand{\laterhat}{\widehat{\later}}
\newcommand{\El}{\mathsf{El}}
\newcommand{\Clock}{\mathsf{Clock}}

\newcommand{\Machine}[1]{\mathsf{Machine}\, {#1}}


% Predomains and EP pairs
\newcommand{\Nat}{\mathsf{Nat}}
\newcommand{\Dyn}{\mathsf{Dyn}}
\newcommand{\ty}[1]{\langle {#1} \rangle}
\newcommand{\li}{L_\mho}
\newcommand{\liclk}[1]{L_\mho [{#1}]}

\newcommand{\ext}[2]{\text{ext}\,{#1}\,{#2}}
\newcommand{\map}[2]{\text{map}\,{#1}\,{#2}}

\newcommand{\ltls}{\ltdyn}
\newcommand{\bisim}{\approx}
\newcommand{\semlt}{\le}
\newcommand{\semltbad}{\lesssim}

%\newcommand{\injarr}{\textsf{Inj}_\to}
%\newcommand{\injnat}{\textsf{Inj}_\mathbb{N}}

\newcommand{\id}{\mathsf{id}}
\newcommand{\ep}{\leadsto}

\newcommand{\emb}[2]{\mathsf{emb}_{#1}({#2})}
\newcommand{\proj}[2]{\mathsf{proj}_{#1}({#2})}

\newcommand{\monto}{\to_m}


% Notation for wait functions
\newcommand{\wre}{w_r^e}
\newcommand{\wle}{w_l^e}
\newcommand{\wrp}{w_r^p}
\newcommand{\wlp}{w_l^p}

\newcommand{\sem}[1]{\llbracket {#1} \rrbracket}
\newcommand{\semgl}[1]{{\sem{#1}}^\text{gl}}


% Denotational model
\newcommand{\errdom}{\textsf{ErrDom}}

\newcommand{\ptb}{\text{ptb}}
\newcommand{\push}{\text{push}}
\newcommand{\pull}{\text{pull}}

\newcommand{\pv}{P^{\mathcal{V}}}
\newcommand{\pe}{P^{\mathcal{E}}}
\newcommand{\ptbv}{\text{ptb}^\mathcal{V}}
\newcommand{\ptbe}{\text{ptb}^\mathcal{E}}

\newcommand{\Ppure}{P}
\newcommand{\Pk}{P^K}
\newcommand{\ptbk}{\text{ptb}^K}


\newcommand{\upf}{\text{up}}
\newcommand{\dnf}{\text{dn}
}
\newcommand{\arr}{\to}
\newcommand{\comp}{}

\newcommand{\ltsq}[2]{\mathrel{\ltdyn^{#1}_{#2}}}
\newcommand{\bisimsq}[2]{\mathrel{\bisim^{#1}_{#2}}}
\newcommand{\ltsqbisim}[2]{\mathrel{{\widetilde{\ltdyn}}^{#1}_{#2}}}
\newcommand{\ltbisim}{\mathrel{\widetilde{\ltdyn}}}
\newcommand{\vf}{\mathcal{V}_f}
\newcommand{\vsq}{\mathcal{V}_{sq}}
\newcommand{\ef}{\mathcal{E}_f}
\newcommand{\esq}{\mathcal{E}_{sq}}

\newcommand{\morbisimid}[1]{\text{Endo}_\bisim({#1})}


\newcommand{\sv}{s_{\mathcal{V}}}
\newcommand{\tv}{t_{\mathcal{V}}}
\newcommand{\rv}{r_{\mathcal{V}}}
\newcommand{\se}{s_{\mathcal{E}}}
\newcommand{\te}{t_{\mathcal{E}}}
\newcommand{\re}{r_{\mathcal{E}}}

\newcommand{\Ff}{F_f}
\newcommand{\Fsq}{F_{sq}}
\newcommand{\Uf}{U_f}
\newcommand{\Usq}{U_{sq}}
\newcommand{\arrf}{\arr_f}
\newcommand{\arrsq}{\arr_{sq}}

\newcommand{\vsim}{\mathcal{V}_{\bisim}}
\newcommand{\esim}{\mathcal{E}_{\bisim}}
\newcommand{\svsim}{s_{\mathcal{V}}^\bisim}
\newcommand{\tvsim}{t_{\mathcal{V}}^\bisim}
\newcommand{\rvsim}{r_{\mathcal{V}}^\bisim}
\newcommand{\sesim}{s_{\mathcal{E}}^\bisim}
\newcommand{\tesim}{t_{\mathcal{E}}^\bisim}
\newcommand{\resim}{r_{\mathcal{E}}^\bisim}



\newcommand{\ve}{\mathcal{V}_e}
\newcommand{\ee}{\mathcal{E}_e}

\newcommand{\vr}{\mathcal{V}_r}
\newcommand{\er}{\mathcal{E}_r}

\newcommand{\relatedin}[3]{{#1}\, {#3}\, {#2}}
\newcommand{\binrel}[1]{\mathbin{#1}}


\newcommand{\da}{\downarrow}

\newcommand{\ret}{\text{ret}}
\newcommand{\bind}[3]{{#1} <- {#2}\, ; \, {#3}}

\newcommand{\piv}{\Pi^\mathcal{V}}
\newcommand{\pie}{\Pi^\mathcal{E}}

\newcommand{\upl}{\textsc{UpL}}
\newcommand{\upr}{\textsc{UpR}}
\newcommand{\dnl}{\textsc{DnL}}
\newcommand{\dnr}{\textsc{DnR}}

\newcommand{\delre}{\delta^{r,e}}
\newcommand{\delle}{\delta^{l,e}}
\newcommand{\delrp}{\delta^{r,p}}
\newcommand{\dellp}{\delta^{l,p}}

\newcommand{\qordeq}{\bisim}

\newcommand{\inat}{\text{Inj}_\mathbb{N}}
\newcommand{\itimes}{\text{Inj}_\times}
\newcommand{\iarr}{\text{Inj}_\to}

\newcommand{\tnat}{\mathsf{nat}}
\newcommand{\ttimes}{\mathsf{times}}
\newcommand{\tfun}{\mathsf{fun}}

\newcommand{\cased}[7]{
    \begin{align*}
    \text{Case}_D ({#1}) &\text{ of } \{ \\
        &\alt \text{nat}({#2}) \to {#3}  \\
        &\alt \text{times}({#4}) \to {#5}  \\
        &\alt \text{fun}({#6}) \to {#7} \}
    \end{align*}
}

\newcommand{\inone}{\mathsf{in}_1}
\newcommand{\intwo}{\mathsf{in}_2}
\newcommand{\inthree}{\mathsf{in}_3}
\newcommand{\infour}{\mathsf{in}_4}
\newcommand{\infive}{\mathsf{in}_5}


\newcommand{\delay}{\text{Delay}}
\newcommand{\ledelay}{\le^{\text{Del}}}
\newcommand{\bisimdelay}{\bisim^{\text{Del}}}
\newcommand{\tnow}{\mathsf{now}}
\newcommand{\tlater}{\mathsf{later}}
\newcommand{\Da}{\Downarrow}
%DIF PREAMBLE EXTENSION ADDED BY LATEXDIFF
%DIF UNDERLINE PREAMBLE %DIF PREAMBLE
\RequirePackage[normalem]{ulem} %DIF PREAMBLE
\RequirePackage{color}\definecolor{RED}{rgb}{1,0,0}\definecolor{BLUE}{rgb}{0,0,1} %DIF PREAMBLE
\providecommand{\DIFadd}[1]{{\protect\color{blue}\uwave{#1}}} %DIF PREAMBLE
\providecommand{\DIFdel}[1]{{\protect\color{red}\sout{#1}}}                      %DIF PREAMBLE
%DIF SAFE PREAMBLE %DIF PREAMBLE
\providecommand{\DIFaddbegin}{} %DIF PREAMBLE
\providecommand{\DIFaddend}{} %DIF PREAMBLE
\providecommand{\DIFdelbegin}{} %DIF PREAMBLE
\providecommand{\DIFdelend}{} %DIF PREAMBLE
\providecommand{\DIFmodbegin}{} %DIF PREAMBLE
\providecommand{\DIFmodend}{} %DIF PREAMBLE
%DIF FLOATSAFE PREAMBLE %DIF PREAMBLE
\providecommand{\DIFaddFL}[1]{\DIFadd{#1}} %DIF PREAMBLE
\providecommand{\DIFdelFL}[1]{\DIFdel{#1}} %DIF PREAMBLE
\providecommand{\DIFaddbeginFL}{} %DIF PREAMBLE
\providecommand{\DIFaddendFL}{} %DIF PREAMBLE
\providecommand{\DIFdelbeginFL}{} %DIF PREAMBLE
\providecommand{\DIFdelendFL}{} %DIF PREAMBLE
%DIF COLORLISTINGS PREAMBLE %DIF PREAMBLE
\RequirePackage{listings} %DIF PREAMBLE
\RequirePackage{color} %DIF PREAMBLE
\lstdefinelanguage{DIFcode}{ %DIF PREAMBLE
%DIF DIFCODE_UNDERLINE %DIF PREAMBLE
  moredelim=[il][\color{red}\sout]{\%DIF\ <\ }, %DIF PREAMBLE
  moredelim=[il][\color{blue}\uwave]{\%DIF\ >\ } %DIF PREAMBLE
} %DIF PREAMBLE
\lstdefinestyle{DIFverbatimstyle}{ %DIF PREAMBLE
	language=DIFcode, %DIF PREAMBLE
	basicstyle=\ttfamily, %DIF PREAMBLE
	columns=fullflexible, %DIF PREAMBLE
	keepspaces=true %DIF PREAMBLE
} %DIF PREAMBLE
\lstnewenvironment{DIFverbatim}{\lstset{style=DIFverbatimstyle}}{} %DIF PREAMBLE
\lstnewenvironment{DIFverbatim*}{\lstset{style=DIFverbatimstyle,showspaces=true}}{} %DIF PREAMBLE
%DIF END PREAMBLE EXTENSION ADDED BY LATEXDIFF

\begin{document}

\title{Denotational Semantics of Gradual Typing using Synthetic Guarded Domain Theory}
\author{Eric Giovannini}
\affiliation{
  \department{Electrical Engineering and Computer Science}
  \institution{University of Michigan}
  \country{USA}
}
\email{ericgio@umich.edu}

\author{Tingting Ding}
\affiliation{
  \department{Electrical Engineering and Computer Science}
  \institution{University of Michigan}
  \country{USA}
}
\email{tingtind@umich.edu}

\author{Max S. New}
\affiliation{
  \department{Electrical Engineering and Computer Science}
  \institution{University of Michigan}
  \country{USA}
}
\email{maxsnew@umich.edu}

\begin{abstract}
  Gradually typed programming languages, which allow for soundly
  mixing static and dynamically typed programming styles, present a
  strong challenge for metatheorists. Even the simplest sound
  gradually typed languages feature at least recursion and errors,
  with realistic languages featuring furthermore runtime allocation of
  memory locations and dynamic type tags. Further, the desired
  metatheoretic properties of gradually typed languages have become
  increasingly sophisticated: validity of \DIFdelbegin \DIFdel{type based }\DIFdelend \DIFaddbegin \DIFadd{type-based }\DIFaddend equational
  reasoning as well as the relational property known as
  graduality. Many recent works have tackled verifying these
  properties, but the resulting mathematical developments are highly
  repetitive and tedious, with few reusable theorems persisting across
  different developments.

  In this work, we present a new denotational semantics for gradual
  typing developed using guarded domain theory. Guarded domain theory
  combines the generality of step-indexed logical relations for
  modeling advanced programming features with the modularity and
  reusability of denotational semantics. We demonstrate the
  feasibility of this approach with a model of \DIFaddbegin \DIFadd{a simple }\DIFaddend gradually
  typed lambda calculus and prove the validity of beta-eta equality
  and the graduality theorem for the denotational model. This model
  should provide the basis for a reusable mathematical theory of
  gradually typed program semantics. Finally, \DIFdelbegin \DIFdel{a recent extension to Agda allows
  synthetic guarded domain theory to be readily mechanized in a proof
  assistant, and we formalize many }\DIFdelend \DIFaddbegin \DIFadd{we have mechanized most
  }\DIFaddend of the core theorems of our development \DIFaddbegin \DIFadd{in Guarded Cubical Agda, a
  recent extension of Agda with support for the guarded recursive
  constructions we use}\DIFaddend .
%%     We develop a denotational semantics for a simple gradually typed language
%%     that is adequate and proves the graduality theorem.
%%     %
%%     The denotational semantics is constructed using \emph{synthetic
%%     guarded domain theory} working in a type theory with a later
%%     modality and clock quantification.
%%     %
%%     This provides a remarkably simple presentation of the semantics,
%%     where gradual types are interpreted as ordinary types in our ambient
%%     type theory equipped with an ordinary preorder structure to model
%%     the error ordering.
%%     %
%%     This avoids the complexities of classical domain-theoretic models
%%     (New and Licata) or logical relations models using explicit
%%     step-indexing (New and Ahmed).
%%     %
%%     In particular, we avoid a major technical complexity of New and
%%     Ahmed that requires two logical relations to prove the graduality
%%     theorem.

%%     By working synthetically we can treat the domains in which gradual
%%     types are interpreted as if they were ordinary sets. This allows us
%%     to give a ``na\"ive'' presentation of gradual typing where each
%%     gradual type is modeled as a well-behaved subset of the universal
%%     domain used to model the dynamic type, and type precision is modeled
%%     as simply a subset relation.
%%     %
\end{abstract}

\maketitle

% Outline

% 1. Intro: What do we want out of gradually typed languages and why
%    is it hard to prove? Explanation: Gradual Typing inherently
%    involves recursive types, multiple effects, relational properties.
%    We argue that the increasing complexity of the metatheory of
%    gradual typing makes it a good candidate for 

% 2. Extensional Dream Semantics: double categorical

% 3. The Problem with Step-indexing

% 4. A Compromise

% 5. Formalization in Guarded Cubical Agda

\newif\ifdraft
\drafttrue
\renewcommand{\max}[1]{\ifdraft{\color{blue}[{\bf Max}: #1]}\fi}
\newcommand{\eric}[1]{\ifdraft{\color{orange}[{\bf Eric}: #1]}\fi}
\newcommand{\tingting}[1]{\ifdraft{\color{red}[{\bf Tingting}: #1]}\fi}

\section{Introduction}

% gradual typing, graduality
\subsection{Gradual Typing and Graduality}
In programming language design, there is a tension between \emph{static} typing
and \emph{dynamic} typing disciplines. With static typing, the code is
type-checked at compile time, while in dynamic typing, the type checking is
deferred to run-time. Both approaches have benefits \DIFdelbegin \DIFdel{: }\DIFdelend \DIFaddbegin \DIFadd{and excel in different
scenarios, }\DIFaddend with static typing \DIFdelbegin \DIFdel{, the programmer is assured that if the program passes the type-checker, their
program is free of type errors, and moreover, soundness of the equational theory implies
that program optimizations are valid.
Meanwhile, dynamic typing allows the programmer to rapidly prototype
their application code without needing to commit to }\DIFdelend \DIFaddbegin \DIFadd{offering compile-time assurance of a program's
type safety and type-based reasoning principles that justify program
optimizations, and dynamic typing allowing for rapid prototyping of a codebase
without committing to }\DIFaddend fixed type signatures\DIFdelbegin \DIFdel{for their
functions.
}\DIFdelend \DIFaddbegin \DIFadd{.
%DIF > 
}\DIFaddend Most languages choose between static or dynamic typing and as a result,
programmers that initially write their code in a dynamically typed language \DIFdelbegin \DIFdel{in order to benefit from faster prototyping and development time }\DIFdelend need
to rewrite some or all of their codebase in a static language if they would like
to receive the benefits of static typing once their codebase has matured.

\emph{\DIFdelbegin \DIFdel{Gradually-typed }\DIFdelend \DIFaddbegin \DIFadd{Gradually typed }\DIFaddend languages} \cite{siek-taha06, tobin-hochstadt06} seek to
resolve this tension by allowing for both static and dynamic typing disciplines
to be used in the same codebase\DIFdelbegin \DIFdel{,
and by supporting }\DIFdelend \DIFaddbegin \DIFadd{. These languages support }\DIFaddend smooth interoperability
between statically-typed and dynamically-typed \DIFdelbegin \DIFdel{code.
This flexibility allows programmers to begin their projects in a dynamic style and
enjoy the benefits of dynamic typing related to rapid prototyping and easy modification
while the codebase ``solidifies''. Over time, as parts of the code become more mature
and the programmer is more certain of what the
types should be, the code can be
}\DIFdelend \DIFaddbegin \DIFadd{styles, allowing the programmer to
begin with fully dynamically-typed code and }\DIFaddend \emph{gradually} \DIFdelbegin \DIFdel{migrated }\DIFdelend \DIFaddbegin \DIFadd{migrate portions of the
codebase }\DIFaddend to a statically typed style without needing to rewrite the project in a
completely different language.

%Gradually-typed languages should satisfy two intuitive properties.
% The following two properties have been identified as useful for gradually typed languages.

\DIFdelbegin \DIFdel{In order for this to work as expected, gradually-typed languages should allow for
different parts of the codebase to be in different places along the spectrum from
dynamic to static, and allow for those different parts to interact with one another.
Moreover, gradually-typed languages should support the smooth migration from
dynamic typing to static typing, in that the programmer can initially leave off the
typing annotations and provide them later without altering the meaning of the program.
%DIF <  Sound gradual typing
Furthermore, the parts of the program that are written in a dynamic
style should soundly interoperate with the parts that are written in a
static style.  That is, the interaction between the static and dynamic
components of the codebase should preserve, to the extent possible,
the guarantees made by the static types.  In particular, while
statically-typed code can error at runtime in a gradually-typed
language, such an error can always be traced back to a
dynamically-typed term that violated the typing contract imposed by
statically typed code. Further, static type assertions are sound in
the static portion, and should enable type-based reasoning and
optimization.
}\DIFdelend %DIF >  \eric{This paragraph could be deleted}
%DIF > % In order for this to work as expected, gradually-typed languages should allow for
%DIF > % different parts of the codebase to be in different places along the spectrum from
%DIF > % dynamic to static, and allow for those different parts to interact with one another.
%DIF > % Moreover, gradually-typed languages should support the smooth migration from
%DIF > % dynamic typing to static typing, in that the programmer can initially leave off the
%DIF > % typing annotations and provide them later without altering the meaning of the program.
%DIF > % % Sound gradual typing
%DIF > % Furthermore, the parts of the program that are written in a dynamic
%DIF > % style should soundly interoperate with the parts that are written in a
%DIF > % static style.  That is, the interaction between the static and dynamic
%DIF > % components of the codebase should preserve, to the extent possible,
%DIF > % the guarantees made by the static types.  In particular, while
%DIF > % statically-typed code can error at runtime in a gradually-typed
%DIF > % language, such an error can always be traced back to a
%DIF > % dynamically-typed term that violated the typing contract imposed by
%DIF > % statically typed code. Further, static type assertions are sound in
%DIF > % the static portion, and should enable type-based reasoning and
%DIF > % optimization.

% Moreover, gradually-typed languages should allow for
% different parts of the codebase to be in different places along the spectrum from
% dynamic to static, and allow for those different parts to interact with one another.
% In a \emph{sound} gradually-typed language,
% this interaction should respect the guarantees made by the static types.

% Graduality property
One of the fundamental theorems for gradually typed languages is
\emph{graduality}, also known as the \emph{dynamic gradual guarantee},
originally defined by Siek, Vitousek, Cimini, and Boyland
\cite{siek_et_al:LIPIcs:2015:5031, new-ahmed2018}.
%
Informally, graduality says that \DIFdelbegin \DIFdel{going from a }\DIFdelend \DIFaddbegin \DIFadd{migrating code from }\DIFaddend dynamic to
static \DIFdelbegin \DIFdel{style }\DIFdelend \DIFaddbegin \DIFadd{typing }\DIFaddend should only allow for the introduction of static or
dynamic type errors, and not otherwise change the behavior of the
program.
%
This is a way to capture programmer intuition that increasing type
precision corresponds to a generalized form of runtime assertions in
that there are no observable behavioral changes up to the point of the
first dynamic type error\footnote{once a dynamic type error is raised,
in languages where the type error can be caught, program behavior may
then further diverge, but this is typically not modeled in gradual
calculi.}.
%
Fundamentally, this property comes down to the behavior of
\emph{runtime type casts}, which implement these generalized runtime
assertions.

Additionally, gradually typed languages should offer some of the
benefits of static typing. While \DIFdelbegin \DIFdel{classical }\DIFdelend \DIFaddbegin \DIFadd{standard }\DIFaddend type soundness, that
well-typed programs are free from runtime errors, is not compatible
with runtime type errors, \DIFdelbegin \DIFdel{we can instead focus on }\DIFdelend \DIFaddbegin \DIFadd{it is possible instead to prove that
gradually typed languages validate }\DIFaddend \emph{type-based reasoning}. For
instance, while dynamically typed $\lambda$ calculi only satisfy
$\beta$ equality for their type formers, statically typed $\lambda$
calculi additionally satisfy type-dependent $\eta$ properties that
ensure that functions are determined by their behavior under
application and that pattern matching on data types is safe and
exhaustive. \DIFaddbegin \DIFadd{A gradually typed calculus that validates these
type-dependent $\eta$ laws then provides some of the type-based
reasoning that dynamic languages lack.
}\DIFaddend 

\DIFdelbegin \DIFdel{More concretely, consider a gradually typed language whose only
effects are gradual type errors and divergence. Then if we fix a
result type of natural numbers, a whole program semantics is a partial
function from closed programs to either natural numbers or errors:
}\[ \DIFdel{-\Downarrow : \{M \,|\, \cdot \vdash M : \nat \} \rightharpoonup \mathbb{N} \cup \{\mho\} }\]%DIFAUXCMD
\DIFdel{where $\mho$ is notation for a runtime type error. We write $M
\Downarrow n$ and $M\Downarrow \mho$ to mean this semantics is defined
as a number or error, and $M\Uparrow$ to mean the semantics is
undefined, representing divergence.
%DIF < 
A well-behaved semantics should then satisfy several properties. First, it
should be }\emph{\DIFdel{adequate}}%DIFAUXCMD
\DIFdel{: natural number constants should step to
themselves. Second it should validate type based reasoning. To
formalize type based reasoning we give a typed equational theory for
terms of the language $M \equiv N$ for when two terms should be
considered equivalent. Then we want to verify that the big step
semantics respects this equational theory: if closed programs $M \equiv
N$ are equivalent in the equational theory then they have the same
semantics, $M \Downarrow n \iff N \downarrow n$ and $M\Uparrow \iff N
\Uparrow$ and $M \Downarrow \mho \iff N \Downarrow \mho$.
%DIF < 
Lastly, the graduality property is defined by giving an
}\emph{\DIFdel{inequational}} %DIFAUXCMD
\DIFdel{theory called term precision, where $M \ltdyn N$
roughly means that $M$ and $N$ have the same type erasure and $M$ has
at each point in the program a more precise/static type than $N$.
%DIF < 
Then, the graduality property states that if $M \ltdyn N$ are whole
programs then $M$ must either have the same behavior as $N$ or error:
Either $M\Downarrow \mho$ or $M \Downarrow n $ and $N \Downarrow n$ or
$M \Uparrow $ and $N \Uparrow$}\footnote{\DIFdel{we use a slightly more complex
definition of this relation in our technical development below that is
classically equivalent but constructively weaker}}%DIFAUXCMD
\addtocounter{footnote}{-1}%DIFAUXCMD
\DIFdel{.
}\DIFdelend %DIF >  moved the remaining paragraphs to a new section later in the intro

\subsection{Denotational Semantics in Guarded Domain Theory}

Our goal in this work is to provide an \emph{expressive},
\emph{reusable}, \emph{compositional} semantic framework for defining
such well-behaved semantics of \DIFaddbegin \DIFadd{gradually typed }\DIFaddend programs.
%
Our approach to achieving this goal is to provide a compositional
\emph{denotational semantics}, mapping types to a kind of semantic
domain, terms to functions and relations such as term precision to
proofs of semantic relations between the denoted functions.
%
Since the denotational constructions are all syntax-independent, the
constructions we provide may be reused for similar languages. Since it
is compositional, components can be mixed and matched depending on
what source language features are present.
%DIF < 
\DIFdelbegin \DIFdel{Providing this }\DIFdelend \DIFaddbegin 

\DIFadd{Providing a }\DIFaddend semantics for gradual typing is inherently complicated in
that it involves: (1) recursion and recursive types through the
presence of dynamic types, (2) effects in the form of divergence and
errors (3) relational models in capturing the graduality
property. Recursion and recursive types must be handled using some
flavor of domain theory. Effects can be modeled using monads in the
style of Moggi, or adjunctions in the style of
Levy\DIFaddbegin \DIFadd{\mbox{%DIFAUXCMD
\cite{MOGGI199155,39155,levy99,levystacks}}\hskip0pt%DIFAUXCMD
}\DIFaddend . Relational properties and their verification
lead naturally to the use of reflexive graph categories or double
categories\DIFaddbegin \DIFadd{\mbox{%DIFAUXCMD
\cite{dunphythesis,new-licata18}}\hskip0pt%DIFAUXCMD
}\DIFaddend .

The only prior denotational semantics for gradual typing was given by
New and Licata and is based on a classical Scott-style \emph{domain
theory} \cite{new-licata18}. The fundamental idea is to equip
$\omega$-CPOs with an additional ``error ordering'' $\ltdyn$ which
models the graduality ordering, and for casts to arise from
\emph{embedding-projection pairs}. Then the graduality property
follows as long as all language constructs can be interpreted using
constructions that are monotone with respect to the error ordering.
%
This framework has the benefit of being compositional, and was
expressive enough to be extended to model dependently typed gradual
typing \DIFdelbegin \DIFdel{\mbox{%DIFAUXCMD
\cite{10.1145/3495528}}\hskip0pt%DIFAUXCMD
}\DIFdelend \DIFaddbegin \DIFadd{\mbox{%DIFAUXCMD
\cite{gradualizing-cic}}\hskip0pt%DIFAUXCMD
}\DIFaddend .
%
However, an approach based on classical domain theory has fundamental
limitations: domain theory \DIFdelbegin \DIFdel{is }\DIFdelend \DIFaddbegin \DIFadd{appears }\DIFaddend incapable of modeling certain perversely
recursive features of programming languages such as dynamic type tag
generation and higher-order references \DIFaddbegin \DIFadd{\mbox{%DIFAUXCMD
\cite{Birkedal-Stovring-Thamsborg-2009}}\hskip0pt%DIFAUXCMD
}\DIFaddend , which are commonplace in
real-world gradually typed systems as well as gradual calculi.
%DIF > 
Our long-term goal is to develop a \DIFdelbegin \DIFdel{semantics that will }\DIFdelend \DIFaddbegin \DIFadd{denotational approach that can
}\DIFaddend scale up to these advanced features, and so we \DIFdelbegin \DIFdel{begin with this work by departing from
}\DIFdelend \DIFaddbegin \DIFadd{must abandon }\DIFaddend classical
domain theory \DIFdelbegin \DIFdel{.
}\DIFdelend \DIFaddbegin \DIFadd{as the foundation in order to make progress. In this
work, we provide first steps towards this goal by adapting work on a
simple gradually typed lambda calculus to }\emph{\DIFadd{guarded domain
theory}}\DIFadd{, a task that already involves nontrivial issues.
Our goal is for this model to scale to dynamic type tag
generation and higher-order references in future work.
}\DIFaddend 

\DIFdelbegin \DIFdel{The }\DIFdelend \DIFaddbegin \emph{\DIFadd{Guarded}} \DIFadd{domain theory is the }\DIFaddend main denotational alternative to
classical domain theory that can successfully model these advanced
features\DIFdelbegin \DIFdel{is }\emph{\DIFdel{guarded domain
theory}}%DIFAUXCMD
\DIFdel{, which we adopt in this work}\DIFdelend . While classical domain theory is based on modeling types as
ordered sets with certain joins, guarded domain theory is based on an
entirely different foundations, sometimes (ultra)metric spaces but
more commonly as ``step-indexed sets'', i.e., objects in the
\emph{topos of trees}
\DIFdelbegin \DIFdel{, i.e., }\DIFdelend \DIFaddbegin \DIFadd{\mbox{%DIFAUXCMD
\cite{birkedal-mogelberg-schwinghammer-stovring2011}}\hskip0pt%DIFAUXCMD
.  Such an object
consists of a family $\{X_n\}_{n \in \mathbb{N}}$ of sets along with
restriction functions $r_n : X_{n+1} \to X_n$ for all $n$.  (in
category theoretic terminology, these are }\DIFaddend presheaves on the poset of
natural numbers\DIFdelbegin \DIFdel{\mbox{%DIFAUXCMD
\cite{birkedal-mogelberg-schwinghammer-stovring2011}}\hskip0pt%DIFAUXCMD
.}\DIFdelend \DIFaddbegin \DIFadd{.)  }\DIFaddend We think of \DIFdelbegin \DIFdel{such }\DIFdelend a $\mathbb{N}$-indexed set as \DIFdelbegin \DIFdel{modeling a kind of ``sequence of approximations'' }\DIFdelend \DIFaddbegin \DIFadd{an
infinite sequence of increasingly precise approximations }\DIFaddend to the true
type being modeled.
%
Key to guarded domain theory is that there is \DIFdelbegin \DIFdel{a ``later'' }\DIFdelend \DIFaddbegin \DIFadd{an }\DIFaddend operator
$\triangleright$ on \DIFdelbegin \DIFdel{types. Thinking of
types as a sequence }\DIFdelend \DIFaddbegin \DIFadd{step-indexed sets called ``later''. In terms of
sequences }\DIFaddend of approximations, the later operator delays the
approximation by one step. Then the crucial axiom of guarded domain
theory is that any guarded \DIFaddbegin \DIFadd{domain }\DIFaddend equation $X \cong F(\triangleright
X)$ has a unique solution. This allows guarded domain theory to model
essentially \emph{any} recursive concept, with the caveat that \DIFdelbegin \DIFdel{it is guarded by a later.
%DIF < 
This caveat is the main source of difficulty in adapting a semantic
approach based on classical domain theory to guarded domain theory:
classical domain theory has limitations in what it can model, but it
provides }\emph{\DIFdel{exact}} %DIFAUXCMD
\DIFdel{solutions to domain equations when it
applies. When adapting the New-Licata approach to guarded domain
theory, this means we must work with an }\emph{\DIFdel{intensional}} %DIFAUXCMD
\DIFdel{semantics,
one where unfolding the dynamic type requires an observable
computational step: e.g., a function that pattern matches on an
element of the dynamic type and then returns it unchanged is
}\emph{\DIFdel{not}} %DIFAUXCMD
\DIFdel{equal to the identity function, because it ``costs'' a step
to perform the pattern match.
%DIF < 
This leads to the main departure of our semantics from the New-Licata
approach: while the New-Licata semantics equips $\omega$-CPOs with an
error ordering, we work with types equipped with not just an error
ordering but also a }\emph{\DIFdel{bisimilarity}} %DIFAUXCMD
\DIFdel{relation that reflects when
elements differ only in the number of abstract computational steps
that they take.
%DIF < 
Because bisimilarity, unlike equality, is not transitive this
necessitates further changes to the core concepts of our semantics.
}\DIFdelend \DIFaddbegin \DIFadd{the
recursion is }\emph{\DIFadd{guarded}} \DIFadd{by a use of the later operator.
}\DIFaddend 

\DIFdelbegin \subsection{\DIFdel{Synthetic Guarded Domain Theory}}
%DIFAUXCMD
\addtocounter{subsection}{-1}%DIFAUXCMD
\DIFdelend %DIF >  The theorems we want are about the whole sequence of approximations
%DIF >  whereas working with presheaves is about each finite approximation
%DIF >  The statement of graduality is in terms of global elements
%DIF >  \eric{This next paragraph could be moved to the section on adequacy, but since
%DIF >  the example is in the analytic setting it also makes sense to keep it here. I'm
%DIF >  not sure which option makes more sense.} While the definitions in guarded domain
%DIF >  theory are constructed as sequences of approximations, results about semantics
%DIF >  of programs, e.g., graduality, are actually statements concerning entire
%DIF >  \emph{sequences} of approximations. As a simple example, consider the set of
%DIF >  programs that may take a computational step or return unit.\footnote{This
%DIF >  example is adapted from \cite{mogelberg-paviotti2016}.} In the topos of trees,
%DIF >  we can model the set of these programs as the indexed family of sets $\{X_n\}_{n
%DIF >  \in \mathbb{N}}$ where for each $n \ge 1$ we define $X_n = \{0,1,\dots,n-1,
%DIF >  \bot\}$. Here, $i$ denotes a program that steps $i$ times and then returns, and
%DIF >  $\bot$ denotes a program that fails to terminate in $n-1$ steps or fewer. For
%DIF >  any fixed $n$, the set $X_n$ fails to distinguish a program that takes $n$ steps
%DIF >  and then terminates from a program that never terminates. On the other hand, if
%DIF >  we define the denotation of a program to be a \emph{sequence} of elements $x_i
%DIF >  \in X_i$ for all $i \in \mathbb{N}$, then the denotation of the diverging
%DIF >  program will be distinct from that of a program that terminates, regardless of
%DIF >  the number of steps it takes.

While guarded domain theory can be presented analytically using
ultrametric spaces or the topos of trees, in practice it is
considerably simpler to work \emph{synthetically} by working in a
non-standard \DIFdelbegin \DIFdel{foundation such as }\DIFdelend \DIFaddbegin \DIFadd{foundational system such as }\emph{\DIFadd{guarded type
theory}}\DIFadd{. In }\DIFaddend guarded type theory \DIFdelbegin \DIFdel{where the later modality }\DIFdelend \DIFaddbegin \DIFadd{later }\DIFaddend is taken as a primitive
operation on types, and we take as an axiom that guarded domain
equations have a (necessarily unique) solution. \DIFdelbegin \DIFdel{This has the added benefit that providing a denotational
semantics of gradually typed terms is as simple as a semantics of an effectful language using a
monad on the category }\DIFdelend \DIFaddbegin \DIFadd{The benefit }\DIFaddend of \DIFdelbegin \DIFdel{sets}\DIFdelend \DIFaddbegin \DIFadd{this
synthetic approach is that when working in the non-standard
foundation, we don't need to model an object language type as a
step-indexed set, but instead simply as a set, and object-language
terms can be modeled in the Kleisli category of a simple monad defined
using guarded recursion}\DIFaddend . Not only does this make on-paper reasoning
about guarded domain theory easier, it also enables a simpler avenue
to verification in a proof assistant. Whereas formalizing analytic
guarded domain theory would require \DIFaddbegin \DIFadd{a }\DIFaddend significant theory of presheaves
and making sure that all constructions are functors on categories of
presheaves, formalizing synthetic guarded domain theory can be done by
directly adding the later modality and the guarded fixed point
property axiomatically.
%
% \max{TODO: more here, specifically this is where we should talk about adequacy I think}
%DIF > 
\DIFaddbegin \DIFadd{In this paper, we work informally in a guarded type theory which is described in
more detail in Section \ref{sec:guarded-type-theory}.
}\DIFaddend 

\DIFdelbegin \DIFdel{Because the solutions we obtain working in guarded domain theory are constructed
as a sequence of increasingly better approximations,
when we want to establish a property that holds in the limit,
we must reason about all finite
approximations. In the setting of the topos of trees, this manifests as the need
to consider }\emph{\DIFdel{global elements}} %DIFAUXCMD
\DIFdel{of }\DIFdelend %DIF > 
%DIF >  As in the setting of analytic guarded domain theory, in the synthetic setting it
%DIF >  is also possible to construct global solutions. The specific approach we use in
%DIF >  this work is that of \emph{clocks} and \emph{clock-quantification}
%DIF >  \cite{atkey-mcbride2013}.
\DIFaddbegin 

%DIF >  \max{shouldn't we say we use clock quantification?}
%DIF >  In the setting of synthetic guarded domain theory, there is an analogue of
%DIF >  constructing global solutions. To construct a global solution we must use an
%DIF >  additional construct, e.g., the $\square$ modality
%DIF >  \cite{10.1007/978-3-319-08918-8_8}, whose semantics in the topos of trees model
%DIF >  is to compute the global elements of a presheaf. A related approach involves
%DIF >  objects known as \emph{clocks}, whereby \emph{clock-quantification}
%DIF >  \cite{atkey-mcbride2013} provides a means to obtain a global solution.


\subsection{\DIFadd{Adapting the New-Licata Model to the Guarded Setting}}

%DIF >  Working in a step-counting model, but graduality is independent of steps
%DIF >  Casts take stpes, so we need to reason up to weak bisimilarity, but it is not transitive
%DIF >  New-Licata freely uses transitivity, but we can't adapt their proof methodology because it uses transitivity pervasively
%DIF >  Thus the old proofs from New-Licata that use transitivity don't apply.
%DIF >  Work with a combination of weak bisimilarity and step-counting reasoning
%DIF >  Need a tiny bit of weak bisimilarity which are perturbations, explicit synchronizations that we manipulate syntactically

\DIFadd{Since guarded domain theory only provides solutions to guarded domain equations,
there is no systematic way to convert a classical domain-theoretic semantics to
a guarded one.  Classical domain theory has limitations in what it can model,
but it provides }\emph{\DIFadd{exact}} \DIFadd{solutions to domain equations when it applies. When
adapting the New-Licata approach to guarded domain theory, }\DIFaddend the \DIFdelbegin \DIFdel{presheaves, i.e., a family of elements
$x_i \in X_i$ compatible with the restriction maps}\DIFdelend \DIFaddbegin \DIFadd{presence of later
in the semantics makes it }\emph{\DIFadd{intensional}}\DIFadd{: unfolding the dynamic type
requires an observable computational step}\DIFaddend . For example, \DIFdelbegin \DIFdel{consider the
set of programs that may take a computational step or return unit.
In the topos
of
trees, we can model this set as the indexed family of sets $X_n =
\{0,1,\dots,n-1, \bot\}$ for $n \ge 1$, where $i$ denotes a program that steps
$i$ times and then returns, and $\bot$ denotes a program that fails to terminate
in $n-1$ stepsor fewer. For any fixed $n$, the set $X_n$ fails to distinguish a
program that takes $n$ steps and then terminates from a program that never
terminates. On the other hand, if we instead consider the denotation of a
program to be a global element of
}\DIFdelend \DIFaddbegin \DIFadd{a function that pattern
matches on an element of the dynamic type and then returns it unchanged is
}\emph{\DIFadd{not}} \DIFadd{equal to the identity function, because it ``costs'' a step to
perform the pattern match.
}

%DIF >  - The essential role of transitivity
%DIF >  - We can't have transitivity + extensional reasoning
%DIF >  - Solution: split the ordering into two, and introduce perturbations
%DIF >    to recover some amount of transitive reasoning

\DIFadd{The intensional nature of guarded domain theory leads to the main departure of
our semantics from the New-Licata approach. Because casts involve computational
steps, }\DIFaddend the \DIFdelbegin \DIFdel{presheaf $\{X_n\}$ -- a
family of elements $x_i \in X_i$ for all $i \in \mathbb{N}$ -- then the denotation of the diverging program will be distinct from that of a program that terminates, regardless of the number of steps it takes}\DIFdelend \DIFaddbegin \DIFadd{graduality property must be insensitive to the steps taken by terms.
This means that the model must allow for reasoning }\emph{\DIFadd{up to weak
bisimilarity}}\DIFadd{, where two terms are weakly bisimilar if they differ only in their
number of computational steps}\DIFaddend . \DIFaddbegin \DIFadd{However, the weak bisimilarity relation is not
transitive in the guarded setting, which follows from a no-go theorem we
establish (Theorem \ref{thm:no-go}). The New-Licata proofs freely use
transitivity, and we argue in Section \ref{sec:towards-relational-model} that
some amount of transitive reasoning is nessecary for defining a
syntax-independent model of gradual typing. As a result, the lack of
transitivity of weak bisimilarity presents a challenge in adapting the
New-Licata model to the guarded setting. Our solution is to work with a
combination of weak bisimilarity and step-sensitive reasoning, where the latter
notion }\emph{\DIFadd{is}} \DIFadd{transitive. To deal with the fact that casts take computational
steps, we introduce the novel concept of }\emph{\DIFadd{syntactic perturbations}}\DIFadd{, which
are explicit synchronizations that we manipulate syntactically. The combination
of lock-step reasoning with perturbations is transitive and is able to account
for the step-insensitive nature of the graduality property.
}\DIFaddend 

\DIFdelbegin \DIFdel{In the setting of synthetic guarded domain theory, this same idea applies. To
construct a global solution we must use an additional construct, }\DIFdelend %DIF >  These can then be used to explicitly \emph{synchronize} elements to ensure
%DIF >  that they are in lock-step.
%DIF >  %
%DIF >  In the end, we have come to a resolution of the issue: the lock-step error
%DIF >  ordering where we employ syntactic perturbations to synchronize elements is
%DIF >  still transitive, while simultaneously being able to relate terms that involve
%DIF >  an explicit, known pattern of computational steps.
\DIFaddbegin 


\subsection{\DIFadd{Adequacy}}

%DIF >  Our approach to proving graduality by constructing a denotational model in
%DIF >  guarded type theory consists of two main steps. First, there is the construction
%DIF >  of the model itself. This gives us an interpretation of the terms as well as the
%DIF >  axioms that model the graduality property. 

\DIFadd{Having defined the model of gradual typing in guarded type theory, we need to
show that the guarded model induces a sensible set-theoretic semantics in
ordinary mathematics. We cannot hope to derive such a semantics for }\emph{\DIFadd{all}}
\DIFadd{types (}\DIFaddend e.g., the \DIFdelbegin \DIFdel{$\square$ modality \mbox{%DIFAUXCMD
\cite{10.1007/978-3-319-08918-8_8}}\hskip0pt%DIFAUXCMD
, whose semantics in the topos of trees model is to compute the global elements of a presheaf. A related
approach involves objects known as }\emph{\DIFdel{clocks}}%DIFAUXCMD
\DIFdel{, whereby
}\emph{\DIFdel{clock-quantification}} %DIFAUXCMD
\DIFdel{\mbox{%DIFAUXCMD
\cite{atkey-mcbride2013} }\hskip0pt%DIFAUXCMD
provides a
means to obtain
a global solution. }\DIFdelend \DIFaddbegin \DIFadd{dynamic type), but for the subset of closed terms of base type,
we should be able to extract a well-behaved semantics for which the graduality
property holds.
}\DIFaddend 

\DIFdelbegin \DIFdel{This has important ramifications for defining an adequate denotational }\DIFdelend \DIFaddbegin \DIFadd{More concretely, consider a gradually typed language whose only effects are
gradual type errors and divergence (errors arise from failing casts; divergence
arises because we can use the dynamic type to encode untyped lambda calculus;
see Section \ref{sec:GTLC}). If we fix a result type of natural numbers, a
``big-step'' semantics is a partial function from closed programs to either natural
numbers or errors:
%DIF > 
}\[ \DIFadd{-\Downarrow : \{M \,|\, \cdot \vdash M : \nat \} \rightharpoonup \mathbb{N}
\cup \{\mho\} }\] 
%DIF > 
\DIFadd{where $\mho$ is notation for a runtime type error. We write $M \Downarrow n$ and
$M\Downarrow \mho$ to mean this semantics is defined as a number or error, and
$M\Uparrow$ to mean the semantics is undefined, representing divergence.
%DIF > 
A well-behaved semantics should then satisfy several properties. First, it
should be }\emph{\DIFadd{adequate}}\DIFadd{: natural number constants should evaluate to
themselves: $n \Downarrow n$. Second, it should validate type based reasoning.
To formalize type based reasoning, languages typically have an equational theory
$M \equiv N$ specifying when two terms should be considered equivalent.
Then we want to verify that the big step semantics respects this equational
theory: if closed programs $M \equiv N$ are equivalent in the equational theory
then they have the same semantics, $M \Downarrow n \iff N \Downarrow n$ and
$M\Uparrow \iff N \Uparrow$ and $M \Downarrow \mho \iff N \Downarrow \mho$.
}


%DIF >  The move to synthetic setting is useful, but now you're stuck in the guarded setting
\DIFadd{We have to take care when defining such a big-step }\DIFaddend semantics in
guarded \DIFdelbegin \DIFdel{domain }\DIFdelend \DIFaddbegin \DIFadd{type }\DIFaddend theory, as \DIFdelbegin \DIFdel{we seek to do in this paper: we do not want our denotational semantics to conflate a diverging program with a program that fails
to terminate in sufficiently few steps.
Thus, establishing adequacy of a
semantics in guarded domain theory will require a means of globalizing, a point we will return to in Section \ref{sec:big-step-term-semantics}}\DIFdelend \DIFaddbegin \DIFadd{a naive definition of termination is
incorrect: we can prove using guarded recursion that an infinite loop
terminates! This is because the definition of termination inside
guarded logic has a different meaning, essentially meaning that for
every finite number $n$, either the program terminates in $< n$ steps
or it takes $n+1$ steps, which is equivalent to a statement of
progress rather than termination. For this reason, we need a kind of
``escape hatch'' in our ambient type theory to leave the guarded
setting and make definitions that are interpreted as constructions in
ordinary mathematics. We incorporate this by using
}\emph{\DIFadd{clock-quantification}} \DIFadd{\mbox{%DIFAUXCMD
\cite{atkey-mcbride2013, kristensen-mogelberg-vezzosi2022}
}\hskip0pt%DIFAUXCMD
%DIF >  \max{Eric should provide a citation for Clocked cubical here since that's close to our actual approach}
in our type theory. In particular, the guarded type theory in which we work has
a notion of clocks that index all of our guarded constructions. This is
discussed in more detail in Section \ref{sec:big-step-term-semantics}.
}

\DIFadd{To define the big-step semantics, we apply clock quantification to obtain
an interpretation of closed programs of base type as partial functions.
We describe the process in more detail in Section
\ref{sec:big-step-term-semantics}. The resulting term semantics is adequate, and
furthemore it validates the equational theory, since equivalent terms in the
equational theory denote equal terms in the big-step semantics.
}

\DIFadd{Lastly, the semantics should be adequate for the graduality property. Graduality
is typically axiomatized by giving an }\emph{\DIFadd{inequational}} \DIFadd{theory called term
precision, where $M \ltdyn N$ roughly means that $M$ and $N$ have the same type
erasure and $M$ has at each point in the program a more precise/static type than
$N$. Then adequacy says that if $M$ and $N$ are whole programs returning type
$\nat$ and $M \ltdyn N$, then the denotations of $M$ and $N$ as big-step terms
should be related in the expected way: either $M$ errors, or $M$ and $N$ have
the same extensional behavior. That is, either $M\Downarrow \mho$ or $M
\Downarrow n $ and $N \Downarrow n$ or $M \Uparrow $ and $N
\Uparrow$}\footnote{\DIFadd{we use a slightly more complex definition of this relation in
our technical development below that is classically equivalent but
constructively weaker}}\DIFaddend .
%DIF > 
\DIFaddbegin \DIFadd{To prove that the semantics is adequate for graduality, we again apply clock
quantification, this time to the relation that denotes the term precision
ordering. We show that a term precision ordering $M \ltdyn N$ implies the
corresponding ordering on the partial functions denoted by $M$ and $N$. The
details of the proof are given in Section \ref{sec:adequacy}.
}\DIFaddend 


% For example, when establishing the adequacy of a semantics defined using
% guarded type theory, we must reason about all finite approximations.

% In this paper, we develop an adequate denotational semantics that satisfies
% graduality and soundness of the equational theory of cast calculi using
% synthetic guarded domain theory.  


\subsection{Contributions \DIFaddbegin \DIFadd{and Outline}\DIFaddend }

The main contribution of this work is a compositional denotational semantics \DIFdelbegin \DIFdel{for gradually typed languages }\DIFdelend \DIFaddbegin \DIFadd{in
guarded type theory for a simple gradually typed language }\DIFaddend that validates
$\beta\eta$ equality and satisfies a graduality theorem.
\DIFdelbegin \DIFdel{A great deal of the }\DIFdelend %DIF > 
%DIF >  In particular, the notion of syntactic perturbations stands out as a key
%DIF >  technical contribution and is the most significant and novel feature of our
%DIF >  model that allowed us to successfully adapt the denotational approach of New and
%DIF >  Licata to the guarded setting.
%DIF >  In particular, the notion of syntactic perturbations stands out as a key feature
%DIF >  of our model, and is our most significant technical contribution allowing us to
%DIF >  successfully adapt the denotational approach of New and Licata to the guarded
%DIF >  setting.
%DIF > 
\DIFaddbegin \DIFadd{Within our semantics, the notion of syntactic perturbation is our most
significant technical contribution, allowing us to successfully adapt the
denotational approach of New and Licata to the guarded setting.
%DIF > 
Most of our }\DIFaddend work has further been verified in Guarded Cubical Agda
\cite{veltri-vezzosi2020}, demonstrating that the semantics is readily
mechanizable. \DIFaddbegin \DIFadd{We provide an overview of the mechanization effort, and what
remains to be formalized, in Section \ref{sec:mechanization}.
%DIF > 
}\DIFaddend 


\DIFdelbegin %DIFDELCMD < \begin{enumerate}
\begin{enumerate}%DIFAUXCMD
%DIFDELCMD < \item %%%
\item%DIFAUXCMD
\DIFdel{First, we give a simple concrete term semantics where we show
  how to model the dynamic type as a solution to a guarded domain equation.
}%DIFDELCMD < \item %%%
\item%DIFAUXCMD
\DIFdel{Next, we identify where prior work on classical domain theoretic
  semantics of gradual typing breaks down when using guarded semantics
  of recursive types.
}%DIFDELCMD < \item %%%
\item%DIFAUXCMD
\DIFdel{We develop a key new concept of }\emph{\DIFdel{syntactic perturbations}}%DIFAUXCMD
\DIFdel{,
  which allow us to recover enough extensional reasoning to model the
  graduality property compositionally.
}%DIFDELCMD < \item %%%
\item%DIFAUXCMD
\DIFdel{We combine this insight together with an abstract categorical
  model of gradual typing using reflexive graph categories and
  call-by-push-value to give a compositional construction of our
  denotational model.
}%DIFDELCMD < \item %%%
\item%DIFAUXCMD
\DIFdel{We prove that the resulting denotational model provides a
  well-behaved semantics as defined above by proving }\emph{\DIFdel{adequacy}}%DIFAUXCMD
\DIFdel{, respect for an equational theory and the graduality property.
}
\end{enumerate}%DIFAUXCMD
%DIFDELCMD < \end{enumerate}
%DIFDELCMD < %%%
\DIFdelend %DIF >  Syntactic perturbations 
%DIF > 
%DIF >  The notion of \emph{syntactic perturbations} allow us to encode the steps
%DIF >  taken by a term in a form that we can manipulate. This allows us to recover
%DIF >  enough extensional reasoning to model the graduality property compositionally.

%DIF >  Syntactic perturbations give us a type-directed way of imposing delays on terms.
\DIFaddbegin 

\begin{comment}
\begin{enumerate}
\item First, we give a simple concrete term semantics where we show
  how to model the dynamic type as a solution to a guarded domain equation.
\item Next, we identify where prior work on classical domain theoretic
  semantics of gradual typing breaks down when using guarded semantics
  of recursive types.
\item We develop a key new concept of \emph{syntactic perturbations},
  which allow us to recover enough extensional reasoning to model the
  graduality property compositionally.
\item We combine this insight together with an abstract categorical
  model of gradual typing using reflexive graph categories and
  call-by-push-value to give a compositional construction of our
  denotational model.
\item We prove that the resulting denotational model provides a
  well-behaved semantics as defined above by proving \emph{adequacy},
  respect for an equational theory and the graduality property.
\end{enumerate}
\end{comment}

\DIFaddend The paper is laid out as follows:
%DIF >  \max{I think there's too much detail in this outline, especially parts 4 and 5. Can we instead expand those into the Intro?}
\begin{enumerate}
\item In Section \ref{sec:GTLC} we fix our input language, a fairly
  typical gradually typed cast calculus.
\item In Section \ref{sec:concrete-term-model} we develop a
  \DIFdelbegin \DIFdel{first
  }\DIFdelend denotational semantics in synthetic guarded domain theory \DIFdelbegin \DIFdel{that
  satisfies adequacy and respects }\DIFdelend \DIFaddbegin \DIFadd{for the
  }\emph{\DIFadd{terms}} \DIFadd{of the gradual lambda calculus.  The model is adeqauate
  and validates }\DIFaddend the equational theory, but \DIFdelbegin \DIFdel{does not
  validate }\DIFdelend \DIFaddbegin \DIFadd{it does not satisfy
  }\DIFaddend graduality. We use this to introduce some of our main technical
  tools: modeling recursive types in guarded type theory and modeling
  effects using call-by-push-value.
\item In Section \ref{sec:towards-relational-model} we show where the New-Licata
  classical domain theoretic approach fails to adapt cleanly to the guarded
  setting and explore the difficulties of proving graduality in an intensional
  model. \DIFaddbegin \DIFadd{We prove a no-go theorem about extensional, transitive relations,
  introduce the lock-step error ordering and weak bisimilarity relation, and
  motivate the need for perturbations.
}\DIFaddend \item In Section \ref{sec:concrete-relational-model} we describe the
  \DIFdelbegin \DIFdel{resulting
  }\DIFdelend \DIFaddbegin \DIFadd{construction of the }\DIFaddend model in detail\DIFdelbegin \DIFdel{. We then }\DIFdelend \DIFaddbegin \DIFadd{, and discuss the
  translation of the syntax and axioms of the gradually typed cast calculus into the model.
  Lastly, we }\DIFaddend prove that the \DIFdelbegin \DIFdel{resulting denotational model is }\emph{\DIFdel{adequate}} %DIFAUXCMD
\DIFdelend \DIFaddbegin \DIFadd{model is adequate }\DIFaddend for the graduality property\DIFdelbegin \DIFdel{: a closed term precision $M \ltdyn
  N : \nat$ has the expected semantics, i.e., that $M$ errors or $M$ and $N$ have the
  same extensional behavior.
  }\DIFdelend \DIFaddbegin \DIFadd{.
  %DIF >  The resulting model validates the axioms for type and term precision specified
  %DIF >  in Section \ref{sec:GTLC}. The final step is to prove that the model is
  %DIF >  \emph{adequate} for the graduality property: a closed term precision $M \ltdyn
  %DIF >  N : \nat$ has the expected semantics, i.e., that $M$ errors or $M$ and $N$
  %DIF >  have the same extensional behavior. We do so by extending the globalization
  %DIF >  techniques used in defining the big-step term model to account for the
  %DIF >  interpretation of term precision.
}\DIFaddend \item In Section \ref{sec:discussion} we discuss prior work on proving
  graduality, \DIFdelbegin \DIFdel{our }\DIFdelend \DIFaddbegin \DIFadd{the }\DIFaddend partial mechanization of \DIFdelbegin \DIFdel{these results and discuss
  }\DIFdelend \DIFaddbegin \DIFadd{our results in Agda, and
  }\DIFaddend future directions for denotational semantics of gradual typing.
\end{enumerate}


%\input{technical-background}

\section{Syntactic Theory of Gradually Typed Lambda Calculus}\label{sec:GTLC}

Here we give an overview of a fairly standard cast calculus for
gradual typing along with its (in-)equational theory that capture our
desired notion of type-based reasoning and graduality. The main
departure from prior work is our explicit treatment of type precision
derivations and an equational theory of those derivations.

We give the basic syntax and select typing rules in
Figure~\ref{fig:gtlc-syntax}. We include a dynamic type, a type of
numbers, the call-by-value function type $A \ra A'$ and products.
%
We include a syntax for \emph{type precision} derivations $c : A
\ltdyn A'$; the typing is given in Figure~\ref{fig:gtlc-syntax}.
%
Any type precision derivation $c : A \ltdyn A'$ induces a pair of
casts, the upcast $\upc c : A \ra A'$ and the downcast $\dnc c : A' \ra
A$.
%
The syntactic intuition is that $c$ is a proof that $A$ is ``less
dynamic'' than $A'$. Semantically, this gives us coercions back and
forth where the upcast is (to a first-order) a pure function whereas
the downcast can fail.
%
These casts are inserted automatically in an elaboration from a
surface language. In this work, we are focused on semantic aspects and
so elide these standard details.
%
The syntax of precision derivations includes reflexivity $r(A)$ and
transitivity $cc'$ as well as monotonicity $c \ra c'$ and $c \times
c'$ that are \emph{covariant} in all arguments and finally generators
$\inat,\iarr,\itimes$ that correspond to the type tags of our dynamic
type.
%
We additionally impose an equational theory $c \equiv c'$ on the
derivations that implies that the corresponding casts are weakly
bisimilar in the semantics.
%
We impose category axioms for the reflexivity and
transitivity and functoriality for the monotonicity rules.
%
We note the following two admissible principles: any two derivations
$c,c' : A \ltdyn A'$ of the same fact are equivalent $c \equiv c'$ and
for any $A$, there is a derivation $\textrm{dyn}(A): A \ltdyn\dyn$. That is, $\dyn$ is the ``most dynamic'' type.

\DIFaddbegin \DIFadd{There is a more common set of rules for type precision where reflexivity and
transitivity are admissible, and whenever $A \ltdyn A'$, there is a unique
precision derivation witnessing this. These rules are shown in Section
\ref{sec:appendix-gtlc-syntax} in the Appendix. The reason for choosing our
system rather than that one is that in our semantics, equivalent type precision
derivations do not denote }\emph{\DIFadd{equal}} \DIFadd{relations. Instead, they denote
relations that are }\emph{\DIFadd{quasi-equivalent}}\DIFadd{, i.e., if two terms are related by
one then they are related also by the other up to insertion of delays (see
Definition \ref{def:quasi-equivalent} for the details).
%DIF > 
However, because all type precision derivations in our system are equivalent, it
is straightforward to define a translation from the more standard system of type
and term precision into ours, so ultimately our graduality proofs can still be
applied to the standard formulation.
}

\DIFaddend \begin{figure}
  \begin{mathpar}
    \begin{array}{rcl}
    \text{Types } A &::=& \nat \alt \,\dyn \alt A \ra A' \alt A \times A'\\
    \text{Type Precision } c &::=& r(A) \alt c c' \alt \iarr \alt \inat \alt \itimes \alt c \ra c' \alt c \times c'\\
    \text{Values } V &::=& x \alt \upc c V \alt \zro \alt \suc\, V \alt \lda{x}{M} \alt (V,V') \\ 
    \text{Terms } M,N &::=& \err\alt \upc c M \alt \dnc c M \alt \zro \alt \suc\, M \alt \lda{x}{M} \\ 
     &&\alt M\, N \alt (M,N) \alt \textrm{let } (x,y) = M \textrm{ in } N\\
    \text{Contexts } \Gamma &::= &\cdot \alt \Gamma, x : A \\
    \text{Ctx Precision } \Delta &::=& \cdot\alt \Delta,x:c
  \end{array}

  \inferrule
  {\Gamma \vdash M : A \and c : A \ltdyn A'}
  {\Gamma \vdash \upc c M : A'}

  \inferrule
  {\Gamma \vdash N : A' \and c : A \ltdyn A'}
  {\Gamma \vdash \dnc c N : A}

  \inferrule{}{\Gamma \vdash \mho : A}
  \end{mathpar}
  \begin{mathpar}
    \inferrule{}{r(A) : A \ltdyn A}\and
    \inferrule{c : A \ltdyn A' \and c' : A' \ltdyn A''}{cc' : A \ltdyn A''}\and
    \inferrule{}{\iarr \colon \dyn \ra \dyn \ltdyn \dyn}\and
    \inferrule{}{\inat \colon \nat \ltdyn \dyn}\and
    \inferrule{}{\itimes \colon \dyn \times \dyn \ltdyn \dyn}\and
    \inferrule{c_i : A_i \ltdyn A_i' \and c_o : A_o \ltdyn A_o'}{c_i \ra c_o : (A_i \ra A_o) \ltdyn (A_i' \ra A_o')}\and
    \inferrule{c_1 : A_1 \ltdyn A_1' \and c_2 : A_2 \ltdyn A_2'}{c_1 \times c_2 : (A_1 \times A_2) \ltdyn (A_1' \times A_2')}\and
     r(A)c \equiv c\and
     c \equiv cr(A')\and
     c(c'c'') \equiv (cc')c''\and
     r(A_i \ra A_o) \equiv r(A_i) \ra r(A_o)\and
     r(A_1\times A_2) \equiv r(A_1) \times r(A_2)\and
     (c_i \ra c_o)(c_i' \ra c_o')\equiv (c_ic_i' \ra c_oc_o') \and
     (c_1\times c_2)(c_1'\times c_2')\equiv (c_1c_1' \times c_2c_2')
  \end{mathpar}
  \caption{GTLC Cast Calculus Syntax, Type Precision Derivations and Precision Equivalence}
  \label{fig:gtlc-syntax}
\end{figure}

Next, we consider the axiomatic (in)equational reasoning principles
for terms: $\beta\eta$ equality and term precision in
Figure~\ref{fig:term-prec}.
%
We include standard CBV $\beta\eta$ rules for function and product
types, as well as equations stating that casts are given functorially.
%
Next, we have \emph{term} precision, an extension of
type precision to terms.
%
The form of the term precision rule is $\Delta \vdash M \ltdyn M' : c$
where $\Delta$ is a context where variables are assigned to type
precision derivations.
%
The judgment is only well formed when every use of $x : c$ for $c : A
\ltdyn A'$ is used with type $A$ in $M$ and $A'$ in $M'$ and similarly
the output types match $c$.
%
We elide the congruence rules for every type constructor, e.g., that
$M \ltdyn M'$ and $N \ltdyn N'$ that $M\,N \ltdyn M'\,N'$.
%
With such congruence rules, reflexivity $M \ltdyn M$ is
admissible. Transitivity, on the other hand, is intentionally not
taken as a primitive rule, matching the original formulation of the
dynamic gradual guarantee \cite{siek_et_al:LIPIcs:2015:5031}.
%
We include a rule that says that equivalent type precision derivations
$c \equiv c'$ are equivalent for the purposes of term precision.
%
% Removed retraction
%
% The next rule is the \emph{retraction} principle, which states that a
% downcast after an upcast is equivalent to doing nothing at all, since
% intuitively the upcasted value should already satisfy the type. Here
% $\equidyn$ means we require each is $\ltdyn$ the other, with
% reflexivity precision derivations.
%
Finally, we include 4 rules for reasoning about casts. Intuitively
these say that the upcast is a kind of \emph{least upper bound} and
dually that the downcast is a \emph{greatest lower bound}.

As a higher-order gradually typed language\DIFaddbegin \DIFadd{, }\DIFaddend we inherently have to deal
with two effects: errors and divergence. Errors arise from failing
casts, e.g. casting a number to dynamic to a function
$\dnc{\iarr}\upc{\inat} x$. Divergence arises because our dynamic type
allows us to encode untyped lambda calculus, and so we can encode the
$\Omega$ term with the help of casts $\Omega = (\lambda
x:\dyn. (\dnc{\iarr} x)x)(\upc{\iarr}(\lambda x:\dyn. (\dnc{\iarr}
x)x))$.

\begin{figure}
  \begin{mathpar}
  (\lambda x. M)(V) = M[V/x] \and (V : A \ra A') = \lambda x. V\,x\\

   \textrm{let } (x,y) = (V,V') \textrm{ in } N = N[V/x,V'/y] \and
   M[V:A\times A'/p] = \textrm{let } (x,y) = V \textrm{ in } M[(x,y)/p]

  \DIFdelbeginFL \DIFdelFL{\upc{(r(A))}M = M }%DIFDELCMD < \and
%DIFDELCMD <   %%%
\DIFdelFL{\upc{c'}\upc{c}M = \upc{cc'}M }%DIFDELCMD < \and
%DIFDELCMD <   %%%
\DIFdelFL{\dnc{(r(A))}M = M }%DIFDELCMD < \and
%DIFDELCMD <   %%%
\DIFdelFL{\dnc{c}\dnc{c'}M = \dnc{cc'}M
}\DIFdelendFL %DIF >  Removed these as they are not needed. The next rule implies that they are quasi-equivalent.
  %DIF >  \upc{(r(A))}M = M \and
  %DIF >  \upc{c'}\upc{c}M = \upc{cc'}M \and
  %DIF >  \dnc{(r(A))}M = M \and
  %DIF >  \dnc{c}\dnc{c'}M = \dnc{cc'}M

  \DIFdelbeginFL %DIFDELCMD < \inferrule
%DIFDELCMD <   %%%
\DIFdelendFL \DIFaddbeginFL \inferrule*[right=EquivTyPrec]
  \DIFaddendFL {\Delta\vdash M \ltdyn M' : c \and c \equiv c'}
  {\Delta\vdash M \ltdyn M' : c'}

  \DIFdelbeginFL %DIFDELCMD < \inferrule
%DIFDELCMD <   %%%
\DIFdelendFL \DIFaddbeginFL \inferrule*[right=ErrBot]
  \DIFaddendFL {}
  {\Delta \vdash \mho \ltdyn M : c}

  % Removed retraction
  % \inferrule
  % {}
  % {\dnc {c} \upc {c} M \equidyn M}

  \DIFdelbeginFL %DIFDELCMD < \inferrule*[Right=UpL]
%DIFDELCMD <   %%%
\DIFdelendFL \DIFaddbeginFL \inferrule*[right=UpL]
  \DIFaddendFL {M \ltdyn M' : cc_r}
  {\upc {c} M \ltdyn M' : c_r}

  \DIFdelbeginFL %DIFDELCMD < \inferrule*[Right=UpR]
%DIFDELCMD <   %%%
\DIFdelendFL \DIFaddbeginFL \inferrule*[right=UpR]
  \DIFaddendFL {M \ltdyn M' : c_l}
  {M \ltdyn \upc {c} M' : c_lc}

  \DIFdelbeginFL %DIFDELCMD < \inferrule*[Right=DnL]
%DIFDELCMD <   %%%
\DIFdelendFL \DIFaddbeginFL \inferrule*[right=DnL]
  \DIFaddendFL {M \ltdyn M' : c_r}
  {\dnc {c} M \ltdyn M' : cc_r}

  \DIFdelbeginFL %DIFDELCMD < \inferrule*[Right=DnR]
%DIFDELCMD <   %%%
\DIFdelendFL \DIFaddbeginFL \inferrule*[right=DnR]
  \DIFaddendFL {M \ltdyn M' : c_lc}
  {M \ltdyn \dnc {c} M' : c_l}
  \end{mathpar}
  \caption{Equality and Term Precision Rules (Selected)}
  \label{fig:term-prec}
\end{figure}

Our goal in the remainder of this work is to develop compositional
denotational semantics of types, terms, type and term precision from
which we can easily extract a big step semantics that satisfies
graduality and respects the equational theory of the calculus.

%% Here we describe the syntax and typing for the gradually-typed lambda calculus.
%% We also give the rules for syntactic type and term precision.
%% % We define four separate calculi: the normal gradually-typed lambda calculus, which we
%% % call the extensional or \emph{step-insensitive} lambda calculus ($\extlc$),
%% % as well as an \emph{intensional} lambda calculus
%% % ($\intlc$) whose syntax makes explicit the steps taken by a program.

%% Before diving into the details, let us give a brief overview of what we will define.
%% We begin with a gradually-typed lambda calculus $(\extlc)$, which is similar to
%% the normal call-by-value gradually-typed lambda calculus, but differs in that it
%% is actually a fragment of call-by-push-value specialized such that there are no
%% non-trivial computation types. We do this for convenience, as either way
%% we would need a distinction between values and effectful terms; the framework of
%% of call-by-push-value gives us a convenient language to define what we need.

%% We then show that composition of type precision derivations is admissible, as is
%% heterogeneous transitivity for term precision, so it will suffice to consider a new
%% language ($\extlcm$) in which we don't have composition of type precision derivations
%% or heterogeneous transitivity of term precision.

%% We then observe that all casts, except those between $\nat$ and $\dyn$
%% and between $\dyn \ra \dyn$ and $\dyn$, are admissible.
%% % (we can define the cast of a function type functorially using the casts for its domain and codomain).
%% This means it will be sufficient to consider a new language ($\extlcmm$) in which
%% instead of having arbitrary casts, we have injections from $\nat$ and
%% $\dyn \ra \dyn$ into $\dyn$, and case inspections from $\dyn$ to $\nat$ and
%% $\dyn$ to $\dyn \ra \dyn$.

%% From here, we define a \emph{step-sensitive} (also called \emph{intensional}) GSTLC,
%% so-named because it makes the intensional stepping behavior of programs explicit in the syntax.
%% This is accomplished by adding a syntactic ``later'' type and a
%% syntactic $\theta$ that maps terms of type later $A$ to terms of type $A$.
%% Finally, we define a \emph{quotiented} version of the step-sensitive language where
%% we add a rule that equates terms that are the same up to their stepping behavior.

%% % ---------------------------------------------------------------------------------------
%% % ---------------------------------------------------------------------------------------

%% \subsection{Syntax}

%% The language is based on Call-By-Push-Value \cite{levy01:phd}, and as such it has two kinds of types:
%% \emph{value types}, representing pure values, and \emph{computation types}, representing
%% potentially effectful computations.
%% In the language, all computation types have the form $\Ret A$ for some value type $A$.
%% Given a value $V$ of type $A$, the term $\ret V$ views $V$ as a term of computation type $\Ret A$.
%% Given a term $M$ of computation type $B$, the term $\bind{x}{M}{N}$ should be thought of as
%% running $M$ to a value $V$ and then continuing as $N$, with $V$ in place of $x$.


%% We also have value contexts and computation contexts, where the latter can be viewed
%% as a pair consisting of (1) a stoup $\Sigma$, which is either empty or a hole of type $B$,
%% and (2) a (potentially empty) value context $\Gamma$.

%% \begin{align*} % TODO is hole a term?
%%   &\text{Value Types } A := \nat \alt \,\dyn \alt (A \ra A') \\
%%   &\text{Computation Types } B := \Ret A \\
%%   &\text{Value Contexts } \Gamma := \cdot \alt (\Gamma, x : A) \\
%%   &\text{Computation Contexts } \Delta := \cdot \alt \hole B \alt \Delta , x : A \\
%%   &\text{Values } V :=  \zro \alt \suc\, V \alt \lda{x}{M} \alt \up{A}{B} V \\ 
%%   &\text{Terms } M, N := \err_B \alt \matchnat {V} {M} {n} {M'} \\ 
%%   &\quad\quad \alt \ret {V} \alt \bind{x}{M}{N} \alt V_f\, V_x \alt \dn{A}{B} M 
%% \end{align*}

%% The value typing judgment is written $\hasty{\Gamma}{V}{A}$ and 
%% the computation typing judgment is written $\hasty{\Delta}{M}{B}$.

%% \begin{comment}
%% We define substitution for value contexts by the following rules:

%% \begin{mathpar}
%%   \inferrule*
%%   { \gamma : \Gamma' \to \Gamma \and 
%%     \hasty{\Gamma'}{V}{A}}
%%   { (\gamma , V/x ) \colon \Gamma' \to \Gamma , x : A }

%%   \inferrule*
%%   {}
%%   {\cdot \colon \cdot \to \cdot}
%% \end{mathpar}

%% We define substitution for computation contexts by the following rules:

%% \begin{mathpar}
%%     \inferrule*
%%     { \delta : \Delta' \to \Delta \and 
%%       \hasty{\Delta'|_V}{V}{A}}
%%     { (\delta , V/x ) \colon \Delta' \to \Delta , x : A }

%%     \inferrule*
%%     {}
%%     {\cdot \colon \cdot \to \cdot}

%%     \inferrule*
%%     {\hasty{\Delta'}{M}{B}}
%%     {M \colon \Delta' \to \hole{B}}
%% \end{mathpar}
%% \end{comment}

%% The typing rules are as expected, with a cast between $A$ to $B$ allowed only when $A \ltdyn B$.
%% Notice that the upcast of a value is a value, since it always succeeds, while the downcast
%% of a value is a computation, since it may fail.

%% \begin{mathpar}
%%     % Var
%%     \inferrule*{ }{\hasty {\cdot, \Gamma, x : A, \Gamma'} x A}

%%     % Err
%%     \inferrule*{ }{\hasty {\cdot, \Gamma} {\err_B} B} 

%%     % Zero and suc
%%     \inferrule*{ }{\hasty \Gamma \zro \nat}

%%     \inferrule*{\hasty \Gamma V \nat} {\hasty \Gamma {\suc\, V} \nat}

%%     % Match-nat
%%     \inferrule*
%%     {\hasty \Gamma V \nat \and 
%%      \hasty \Delta M B \and \hasty {\Delta, n : \nat} {M'} B}
%%     {\hasty \Delta {\matchnat {V} {M} {n} {M'}} B}

%%     % Lambda
%%     \inferrule* 
%%     {\hasty {\cdot, \Gamma, x : A} M {\Ret A'}} 
%%     {\hasty \Gamma {\lda x M} {A \ra A'}}

%%     % App
%%     \inferrule*
%%     {\hasty \Gamma {V_f} {A \ra A'} \and \hasty \Gamma {V_x} A}
%%     {\hasty {\cdot , \Gamma} {V_f \, V_x} {\Ret A'}}

%%     % Ret
%%     \inferrule*
%%     {\hasty \Gamma V A}
%%     {\hasty {\cdot , \Gamma} {\ret\, V} {\Ret A}}
%%     % TODO should this involve a Delta?

%%     % Bind
%%     \inferrule*
%%     {\hasty \Delta M {\Ret A} \and \hasty{\cdot , \Delta|_V , x : A}{N}{B} } % Need x : A in context
%%     {\hasty {\Delta} {\bind{x}{M}{N}} {B}}

%%     % Upcast
%%     \inferrule*
%%     {A \ltdyn A' \and \hasty \Gamma V A}
%%     {\hasty \Gamma {\up A {A'} V} {A'} }

%%     % Downcast
%%     % \inferrule*
%%     % {A \ltdyn A' \and \hasty {\Gamma} V {A'}}
%%     % {\hasty {\cdot, \Gamma} {\dn A {A'} V} {\Ret A}}

%%     \inferrule* % TODO is this correct?
%%     {B \ltdyn B' \and \hasty {\Delta} {M} {B'}}
%%     {\hasty {\Delta} {\dn B {B'} M} {B}}

%% \end{mathpar}


%% In the equational theory, we have $\beta$ and $\eta$ laws for function type,
%% as well a $\beta$ and $\eta$ law for $\Ret A$.

%% % TODO do we need to add a substitution rule here?
%% \begin{mathpar}
%%   % Function Beta and Eta
%%   \inferrule*
%%   {\hasty {\cdot, \Gamma, x : A} M {\Ret A'} \and \hasty \Gamma V A}
%%   {(\lda x M)\, V = M[V/x]}

%%   \inferrule*
%%   {\hasty \Gamma V {A \ra A}}
%%   {\Gamma \vdash V = \lda x {V\, x}}

%%   % Ret Beta and Eta
%%   \inferrule*
%%   {}
%%   {(\bind{x}{\ret\, V}{N}) = N[V/x]}

%%   \inferrule*
%%   {\hasty {\hole{\Ret A} , \Gamma} {M} {B}}
%%   {\hole{\Ret A}, \Gamma \vdash M = (\bind{x}{\bullet}{M[\ret\, x]})}

%%   % Match-nat Beta
%%   \inferrule*
%%   {\hasty \Delta M B \and \hasty {\Delta, n : \nat} {M'} B}
%%   {\matchnat{\zro}{M}{n}{M'} = M}

%%   \inferrule*
%%   {\hasty \Gamma V \nat \and 
%%    \hasty \Delta M B \and \hasty {\Delta, n : \nat} {M'} B}
%%   {\matchnat{\suc\, V}{M}{n}{M'} = M'}

%%   % Match-nat Eta
%%   % This doesn't build in substitution
%%   \inferrule*
%%   {\hasty {\Delta , x : \nat} M A}
%%   {M = \matchnat{x} {M[\zro / x]} {n} {M[(\suc\, n) / x]}}



%% \end{mathpar}

%% % ---------------------------------------------------------------------------------------
%% % ---------------------------------------------------------------------------------------

%% \subsection{Type Precision}

%% The type precision rules specify what it means for a type $A$ to be more precise than $A'$.
%% We have reflexivity rules for $\dyn$ and $\nat$, as well as rules that $\nat$ is more precise than $\dyn$
%% and $\dyn \ra \dyn$ is more precise than $\dyn$.
%% We also have a transitivity rule for composition of type precision,
%% and also a rule for function types stating that given $A_i \ltdyn A'_i$ and $A_o \ltdyn A'_o$, we can prove
%% $A_i \ra A_o \ltdyn A'_i \ra A'_o$.
%% Finally, we can lift a relation on value types $A \ltdyn A'$ to a relation $\Ret A \ltdyn \Ret A'$ on
%% computation types.

%% \begin{mathpar}
%%   \inferrule*[right = \dyn]
%%     { }{\dyn \ltdyn\, \dyn}

%%   \inferrule*[right = \nat]
%%     { }{\nat \ltdyn \nat}

%%   \inferrule*[right = $\ra$]
%%     {A_i \ltdyn A'_i \and A_o \ltdyn A'_o }
%%     {(A_i \ra A_o) \ltdyn (A'_i \ra A'_o)}

%%   \inferrule*[right = $\textsf{Inj}_\nat$]
%%     { }{\nat \ltdyn\, \dyn}

%%   \inferrule*[right=$\textsf{Inj}_{\ra}$]
%%     { }
%%     {(\dyn \ra \dyn) \ltdyn\, \dyn}

%%   \inferrule*[right=ValTrans]
%%     {A \ltdyn A' \and A' \ltdyn A''}
%%     {A \ltdyn A''}

%%   \inferrule*[right=CompTrans]
%%     {B \ltdyn B' \and B' \ltdyn B''}
%%     {B \ltdyn B''}

%%   \inferrule*[right=$\Ret{}$]
%%     {A \ltdyn A'}
%%     {\Ret {A} \ltdyn \Ret {A'}}

%%     % TODO are there other rules needed for computation types?

  
%% \end{mathpar}

%% % Type precision derivations
%% Note that as a consequence of this presentation of the type precision rules, we
%% have that if $A \ltdyn A'$, there is a unique precision derivation that witnesses this.
%% As in previous work, we go a step farther and make these derivations first-class objects,
%% known as \emph{type precision derivations}.
%% Specifically, for every $A \ltdyn A'$, we have a derivation $c : A \ltdyn A'$ that is constructed
%% using the rules above. For instance, there is a derivation $\dyn : \dyn \ltdyn \dyn$, and a derivation
%% $\nat : \nat \ltdyn \nat$, and if $c_i : A_i \ltdyn A_i$ and $c_o : A_o \ltdyn A'_o$, then
%% there is a derivation $c_i \ra c_o : (A_i \ra A_o) \ltdyn (A'_i \ra A'_o)$. Likewise for
%% the remaining rules. The benefit to making these derivations explicit in the syntax is that we
%% can perform induction over them.
%% Note also that for any type $A$, we use $A$ to denote the reflexivity derivation that $A \ltdyn A$,
%% i.e., $A : A \ltdyn A$.
%% Finally, observe that for type precision derivations $c : A \ltdyn A'$ and $c' : A' \ltdyn A''$, we
%% can define (via the rule ValComp) their composition $c \relcomp c' : A \ltdyn A''$.
%% The same holds for computation type precision derivations.
%% This notion will be used below in the statement of transitivity of the term precision relation.

%% % ---------------------------------------------------------------------------------------
%% % ---------------------------------------------------------------------------------------

%% \subsection{Term Precision}

%% We allow for a \emph{heterogeneous} term precision judgment on terms values $V$ of type
%% $A$ and $V'$ of type $A'$ provided that $A \ltdyn A'$ holds. Likewise, for computation
%% types $B \ltdyn B'$, if $M$ has type $B$ and $M'$ has type $B'$, we can form the judgment
%% that $M \ltdyn M'$.

%% % Type precision contexts
%% % TODO should we include the formal definitions of value and computation type precision contexts?
%% In order to deal with open terms, we will need the notion of a type precision \emph{context}, which we denote
%% $\gamlt$. This is similar to a normal context but instead of mapping variables to types,
%% it maps variables $x$ to related types $A \ltdyn A'$, where $x$ has type $A$ in the left-hand term
%% and $B$ in the right-hand term. We may also write $x : d$ where $d : A \ltdyn A'$ to indicate this.
%% Similarly, we have computation type precision contexts $\Delta^\ltdyn$. Similar to ``normal'' computation
%% type precision contexts $\Delta$, these consist of (1) a stoup $\Sigma$ which is either empty or
%% has a hole $\hole{d}$ for some computation type precision derivation $d$, and (2) a value type precision context
%% $\Gamma^\ltdyn$.

%% % An equivalent way of thinking of type precision contexts is as a pair of ``normal" typing
%% % contexts $\Gamma, \Gamma'$ with the same domain such that $\Gamma(x) \ltdyn \Gamma'(x)$ for
%% % each $x$ in the domain.
%% % We will write $\gamlt : \Gamma \ltdyn \Gamma'$ when we want to emphasize the pair of contexts.
%% % Conversely, if we are given $\gamlt$, we write $\gamlt_l$ and $\gamlt_r$ for the normal typing contexts on each side.

%% An equivalent way of thinking of a type precision context $\gamlt$ is as a
%% pair of ``normal" typing contexts, $\gamlt_l$ and $\gamlt_r$, with the same
%% domain and such that $\gamlt_l(x) \ltdyn \gamlt_r(x)$ for each $x$ in the domain.
%% We will write $\gamlt : \gamlt_l \ltdyn \gamlt_r$ when we want to emphasize the pair of contexts.

%% As with type precision derivations, we write $\Gamma$ to mean the ``reflexivity" type precision context
%% $\Gamma : \Gamma \ltdyn \Gamma$.
%% Concretely, this consists of reflexivity type precision derivations $\Gamma(x) \ltdyn \Gamma(x)$ for
%% each $x$ in the domain of $\Gamma$.
%% Similarly, we also have reflexivity for computation type precision contexts.
%% %
%% Furthermore, we write $\gamlt_1 \relcomp \gamlt_2$ to denote the ``composition'' of $\gamlt_1$ and $\gamlt_2$
%% --- that is, the precision context whose value at $x$ is the type precision derivation
%% $\gamlt_1(x) \relcomp \gamlt_2(x)$. This of course assumes that each of the type precision
%% derivations is composable, i.e., that the RHS of $\gamlt_1(x)$ is the same as the left-hand side of $\gamlt_2(x)$.
%% We define the same for computation type precision contexts $\deltalt_1$ and $\deltalt_2$,
%% provided that both the computation type precision contexts have the same ``shape'', which is defined as
%% (1) either the stoup is empty in both, or the stoup has a hole in both, say $\hole{d}$ and $\hole{d'}$
%% where $d$ and $d'$ are composable, and (2) their value type precision contexts are composable as described above.

%% The rules for term precision come in two forms. We first have the \emph{congruence} rules,
%% one for each term constructor. These assert that the term constructors respect term precision.
%% The congruence rules are as follows:

%% \begin{mathpar}

%%   \inferrule*[right = Var]
%%     { c : A \ltdyn B \and \gamlt(x) = (A, B) } 
%%     { \etmprec {\gamlt} x x c }

%%   \inferrule*[right = Zro]
%%     { } {\etmprec \gamlt \zro \zro \nat }

%%   \inferrule*[right = Suc]
%%     { \etmprec \gamlt V {V'} \nat } {\etmprec \gamlt {\suc\, V} {\suc\, V'} \nat}

%%   \inferrule*[right = MatchNat]
%%   {\etmprec \gamlt V {V'} \nat \and 
%%     \etmprec \deltalt M {M'} d \and \etmprec {\deltalt, n : \nat} {N} {N'} d}
%%   {\etmprec \deltalt {\matchnat {V} {M} {n} {N}} {\matchnat {V'} {M'} {n} {N'}} d}

%%   \inferrule*[right = Lambda]
%%     { c_i : A_i \ltdyn A'_i \and 
%%       c_o : A_o \ltdyn A'_o \and 
%%       \etmprec {\cdot , \gamlt , x : c_i} {M} {M'} {\Ret c_o} } 
%%     { \etmprec \gamlt {\lda x M} {\lda x {M'}} {(c_i \ra c_o)} }

%%   \inferrule*[right = App]
%%     { c_i : A_i \ltdyn A'_i \and
%%       c_o : A_o \ltdyn A'_o \\\\
%%       \etmprec \gamlt {V_f} {V_f'} {(c_i \ra c_o)} \and
%%       \etmprec \gamlt {V_x} {V_x'} {c_i}
%%     } 
%%     { \etmprec {\cdot , \gamlt} {V_f\, V_x} {V_f'\, V_x'} {\Ret {c_o}}}

%%   \inferrule*[right = Ret]
%%     {\etmprec {\gamlt} V {V'} c}
%%     {\etmprec {\cdot , \gamlt} {\ret\, V} {\ret\, V'} {\Ret c}}

%%   \inferrule*[right = Bind]
%%     {\etmprec {\deltalt} {M} {M'} {\Ret c} \and 
%%      \etmprec {\cdot , \deltalt|_V , x : c} {N} {N'} {d} }
%%     {\etmprec {\deltalt} {\bind {x} {M} {N}} {\bind {x} {M'} {N'}} {d}}
%% \end{mathpar}

%% We then have additional equational axioms, including transitivity, $\beta$ and $\eta$ laws, and
%% rules characterizing upcasts as least upper bounds, and downcasts as greatest lower bounds.

%% We write $M \equidyn N$ to mean that both $M \ltdyn N$ and $N \ltdyn M$.

%% % TODO adapt these for value/computation distinction
%% % TODO substitution rules for values and terms?
%% \begin{mathpar}
%%   \inferrule*[right = $\err$]
%%     { \hasty {\deltalt_l} M B }
%%     {\etmprec {\Delta} {\err_B} M B}

%%   \inferrule*[right = Transitivity]
%%     { d : B \ltdyn B' \and d' : B' \ltdyn B'' \\\\
%%      \etmprec {\deltalt_1} {M} {M'} {d} \and
%%      \etmprec {\deltalt_2} {M'} {M''} {d'} } 
%%     {\etmprec {\deltalt_1 \relcomp \deltalt_2} {M} {M''} {d \relcomp d'} }


%%   \inferrule*[right = $\beta$-fun]
%%     { \hasty {\cdot, \Gamma, x : A_i} M {\Ret A_o} \and
%%       \hasty {\Gamma} V {A_i} } 
%%     { \etmequidyn {\cdot, \Gamma} {(\lda x M)\, V} {M[V/x]} {\Ret A_o} }

%%   \inferrule*[right = $\eta$-fun]
%%     { \hasty {\Gamma} {V} {A_i \ra A_o} } 
%%     { \etmequidyn \Gamma {\lda x (V\, x)} V {A_i \ra A_o} }

%%   % Match-nat beta and eta



%%   \inferrule*[right = $\beta$-ret]
%%     {}
%%     {\bind{x}{\ret\, V}{N} \equidyn N[V/x]}

%%   \inferrule*[right = $\eta$-ret]
%%     {\hasty {\hole{\Ret A} , \Gamma} {M} {B}}
%%     {\hole{\Ret A}, \Gamma \vdash M \equidyn \bind{x}{\bullet}{M[\ret\, x]}}


%%   % Could specify \gamlt : \Gamma \ltdyn \Gamma'
%%   % and then we wouldn't need to say l and r

%%   \inferrule*[right = UpR]
%%     { d : A \ltdyn A' \and 
%%       \hasty {\Delta} {M} {A} } 
%%     { \etmprec {\Delta} {M} {\up {A} {A'} M} {d}  }

%%   \inferrule*[right = UpL]
%%     { d : A \ltdyn A' \and
%%       \etmprec {\deltalt} {M} {N} {d} } 
%%     { \etmprec {\deltalt} {\up {A} {A'} M} {N} {A'} }

%%   \inferrule*[right = DnL]
%%     { d : B \ltdyn B' \and 
%%       \hasty {\Delta} {M} {B'} } 
%%     { \etmprec {\Delta} {\dn {B} {B'} M} {M} {d} }

%%   \inferrule*[right = DnR]
%%     { d : B \ltdyn B' \and
%%       \etmprec {\deltalt} {M} {N} {d} } 
%%     { \etmprec {\deltalt} {M} {\dn {B} {B'} N} {B} }
%% \end{mathpar}

%% % TODO explain the least upper bound/greatest lower bound rules
%% The rules UpR, UpL, DnL, and DnR were introduced in \cite{new-licata18} as a means
%% of cleanly axiomatizing the intended behavior of casts in a way that
%% doesn't depend on the specific constructs of the language.
%% Intuitively, rule UpR says that the upcast of $M$ is an upper bound for $M$
%% in that $M$ may error more, and UpL says that the upcast is the \emph{least}
%% such upper bound, in that it errors more than any other upper bound for $M$.
%% Conversely, DnL says that the downcast of $M$ is a lower bound, and DnR says
%% that it is the \emph{greatest} lower bound.
%% % These rules provide a clean axiomatization of the behavior of casts that doesn't
%% % depend on the specific constructs of the language.

%% % ---------------------------------------------------------------------------------------
%% % ---------------------------------------------------------------------------------------
%% \subsection{Removing Transitivity as a Primitive}

%% The first observation we make is that transitivity of type precision, and heterogeneous
%% transitivity of term precision, are admissible. That is, consider a related language which
%% is the same as $\extlc$ except that we have removed the composition rule for type precision and
%% the heterogeneous transitivity rule for type precision. Denote this language by $\extlcm$.
%% We claim that in this new language, the rules we removed are derivable from the remaining rules.

%% To see this, suppose $\gamlt : \Gamma \ltdyn \Gamma'$ and $d : A \ltdyn A'$, and that
%%  $\etmprec {\gamlt} {V} {V'} {d}$, as shown in the diagram below:

%% % https://q.uiver.app/?q=WzAsNCxbMCwwLCJcXEdhbW1hIl0sWzAsMSwiXFxHYW1tYSciXSxbMSwwLCJBIl0sWzEsMSwiQSciXSxbMCwxLCJcXGx0ZHluIiwzLHsic3R5bGUiOnsiYm9keSI6eyJuYW1lIjoibm9uZSJ9LCJoZWFkIjp7Im5hbWUiOiJub25lIn19fV0sWzIsMywiXFxsdGR5biIsMyx7InN0eWxlIjp7ImJvZHkiOnsibmFtZSI6Im5vbmUifSwiaGVhZCI6eyJuYW1lIjoibm9uZSJ9fX1dLFswLDIsIlYiXSxbMSwzLCJWJyJdLFs2LDcsIlxcbHRkeW4iLDMseyJzaG9ydGVuIjp7InNvdXJjZSI6MjAsInRhcmdldCI6MjB9LCJzdHlsZSI6eyJib2R5Ijp7Im5hbWUiOiJub25lIn0sImhlYWQiOnsibmFtZSI6Im5vbmUifX19XV0=
%% \[\begin{tikzcd}[ampersand replacement=\&]
%% 	\Gamma \& A \\
%% 	{\Gamma'} \& {A'}
%% 	\arrow["\ltdyn"{marking}, draw=none, from=1-1, to=2-1]
%% 	\arrow["\ltdyn"{marking}, draw=none, from=1-2, to=2-2]
%% 	\arrow[""{name=0, anchor=center, inner sep=0}, "V", from=1-1, to=1-2]
%% 	\arrow[""{name=1, anchor=center, inner sep=0}, "{V'}", from=2-1, to=2-2]
%% 	\arrow["\ltdyn"{marking}, draw=none, from=0, to=1]
%% \end{tikzcd}\]

%% Now note that this is equivalent, by the cast rule UpL, to
%% $\etmprec {\Gamma'} {\up{A}{A'} V} {V'} {A'}$,
%% where as noted above, $\Gamma'$ refers to the context $\Gamma'$ viewed as a reflexivity
%% precision context and likewise the $A'$ at the end refers to the reflexivity derivation $A' \ltdyn A'$.

%% % https://q.uiver.app/?q=WzAsMixbMCwwLCJcXEdhbW1hJyJdLFsxLDAsIkEnIl0sWzAsMSwiXFx1cCB7QX0ge0EnfSBWIiwwLHsiY3VydmUiOi0yfV0sWzAsMSwiViciLDIseyJjdXJ2ZSI6Mn1dLFsyLDMsIlxcbHRkeW4iLDMseyJzaG9ydGVuIjp7InNvdXJjZSI6MjAsInRhcmdldCI6MjB9LCJzdHlsZSI6eyJib2R5Ijp7Im5hbWUiOiJub25lIn0sImhlYWQiOnsibmFtZSI6Im5vbmUifX19XV0=
%% \[\begin{tikzcd}[ampersand replacement=\&]
%% 	{\Gamma'} \& {A'}
%% 	\arrow[""{name=0, anchor=center, inner sep=0}, "{\up {A} {A'} V}", curve={height=-12pt}, from=1-1, to=1-2]
%% 	\arrow[""{name=1, anchor=center, inner sep=0}, "{V'}"', curve={height=12pt}, from=1-1, to=1-2]
%% 	\arrow["\ltdyn"{marking}, draw=none, from=0, to=1]
%% \end{tikzcd}\]

%% Now consider the situation shown below:

%% % https://q.uiver.app/?q=WzAsNixbMCwwLCJcXEdhbW1hIl0sWzAsMSwiXFxHYW1tYSciXSxbMCwyLCJcXEdhbW1hJyciXSxbMiwwLCJBIl0sWzIsMSwiQSciXSxbMiwyLCJBJyciXSxbMiw1LCJWJyciXSxbMSw0LCJWJyJdLFswLDMsIlYiXSxbMyw0LCJcXGx0ZHluIiwzLHsic3R5bGUiOnsiYm9keSI6eyJuYW1lIjoibm9uZSJ9LCJoZWFkIjp7Im5hbWUiOiJub25lIn19fV0sWzQsNSwiXFxsdGR5biIsMyx7InN0eWxlIjp7ImJvZHkiOnsibmFtZSI6Im5vbmUifSwiaGVhZCI6eyJuYW1lIjoibm9uZSJ9fX1dLFswLDEsIlxcbHRkeW4iLDMseyJzdHlsZSI6eyJib2R5Ijp7Im5hbWUiOiJub25lIn0sImhlYWQiOnsibmFtZSI6Im5vbmUifX19XSxbMSwyLCIiLDEseyJzdHlsZSI6eyJib2R5Ijp7Im5hbWUiOiJub25lIn0sImhlYWQiOnsibmFtZSI6Im5vbmUifX19XSxbMSwyLCJcXGx0ZHluIiwzLHsic3R5bGUiOnsiYm9keSI6eyJuYW1lIjoibm9uZSJ9LCJoZWFkIjp7Im5hbWUiOiJub25lIn19fV0sWzgsNywiXFxsdGR5biIsMyx7InNob3J0ZW4iOnsic291cmNlIjoyMCwidGFyZ2V0IjoyMH0sInN0eWxlIjp7ImJvZHkiOnsibmFtZSI6Im5vbmUifSwiaGVhZCI6eyJuYW1lIjoibm9uZSJ9fX1dLFs3LDYsIlxcbHRkeW4iLDMseyJzaG9ydGVuIjp7InNvdXJjZSI6MjAsInRhcmdldCI6MjB9LCJzdHlsZSI6eyJib2R5Ijp7Im5hbWUiOiJub25lIn0sImhlYWQiOnsibmFtZSI6Im5vbmUifX19XV0=
%% \[\begin{tikzcd}[ampersand replacement=\&]
%% 	\Gamma \&\& A \\
%% 	{\Gamma'} \&\& {A'} \\
%% 	{\Gamma''} \&\& {A''}
%% 	\arrow[""{name=0, anchor=center, inner sep=0}, "{V''}", from=3-1, to=3-3]
%% 	\arrow[""{name=1, anchor=center, inner sep=0}, "{V'}", from=2-1, to=2-3]
%% 	\arrow[""{name=2, anchor=center, inner sep=0}, "V", from=1-1, to=1-3]
%% 	\arrow["\ltdyn"{marking}, draw=none, from=1-3, to=2-3]
%% 	\arrow["\ltdyn"{marking}, draw=none, from=2-3, to=3-3]
%% 	\arrow["\ltdyn"{marking}, draw=none, from=1-1, to=2-1]
%% 	\arrow[draw=none, from=2-1, to=3-1]
%% 	\arrow["\ltdyn"{marking}, draw=none, from=2-1, to=3-1]
%% 	\arrow["\ltdyn"{marking}, draw=none, from=2, to=1]
%% 	\arrow["\ltdyn"{marking}, draw=none, from=1, to=0]
%% \end{tikzcd}\]


%% Using the above observation, we have that the above is equivalent to

%% % https://q.uiver.app/?q=WzAsNCxbMCwwLCJcXEdhbW1hJyJdLFswLDEsIlxcR2FtbWEnJyJdLFsyLDAsIkEnIl0sWzIsMSwiQScnIl0sWzAsMiwiXFx1cCB7QX0ge0EnfSBWIiwwLHsiY3VydmUiOi0yfV0sWzAsMiwiViciLDIseyJjdXJ2ZSI6Mn1dLFsxLDMsIlYnJyIsMix7ImN1cnZlIjoyfV0sWzAsMSwiXFxsdGR5biIsMyx7InN0eWxlIjp7ImJvZHkiOnsibmFtZSI6Im5vbmUifSwiaGVhZCI6eyJuYW1lIjoibm9uZSJ9fX1dLFsyLDMsIlxcbHRkeW4iLDMseyJzdHlsZSI6eyJib2R5Ijp7Im5hbWUiOiJub25lIn0sImhlYWQiOnsibmFtZSI6Im5vbmUifX19XSxbNCw1LCJcXGx0ZHluIiwzLHsic2hvcnRlbiI6eyJzb3VyY2UiOjIwLCJ0YXJnZXQiOjIwfSwic3R5bGUiOnsiYm9keSI6eyJuYW1lIjoibm9uZSJ9LCJoZWFkIjp7Im5hbWUiOiJub25lIn19fV0sWzUsNiwiXFxsdGR5biIsMyx7InNob3J0ZW4iOnsic291cmNlIjoyMCwidGFyZ2V0IjoyMH0sInN0eWxlIjp7ImJvZHkiOnsibmFtZSI6Im5vbmUifSwiaGVhZCI6eyJuYW1lIjoibm9uZSJ9fX1dXQ==
%% \[\begin{tikzcd}[ampersand replacement=\&]
%% 	{\Gamma'} \&\& {A'} \\
%% 	{\Gamma''} \&\& {A''}
%% 	\arrow[""{name=0, anchor=center, inner sep=0}, "{\up {A} {A'} V}", curve={height=-12pt}, from=1-1, to=1-3]
%% 	\arrow[""{name=1, anchor=center, inner sep=0}, "{V'}"', curve={height=12pt}, from=1-1, to=1-3]
%% 	\arrow[""{name=2, anchor=center, inner sep=0}, "{V''}"', curve={height=12pt}, from=2-1, to=2-3]
%% 	\arrow["\ltdyn"{marking}, draw=none, from=1-1, to=2-1]
%% 	\arrow["\ltdyn"{marking}, draw=none, from=1-3, to=2-3]
%% 	\arrow["\ltdyn"{marking}, draw=none, from=0, to=1]
%% 	\arrow["\ltdyn"{marking}, draw=none, from=1, to=2]
%% \end{tikzcd}\]

%% % TODO finish the explanation


%% % ---------------------------------------------------------------------------------------
%% % ---------------------------------------------------------------------------------------

%% \subsection{Removing Casts as Primitives}

%% % We now observe that all casts, except those between $\nat$ and $\dyn$
%% % and between $\dyn \ra \dyn$ and $\dyn$, are admissible, in the sense that
%% % we can start from $\extlcm$, remove casts except the aforementioned ones,
%% % and in the resulting language we will be able to derive the other casts.

%% We now observe that all casts, except those between $\nat$ and $\dyn$
%% and between $\dyn \ra \dyn$ and $\dyn$, are admissible.
%% That is, consider a new language ($\extlcmm$) in which
%% instead of having arbitrary casts, we have injections from $\nat$ and
%% $\dyn \ra \dyn$ into $\dyn$, and case inspections from $\dyn$ to $\nat$ and
%% $\dyn$ to $\dyn \ra \dyn$. We claim that in $\extlcmm$, all of the casts
%% present in $\extlcm$ are derivable.
%% It will suffice to verify that casts for function type are derivable.
%% This holds because function casts are constructed inductively from the casts
%% of their domain and codomain. The base case is one of the casts involving $\nat$
%% or $\dyn \ra \dyn$ which are present in $\extlcmm$ as injections and case inspections.


%% The resulting calculus $\extlcmm$ now lacks transitivity of type precision,
%% heterogeneous transitivity of term precision, and arbitrary casts as primitive
%% notions.

%% \begin{align*}
%%   &\text{Value Types } A := \nat \alt \dyn \alt (A \ra A') \\
%%   &\text{Computation Types } B := \Ret A \\
%%   &\text{Value Contexts } \Gamma := \cdot \alt (\Gamma, x : A) \\
%%   &\text{Computation Contexts } \Delta := \cdot \alt \hole B \alt \Delta , x : A \\
%%   &\text{Values } V :=  \zro \alt \suc\, V \alt \lda{x}{M} \alt \injnat V \alt \injarr V \\ 
%%   &\text{Terms } M, N := \err_B \alt \ret {V} \alt \bind{x}{M}{N}
%%     \alt V_f\, V_x \alt
%%     \\ & \quad\quad \casenat{V}{M_{no}}{n}{M_{yes}} 
%%     \alt \casearr{V}{M_{no}}{f}{M_{yes}}
%% \end{align*}

%% In this setting, rather than type precision, it makes more sense to
%% speak of arbitrary \emph{monotone relations} on types, which we denote by $A \rel A'$.
%% We have relations on value types, as well as on computation types. We also have
%% value relation contexts and computation relation contexts, analogous to the value type
%% precision contexts and computation type precision contexts from before.

%% \begin{align*}
%%   &\text{Value Relations } R := \nat \alt \dyn \alt (R \ra R) \alt\, \dyn\, R(V_1, V_2)\\
%%   &\text{Computation Relations } S := \li R \\
%%   &\text{Value Relation Contexts } \Gamma^{\rel} := \cdot \alt \Gamma^{\rel} , A^{\rel} (x_l : A_l , x_r : A_r)\\
%%   &\text{Computation Relation Contexts } \Delta^{\rel} := \cdot \alt \hole{B^{\rel}} \alt 
%%     \Delta^{\rel} , A^{\rel} (x_l : A_l , x_r : A_r)   \\
%% \end{align*}

%% % TODO rules for relations
%% The forms for relations are as follows:

%% \begin{align*}
%%   A^{\rel}      &\colon A_l      \rel A_r \\
%%   \Gamma^{\rel} &\colon \Gamma_l \rel \Gamma_r \\
%%   B^{\rel}      &\colon B_l      \rel B_r \\
%%   \Delta^{\rel} &\colon \Delta_l \rel \Delta_r
%% \end{align*}



%% Figure \ref{fig:relation-rules} shows the rules for relations. We show only those for value types;
%% the corresponding computation type relation rules are analogous.
%% The rules for relations are as follows. First, we require relations to be reflexive.
%% We also require that they are \emph{profunctorial}, in the sense that a relation between
%% $A$ and $A'$ is closed under the ``homogeneous'' relations on both sides.
%% We also require that they satisfy a substitution principle.

%% \begin{figure}
%%   \begin{mathpar}
%%     \inferrule*[right = Reflexivity]
%%     {\hasty \Gamma V A}
%%     {\refl(\Gamma) \vdash \refl(A)(V, V)}

%%     \inferrule*[right = Profunctoriality]
%%     { \refl(\Gamma^{\rel}_l) \vdash  \refl(A^{\rel}_l) (V_l' , V_l) \\\\ 
%%         \Gamma^{\rel}    \vdash    A^{\rel}    (V_l  , V_r) \\\\
%%       \refl(\Gamma^{\rel}_r) \vdash  \refl(A^{\rel}_r) (V_r  , V_r')
%%     }
%%     {\Gamma^{\rel} \vdash A^{\rel} (V_l', V_r')}

%%     \inferrule*[right = Subst]
%%     { \Gamma'^{\rel} \vdash \Gamma^{\rel} (\gamma_l, \gamma_r) \\\\
%%       \Gamma^{\rel} A^{\rel} (V_l, V_r)
%%     }
%%     {\Gamma'^{\rel} \vdash A^{\rel} (V_l[\gamma_l] , V_r[\gamma_r]) }

%%     % \inferrule*[right = TermSubst]
%%     % { \Delta'^{\rel} \vdash \Delta^{\rel} (\delta_l, \delta_r) \\\\
%%     %   \Delta^{\rel} B^{\rel} (M_l, M_r)
%%     % }
%%     % {\Delta'^{\rel} \vdash B^{\rel} (M_l[\delta_l] , M_r[\delta_r]) }

%%   \end{mathpar}
%%   \caption{Rules for value type relations. The rules for computation type relations are analogous.}
%%   \label{fig:relation-rules}
%% \end{figure}

%% We also have a rule for the restriction of a relation along a function,
%% and we have a rule characterizing relation at function type. The latter states that
%% if under the assumption that $x$ is related to $x'$ by $A^{\rel}$, we can show that $M$
%% is related to $M'$ by $\li A'^{\rel}$, then we have that $\lda{x}{M}$ is related to
%% $\lda{x'}{M'}$ by $A^{\rel} \ra A'^{\rel}$.

%% \begin{mathpar}
%%   \mprset{fraction={===}}

%%   % \inferrule*[]
%%   % { A^{\rel}  (x_l, x_r) \vdash A^{\rel} (V_l, V_r) }
%%   % { A'^{\rel} (x_l, x_r) \vdash A^{\rel} (V_l, V_r)(x_l, x_r) }

%%   \inferrule*[right = Restriction]
%%   { \Gamma^{\rel} \vdash A^{\rel} (V_l (V_l'), V_r (V_r')) }
%%   { \Gamma^{\rel} \vdash (A^{\rel} (V_l, V_r)) (V_l', V_r') }

%%   \inferrule*[right = $\text{Rel}_\ra$]
%%   { A^{\rel} (x, x') \vdash (\li A'^{\rel})(M , M') }
%%   {  \vdash (A^{\rel} \ra A'^{\rel}) (\lda{x}{M}) , (\lda{x'}{M'})}

%% \end{mathpar}



%% % New rules
%% Figure \ref{fig:extlc-minus-minus-typing} shows the new typing rules,
%% and Figure \ref{fig:extlc-minus-minus-eqns} shows the equational rules
%% for case-nat (the rules for case-arrow are analogous).

%% \begin{figure}
%%   \begin{mathpar}
%%       % inj-nat
%%       \inferrule*
%%       {\hasty \Gamma M \nat}
%%       {\hasty \Gamma {\injnat M} \dyn}

%%       % inj-arr 
%%       \inferrule*
%%       {\hasty \Gamma M (\dyn \ra \dyn)}
%%       {\hasty \Gamma {\injarr M} \dyn}

%%       % Case nat
%%       \inferrule*
%%       {\hasty{\Delta|_V}{V}{\dyn} \and 
%%         \hasty{\Delta , x : \nat }{M_{yes}}{B} \and 
%%         \hasty{\Delta}{M_{no}}{B}}
%%       {\hasty {\Delta} {\casenat{V}{M_{no}}{n}{M_{yes}}} {B}}

%%       % Case arr
%%       \inferrule*
%%       {\hasty{\Delta|_V}{V}{\dyn} \and 
%%         \hasty{\Delta , x : (\dyn \ra \dyn) }{M_{yes}}{B} \and 
%%         \hasty{\Delta}{M_{no}}{B}}
%%       {\hasty {\Delta} {\casearr{V}{M_{no}}{f}{M_{yes}}} {B}}
%%   \end{mathpar}
%%   \caption{New typing rules for $\extlcmm$.}
%%   \label{fig:extlc-minus-minus-typing}
%% \end{figure}


%% \begin{figure}
%%   \begin{mathpar}
%%      % Case-nat Beta
%%      \inferrule*
%%      {\hasty \Gamma V \nat}
%%      {\casenat {\injnat {V}} {M_{no}} {n} {M_{yes}} = M_{yes}[V/n]}

%%      \inferrule*
%%      {\hasty \Gamma V {\dyn \ra \dyn} }
%%      {\casenat {\injarr {V}} {M_{no}} {n} {M_{yes}} = M_{no}}

%%      % Case-nat Eta
%%      \inferrule*
%%      {}
%%      {\Gamma , x :\, \dyn \vdash M = \casenat{x}{M}{n}{M[(\injnat{n}) / x]} }


%%      % Case-arr Beta


%%      % Case-arr Eta


%%   \end{mathpar}
%%   \caption{New equational rules for $\extlcmm$ (rules for case-arrow are analogous
%%            and hence are omitted).}
%%   \label{fig:extlc-minus-minus-eqns}
%% \end{figure}



%% % TODO : Updated term precision rules



%% \subsection{The Step-Sensitive Lambda Calculus}\label{sec:step-sensitive-lc}

%% % \textbf{TODO: Subject to change!}

%% Rather than give a semantics to $\extlcmm$ directly, we first introduce another intermediary
%% language, a \emph{step-sensitive} (also called \emph{intensional}) calculus.
%% As mentioned, this language makes the intensional stepping behavior of programs
%% explicit in the syntax. We do this by adding a syntactic ``later'' type and a
%% syntactic $\theta$ that takes terms of type later $A$ to terms of type $A$.

%% % In the step-sensitive syntax, we add a type constructor for later, as well as a
%% % syntactic $\theta$ term and a syntactic $\nxt$ term.
%% We add rules for these new constructs, and also modify the rules for inj-arr and
%% case-arrow, since now the function is not $\Dyn \ra \Dyn$ but rather $\later (\Dyn \ra \Dyn)$.
%% We also add congruence relations for $\later$ and $\nxt$.

%% % TODO show changes

%% \noindent Modified syntax:
%% \begin{align*}
%%   &\text{Value Types } A := \nat \alt \dyn \alt (A \ra A') \alt {\color{red} \later A} \\
%%   &\text{Values } V :=  \zro \alt \suc\, V \alt \lda{x}{M} \alt \injnat V \alt \injarr V 
%%     \alt {\color{red} \nxt\, V} \alt {\color{red} \mathbf{\theta}}
%% \end{align*}

%% \noindent Additional typing rules:
%% \begin{mathpar}
%%   \inferrule
%%   {\hasty \Gamma V A}
%%   {\hasty \Gamma {\nxt\, V} {\later A}}

%%   \inferrule
%%   {}
%%   {\hasty \Gamma \theta {\later A \ra A}}

%%   % \theta(\nxt x) = \theta(y); \texttt{ret}\, x
%% \end{mathpar}

%% \noindent Modified typing rules:
%% \begin{mathpar}

%%   % inj-arr 
%%   \inferrule*
%%   {\hasty \Gamma M {\color{red} \later (\dyn \ra \dyn)}}
%%   {\hasty \Gamma {\injarr M} \dyn}

%%   % Case arr
%%   % TODO if the extensional version is incorrect and needs to change, make
%%   % sure to change this one accordingly
%%   \inferrule*
%%   {\hasty{\Delta|_V}{V}{\dyn} \and 
%%     \hasty{\Delta , x \colon {\color{red} \later (\dyn \ra \dyn)} }{M_{yes}}{B} \and 
%%     \hasty{\Delta}{M_{no}}{B}}
%%   {\hasty {\Delta} {\casearr{V}{M_{no}}{\tilde{f}}{M_{yes}}} {B}}  
%% \end{mathpar}

%% \noindent Additional relations:
%% \begin{mathpar}
%%   \inferrule*[]
%%   {A^{\rel} : A_l \rel A_r}
%%   {\later A^{\rel} : \later A_l \rel \later A_r}

%%   \inferrule*[]
%%   {A^{\rel} (V_l, V_r)}
%%   {\later A^{\rel} (\nxt\, V_l, \nxt\, V_r)}

%% \end{mathpar}

%% % TODO what about the relation for theta? Or is that automatic since it's a function symbol?

%% % TODO beta rule for theta

%% We define the term $\delta$ to be the function $\lda {x} {\theta\, (\nxt\, x)}$.

%% % We define an erasure function from step-sensitive syntax to step-insensitive syntax
%% % by induction on the step-sensitive types and terms.
%% % The basic idea is that the syntactic type $\later A$ erases to $A$,
%% % and $\nxt$ and $\theta$ erase to the identity.



%% \subsection{Quotienting by Syntactic Bisimilarity}

%% We now define a quotiented variant of the above step-sensitive calculus,
%% which we denote by $\intlcbisim$.
%% In this syntax, we add a rule saying, roughly speaking, that 
%% $\theta \circ \nxt$ is the identity. This causes terms that differ only in
%% their intensional behavior to become equal.
%% Note that a priori, this is not the same language as the step-insensitive
%% calculus on which we based the insensitive calculus.

%% Formally, the equational theory for the quotiented syntax is the same as
%% that of the original step-sensitive language, with the addition of the following
%% rule:

%% % TODO is this correct?
%% \begin{mathpar}
%%   \inferrule*
%%   { }
%%   { \theta\, (\nxt\, x) = \bind{y}{(\theta\, V')}{\ret\, x}  }
%% \end{mathpar}

%% This states that the application of $\theta$ to $\nxt\, x$ is equivalent to
%% the computation that applies $\theta$ to $V'$ to obtain a variable $y$, and
%% then simply returns $x$.




% \input{extensional-model}

\section{A Denotational Semantics for Types and Terms}\label{sec:concrete-term-model}

As a first step to giving a well-behaved semantics of GTLC, in this
section we define a simpler semantics that validates all of our
desired properties except for graduality: that is\DIFaddbegin \DIFadd{, }\DIFaddend it is adequate and
validates $\beta\eta$ equality. Our relational semantics will then be
a refinement of this simpler design. In the process we will introduce
some of the technical tools we will use in the refined semantics:
guarded type theory and call-by-push-value. While these tools can be
somewhat technical we emphasize that due to the use of synthetic
guarded domain theory, this initial semantics is essentially just a
``na\"\i ve'' denotational semantics of an effectful language arising
from a monad on sets. But since we are working in the non-standard
foundation of guarded type theory our notion of sets allows for the
construction of more sophisticated sets and monads than classical
logic allows.

% In Section \ref{sec:towards-relational-model}, we will discuss how to extend the
% denotational semantics to accommodate the type and term precision orderings. 
%
%necessitate the abstract components that are found in the definition of the model.

\subsection{Guarded Type Theory}\DIFaddbegin \label{sec:guarded-type-theory}
\DIFaddend 

% \max{TODO: discuss this wording with Eric. Are we working in guarded
%   type theory, SGDT or Ticked Cubical Type Theory in the informal work
%   in the paper?}
We begin with a brief overview of the background theory that we use in this work.
\emph{Synthetic guarded domain theory}, or SGDT for
short, is an axiomatic framework in which we can reason about non-well-founded
recursive constructions while abstracting away the specific details of
step-indexing that we would need to track if we were working analytically. This
allows us to avoid the tedious reasoning associated with traditional
step-indexing techniques. We provide a brief overview here; more details can be
found in \cite{birkedal-mogelberg-schwinghammer-stovring2011}.

SGDT offers a synthetic approach to domain theory that allows for guarded
recursion to be expressed syntactically via a type constructor $\later \colon
\type \to \type$ (pronounced ``later''). The use of a modality to express
guarded recursion was introduced by Nakano \cite{Nakano2000}.
%
Given a type $A$, the type $\later A$ represents an element of type $A$ that is
available one time step later. There is an operator $\nxt : A \to\, \later A$
that ``delays'' an element available now to make it available later. We will use
a tilde to denote a term of type $\later A$, e.g., $\tilde{M}$.

% TODO later is an applicative functor, but not a monad

There is a \emph{guarded fixpoint} operator
\[ \fix : \forall T, (\later T \to T) \to T. \]
That is, to construct a term of type $T$, it suffices to assume that we have
access to such a term ``later'' and use that to help us build a term ``now''.
This operator satisfies the axiom that $\fix f = f (\nxt (\fix f))$.
%DIF > 
In particular, \DIFdelbegin \DIFdel{this axiom applies to propositions }\DIFdelend \DIFaddbegin \DIFadd{$\fix$ applies when type $T$ is instantiated to a proposition 
}\DIFaddend $P : \texttt{Prop}$; \DIFdelbegin \DIFdel{proving a
statement in this manner is known as }\DIFdelend \DIFaddbegin \DIFadd{in that case, it corresponds to }\DIFaddend $\lob$-induction.

% Clocked Cubical Type Theory
\subsubsection{Ticked Cubical Type Theory}
% TODO motivation for Clocked Cubical Type Theory, e.g., delayed substitutions?

Ticked Cubical Type Theory \cite{mogelberg-veltri2019} is an extension of
Cubical Type Theory \cite{CohenCoquandHuberMortberg2017} that has an additional
sort called \emph{ticks}. Ticks were originally introduced in
\cite{bahr-grathwohl-bugge-mogelberg2017}. A tick $t : \tick$ serves as evidence
that one unit of time has passed. In Ticked Cubical Type Theory, the type
$\later A$ is the type of dependent functions from ticks to $A$. The type $A$ is
allowed to depend on $t$, in which case we write $\later_t A$ to emphasize the
dependence.

% TODO next as a function that ignores its input tick argument?

% TODO include a figure with some of the rules for ticks

The rules for tick abstraction and application are similar to those of ordinary
$\Pi$ types. A context now consists of ordinary variables $x : A$ as well as
tick variables $t : \tick$. The presence of the tick variable $t$ in context
$\Gamma, (t : \tick), \Gamma'$ intuitively means that the values of the
variables in $\Gamma$ arrive ``first'', then one time step occurs, and then the
values of the variables in $\Gamma'$ arrive.
%
The abstraction rule for ticks states that if in context $\Gamma, t : \tick$
the term $M$ has type $A$, then in context $\Gamma$ the term $\lambda t.M$ has
type $\later_t A$.
%
Conversely, if we have a term $M$ of type $\later A$, and we have available in
the context a tick $t' : \tick$, then we can apply the tick to $M$ to get a
term $M[t'] : A[t'/t]$. However, there is an important restriction on when we
are allowed to apply ticks\DIFdelbegin \DIFdel{. In order to }\DIFdelend \DIFaddbegin \DIFadd{: To }\DIFaddend apply $M$ to tick $t$, \DIFaddbegin \DIFadd{we require }\DIFaddend $M$ \DIFdelbegin \DIFdel{must }\DIFdelend \DIFaddbegin \DIFadd{to }\DIFaddend be
well-typed in the prefix of the context occurring before the tick $t$. That is,
all variables mentioned in $M$ must be available before $t$. This ensures that
we cannot, for example, define a term of type $\later \laterhs A \to\, \laterhs
A$ via repeated tick application.
% TODO restriction on tick application
%
For the sake of brevity, we will also write tick application as $M_t$.
The constructions we carry out in this paper will take place in a guarded type
theory with ticks.
\subsection{A Call-by-push-value Model}

% TODO there aren't actually computation types in the syntax. They
% arise because we are embedding call-by-value into CBPV.

Our denotational model of GTLC must account for the fact that the
language is call-by-value with non-trivial effects (divergence and
errors).  To do this, we formulate the model using the categorical
structures of Levy's \emph{call-by-push-value}
\cite{levy99}. Call-by-push-value is a categorical model of effects
that is a refinement of Moggi's monadic semantics
\cite{MOGGI199155}. While Moggi works with a (strong) monad $T$ on a
category $\mathcal V$ of pure morphisms, call-by-push-value works
instead with a decomposition of such a monad into an adjunction
between a category $\mathcal V$ of ``value types'' and pure morphisms
and a category $\mathcal E$ of ``computation types'' and linear
morphisms. The monad $T$ on $\calV$ is decomposed into a left adjoint
$F : \calV \to \calE$ constructing the type $F A$ of computations that
return $A$ values and and right adjoint $U : \calE \to \calV$
constructing the type $U B$ of values that are first-class suspended
computations of type $B$. The monad $T$ can then be reconstituted as
$T = UF$ but the main advantage is that many constructions on
effectful programs naturally decompose into constructions on
value/computation types. Most strikingly, the CBV function type $A
\rightharpoonup A'$ decomposes into the composition of three
constructions $U(A \to F A')$ where $\arr : \calV^{op} \times \calE
\to \calE$ is a kind of mixed kinded function type. This decomposition
considerably simplifies treatment of the CBV function
type. Additionally, prior work on gradual typing semantics has shown
that the casts of gradual typing are naturally formulated in terms of
pairs of a pure morphism and a linear morphism, as we explain below
\cite{new-licata-ahmed2019}.

For our simple denotational semantics, the category $\calV$ is simply
the category $\Set$ of sets and functions. Then because GTLC is
call-by-value, all types will be interpreted as objects of this
category, i.e., sets. More interesting are the computation
types. These will be interpreted in a category of sets equipped with
\DIFdelbegin \DIFdel{``algebraic''}\DIFdelend \DIFaddbegin \DIFadd{algebraic}\DIFaddend \footnote{\DIFdelbegin \DIFdel{the }\DIFdelend \DIFaddbegin \DIFadd{The }\DIFaddend $\theta$ structure is not algebraic in the
strictest sense since it does not have finite arity\DIFdelbegin \DIFdel{, but }\DIFdelend \DIFaddbegin \DIFadd{.}\DIFaddend } structure with
morphisms being functions that are homomorphic in this structure.
This algebraic structure needs to model the effects present in our
language, in this case these are errors as well as taking
computational steps (potentially diverging by taking computational
steps forever).
%
We call these algebraic structures simple error domains to distinguish
them from the error domains we define later to model graduality:
\begin{definition}[Simple Error Domains]
  A (simple) error domain $B$ consists of
  \begin{enumerate}
  \item A carrier set $UB$
  \item An element $\mho_{B} : UB$ representing error
  \item A function $\theta_B : \laterhs UB \to UB$ modeling a computational step
  \end{enumerate}
  A homomorphism of error domains $\phi : B \multimap B'$ is a
  function $\phi : UB \to UB'$ that preserves \DIFdelbegin \DIFdel{these in that
  }\DIFdelend \DIFaddbegin \DIFadd{$\mho$ and $\theta$, i.e.,
  }\DIFaddend $\phi(\mho_B) = \mho_{B'}$ and $\phi(\theta_{B}(\tilde{x})) =
  \theta_{B'}(\later (\phi(\tilde{x})))$.  Simple error domains and
  homomorphisms assemble into a category $\ErrDom$ with a forgetful
  functor $U : \ErrDom \to \Set$ taking the underlying set and
  function.
\end{definition}
Here is our first point where we utilize guarded type theory: rather
than simply being a function $UB \to UB$, the ``think'' map $\theta$
takes an element \emph{later}. This makes a major difference, because
the structure of a think map combined with the guarded fixpoint
operator allows us to define recursive elements of $UB$ in that any
function $f : UB \to UB$ has a ``quasi-fixed point'' $\textrm{qfix}(f)
= \fix(f \circ \theta_B)$ satisfying the quasi-fixed point property:
\[ \qfix(f) = f(\delta_B(\qfix f)) \]
where $\delta_B = \theta_B \circ \nxt$ is a map we call the ``delay''
map which trivially delays an element now to be available later.  As
an example, we can define $\Omega_B = \qfix(\id) : UB$, the
``diverging element'' that ``thinks forever'' in that \DIFdelbegin \DIFdel{$\Omega_B =
\delta(\Omega_B)$}\DIFdelend \DIFaddbegin \DIFadd{$\Omega_B =
\delta_B(\Omega_B)$}\DIFaddend . We call this a ``quasi'' fixed point because it is
\emph{nearly} a fixed point except for the presence of the delay map
$\delta_B$, which is irrelevant from an extensional point of view
where we would prefer to ignore differences in the number of steps
that computations take.

% \max{move this}
% The functor $\arr : \Set^{op} \times \errdom \to \errdom$ is defined on objects
% as follows. Given a set $A$ and error domain $B$, $A \arr B$ is the error domain
% whose underlying set is the set of functions $A \to UB$. The error element is
% the constant error function, i.e., $\lambda x. \mho_B$. The map $\theta :
% (\laterhs (A \to UB)) \to (A \to UB)$ is defined by 
% $\theta(\tilde{f}) = \lambda x . \theta_B (\tilde{f}_t)$.
% %
% Given a morphism $f : A_o \to A_i$ and $\phi : B_i \to B_o$, the morphism $f
% \arr \phi : (A_i \arr B_i) \to (A_o \arr B_o)$ is given by pre-composing by $f$
% and post-composing by $\phi$.

Then to model effectful programs, we need to construct a left adjoint
$\li : \Set \to \ErrDom$ to $U$ which constructs the ``free error
domain'' on a set, which models an effectful CBV computation.
%% \subsection{The Lift + Error Monad}\label{sec:lift-monad}
We base the construction on the \emph{guarded
lift monad} described in \cite{mogelberg-paviotti2016}. Here, we augment the
guarded lift monad to accommodate the additional effect of failure.
\begin{definition}[Free error domain]
  For a set $A$, we define the (carrier of) \emph{free error domain} $U(\li A)$ as the unique solution to the guarded domain equation:
  \[ U(\li A) \cong A + 1\, + \laterhs U(\li A). \]
  For which we use the following notation for the three constructors:
  \begin{enumerate}
  \item $\eta \colon A \to U(\li A)$
  \item $\mho \colon U(\li A)$
  \item $\theta \colon \laterhs U(\li A) \to U(\li A)$
  \end{enumerate}
  Here $\mho, \theta$ provide the error domain structure for $\li A$,
  and this has the universal property of being free in that any
  function $f : U(\li A) \to UB$ uniquely extends to a homomorphism
  $f^\dagger : \li A \multimap B$ satisfying $Uf^\dagger \circ \eta =
  f$. In other words, the free error monad is left adjoint $\li \dashv
  U$ to the forgetful functor, with $\eta$ the unit of the adjunction.
  %DIF > 
  \DIFaddbegin \DIFadd{We will write $\theta_t(\dots)$ to mean $\theta (\lambda t. \dots)$.
}\DIFaddend \end{definition}
%% %
%% Unless otherwise mentioned, all constructs involving $\later$ or $\fix$ are
%% understood to be with respect to a fixed clock $k$. So for the above, we really
%% have for each clock $k$ a type $\li^k A$ with respect to that clock.
%% %
%% Formally, the lift monad $\li A$ is defined as the solution to the guarded
%% recursive type equation
%% %
%
An element of $\li A$ should be viewed as a computation that can either (1)
return a value (via $\eta$), (2) raise an error and stop (via $\mho$), or (3)
take a computational step (via $\theta$).
%
The operation $f^\dagger$ can be defined by guarded recursion in the
apparent way, and proven to satisfy the freeness property by L\"ob
induction.
% It is instructive to give at least one example of a use of guarded recursion, so
% we show below how to define extend:
% TODO
%
%
%% Verifying that the monadic laws hold uses \lob-induction and is straightforward.

%% % TODO mention that error domains are algebras of the lift monad?
%% We observe that the guarded lift monad applied to a set $A$ is the underlying
%% set of the free error domain on $A$ $A$, i.e., we have $\li = UF$ where $F :
%% \Set \to \errdom$ and $U : \errdom \to \Set$ and $F \dashv U$. Given a set $A$
%% and an error domain $B$, the set of functions $A \to UB$ is naturally isomorphic
%% to the set of error domain morphisms $FA \to B$.

%% Lastly, we define the function $\delta : \li A \to \li A$ by $\delta = \theta_A
%% \circ \nxt$.

% The function type A \ra A' will be modeled as \sem{A} \to \li \sem{A'}

The free error domain is essential to modeling CBV computation: a term
$\Gamma \vdash M : A$ in GTLC will be modeled as a function $M :
\sem{\Gamma} \to U\li \sem{A}$ from the set denoted by $\Gamma$ to the
underlying set of the free error domain on the set denoted by $A$,
just as in a monadic semantics with monad $U\li$.

\subsection{Modeling the type structure}\label{sec:dynamic-type}
Much of the type structure of GTLC is simple to model: $\nat$ denotes
the set of natural numbers and $\times$ denotes the cartesian product
of sets. To model the CBV function type $A \rightharpoonup A'$ we use
the CBPV decomposition $U(\sem{A} \to \li \sem{A'})$ where $A \to B$
is the obvious point-wise algebraic structure on the set of functions
from $A$ to $UB$.

The final type to model is the the dynamic type $\dyn$.  In the style
of Scott, a classical domain theoretic model for $\dyn$ would be given
by solving a domain equation
%
\[ D \cong \Nat + (D \times D) + U(D \to \li D), \]
%
representing the fact that dynamically typed values can be numbers,
pairs or CBV functions. This equation does not have inductive or
coinductive solutions due to the negative occurrence of $D$ in the
domain of the function type. In classical domain theory this is
avoided by instead taking the initial algebra in a category of
embedding-projection pairs, by which we obtain an exact solution to
this equation. However, in guarded domain theory we instead replace this domain equation with a guarded domain equation.
Consider the following similar-looking equation:
%
\begin{equation}\label{eq:dyn}
D \cong \Nat + (D \times D)\, + \laterhs U(D \to \li D).
\end{equation}
%
Since the negative occurrence of $D$ is guarded under a later, we can
solve this equation by a mixture of inductive types and guarded fixed
points. Specifically, consider the parameterized inductive type
%
\[ D'\, X := \mu T. \Nat + (T \times T) + X. \]
%
Then we can construct $D$ as the unique solution to the guarded domain equation
\[ D \cong D'(\laterhs U(D \to \li D)) \]
which is a guarded equation in that all occurrences of $D$ on the right hand side are guarded by the $\laterhs$.
%
Expanding out the guarded fixed point and least fixed point property
gives us that $D$ satisfies Equation $\ref{eq:dyn}$. 
% \max{give that equation a number}.
%
We then refer to the three injections for numbers, pairs and functions
as $\inat, \itimes, \iarr$ respectively.

\subsection{Term Semantics}\label{sec:term-interpretation}

We now extend our semantics of types to a semantics of terms of
GTLC. As mentioned above, terms $x:A_1,\ldots \vdash M : A$ are
modeled as effectful functions $\sem{M}: \sem{A_1}\times \cdots \to U\li\sem{A}$
and values $x:A_1,\ldots \vdash V : A$ are also \DIFdelbegin \DIFdel{be }\DIFdelend modeled as pure
functions $\sem{V} : \sem{A_1}\times\cdots \to \sem{A}$.
%
The interpretation of the typical CBV features of GTLC into a \DIFaddbegin \DIFadd{model of }\DIFaddend CBPV is
standard \cite{levy99}\DIFdelbegin \DIFdel{, }\DIFdelend \DIFaddbegin \DIFadd{; }\DIFaddend the only additional part we need to account for
is the semantics of casts, shown in Figure \ref{fig:term-semantics}.
%
%% Much of the semantics is similar to a normal call-by-value denotational
%% semantics, so we focus only on the cast semantics. We interpret types as sets;
%% the interpretation of types is given in Figure \ref{fig:type-interpretation}.
%% Contexts $\Gamma = x_1 \colon A_1, \dots, x_n \colon A_n$ are interpreted as the
%% product $\sem{A_1} \times \cdots \times \sem{A_n}$. The semantics of the dynamic
%% type $\dyn$ is the set $D$ introduced in Section \ref{sec:dynamic-type}. The
%% product type $A \times A'$ is interpreted as the Cartesian product of the
%% denotations of $A$ and $A'$. The function type $A \ra A'$ is interpreted as the
%% set of functions from $\sem{A}$ to $\li (\sem {A'})$.
%
%% The interpretation of a value $\hasty {\Gamma} V A$ is a function from
%% $\sem{\Gamma}$ to $\sem{A}$. Likewise, a term $\hasty {\Gamma} M {{A}}$ is
%% interpreted as a function from $\sem{\Gamma}$ to $\li \sem{A}$.
Following the CBPV treatment of gradual typing given by
New-Licata-Ahmed, an upcast $\upc c$ where $c : A \ltdyn A'$ is
interpreted as a pure function $\sem{\upc c} :\sem{A} \to \sem{A'}$ whereas a
downcast $\dnc c$ is interpreted as a homomorphism $\sem{\dnc c} :
\li\sem{A'} \multimap \li\sem{A}$.
%% By the universal property of $\li$, this is determined by a function $\sem{A'} \to U\li\sem{A}$ and
%% We use $\text{ext}$ to go from the former to the latter in the definitions of the downcasts.
We define the semantics of upcasts and downcasts simultaneously by
recursion over type precision derivations.
%
The reflexivity casts are given by identities, and the transitivity
casts by composition. The upcasts for products, and the downcasts for
functions are given by the functorial actions of the type constructors
on the casts for the subformulae.
%
The \emph{up}casts for functions and similarly the \emph{down}casts
for products, on the other hand, are more complex.
%
For these, we use not the ordinary functorial action of type
constructors, but an action we call the \emph{Kleisli} action.

The Kleisli action of type constructors can be defined in any CBPV
model. First, we can define the Kleisli category of value types
$\calV_k$ to have value types as objects, but have morphisms
$\calV_k(A,A') = \calE(FA,FA')$. Similarly, we can define a Kleisli
category of computation types $\calE_k$ to have computation types as
objects, but have morphisms $\calE_k(B,B') = \calV(UB,UB')$. The
intuition is that these are the ``effectful'' morphisms between
value/computation types, whereas the morphisms of $\calV/\calE$ are
``pure''/``linear''. Then every CBPV type constructor $F,U,\times,\to$
has in addition to its usual functorial action on pure/linear
morphisms, an action in each argument on Kleisli morphisms. So for
instance, the function construction is functorial in pure/linear
morphisms $\to : \op{\calV} \times \calE \to \calE$ but on Kleisli
morphisms what we have is for each computation type $B$ a functorial
action $-\to B : \calV_k^{op} \to \calE_k$ given by $-\tok B$ on
objects and for each value type a functorial action $A \tok - :
\calE_k \to \calE_k$. Similarly we get two actions $A \timesk -$ and
$-\timesk A'$ on Kleisli morphisms between value types. In these two
cases we do not get a bifunctorial action because the two Kleisli
actions do not commute past each other, i.e., in general $(A_l'
\timesk f_r \circ f_l \timesk A_r) \neq (f_l \timesk A_r' \circ A_l
\timesk f_r)$. The full definitions of the Kleisli actions are
included in Appendix \ref{sec:kleisli-actions}.

Lastly, we have the upcasts and downcasts for the injections into the dynamic type, which
are the core primitive casts. The upcasts are simply the injections
themselves, except the function case which must include a $\nxt$ to
account for the fact that the functions are under a later in the
dynamic type. The downcasts are similar in that on values
they pattern match on the input and return it if it is of the correct
type, otherwise erroring. Again, the function case is slightly
different in that if the input is in the function case, then it is
actually only a function available later, and so we must insert a
``think'' in order to return it.

% The upcast for $c_i \ra c_o$ and the downcast for $c_1 \times c_2$ make use of
% \emph{Kleisli functors} $\tok$ and $\timesk$.

%% \eric{Introduce Kleisli categories and functors}

%% Recall the definition of type precision derivations given in Figure
%% \ref{fig:typrec}. The semantics of the up- and downcasts for a type precision
%% derivation $c$ are defined by induction on $c$.
%% %
%% For $r(A)$ the up- and downcasts are simply the identity function. 
%% %
%% For the composition $c \comp c'$ we compose the cast for $c$ and the cast for
%% $c'$ as functions.
%% %
%% For $c_i \ra c_o$ we apply the casts for $c_i$ and $c_o$ in the domain and
%% codomain respectively. More concretely, for the upcast we are given a function
%% $V_f : A_i \to \li A_o$ and we must define a function $A_i' \to \li A_o'$. The
%% function first downcasts its argument according to $c_i$, resulting in an
%% element of $\li A_i$. It then applies $\ext{V_f}{}$ to this value to obtain an
%% element of $\li A_o$. Finally, it applies the upcast of $c_o$ via the functorial
%% action of the lift monad, obtaining an element of $\li A_o'$ as needed. For the
%% downcast from $A_i' \to \li A_o'$ to $A_i \to \li A_o$, the resulting function
%% first applies the upcast by $c_i$ to its argument, then applies the original
%% function, and finally applies the downcast by $c_o$ to the result.
%% %
%% Likewise, in the upcast for $c_1 \times c_2$, we apply the cast for $c_1$ on the
%% left and the cast for $c_2$ on the right.

%% The ``base cases'' for the casts are $\inat$, $\itimes$, and $\iarr$. Recall
%% that $D$ is isomorphic to $\Nat\, + (D \times D)\, + \laterhs (D \to \li D)$
%% (say that the sum is left-associative).
%% %
%% For $\inat$, the upcast is simply $\inl \circ \inl$. The downcast has type $D
%% \to \li D$ and performs a case inspection on the sum type. If it is a natural
%% number $n$, then we return $\eta\, n$. Otherwise, we return $\mho$, modelling
%% the fact that a run-time error occurs.
%% %
%% For $\itimes$, the upcast is $\inl \circ \inr$, and the downcast performs a case
%% inspection analogously to the natural number downcast.
%% %
%% % TODO explain this further
%% For $\iarr$, the upcast is given by $\inr \circ \nxt$. The downcast performs a
%% case inspection, and in the $\inr$ case, it uses a $\theta$ in order to gain
%% access to the function under the $\later$.


%
% TODO write this out explicitly?

% TODO should we include function casts?
%% \begin{figure*}
%%   \begin{align*}
%%     \sem{\nat} &= \Nat \\
%%     \sem{\dyn} &= D \\
%%     \sem{A \times A'} &= \sem{A} \times \sem{A'} \\ 
%%     \sem{A \ra A'} &= \sem{A} \To \li \sem{A'} \\
%%   \end{align*}
%%   \caption{Interpretation of types}
%%   \label{fig:type-interpretation}
%% \end{figure*}
\begin{figure*}
  \begin{minipage}{0.3\textwidth}
    \begin{footnotesize}
  \begin{align*}
    % upcasts
    \sem{\upc{r(A)}} &= \id\DIFdelbeginFL \DIFdelFL{_A }\DIFdelendFL \DIFaddbeginFL \DIFaddFL{_{\sem{A}} }\DIFaddendFL \\
    \sem{\upc{(c \comp c')}} &= \sem{\upc{c'}} \circ \sem{\upc{c}} \\
    \sem{\upc{(c_i \ra c_o)}} &= (\sem{\dnc {c_i}} \tok \li\sem{A_o'})\circ (\sem{A_i} \tok U\li(\sem{\upc {c_o}}))\\
    \sem{\upc{(c_1 \times c_2)}} &= \sem{\upc{c_1}} \times \sem{\upc{c_2}}\\
    \sem{\upc{\inat}} &= \inat\\
    \sem{\upc{\itimes}} &= \itimes\\
    \sem{\upc{\iarr}} &= \iarr \circ \nxt\\
    % \sem{\zro}         &= \lambda \gamma . 0 \\
    % \sem{\suc\, V}     &= \lambda \gamma . (\sem{V}\, \gamma) + 1 \\
    % \sem{x \in \Gamma} &= \lambda \gamma . \gamma(x) \\
    % \sem{\lda{x}{M}}   &= \lambda \gamma . \lambda a . \sem{M}\, (*,\, (\gamma , a))  \\
    % \sem{V_f\, V_x}    &= \lambda \gamma . {({(\sem{V_f}\, \gamma)} \, {(\sem{V_x}\, \gamma)})} \\
    % \sem{\err_B}       &= \lambda \gamma . \mho \\
    %
    % downcasts
    \sem{\dnc{r(A)}} &= \id\DIFdelbeginFL \DIFdelFL{_{FA} }\DIFdelendFL \DIFaddbeginFL \DIFaddFL{_{\li \sem{A}} }\DIFaddendFL \\
    \sem{\dnc{(c \comp c')}} &= \sem{\dnc{c}} \circ \sem{\dnc{c'}} \\
    \sem{\dnc{(c_i \ra c_o)}} &= U(\sem{\upc {c_i}} \to \sem{\dnc {c_o}})\\
    \sem{\dnc{(c_1 \times c_2)}} &= (\dnc{c_1} \timesk A_2) \circ (A_1' \timesk \dnc{c_2}) \\
    \sem{\dnc{\inat}} &= (
      \lambda V_d . \text{case $({V_d})$ of}
      \{ \inat\,n \to \eta\, n
         \alt \text{otherwise} \to \mho \})^\dagger \\
    \sem{\dnc{\itimes}} &= (
      \lambda V_d . \text{case $({V_d})$ of}
      \{ \itimes(d_1, d_2) \to \eta\, (d_1, d_2)
         \alt \text{otherwise} \to \mho \})^\dagger \\
    \sem{\dnc{\iarr}} &= (
      \lambda V_d . \text{case $({V_d})$ of}
      \{ \iarr \tilde{f} \to \theta_t\, (\eta (\tilde{f}_t))
         \alt \text{otherwise} \to \mho \})^\dagger \\
    % \sem{\ret\, V}       &= \lambda \gamma . \eta\, \sem{V} \\
    % \sem{\bind{x}{M}{N}} &= \lambda \delta . \ext {(\lambda x . \sem{N}\, (\delta, x))} {\sem{M}\, \delta} \\
  \end{align*}
    \end{footnotesize}
  \end{minipage}
  \caption{Semantics of casts}
  \label{fig:term-semantics}
\end{figure*}

\subsection{Extracting a well-behaved big-step semantics}\label{sec:big-step-term-semantics}

% \max{TODO: Eric needs to add a discussion of adequacy/clock stuff here}
% \eric{Check this}

We now define from the above term model a ``big-step'' semantics. More
concretely, our goal is to define a \DIFaddbegin \DIFadd{partial }\DIFaddend function 
$-\Downarrow : \{M \,|\, \cdot \vdash M : \nat \} \rightharpoonup \mathbb{N} + {\mho}$.

To begin, we must first discuss another aspect of guarded type theory,
\emph{clocks}, originally introduced in \cite{atkey-mcbride2013}. The
constructions of guarded type theory ($\later$, $\nxt$, and $\fix$) we have been
using may be indexed by objects called \emph{clocks}. A clock intuitively serves
as a reference relative to which steps are counted. For instance, given a clock
$k$ and type $X$, the type $\later^k X$ represents a value of type $X$ one unit
of time in the future according to clock $k$.
%
Up to this point, all uses of guarded constructs have been indexed by a single
clock $k$. In particular, we have for each clock $k$ a type $\li^k X$. The
notion of \emph{clock quantification} allows us to pass from the guarded setting
to the usual set-theoretic world. Clock quantification enables us to encode
coinductive types using guarded recursion. That is, given a type $A$ with a free
variable $k : \Clock$, we can form the type $A^{gl} := \forall k. A$ \DIFaddbegin \DIFadd{(here, $gl$ stands for ``global'')}\DIFaddend . Given a
functor $F$ that satisfies a technical condition, if $A$ is a guarded recursive
type satisfying $A \cong F(\later^k A)$, then the type $A^{gl}$ has a final
$F$-coalgebra structure (see Theorem 4.3 of
\cite{kristensen-mogelberg-vezzosi2022}). 

Consider the functor $F(X) = \mathbb{N} + {\mho} + X$. By the definition of the
guarded lift, we have that
% 
$\li^k \mathbb{N} 
  \cong \mathbb{N} + {\mho} + \later^k (\li \mathbb{N}) 
  \cong F(\later \li^k \mathbb{N})$.
%
Furthermore, it can be shown that the functor $F$ satisfies the technical
condition in the statement of the above-mentioned theorem. Thus, the theorem
implies that the type $\li^{gl} \mathbb{N} := \forall k. \li^k \mathbb{N}$ is a
final $F$-coalgebra. In fact, $\li^{gl}$ is isomorphic to a more familiar coninductive type.
%
Given a type $Y$, Capretta's \emph{delay monad} \cite{lmcs:2265},
which we write as $\text{Delay}(Y)$, is defined as the coinductive type
generated by $\tnow : Y \to \delay(Y)$ and $\tlater : \delay(Y) \to \delay(X)$.
We have that $\delay (\mathbb{N} + {\mho})$ is a final coalgebra for the functor
$F$ defined above. This means $\li^{gl} \mathbb{N}$ is isomorphic to
$\delay(\mathbb{N} + {\mho})$. \DIFaddbegin \DIFadd{Denote this isomorphism by $\alpha$.
}\DIFaddend 

\DIFdelbegin \DIFdel{We can }\DIFdelend \DIFaddbegin \DIFadd{Given $d : \delay(\mathbb{N} + {\mho})$, we }\DIFaddend define a notion of termination in
$i$ steps\DIFdelbegin \DIFdel{for the Delay monad $d \da^i
n_?$ where $d : \delay(\mathbb{N} + {\mho})$ and }\DIFdelend \DIFaddbegin \DIFadd{, written $d \da^i$. We define $d \da^i$ inductively by the rules
$(\tnow\, n_?) \da^0$ and $(\tlater\, d) \da^{i+1}$ if $d \da^i$ (here, }\DIFaddend $n_? :
\mathbb{N} + {\mho}$\DIFaddbegin \DIFadd{)}\DIFaddend . From this \DIFdelbegin \DIFdel{, can }\DIFdelend \DIFaddbegin \DIFadd{definition, it is clear that if there is an $i$
such that $d \da^i$, then this $i$ is unique. 
%DIF >  Furthermore, from the proof that
%DIF >  $d$ terminates in $i$ steps, we can extract the result with which $d$
%DIF >  terminates. This is defined by a straightforward induction on the proof that $d \da^i$.
Using this definition of termination, we }\DIFaddend define a partial function 
\DIFdelbegin \DIFdel{$\delay (\mathbb{N} + {\mho})
\rightharpoonup \mathbb{N} + {\mho}$ }\DIFdelend %DIF > 
\DIFaddbegin \DIFadd{$f \colon \delay (\mathbb{N} + {\mho}) \rightharpoonup \mathbb{N} + {\mho}$
%DIF > 
}\DIFaddend by existentially quantifying over $i$ for which $d \da^i n_?$. \DIFaddbegin \DIFadd{That is, we
define $f(d) := n_?$ if $d \da^i n_?$ for some (necessarily unique) $i$, and $f$
is undefined otherwise.
}\DIFaddend 

Now consider the \DIFdelbegin \DIFdel{function $\{M \,|\, \cdot \vdash M : \nat \} \to \li^{gl}
\mathbb{N}$ defined by
$\lambda M. \lambda k .\sem{M}$}\DIFdelend \DIFaddbegin \DIFadd{global version of the term semantics, $\sem{\cdot}^{gl} \colon
\{M \,|\, \cdot \vdash M : \nat \} \to \li^{gl} \mathbb{N}$, defined by
$\sem{M}^{gl} := \lambda (k : \Clock) .\sem{M}$}\DIFaddend . Composing this with the
\DIFdelbegin \DIFdel{above partial function }\DIFdelend \DIFaddbegin \DIFadd{isomorphism $\alpha$ and the partial function $f$ }\DIFaddend yields the desired partial
function
%DIF > 
$-\Downarrow \colon \{M \,|\, \cdot \vdash M : \nat \} \rightharpoonup \mathbb{N} + {\mho}$.
%DIF > 
It is then straightforward to see that this big-step semantics respects the
equational theory and that the denotations of the diverging term $\Omega$, the
error term $\mho$, and a syntactic natural number value $n$ are all distinct.


\section{Challenges to a Guarded Model of Graduality}\label{sec:towards-relational-model}

% In the previous section, we gave a denotational semantics to the terms of the
% gradually-typed lambda calculus using the tools of SGDT. We defined a big-step
% semantics for closed terms of type $\nat$, as functions from $\mathbb{N} \to
% ((\mathbb{N} + 1) + 1)$.

Now that we have seen how to model the terms without regards to
graduality, we next turn to how to enhance this model to provide a
compositional semantics that additionally satisfies graduality. By a
compositional semantics, we mean that we want to provide a
compositional interpretation of \emph{type} and \emph{term precision}
such that we can extract the graduality relation from our model of
term precision.

\subsection{A first attempt: Modeling Term Precision with Posets}

In prior work, New and Licata model the term precision relation by
equipping $\omega$-CPOs with a second ordering $\ltdyn$ they call the
``error ordering'' such that a term precision relationship at a fixed
type $M \ltdyn M' : r(A)$ implies the error ordering $\sem{M} \ltdyn
\sem{M'}$ in the semantics. Then the error ordering on the free error
domain $\li A$ is based on the graduality property: essentially either
$\sem{M}$ errors or both terms diverge or both terms terminate with
values related by the ordering relation on $A$.
%
The obvious first try at a guarded semantics that models graduality
would be to similarly enhance our value types and computation types to
additionally carry a poset structure. For most type formers this
ordering can be defined as lifting poset structures on the
subformulae; the key design choice is in the ordering on the free
error domain $\li A$.

The ordering on $\li A$ should model the graduality property, but how
to interpret this is not so straightforward. The graduality property
as stated mentions divergence, but we only work indirectly with
diverging computations in the form of the think maps. The closest
direct encoding of the graduality property is an ordering we call the
\emph{step-insensitive error ordering}, defined by guarded recursion
in Figure~\ref{fig:step-insensitive-error-ordering}.
%
Firstly, to model graduality, $\mho$ is the least element. Next, two
returned values are related just when they are in the ordering $\ltdyn$
on $A$.
%
Two thinking elements are related just when they are related later.
%
The interesting cases are then those where one side is thinking and
the other has completed to either an error or a value.
%
Since the graduality property is \emph{extensional}, i.e., oblivious
to the number of steps taken, it is sensible to say that if one side
has terminated and the other side is thinking, that we must require
the thinking side to eventually terminate with a related behavior,
which is the content of the final \DIFdelbegin \DIFdel{$4$ }\DIFdelend \DIFaddbegin \DIFadd{three }\DIFaddend cases of the definition.
\begin{figure*}
      \begin{align*}
        \mho \semltbad l &\text{ iff } \top \\
        %
        \eta\, x \semltbad \eta\, y &\text{ iff } 
            x \DIFdelbeginFL %DIFDELCMD < \semlt %%%
\DIFdelendFL \DIFaddbeginFL \DIFaddFL{\ltdyn }\DIFaddendFL y \\
        %		
        \theta\, \tilde{l} \semltbad \theta\, \tilde{l'} &\text{ iff } 
            \later_t (\tilde{l}_t \semltbad \tilde{l'}_t) \\
        %	
        \theta\, \tilde{l} \semltbad \mho &\text{ iff } \exists n. \theta\, \tilde{l} = \delta^n(\mho) \\
        %	
        \theta\, \tilde{l} \semltbad \eta\, y &\text{ iff } \exists n. \exists x \ltdyn y.
            (\theta\, \tilde{l} = \delta^n(\eta\, x)) \\
        %	
        \eta\, x \semltbad \theta\, \tilde{l'} &\text { iff }
            \exists n. \exists y \DIFdelbeginFL \DIFdelFL{\geq }\DIFdelendFL \DIFaddbeginFL \DIFaddFL{\gtdyn }\DIFaddendFL x. (\theta\, \tilde{l'} = \delta^n (\eta\, y))
    \end{align*}
    \caption{Step-insensitive error ordering}
    \label{fig:step-insensitive-error-ordering}
\end{figure*}

This definition of $\semltbad$ is sufficient to model the graduality
property in that if $l \semltbad l'$ for $l,l' : \li \mathbb{N}$ where
$\mathbb{N}$ is the flat poset whose ordering is equality, then the
big-step semantics we derive for $l, l'$ does in fact satisfy the
intended graduality relation.
%
However, there is a major issue with this definition: it's not a poset
at all: it is not anti-symmetric, but more importantly it is
\emph{not} transitive!
%
To see why, we observe the following undesirable property of any
relation that is transitive and ``step-insensitive'' on one side:
\begin{theorem}\label{thm:no-go}
  Let $R$ be a binary relation on the free error domain $U(\li A)$. If
  $R$ satisfies the following properties
  \begin{enumerate}
  \item Transitivity
  \item $\theta$-congruence: If $\later_t (\tilde{x}_t \binrel{R} \tilde{y}_t)$, then $\theta(\tilde{x}) \binrel{R} \theta(\tilde{y})$.
  \item Right step-insensitivity: If $x \binrel{R} y$ then $x \binrel{R} \delta y$.
  \end{enumerate}
  Then for any $l : U(\li A)$, we have $l \binrel{R} \Omega$. If $R$ is left
  step-insensitive instead then $\Omega \binrel{R} x$.
\end{theorem}
\begin{proof}
  By L\"ob induction\DIFdelbegin \DIFdel{, in the appendix }\DIFdelend \DIFaddbegin \DIFadd{: we assume that $l$ is related to $\Omega$ later, which implies that
  $\theta (\nxt\, l) \binrel{R} \theta (\nxt\, \Omega)$. Then observe that 
  $l \binrel{R} \theta (\nxt\, l) \binrel{R} \theta (\nxt\, \Omega) = \Omega$.
  See the appendix for more details}\DIFaddend .
\end{proof}
\begin{corollary}
  Let $R$ be a binary relation on $U(\li A)$. If $R$ satisfies
  transitivity, $\theta$-congruence and left and right
  step-insensitivity, then $R$ is the total relation: $\forall x, y. x
  \binrel{R} y$.
\end{corollary}
\begin{proof}
  $x \binrel{R} \Omega$ and $\Omega \binrel{R} y$ by the previous
  theorem. Then by transitivity $x \binrel R y$.
\end{proof}

This in turn implies $\semltbad$ is not transitive, since it is a
$\theta$-congruence and step-insensitive on both sides, but not
trivial: e.g., for the flat poset on $\mathbb N$, $\eta 0 \semltbad
\eta 1$ is false.
%
This shows definitively that we cannot provide a compositional
semantic graduality proof based on types as posets if we use this
step-insensitive ordering on $\li A$.

\subsection{Why transitivity is essential}

At this point we might try to weaken \DIFdelbegin \DIFdel{from }\DIFdelend our initial attempt at
providing a poset semantics to a semantics where types are equipped
with merely a reflexive relation, and continue using the
step-insensitive ordering. However\DIFaddbegin \DIFadd{, }\DIFaddend some level of transitivity is
absolutely essential to providing compositional reasoning, and
transitivity is used pervasively in the prior work by New and Licata.
%
The reason that it is essential \DIFdelbegin \DIFdel{has to do with }\DIFdelend \DIFaddbegin \DIFadd{becomes clear by examining }\DIFaddend the details of the
graduality proof, so first we expand on how we prove graduality compositionally.
First\DIFdelbegin \DIFdel{of all}\DIFdelend , we need to extend our compositional semantics of types and terms to \DIFdelbegin \DIFdel{also }\DIFdelend a
compositional semantics of \emph{type precision derivations} and \emph{term
precision}.
%
Sticking to a poset-based semantics, if our value types are \DIFdelbegin \DIFdel{posets, we then }\DIFdelend \DIFaddbegin \DIFadd{interpreted as
posets, then we }\DIFaddend want our precision derivations $c : A \ltdyn A'$ to be
\DIFaddbegin \DIFadd{interpreted as a }\DIFaddend \emph{\DIFdelbegin \DIFdel{poset relations}%DIFDELCMD < \MBLOCKRIGHTBRACE %%%
\DIFdel{on $A,
A'$, that is relations between the underlying sets of $A, A'$ that are
}\DIFdelend \DIFaddbegin \DIFadd{relation}} \DIFadd{between $A$ and $A'$. It is also natural to
ask that the relation $c$ interact with the orderings on $A$ and $A'$. More
specifically, if $x' \ltdyn_{A} x$ and $x \mathrel{c} y$, then $x' \mathrel{c}
y$. Similarly, if $x \mathrel{c} y$ and $y \ltdyn_{A'} y'$, then $x \mathrel{c}
y'$. We summarize this requirement by saying $c$ is }\DIFaddend \emph{downward-closed} in
$A$ and \DIFdelbegin \emph{\DIFdel{upward-closed}} %DIFAUXCMD
\DIFdelend \DIFaddbegin \DIFadd{$\emph{upward-closed}$ }\DIFaddend in $A'$. We \DIFdelbegin \DIFdel{denote }\DIFdelend \DIFaddbegin \DIFadd{will see below why this requirement
is necessary for the graduality proof.
}

%DIF >  ...we then want our precision derivations $c : A \ltdyn A'$ to be \emph{poset relations} on $A,
%DIF >  A'$, that is, relations between the underlying sets of $A, A'$ that are
%DIF >  \emph{downward-closed} in $A$ and \emph{upward-closed} in $A'$.
%DIF >  The former means that if $x' \ltdyn_{A} x$ and $x \mathrel{c} y$, then $x' \mathrel{c} y$,
%DIF >  while the latter says that if $x \mathrel{c} y$ and $y \ltdyn_{A'} y'$, then $x \mathrel{c} y'$.

\DIFadd{We call }\DIFaddend such a relation between $A$ and $A'$ \DIFdelbegin \DIFdel{by }\DIFdelend \DIFaddbegin \DIFadd{a }\emph{\DIFadd{poset relation}}\DIFadd{, and we
write }\DIFaddend $c : A \rel A'$. The paradigmatic example of a poset relation is \DIFdelbegin \DIFdel{the poset ordering itself, $\ltdyn_A : A \rel A$, which }\DIFdelend \DIFaddbegin \DIFadd{given by
the ordering $\ltdyn_A$. We call this the reflexive poset relation and
denote it by $r(A) : A \rel A$. This }\DIFaddend is a poset relation because \DIFdelbegin \DIFdel{it }\DIFdelend \DIFaddbegin \DIFadd{the ordering
$\ltdyn_A$ }\DIFaddend is transitive.
\DIFaddbegin 

%DIF >  The paradigmatic example of a poset relation is the ``reflexive'' relation,
%DIF >  which is given by the ordering $\ltdyn_A$ and denoted by $r(A)$.

\DIFaddend %
Given two poset relations $c : A_1 \rel A_2$ and $c' : A_2 \rel A_3$, we can
form their \emph{composition} in the usual manner, i.e., $x \binrel{c \comp c'} y$ if and
only if there is an element $z : A_2$ with $x \binrel{c} z$ and $z \binrel{c'} y$.
%We write the composition of relations as $c \comp c'$.

% \max{TODO: Eric pick
%   up from here. Go into the squares for casts and how representability
%   gives us a compositional approach to proving them}

% \eric{Picked up here; Max, please read this.}

These relations are then used to model term precision. For closed programs $M
\ltdyn M' : c$ where $c : A \ltdyn A'$, the terms denote elements of the free
error domain $\sem{M},\sem{M'} : U\li \sem{A}$\DIFaddbegin \DIFadd{, }\DIFaddend and we model the term precision
semantically as the relation holding between those denoted elements $\sem{M}
\binrel{(U\li\sem{c})} \sem{M'}$\DIFdelbegin \DIFdel{where
}\DIFdelend \DIFaddbegin \DIFadd{. Here }\DIFaddend $U\li\sem{c}$ is a relational lifting of
$U\circ\li$\DIFaddbegin \DIFadd{, which extends $\sem{c}$ to account for error and computational
steps}\DIFaddend .
%
More generally, $M$ and $M'$ can be \emph{open} terms, in which case
they denote functions, not just elements. We model that as not simply
elements being related but in a typical logical-relations style as
\emph{preserving} relatedness: if they are passed in inputs related by
the domain relations then their output is related by the codomain
relation.
%
To capture this preservation of relations property, we use the notion of a
\emph{square}, which specifies a relation between two monotone functions. The
definition of square is as follows. Let $A_i, A_o, A_i', A_o'$ be
partially-ordered sets. Let $c_i : A_i \rel A_i'$ and $c_o : A_o \rel A_o'$ be
poset relations, and let $f : A_i \to A_o$ and $g : A_i' \to A_o'$ be monotone
functions. We say that $f \ltsq{c_i}{c_o} g$ if for all $x : A_i$ and $y : A_i'$
with $x \binrel{c_i} y$, we have $f(x) \binrel{c_o} g(y)$. We call this a
\emph{square} because we visualize this situation as follows:
%
% https://q.uiver.app/#q=WzAsNCxbMCwwLCJBX2kiXSxbMCwxLCJBX28iXSxbMSwwLCJBX2knIl0sWzEsMSwiQV9vJyJdLFsyLDMsImciXSxbMCwxLCJmIiwyXSxbMCwyLCJjX2kiLDAseyJzdHlsZSI6eyJib2R5Ijp7Im5hbWUiOiJiYXJyZWQifSwiaGVhZCI6eyJuYW1lIjoibm9uZSJ9fX1dLFsxLDMsImNfbyIsMix7InN0eWxlIjp7ImJvZHkiOnsibmFtZSI6ImJhcnJlZCJ9LCJoZWFkIjp7Im5hbWUiOiJub25lIn19fV1d
\[\begin{tikzcd}[ampersand replacement=\&]
  {A_i} \& {A_i'} \\
  {A_o} \& {A_o'}
  \arrow["{c_i}", "\shortmid"{marking}, no head, from=1-1, to=1-2]
  \arrow["f"', from=1-1, to=2-1]
  \arrow["g", from=1-2, to=2-2]
  \arrow["{c_o}"', "\shortmid"{marking}, no head, from=2-1, to=2-2]
\end{tikzcd}\]
%
Squares are then used to model the term precision ordering in that
$\Delta \vdash M \ltdyn N : c$ implies $\sem{M}
\ltsq{\sem{\Delta}}{U\li\sem{c}} \sem{N}$ where here $\sem{\Delta =
  x_1:c_1,\ldots}$ is the action of the product on the relations
$c_1,\ldots$.

%% % https://q.uiver.app/#q=WzAsNCxbMCwwLCJcXEdhbW1hIl0sWzAsMSwiVVxcbGkgQSJdLFsxLDAsIlxcR2FtbWEnIl0sWzEsMSwiVVxcbGkgQSciXSxbMiwzLCJcXHNlbXtOfSJdLFswLDEsIlxcc2Vte019IiwyXSxbMCwyLCJcXEdhbW1hXntcXGx0ZHlufSIsMCx7InN0eWxlIjp7ImJvZHkiOnsibmFtZSI6ImJhcnJlZCJ9LCJoZWFkIjp7Im5hbWUiOiJub25lIn19fV0sWzEsMywiVVxcbGkgYyIsMix7InN0eWxlIjp7ImJvZHkiOnsibmFtZSI6ImJhcnJlZCJ9LCJoZWFkIjp7Im5hbWUiOiJub25lIn19fV1d
%% \[\begin{tikzcd}[ampersand replacement=\&]
%% 	\Gamma \& {\Gamma'} \\
%% 	{U\li A} \& {U\li A'}
%% 	\arrow["{\Gamma^{\ltdyn}}", "\shortmid"{marking}, no head, from=1-1, to=1-2]
%% 	\arrow["{\sem{M}}"', from=1-1, to=2-1]
%% 	\arrow["{\sem{N}}", from=1-2, to=2-2]
%% 	\arrow["{U\li c}"', "\shortmid"{marking}, no head, from=2-1, to=2-2]
%% \end{tikzcd}\]

% \max{This needs to also talk about relational composition and horizontal pasting squares, since that's transitivity!}
We note that for any $c$, there is a ``vertical identity'' square $\id
\ltsq{c}{c} \id$, as can be seen by unfolding the
definition. Similarly, for any \DIFdelbegin \DIFdel{montone }\DIFdelend \DIFaddbegin \DIFadd{monotone }\DIFaddend function $f : A_i \to A_o$,
there is a ``horizontal identity'' square \DIFdelbegin \DIFdel{$f
\ltsq{\le_{A_i}}{\le_{A_o}} f$ }\DIFdelend \DIFaddbegin \DIFadd{$f
\ltsq{r(A_i)}{r(A_o)} f$ }\DIFaddend arising from the fact that $f$ is
monotone.
%
We can \emph{compose} squares vertically, i.e., if we have $f
\ltsq{c_1}{c_2} f'$ and $g \ltsq{c_2}{c_3} g'$ then we obtain a square
$g \circ f \ltsq{c_1}{c_3} g' \circ f'$. Likewise, we can compose
squares horizontally: If $f \ltsq{c_i}{c_o} g$ and $g
\ltsq{c_i'}{c_o'} h$, then we obtain a square $f \ltsq{c_i \comp
  c_i'}{c_o \comp c_o'} h$. Horizontal composition of squares
corresponds to transitivity of term precision. While we do not dwell
on categorical abstractions in this work, we note that this structure
of functions, relations and squares forms a locally thin \emph{double
category} as used in previous work \cite{new-licata18}.

Now that we have an intended model of term precision in the form of
squares, we can see what is required to give a compositional proof of
graduality. Most of the cases of term precision are just congruence
and follow easily. The core of the proof of the graduality property is
in proving the validity of the \emph{cast rules} \DIFdelbegin \DIFdel{UpL, UpR, DnL and
DnR}\DIFdelend \DIFaddbegin \DIFadd{$\upl$, $\upr$, $\dnl$, and
$\dnr$}\DIFaddend .
%
These rules specify a relationship between the semantics of type
precision derivations $c$ and the corresponding casts $\upc c, \dnc c$ (see Section
\ref{sec:GTLC}).
%
To prove graduality compositionally then, we will need to codify in an
extra property on $\sem{c}$ that the relation $\sem{c}$ has some
correspondence with the casts $\sem{\upc c}, \sem{\dnc c}$.
%
What property should we require?

Consider the \DIFdelbegin \DIFdel{UpL }\DIFdelend \DIFaddbegin \DIFadd{$\upl$ }\DIFaddend rule. It states that if
$c : A_1 \ltdyn A_2$ and $c' : A_2 \ltdyn A_3$ and $M : A_1$ is
related to $N : A_3$ via the composition $c \comp c'$, then the upcast
$\upc{c} M$ is related to $N$ via $c'$. Since the upcasts are pure
functions in the semantics this can be deduced from the existence of
the following square, where the relation on the top is $\sem{c} \comp \sem{c'}$:
% \max{the square doesn't make sense because we
%   haven't discussed horizontal composition of relations or that this
%   notation means horizontal composition of relations}
%
% https://q.uiver.app/#q=WzAsNSxbMCwwLCJBXzEiXSxbMSwwLCJBXzIiXSxbMiwwLCJBXzMiXSxbMCwxLCJBXzIiXSxbMiwxLCJBXzMiXSxbMyw0LCJjJyIsMix7InN0eWxlIjp7ImJvZHkiOnsibmFtZSI6ImJhcnJlZCJ9LCJoZWFkIjp7Im5hbWUiOiJub25lIn19fV0sWzAsMywiXFx1cGN7Y30iLDJdLFsyLDQsIlxcaWQiXSxbMCwxLCJjIiwwLHsic3R5bGUiOnsiYm9keSI6eyJuYW1lIjoiYmFycmVkIn0sImhlYWQiOnsibmFtZSI6Im5vbmUifX19XSxbMSwyLCJjJyIsMCx7InN0eWxlIjp7ImJvZHkiOnsibmFtZSI6ImJhcnJlZCJ9LCJoZWFkIjp7Im5hbWUiOiJub25lIn19fV1d
\[\begin{tikzcd}[ampersand replacement=\&]
  {\sem{A_1}} \& {\sem{A_2}} \& {\sem{A_3}} \\
  {\sem{A_2}} \&\& {\sem{A_3}}
  \arrow["\sem{c}", "\shortmid"{marking}, no head, from=1-1, to=1-2]
  \arrow["{\sem{\upc{c}}}"', from=1-1, to=2-1]
  \arrow["{\sem{c'}}", "\shortmid"{marking}, no head, from=1-2, to=1-3]
  \arrow["\id", from=1-3, to=2-3]
  \arrow["{\sem{c'}}"', "\shortmid"{marking}, no head, from=2-1, to=2-3]
\end{tikzcd}\]
While at first look this seems to be specifying a relationship between
$\sem{c}$ and $\sem{\upc c}$, there is a problem: it also quantifies
\DIFdelbegin \DIFdel{involves }\DIFdelend \DIFaddbegin \DIFadd{over }\DIFaddend an arbitrary other relation $\sem{c'}$! This means we cannot
require the existence of this square as part of the definition of a
relation between value types, as it is self-referential. This would
seem to imply that we cannot give a compositional model for
graduality. However, New and Licata observed that in the presence of
transitivity, we can \emph{derive} the above squares from simpler ones
that do not involve composition of relations. Below \DIFdelbegin \DIFdel{is the square
corresponding to UpL:
%DIF <  https://q.uiver.app/#q=WzAsNCxbMCwwLCJBXzEiXSxbMCwxLCJBXzIiXSxbMSwwLCJBXzIiXSxbMSwxLCJBXzIiXSxbMCwyLCJjIiwwLHsic3R5bGUiOnsiYm9keSI6eyJuYW1lIjoiYmFycmVkIn0sImhlYWQiOnsibmFtZSI6Im5vbmUifX19XSxbMSwzLCJcXGxlX3tBXzJ9IiwyLHsic3R5bGUiOnsiYm9keSI6eyJuYW1lIjoiYmFycmVkIn0sImhlYWQiOnsibmFtZSI6Im5vbmUifX19XSxbMCwxLCJcXHVwY3tjfSIsMl0sWzIsMywiXFxpZCJdXQ==
}\[\DIFdel{\begin{tikzcd}[ampersand replacement=\&]
	{A_1} \& {A_2} \\
	{A_2} \& {A_2}
	\arrow["c", "\shortmid"{marking}, no head, from=1-1, to=1-2]
	\arrow["{\upc{c}}"', from=1-1, to=2-1]
	\arrow["\id", from=1-2, to=2-2]
	\arrow["{\le_{A_2}}"', "\shortmid"{marking}, no head, from=2-1, to=2-2]
\end{tikzcd}}\]%DIFAUXCMD
\DIFdelend \DIFaddbegin \DIFadd{are the simpler squares
corresponding to $\upl$, $\upr$, $\dnl$, and $\dnr$:
}\begin{center}
  \begin{tabular}{ c | c | c | c} 
    %DIF >  UpL
    %DIF >  https://q.uiver.app/#q=WzAsNCxbMCwwLCJBXzEiXSxbMCwxLCJBXzIiXSxbMSwwLCJBXzIiXSxbMSwxLCJBXzIiXSxbMCwyLCJjIiwwLHsic3R5bGUiOnsiYm9keSI6eyJuYW1lIjoiYmFycmVkIn0sImhlYWQiOnsibmFtZSI6Im5vbmUifX19XSxbMSwzLCJyKEFfMikiLDIseyJzdHlsZSI6eyJib2R5Ijp7Im5hbWUiOiJiYXJyZWQifSwiaGVhZCI6eyJuYW1lIjoibm9uZSJ9fX1dLFswLDEsIlxcdXBje2N9IiwyXSxbMiwzLCJcXGlkIl0sWzYsNywiXFx0ZXh0e1VwTH0iLDMseyJzaG9ydGVuIjp7InNvdXJjZSI6MjAsInRhcmdldCI6MjB9LCJzdHlsZSI6eyJib2R5Ijp7Im5hbWUiOiJub25lIn0sImhlYWQiOnsibmFtZSI6Im5vbmUifX19XV0=
    \begin{tikzcd}[ampersand replacement=\&]
      {A_1} \& {A_2} \\
      {A_2} \& {A_2}
      \arrow["c", "\shortmid"{marking}, no head, from=1-1, to=1-2]
      \arrow[""{name=0, anchor=center, inner sep=0}, "{\upc{c}}"', from=1-1, to=2-1]
      \arrow[""{name=1, anchor=center, inner sep=0}, "\id", from=1-2, to=2-2]
      \arrow["{r(A_2)}"', "\shortmid"{marking}, no head, from=2-1, to=2-2]
      \arrow["{\text{UpL}}"{marking, allow upside down}, draw=none, from=0, to=1]
  \end{tikzcd} &
    \DIFaddend %
    \DIFdelbegin \DIFdel{This square, with a dual one related to UpR, }\DIFdelend %DIF >  UpR
    %DIF >  https://q.uiver.app/#q=WzAsNCxbMCwwLCJBXzEiXSxbMCwxLCJBXzEiXSxbMSwwLCJBXzEiXSxbMSwxLCJBXzIiXSxbMCwyLCJyKEFfMSkiLDAseyJzdHlsZSI6eyJib2R5Ijp7Im5hbWUiOiJiYXJyZWQifSwiaGVhZCI6eyJuYW1lIjoibm9uZSJ9fX1dLFsxLDMsImMiLDIseyJzdHlsZSI6eyJib2R5Ijp7Im5hbWUiOiJiYXJyZWQifSwiaGVhZCI6eyJuYW1lIjoibm9uZSJ9fX1dLFswLDEsIlxcaWQiLDJdLFsyLDMsIlxcdXBjIGMiXSxbNiw3LCJcXHRleHR7VXBSfSIsMyx7InNob3J0ZW4iOnsic291cmNlIjoyMCwidGFyZ2V0IjoyMH0sInN0eWxlIjp7ImJvZHkiOnsibmFtZSI6Im5vbmUifSwiaGVhZCI6eyJuYW1lIjoibm9uZSJ9fX1dXQ==
    \DIFaddbegin \begin{tikzcd}[ampersand replacement=\&]
      {A_1} \& {A_1} \\
      {A_1} \& {A_2}
      \arrow["{r(A_1)}", "\shortmid"{marking}, no head, from=1-1, to=1-2]
      \arrow[""{name=0, anchor=center, inner sep=0}, "\id"', from=1-1, to=2-1]
      \arrow[""{name=1, anchor=center, inner sep=0}, "{\upc c}", from=1-2, to=2-2]
      \arrow["c"', "\shortmid"{marking}, no head, from=2-1, to=2-2]
      \arrow["{\text{UpR}}"{marking, allow upside down}, draw=none, from=0, to=1]
    \end{tikzcd} &
    %DIF > 
    %DIF > 
    %DIF > 
    %DIF >  DnL
    %DIF >  https://q.uiver.app/#q=WzAsNCxbMCwwLCJBXzIiXSxbMCwxLCJBXzEiXSxbMSwwLCJBXzIiXSxbMSwxLCJBXzIiXSxbMCwyLCJyKEFfMikiLDAseyJzdHlsZSI6eyJib2R5Ijp7Im5hbWUiOiJiYXJyZWQifSwiaGVhZCI6eyJuYW1lIjoibm9uZSJ9fX1dLFsxLDMsImMiLDIseyJzdHlsZSI6eyJib2R5Ijp7Im5hbWUiOiJiYXJyZWQifSwiaGVhZCI6eyJuYW1lIjoibm9uZSJ9fX1dLFswLDEsIlxcZG5jIGMiLDJdLFsyLDMsIlxcaWQiXSxbNiw3LCJcXHRleHR7RG5MfSIsMyx7InNob3J0ZW4iOnsic291cmNlIjoyMCwidGFyZ2V0IjoyMH0sInN0eWxlIjp7ImJvZHkiOnsibmFtZSI6Im5vbmUifSwiaGVhZCI6eyJuYW1lIjoibm9uZSJ9fX1dXQ==
    \begin{tikzcd}[ampersand replacement=\&]
      {A_2} \& {A_2} \\
      {A_1} \& {A_2}
      \arrow["{r(A_2)}", "\shortmid"{marking}, no head, from=1-1, to=1-2]
      \arrow[""{name=0, anchor=center, inner sep=0}, "{\dnc c}"', from=1-1, to=2-1]
      \arrow[""{name=1, anchor=center, inner sep=0}, "\id", from=1-2, to=2-2]
      \arrow["c"', "\shortmid"{marking}, no head, from=2-1, to=2-2]
      \arrow["{\text{DnL}}"{marking, allow upside down}, draw=none, from=0, to=1]
    \end{tikzcd} &
    %DIF > 
    %DIF >  DnR
    %DIF >  https://q.uiver.app/#q=WzAsNCxbMCwwLCJBXzEiXSxbMCwxLCJBXzEiXSxbMSwwLCJBXzIiXSxbMSwxLCJBXzEiXSxbMCwyLCJjIiwwLHsic3R5bGUiOnsiYm9keSI6eyJuYW1lIjoiYmFycmVkIn0sImhlYWQiOnsibmFtZSI6Im5vbmUifX19XSxbMSwzLCJyKEFfMSkiLDIseyJzdHlsZSI6eyJib2R5Ijp7Im5hbWUiOiJiYXJyZWQifSwiaGVhZCI6eyJuYW1lIjoibm9uZSJ9fX1dLFswLDEsIlxcaWQiLDJdLFsyLDMsIlxcZG5jIGMiXSxbNiw3LCJcXHRleHR7RG5SfSIsMyx7InNob3J0ZW4iOnsic291cmNlIjoyMCwidGFyZ2V0IjoyMH0sInN0eWxlIjp7ImJvZHkiOnsibmFtZSI6Im5vbmUifSwiaGVhZCI6eyJuYW1lIjoibm9uZSJ9fX1dXQ==
    \begin{tikzcd}[ampersand replacement=\&]
      {A_1} \& {A_2} \\
      {A_1} \& {A_1}
      \arrow["c", "\shortmid"{marking}, no head, from=1-1, to=1-2]
      \arrow[""{name=0, anchor=center, inner sep=0}, "\id"', from=1-1, to=2-1]
      \arrow[""{name=1, anchor=center, inner sep=0}, "{\dnc c}", from=1-2, to=2-2]
      \arrow["{r(A_1)}"', "\shortmid"{marking}, no head, from=2-1, to=2-2]
      \arrow["{\text{DnR}}"{marking, allow upside down}, draw=none, from=0, to=1]
    \end{tikzcd}
  \end{tabular}
\end{center}
%DIF > 
\DIFadd{These squares }\DIFaddend say that the relation $c$ is \emph{representable} by the upcast
morphism $\upc{c}$, essentially that the relation is a kind of \emph{graph} of
the function in that $x \binrel{\sem c} y$ if and only if $\sem{\upc c}(x)
\leq_{A_2} y$ \cite{shulman2007}.

We can use \DIFdelbegin \DIFdel{this simpler }\DIFdelend \DIFaddbegin \DIFadd{the simpler $\upl$ }\DIFaddend square to derive the original \DIFdelbegin \DIFdel{UpL square as long as we can
horizontally compose }\DIFdelend \DIFaddbegin \DIFadd{$\upl$ square by
horizontally composing }\DIFaddend with the identity square for $c'$:
%
% https://q.uiver.app/#q=WzAsNixbMCwwLCJBXzEiXSxbMSwwLCJBXzIiXSxbMiwwLCJBXzMiXSxbMCwxLCJBXzIiXSxbMiwxLCJBXzMiXSxbMSwxLCJBXzIiXSxbMCwxLCJjIiwwLHsic3R5bGUiOnsiYm9keSI6eyJuYW1lIjoiYmFycmVkIn0sImhlYWQiOnsibmFtZSI6Im5vbmUifX19XSxbMSwyLCJjJyIsMCx7InN0eWxlIjp7ImJvZHkiOnsibmFtZSI6ImJhcnJlZCJ9LCJoZWFkIjp7Im5hbWUiOiJub25lIn19fV0sWzMsNSwiXFxsdGR5bl97QV8yfSIsMix7InN0eWxlIjp7ImJvZHkiOnsibmFtZSI6ImJhcnJlZCJ9LCJoZWFkIjp7Im5hbWUiOiJub25lIn19fV0sWzUsNCwiYyciLDIseyJzdHlsZSI6eyJib2R5Ijp7Im5hbWUiOiJiYXJyZWQifSwiaGVhZCI6eyJuYW1lIjoibm9uZSJ9fX1dLFswLDMsIlxcdXBje2N9IiwyXSxbMiw0LCJcXGlkIl0sWzEsNSwiXFxpZCIsMl1d
\[\begin{tikzcd}[ampersand replacement=\&]
	{A_1} \& {A_2} \& {A_3} \\
	{A_2} \& {A_2} \& {A_3}
	\arrow["c", "\shortmid"{marking}, no head, from=1-1, to=1-2]
	\arrow["{\upc{c}}"', from=1-1, to=2-1]
	\arrow["{c'}", "\shortmid"{marking}, no head, from=1-2, to=1-3]
	\arrow["\id"', from=1-2, to=2-2]
	\arrow["\id", from=1-3, to=2-3]
	\arrow["{\ltdyn_{A_2}}"', "\shortmid"{marking}, no head, from=2-1, to=2-2]
	\arrow["{c'}"', "\shortmid"{marking}, no head, from=2-2, to=2-3]
\end{tikzcd}\]
%
Then to complete the proof, we observe that the composition of $\ltdyn_{A_2}$ with $c'$ is equal to
$c'$, because $c'$ is \emph{downward-closed} under the relation on $A_2$.
\DIFaddbegin \DIFadd{An analogous argument shows that we can derive the original $\upr$ square from the simpler version of $\upr$,
this time using the fact that poset relations are }\emph{\DIFadd{upward-closed}}\DIFadd{.
}\DIFaddend % Thus, we
% obtain the original UpL square by vertically composing the square shown above
% with the square $\id \ltsq{\le_{A_2}\, c'}{c'} \id$ on the bottom.

\DIFdelbegin \DIFdel{The actual place where transitivity of the relations on the value types is
needed in this argument occurs when $A_3 = A_2$ and we take $c'$ to be
$\ltdyn_{A_2}$ (recall that $c'$ is universally quantified, so it can be any
relation). The fact that $\ltdyn_{A_2}$ is downward-closed is equivalent to it
being transitive. In other words,
the relations $\ltdyn_A$ on value types must be
poset relations }\DIFdelend \DIFaddbegin \DIFadd{To recap, we have shown that }\DIFaddend in order to carry out the \DIFdelbegin \DIFdel{above compositional construction.
}%DIFDELCMD < 

%DIFDELCMD < %%%
\DIFdel{So we see that in order to get a sensible condition that can give us a
compositional definition of a well-behaved relation, some amount of
transitivereasoning on squares is required, and }\DIFdelend \DIFaddbegin \DIFadd{graduality proof
compositionally, we need that the relations are closed under the ordering on
each side. In particular, this implies that the ordering relations $\ltdyn_{A}$
must be transitive, }\DIFaddend so we cannot entirely drop posets from our semantics.

%DIF >  The actual place where transitivity of the relations on the value types is
%DIF >  needed in this argument occurs when $A_3 = A_2$ and we take $c'$ to be
%DIF >  $\ltdyn_{A_2}$ (recall that $c'$ is universally quantified, so it can be any
%DIF >  relation). The fact that $\ltdyn_{A_2}$ is downward-closed is equivalent to it
%DIF >  being transitive. In other words, the relations $\ltdyn_A$ on value types must be
%DIF >  poset relations in order to carry out the above compositional construction.
\DIFaddbegin 

%DIF >  So we see that in order to get a sensible condition that can give us a
%DIF >  compositional definition of a well-behaved relation, some amount of
%DIF >  transitive reasoning on squares is required, and so we cannot entirely
%DIF >  drop posets from our semantics.

\DIFaddend \subsection{Resolution: Splitting Error Ordering and Bisimilarity}\label{sec:lock-step-and-weak-bisim}

Given that some amount of transitive reasoning is essential for
modeling type and term precision compositionally, we need to revisit
our ordering relation on $\li A$ to get one that is
transitive. Theorem \ref{thm:no-go} tells us, however, that any non-trivial such
relation must either not be a congruence with respect to $\theta$, or
must not be left- and right-step-insensitive.
\footnote{Technically, the theorem only implies that we need to drop
one of left- or right-step-insensitivity. We might then attempt
to work with two separate ``step-semi-sensitive'' relations on $\li
A$, each closed under delays on the left and right respectively. This
has some similarities to prior work on logical relations models
\cite{new-licata-ahmed2019,new-giovannini-licata-2022}, but we do not take this
approach because the relations still do not seem to be provably
transitive, though we have not been able to prove that they are not
transitive either.}
% \max{Eric check if this footnote is correct}}
%
We cannot forego the property of being a $\theta$-congruence, as
without this we would not be able to prove basic properties of the
relation using \lob-induction, e.g., that the extension $f^\dagger$ is
\emph{monotone}, which is used pervasively in the proof of graduality.
%
Thus, we choose to sacrifice left and right step-insensitivity. That
is, we will define an ordering relation that requires terms to be in
``lock-step''. In order for two computations to be related in this
ordering, they must have the exact same stepping behavior
(unless/until the left-hand side results in an error).

More formally, we define the \emph{lock-step error-ordering}, with the idea
being that two computations $l$ and $l'$ are related if they are in lock-step
with regard to their intensional behavior, up to $l$ erroring. Figure
\ref{fig:lock-step-error-ordering} gives the definition of this relation.

\begin{figure}
  \begin{minipage}{0.5\textwidth}
    \fbox{$l_1 \ltls l_2$}
    \begin{align*}
        &\eta\, x \ltls \eta\, y \text{ if } 
            x \mathbin{\ltdyn_A} y \\
        %		
        &\mho \ltls l' \\
        %
        &\theta\, \tilde{l} \ltls \theta\, \tilde{l'} \text{ if } 
            \later_t (\tilde{l}_t \ltls \tilde{l'}_t)
  \end{align*}\end{minipage}\begin{minipage}{0.5\textwidth}
\fbox{$l_1 \bisim l_2$}
    \begin{align*}
        \mho \bisim \mho &\text{ iff } \top \\
      %
        \eta\, x \bisim \eta\, y &\text{ iff } x \bisim_A y \\
      %		
        \theta\, \tilde{x} \bisim \theta\, \tilde{y} &\text{ iff } \later_t (\tilde{x}_t \bisim \tilde{y}_t) \\
      %	
        \theta\, \tilde{x} \bisim \mho &\text{ iff } \exists n. \theta\, \tilde{x} = \delta^n(\mho)\\
      %	
        \theta\, \tilde{x} \bisim \eta\, y &\text{ iff } \exists n. \exists x \bisim_A y.
          (\theta\, \tilde{x} = \delta^n(\eta\, x))\\
      %
        \mho \bisim \theta\, \tilde{y} &\text{ iff } \exists n. \theta\, \tilde{y} = \delta^n(\mho) \\
      %	
        \eta\, x \bisim \theta\, \tilde{y} &\text { iff } \exists n. \exists y \bisim_A x. (\theta\, \tilde{y} = \delta^n (\eta\, y))
      \end{align*}
  \end{minipage}
    \caption{lock-step error ordering and weak bisimilarity}
    \label{fig:lock-step-error-ordering}
\end{figure}

When both sides are $\eta$, then we check that the returned values are related
in $\ltdyn_A$. The error term $\mho$ is the least element. Lastly, if both sides step
(i.e., are a $\theta$) then we compare the resulting computations one time step
later.
%
It is straightforward to prove using \lob-induction that this relation is
reflexive, transitive and anti-symmetric given that the \DIFdelbegin \DIFdel{underyling relation $R$
}\DIFdelend \DIFaddbegin \DIFadd{underlying relation $\ltdyn_A$
}\DIFaddend has those properties. The lock-step ordering is therefore the partial ordering
we will associate with $\li A$.
%
More generally we can define a heterogeneous version of this ordering that lifts
poset relation $c : A \rel A'$ to a poset relation $\li c : \li A \rel \li A'$.
\DIFaddbegin \DIFadd{The definition is nearly identical, except that in the case when both sides are $\eta$ we
check that the returned values are related in $c$.
}\DIFaddend 

However, we also cannot completely drop step-\emph{insensitivity} from
our model. The graduality property itself is not sensitive to steps,
and in fact there are terms related by term precision that take
differing numbers of steps.
%
The offending step arises precisely from the place where our
definition of $D$ uses a $\later$: the function case.
%
Consider the representability condition corresponding to the \DIFdelbegin \DIFdel{DnL }\DIFdelend \DIFaddbegin \DIFadd{$\dnl$ }\DIFaddend rule
for $\iarr \colon \dyntodyn \ltdyn\, \dyn$. This is the square $\sem{\dnc{\iarr}} \ltsq{r({\li D})}{\li\sem{\iarr}} \id$.
%% % https://q.uiver.app/#q=WzAsNCxbMCwwLCJcXGxpIEQiXSxbMCwxLCJcXGxpIFUoRCBcXGFyciBcXGxpIEQpIl0sWzEsMCwiXFxsaSBEIl0sWzEsMSwiXFxsaSBEIl0sWzAsMiwiXFxsZXFfe1xcbGkgRH0iLDAseyJzdHlsZSI6eyJib2R5Ijp7Im5hbWUiOiJiYXJyZWQifSwiaGVhZCI6eyJuYW1lIjoibm9uZSJ9fX1dLFsxLDMsIlxcbGlcXGlhcnIiLDIseyJzdHlsZSI6eyJib2R5Ijp7Im5hbWUiOiJiYXJyZWQifSwiaGVhZCI6eyJuYW1lIjoibm9uZSJ9fX1dLFswLDEsIlxcZG5je1xcaWFycn0iLDJdLFsyLDMsIlxcaWQiXV0=
%% \[\begin{tikzcd}[ampersand replacement=\&]
%% 	{\li D} \& {\li D} \\
%% 	{\li U(D \arr \li D)} \& {\li D}
%% 	\arrow["{\leq_{\li D}}", "\shortmid"{marking}, no head, from=1-1, to=1-2]
%% 	\arrow["{\sem{\dnc{\iarr}}}"', from=1-1, to=2-1]
%% 	\arrow["\id", from=1-2, to=2-2]
%% 	\arrow["\li\sem{\iarr}"', "\shortmid"{marking}, no head, from=2-1, to=2-2]
%% \end{tikzcd}\]
For this square to be valid with the lock-step ordering the left and
right hand side would need to take the same number of steps, at least
when the left-hand side does not error.
%
However\DIFaddbegin \DIFadd{, }\DIFaddend the right hand side \emph{always} takes $0$ steps, as it is
the identity function, whereas on the left, our definition of
$\sem{\dnc{\iarr}}$ had to take an observable step when its input is a
function: this was inherent to the fact that the function case of the
dynamic type is guarded by a later.

\DIFdelbegin \DIFdel{However, observe }\DIFdelend \DIFaddbegin \DIFadd{Observe }\DIFaddend that we can remedy this particular situation by
replacing the $\id$ on the right hand side by an innocuous function
$\delta^* = (\delta \circ \eta)^\dagger$ that on an input value takes
a single computational step, but is otherwise the identity
function. If we were to ignore computational steps, then this function
\emph{would} be the identity function. We call such a function that is
the identity function except for the introduction of computational
steps a \emph{perturbation}.

We formalize this property of being equivalent ``except for steps''
with a second relation on the free error domain $\li A$: \emph{weak
bisimilarity}, defined in Figure~\ref{fig:lock-step-error-ordering}, which is
parameterized by a binary relation $\bisim_A$ on $A$
\cite{mogelberg-paviotti2016}.
% For a type $X$, we define a relation on $\li X$, called ``weak bisimilarity",
% written $l \bisim l'$. 
Two errors are bisimilar, and when both sides are $\eta$, we ensure
that the underlying values are bisimilar in the underlying
bisimilarity relation on $A$. When both sides are thinking, we ensure
the terms are bisimilar later.  Most importantly, when one side is
thinking but the other terminates at $\eta x$ (i.e., one side steps),
we stipulate that the $\theta$-term runs to $\eta y$ where $x$ is
related to $y$. And similarly, if one side is thinking and the other
errors, we ensure the thinking side eventually errors.

It can be shown (by \lob-induction) that weak bisimilarity is
reflexive and symmetric. Since it is non-trivial, a $\theta$
congruence and step-insensitive on both sides, by
Theorem~\ref{thm:no-go}, we also know that it is \emph{not}
transitive.
%
We will then require our denotations of types to be not just posets,
but posets additionally equipped with a reflexive symmetric relation
$\bisim_A$.

Then we can refine our denotation of term precision $\Delta \vdash M
\ltdyn N : c$ to mean not that $\sem{M}$ and $\sem{N}$ are necessarily
in a lock-step error ordering directly, but that they can be
``synchronized'' to do so, that is that there exist $f \bisim \sem{M}$
and $g \bisim \sem{N}$ such that $f \ltsq{\sem{\Delta}}{U\li\sem{c}}
g$.
%
Note that this in turn implies that $\sem{M}$ and $\sem{N}$ are
related in the original \emph{step-insensitive} error ordering that we
sought to prove!
%
We additionally weaken our representability squares so that our
upcasts and downcasts are not required to be in a lock-step ordering
with an identity but instead to be in a lock-step ordering with some
perturbation, that is a function \emph{bisimilar to} the identity. We
call this weakened form of representability
\emph{quasi}-representability.
%
Then if our relations are quasi-representable we can in fact prove the
validity of our cast rules.
%
The remainder of the proof of graduality then is to show that this
property of being quasi-representable is itself compositional: that
all of the constructions we have on type precision preserve the
property of being quasi-representable.
%
While this is true for representability, it is not quite true for
quasi-representability. The reason is that perturbations are not quite
as well behaved as \emph{actual} identity functions. To solve this
final issue, we attach one final piece of information to our types
$A$: a type of \emph{syntactic} perturbations that represent functions
bisimilar to the identity, but are presented as data so that we can
perform operations such as composition and functorial actions on them.
%
With this notion of syntactic perturbation, and additional
requirements that the relations interact well with perturbations, we
finally achieve our desired result: a compositional,
syntax-independent proof of graduality.


%% Importantly, we note that as with the original ordering $\semltbad$ defined at
%% the beginning of this section, the relation just defined is not transitive
%% (again as a result of the above no-go theorem). It may therefore seem that we
%% have not solved the original issue we faced. In a sense, this is true. The model
%% of gradual typing that we end up with will not support \emph{extensional}
%% compositional reasoning. However, by decomposing the denotation of term
%% precision in the above manner, we can employ \emph{intensional} compositional
%% reasoning by working with the lock-step error ordering. Thus, the tentative plan
%% going forward will be to carry out the proofs compositionally using the
%% lock-step ordering. We will then apply the closure under weak bisimilarity in
%% the denotation of term precision given above to handle the aspects involving
%% stepping.

%% Notice that this delay
%% function is weakly bisimilar to the identity function in that $\delta^*\, x
%% \bisim x$ for all $x$. \footnote{This follows from the fact that $\delta$ is
%% weakly bisimilar to the identity ($\delta\, x \bisim x$ for all $x$) and that
%% $-^\dagger$ preserves weak bisimilarity.} Thus, it will disappear when we apply
%% the bisimilarity-closure construction described above. That is, this delay makes
%% no difference in the extensional setting, but its presence is crucial in the
%% intensional setting.


%% The precise definition of $\iarr$ as a poset relation is not important
%% for this discussion.
%% %

%% \eric{We may want to introduce error domain relations and squares at this point,
%% since the squares for DnL and DnR involve lift.}
%
%% We now recall the definition of the semantics of a downcast given in Section
%% \ref{sec:term-interpretation}. The downcast on the left will insert a $\theta$
%% in the case where the value of type $D$ is a later-function $\tilde{f}$.
%% %
%% Thus, in order for the downcast to be related to the RHS in the lock-step error
%% ordering on $\li (D \to \li D)$, the RHS must be of the form $\theta(\dots)$,
%% and moreover, after one time step, the argument of $\theta$ on the LHS must be
%% related to the argument of $\theta$ on the RHS. As it stands now, this need not
%% be the case, e.g., if the RHS is of the form $\eta(\dots)$. Thus, we conclude
%% that the DnL rule does not in general hold under the lock-step error ordering.


%% \max{resume here}
%% To deal with
%% terms that take differing numbers of steps, we then define a separate
%% relation called \emph{weak bisimilarity} that relates terms that are
%% extensionally equal and may differ by a finite number of delays. Then,
%% the semantics of the error ordering for the guarded lift monad will
%% involve a combination of these two relations: a ``closure'' of the
%% lock-step error ordering under weak bisimilarity on both sides.
%% % Although the combined relation will not be transitive (for the same reason that
%% % $\semltbad$ is not transitive), this...
%% This decomposition has the advantage that we can recover some transitive
%% reasoning and push much of the reasoning about stepping to the margins of the
%% development.

%% % ...to carry out the proofs using the lock-step ordering and handle the stepping
%% % behavior separately via weak bisimilarity. In the end we combine everything
%% % together using the above denotation for term precision


%% \subsection{Extending the Model to Higher-Order Types}
%% \max{this subsection doesn't really say anything concrete. If we want to talk about perturbations, actually talk about them}

%% % At this point, we could carry out the full construction of a concrete relational
%% % model with these additional features. 

%% The above techniques give a semantic interpretation of term precision for closed
%% terms of base type. We now consider how to extend these constructions to
%% potentially-open terms of \emph{all} types. In particular, in order to model
%% higher-order data types (i.e., functions) we need to equip \emph{every} semantic
%% object with not only a partial ordering relation but also a ``bisimilarity''
%% relation that is reflexive and symmetric. We must similarly equip every object
%% with a structured set of delays that can be inserted and manipulated to ensure
%% the appropriate cast rules hold at all higher-order types. For example, the
%% upcast for a derivation $c_i \ra c_o$ involves the downcast corresponding to
%% $c_i$ and the upcast corresponding to $c_o$. We must be able to insert delays in
%% a functorial manner to mimic the structure of the casts. In the next section, we
%% will make the necessary definitions and describe the relevant constructions.

% For the sake of reusability and modularity, rather than carry out this
% construction in the concrete setting developed here, we will instead return to
% the abstract setting and adjust our definition of model to account for these
% requirements. We will break the construction into smaller steps and isolate the
% pieces that require the techniques of SGDT from those that do not. Then with
% this framework at our disposal, we will return to the construction of a concrete
% model in Section \ref{sec:concrete-model}.
% , taking as a starting point the definitions introduced in the current section.

% Thus in the next section we define revised notions of a model of
% gradual typing based on the lessons learned in the present section. 

% The lack of transitivity presents a major barrier towards giving a compositional
% semantics to gradual typing: if our relations are not transitive, then we cannot
% compose squares horizontally. But we have argued above that horizontal
% composition is essential for modelling the axioms of gradual typing in a
% syntax-independent manner.

% In particular, to model the UpL/UpR/DnL/DnR rules
% for casts in the presence of transitivity, it is sufficient to establish the
% validity of simpler rules that do not ``build in'' composition. We can then use
% transitivity to derive the original versions of the rules. On the other hand,
% without transitivity we must instead validate the cast rules in the model ``from
% scratch''. In fact, we cannot even define the semantic interpretation of type
% precision in a syntax-independent manner. The issue is that in the definition of
% relation, the requirement that it be representable now involves quantifying over
% all other relations, which is circular.

% So, although this approach would suffice for proving graduality, it lacks the
% compositionality that we seek in a reusable framework for the semantics of
% gradual typing.

%Because it is convenient to make use of transitive reasoning in proving graduality, ...


% Doing so, we define a \emph{lock-step} error ordering, where roughly speaking,
% in order for computations to be related, they must have the same stepping
% behavior. We then formulate a separate relation, \emph{weak bisimilarity}, that
% relates computations that are extensionally equal and may differ only in their
% stepping behavior. % up to a finite number of \delta's
% Finally, the semantics of term precision will involve a combination of these two
% relations, a sort of closure of the lock-step ordering under weak bisimilarity
% on both sides. This decomposition has the advantage that we can recover some
% transitive reasoning and push the parts involving stepping to the margins of the
% development.



\begin{comment}
\subsection{Modeling Type and Term Precision}
\max{TODO: figure out where this stuff goes}
% We will model type precision $c : A \ltdyn A'$ as a relation between the sets
% $\sem{A}$ and $\sem{A'}$. 

To model term precision, we begin by equipping the denotation of every type with
an ordering relation. Since term precision is reflexive and transitive, and
since we identify terms that are equi-precise, we model value types as
partially-ordered sets and values $\Gamma \vdash V : A$ as \emph{monotone}
functions. Analogously, every error domain is now equipped with a partial
ordering for which the error element is the bottom element, and the map
$\theta_B : \laterhs B \to B$ is now required to be monotone. 
%
Morphisms of error domains are morphisms of the underlying partially-ordered
sets that preserve the error element and $\theta$ map, as was the case in the
previous section.

We model a type precision relation $c : A \ltdyn A'$ as a \emph{monotone
relation}, i.e., a relation $c$ that is upward-closed under the relation on $A'$
and downward-closed under the relation on $A$. We denote such a relation between
$A$ and $A'$ by $c : A \rel A'$. The relation on the poset $A$ is denoted
$r(A)$.
% This extends to products $c_1 \times c_2$ in the obvious way.
Composition of relations on predomains is the usual relational composition
(which is truncated to be propositional).



\subsubsection{Modelling Term Precision}\label{sec:modeling-term-precision}




It remains to define $Fc$, i.e., the action of $F$ on relations.
%
% However, lifting a relation between $A$ and $A'$ to a relation between $\li A$
% and $\li A'$ proves problematic if we allow computations that take different
% numbers of steps to be related. To illustrate the issue, let us define an
% ordering $\semltbad$ between $\li A$ and $\li A'$; we call this the
% \emph{step-insensitive error ordering}. 
%
We note that the graduality property and the axioms of the inequational theory
are independent of the intensional stepping behavior of terms, so our ultimate
notion that interprets term precision will need to be oblivious to stepping as
well.
%
To that end, let us define a \emph{step-insensitive error ordering} $\semltbad$
between $\li A$ and $\li A'$; The ordering is parameterized by an ordering
relation $\le$ between $A$ and $A'$. The definition is by guarded recursion and
is shown in Figure \ref{fig:step-insensitive-error-ordering}. Recall that
$\delta : \li A \to \li A$ is defined by $\delta = \theta_A \circ \nxt$.

Two computations that immediately return $(\eta)$ are related if the returned
values are related in the underlying ordering. The computation that errors
$(\mho)$ is the least term in the ordering. If both sides step (i.e., both sides
are $\theta$), then we allow one time step to pass and compare the resulting
terms (this is where use the relation defined ``later'' and is why we employ
guarded recursion to define the relation).
%
Lastly, if one side steps and the other immediately returns a value, then in
order for these terms to be related, the side that steps must terminate with a
value in some finite number of steps $n$, and that value must be related to the
value returned by the other side. Likewise, if the LHS steps and the RHS
immediately errors, then in order to be related, the LHS must eventually
terminate with error.

\end{comment}


% \input{semantics}

%% \input{ordering}

%% \input{graduality}

% \input{categorical-model}

%\input{denotational-model}
%\input{terms}
\section{A Guarded Model of Graduality}\label{sec:concrete-relational-model}

Next we describe in detail our compositional model of GTLC that
satisfies graduality. The definitions proceed in two phases: in the
first phase we define our objects, functions, relations and squares to
model the error ordering and bisimilarity, and then we add the extra
structure of perturbations and quasi-representability to our objects
and relations to ensure that we can model the cast rules.
%
Once these definitions are complete we show how to interpret the
types, type precision, terms and term precision compositionally in our
model, and finally we conclude by extracting a big step semantics from
this model that satisfies the graduality theorem.

\subsection{Phase One: Predomains and Error Domains}

As discussed in the previous section our model of types must now come
with both an error ordering and a bisimilarity relation.
\begin{definition}
A \textbf{predomain} $A$ consists of a set $A$ along with two relations:
\begin{itemize}
    \item A partial order $\ltdyn_A$.
    \item A reflexive, symmetric ``bisimilarity'' relation $\bisim_A$.
\end{itemize}
A \textbf{morphism of predomains} $f : A \to A'$ is a function between
the underlying sets that is both \emph{monotone} (if
$x \ltdyn_A x'$, then $f(x) \ltdyn_{A'} f(x')$), and
\emph{extensional} (if $x \bisim_A x'$, then $f(x) \bisim_{A'} f(x')$).
\end{definition}

We will generally write $A$ for both a predomain and its underlying set; if we
need to emphasize the difference we will write $|A|$ for the underlying set of
$A$.
%
Given a predomain $A$, we can form the predomain $\later A$. The
underlying set is $\later |A|$ and the ordering is given by $\tilde{x}
\ltdyn_{\later A} \tilde{x'}$ iff $\later_t(\tilde{x}_t \ltdyn_A
\tilde{x'}_t)$\DIFdelbegin \DIFdel{. Likewise }\DIFdelend \DIFaddbegin \DIFadd{, and likewise }\DIFaddend for bisimilarity.

We similarly enhance our error domains with ordering and bisimilarity
relations, but we further need to impose conditions on the algebraic
structure: the error element should be the least element (so that it
models graduality), the think map should be monotone and extensional
and the bisimilarity relation should relate $\delta$ to the identity
(so that it models bisimilarity).
%
\begin{definition}
An \textbf{error domain} $B$ consists of a predomain $B$ along with the following data:
\begin{itemize}
    \item A distinguished ``error'' element $\mho_B \in B$ such that $\mho \ltdyn_B x$ for all $x$
    \item A morphism of predomains $\theta_B \colon \later B \to B$
    \item For all $x : B$ we have $\theta_B(\nxt\, x) \bisim_B x$
\end{itemize}
A \textbf{morphism of error domains} is a morphism of the underlying
predomains that preserves the error element and the $\theta$ map.
\end{definition}
For an error domain $B$, we define the predomain morphism $\delta_B := \theta_B
\circ \nxt$.

A \emph{predomain relation} $c : A \rel A'$ is a \DIFdelbegin \DIFdel{monotone relation }\DIFdelend \DIFaddbegin \DIFadd{relation between $A$ and $A'$
that is upward closed in $A'$ and downward closed in $A$, }\DIFaddend as defined in the
previous section (no condition is required for bisimilarity). The ``reflexive''
predomain relation is written $r(A)$ and is given by the ordering $\ltdyn_A$.
The relational composition of two \DIFdelbegin \DIFdel{predomains is also monotone}\DIFdelend \DIFaddbegin \DIFadd{predomain relations is also a predomain
relation}\DIFaddend , and $r(A)$ is the identity for this composition.

Error domain relations are given by predomain relations that are
additionally a congruence with respect to the algebraic structures
$\theta$ and $\mho$.
\begin{definition}
  An \textbf{error domain relation} $d : B \rel B'$ is a monotone relation $d$ on the underlying predomains satisfying
  \begin{enumerate}
     \item ($\mho$ congruence): $\mho_B \binrel{d} \mho_{B'}$.
     \item ($\theta$ congruence): For all $\tilde{x}$ in $\laterhs B$ and
     $\tilde{y} \in \laterhs B'$, if $\laterhs_t (\tilde{x}_t \binrel{d} \tilde{y}_t)$, then 
     $(\theta_B (\tilde{x}) \binrel{d} \theta_{B'} (\tilde{y}))$. 
  \end{enumerate}
\end{definition}
Again we have a reflexive error domain relation $r(B)$ given by the
ordering. The relational composition of two error domain relations in
general is not an error domain relation ($\theta$ congruence in
particular is not preserved), but we instead use an alternative ``free
composition'' \DIFdelbegin \DIFdel{$d \odot d'$ }\DIFdelend \DIFaddbegin \DIFadd{$d \odot d' : B \rel B''$ }\DIFaddend that is inductively generated from the
relational composition. \DIFdelbegin \DIFdel{Details of this construction are in the
appendix (Definition \ref{def:free-comp-ed-rel}) .
}\DIFdelend \DIFaddbegin \DIFadd{The rules are given in Figure \ref{fig:free-comp-ed-rel}.
%DIF >  the appendix (Definition \ref{def:free-comp-ed-rel}).
}\DIFaddend 

\DIFdelbegin \DIFdel{We also will use }\DIFdelend \DIFaddbegin \begin{figure}
  \begin{mathpar}
    \inferrule
    {\DIFaddFL{x \mathbin{d} y }\and \DIFaddFL{y \mathbin{d'} z}}
    {\DIFaddFL{x \mathbin{(d \relcomp d')} z}}

    \inferrule
    {\DIFaddFL{x' \ltdyn_{B_1} x }\and \DIFaddFL{x \mathbin{(d \relcomp d')} z}}
    {\DIFaddFL{x' \mathbin{(d \relcomp d')} z}}

    \inferrule
    {\DIFaddFL{x \mathbin{(d \relcomp d')} z }\and \DIFaddFL{z \ltdyn_{B_3} z'}}
    {\DIFaddFL{x \mathbin{(d \relcomp d')} z'}}

    \inferrule
    { }
    {\DIFaddFL{\mho_{B_1} \mathbin{(d \relcomp d')} z}}

    \inferrule
    {\later\DIFaddFL{_t[ \tilde{x}_t \mathbin{(d \relcomp d')} \tilde{z}_t ] }}
    {\DIFaddFL{\theta_{B_1}(}\tilde{x}\DIFaddFL{) \mathbin{(d \relcomp d')} \theta_{B_3}(}\tilde{z}\DIFaddFL{) }}
  \end{mathpar}
  \caption{\DIFaddFL{The free composition of error domain relations.}}
  \label{fig:free-comp-ed-rel}
\end{figure}

\DIFadd{As in the previous section, we also have }\DIFaddend the notion of a square $f
\ltsq{c_i}{c_o} g$ between predomain morphisms\DIFdelbegin \DIFdel{as well as }\DIFdelend \DIFaddbegin \DIFadd{, which is defined to mean that
for all $x \binrel{c_i} y$, we have $f(x) \binrel{c_o} g(y)$. Likewise, we have
}\DIFaddend squares $\phi \ltsq{d_i}{d_o} \psi$ between error domain morphisms, which are
defined similarly, with no extra interaction with the algebraic structure.
%DIF > 
\DIFaddbegin \DIFadd{The free composition of error domain relations has the universal property that
$\phi \ltsq{d \relcomp d'}{d''} \phi'$ if and only if for any $x \binrel{d} y$
and $y \binrel{d'} z$, we have $\phi(x) \binrel{d''} \phi'(z)$.
}\DIFaddend 

In our denotational model of GTLC, we map every type $A$ to a
predomain $\sem{A}$. These constructions are given by the same
constructions as in Section~\ref{sec:concrete-term-model}, but now enriched to
act also on orderings and bisimilarity. For products this is the
obvious lifting, and for functions we use the CBPV decomposition $U(A
\to \li A')$ where the only one that is not clear is $\li A'$.
\begin{definition}
  The free error domain $\li A$ has as the underlying set the free
  simple error domain with ordering given by the lock-step error
  ordering and bisimilarity by weak bisimilarity both generated by the
  ordering/bisimilarity on $A$. This is free in the same sense as the
  free simple error domain that it is left adjoint to $U$.
\end{definition}

%% \subsubsection{Functors}\label{sec:free-error-domain}

%% The functors $U$, $\li$, $\times$, and $\arr$ introduced in the previous section
%% act on predomains and error domains and also have actions on their morphisms,
%% relations, and squares. The functor $U$ simply takes the underlying predomain
%% relation of an error domain relation. Given a relation $c_1 : A_1 \rel A_1'$ and
%% $c_2 : A_2 \rel A_2'$, the relation $c_1 \times c_2 : (A_1 \times A_2) \rel
%% (A_1' \times A_2')$ is defined by $((x, y), (x', y')) \in (c_1 \times c_2)$ if
%% and only if $(x, x') \in c_1$ and $(y, y') \in c_2$. Given a relation $c : A
%% \rel A'$ and an error domain relation $d : B \rel B'$, the relation $c \arr d :
%% (A \arr B) \rel (A' \arr B')$ is the error domain relation defined by $(f, g)
%% \in (c \arr d)$ if and only if for all $x : A$ and $x' : A'$ with $(x , x') \in
%% c$ we have $(f(x), g(x')) \in Ud$. The fact that this relation respects error %\max{preserves error?}
%% and preserves $\theta$ follows from the corresponding properties for $d$.

%% % Free error domain
%% The functor $\li$ takes a predomain $A$ to the free error domain on
%% $A$. We first define the predomain $U\li A$. The underlying set of
%% $U\li A$ is defined to be $\li |A|$ (the guarded lift monad applied to
%% $|A|$). The ordering relation is the lock-step error ordering
%% introduced in the previous section , and the bisimilarity relation is
%% the weak bisimilarity relation on $\li |A|$ also defined in the
%% previous section. With these relations, the constructors $\eta$ and
%% $\theta$ of the guarded lift monad are in fact morphisms of
%% predomains, i.e., they are monotone and preserve bisimilarity.
%% %
%% We observe that this predomain structure extends to an error domain structure by
%% noting that the required error element is given by the constructor $\mho$ and
%% the required $\theta$ map is the constructor $\theta$. Lastly, it can be shown
%% that the delay morphism $\delta = \theta \circ \nxt : UB \to UB$ satisfies $x \bisim
%% \delta\, x$ for all $x$ as is required by the definition of error domain.

% The monadic bind operation $-^\dagger$ takes a predomain morphism $f : A \to UB$
% to an error domain morphism $\li A \multimap B$; with this we define the action of $\li$ on
% morphisms: Given $f : A \to A'$ we define $\li f = (\eta \circ f)^\dagger : \li A \to \li A'$.
%\max{what is the point of this sentence?}

Additionally, all of these core constructions $\times, \to, U, \li$
act on heterogeneous relations as well in an analogous way.

The most interesting type to interpret is the dynamic
type. The beauty of guarded domain theory is that once we have define
$\laterhs$ for predomains, we can solve the same domain equation as in
Section~\ref{sec:concrete-term-model} changed only in that we interpret it as a
domain equation on predomains rather than sets:
%
\[ D \cong \mathbb{N}\, + (D \times D)\, + \laterhs U(D \arr \li D) \]
The complete definition of the ordering $\ltdyn_D$ is provided in the appendix
for completeness. \DIFaddbegin \DIFadd{We refer to the three injections as $e_\nat$, $e_\times$,
and $e_\to$, respectively.
}\DIFaddend 

To interpret \DIFdelbegin \DIFdel{term }\DIFdelend \DIFaddbegin \DIFadd{type }\DIFaddend precision, we define three relations involving $D$
as follows. We define $\inat : \mathbb{N} \rel D$ by $(n, d) \in
\inat$ iff $e_\nat(n) \ltdyn_D d$. We similarly define $\itimes : D
\times D \rel D$ by $((d_1, d_2), d) \in \itimes$ iff $e_\times(d_1,
d_2) \ltdyn_D d$, and we define $\text{inj}_\to : U(D \arr \li D) \rel
D$ by $(f, d) \in \iarr$ iff \DIFdelbegin \DIFdel{$e_\to(f) \ltdyn_D d$}\DIFdelend \DIFaddbegin \DIFadd{$(e_\to(f) \circ \nxt) \ltdyn_D d$}\DIFaddend .

\subsection{Phase Two: Perturbations and Quasi-Representable Relations}

We now introduce the additional definitions needed to complete the construction
of our model. Recall from Section \ref{sec:lock-step-and-weak-bisim} that in order
to establish the \DIFdelbegin \DIFdel{DnL }\DIFdelend \DIFaddbegin \DIFadd{$\dnl$ }\DIFaddend rule for the downcast of $\iarr$, we needed to adjust the
rule by inserting a ``delay'' on the right-hand side. A similar adjustment is
needed for the \DIFdelbegin \DIFdel{DnR }\DIFdelend \DIFaddbegin \DIFadd{$\dnr$ }\DIFaddend rule for $\iarr$. Moreover, the need to insert delays impacts
the semantics of the cast rules for \emph{all} relations, because of the
functorial nature of casts. That is, the upcast at a function type $c_i \ra c_o$
involves a downcast in the domain and an upcast in the codomain. This has two
consequences: first, the squares corresponding to the rules \DIFdelbegin \DIFdel{UpL and UpR }\DIFdelend \DIFaddbegin \DIFadd{$\upl$ and $\upr$ }\DIFaddend may also
require the insertion of a delay. Second, we need to be able to insert
``higher-order'' delays in a way that follows the structures of the casts.

To accomplish this, we equip the predomains $A$ in our semantics with
a monoid $M_A$ of \emph{syntactic perturbations}, as well as a means
of \emph{interpreting} these syntactic perturbations as
\emph{semantic} perturbations on $A$, i.e., as endomorphisms that are
bisimilar to the identity. We call such an object a \emph{value
object}\DIFaddbegin \DIFadd{, }\DIFaddend and this will serve as the final definition of a denotation of
a type in our semantics. \DIFdelbegin \DIFdel{We }\DIFdelend \DIFaddbegin \DIFadd{Similarly, we }\DIFaddend define computation objects to \DIFdelbegin \DIFdel{similarly }\DIFdelend be
error domains equipped with a notion of syntactic perturbation:
\begin{definition}[value and computation objects + morphisms]
  A \textbf{value object} consists of a predomain $A$, a monoid $M_A$ and a
  homomorphism of monoids \DIFdelbegin \DIFdel{$M_A \to \morbisimid{A} := \{ f : A \to A \mid f
  \bisim \id_A \}$}\DIFdelend \DIFaddbegin \DIFadd{$i_A \colon M_A \to \morbisimid{A} := \{ f : A \to A \mid f
  \bisim \id_A \}$}\DIFaddend .

  Likewise, a \textbf{computation object} consists of an error domain $B$, a
  monoid $M_B$ and a homomorphism of monoids \DIFdelbegin \DIFdel{$M_B \to \morbisimid{B}$}\DIFdelend \DIFaddbegin \DIFadd{$i_B \colon M_B \to \morbisimid{B}$}\DIFaddend .

  We define a \textbf{value morphism} to be \DIFaddbegin \DIFadd{simply }\DIFaddend a morphism of the underlying
  predomains, and a \textbf{computation morphism} to be a morphism of the
  underlying error domains.
\end{definition}

\DIFaddbegin \begin{remark}
  \DIFadd{Syntactic perturbations are equipped with a monoid structure in
  order to support identity and composition relations on value and
  computation types. That is, the identity relation requires an
  identity perturbation and composition of relations requires a
  composition operation on perturbations.
  %DIF > % The reason why the set of syntactic perturbations needs to carry a monoid
  %DIF > % structure and why the interpretation needs to be a homomorphism of monoids is
  %DIF > % related to the construction of the identity relations and composition of
  %DIF > % relations on value/computation types. As part of this construction, we will
  %DIF > % need to produce a syntactic perturbation corresponding to the identity
  %DIF > % predomain/error domain relation, as well as corresponding to the composition
  %DIF > % of relations. The monoid structure gives us what we need to produce these
  %DIF > % perturbations.
  %DIF >  we will need to produce a syntactic perturbation corresponding to the identity
  %DIF >  relations $r(A)$, as well as one for the composition of relations. The monoid
  %DIF >  structure gives us what we need to produce these elements. Likewise, the
  %DIF >  reason that the interpretations as semantic perturbations must be
  %DIF >  \emph{homomorphisms} of monoids is that we need a syntactic perturbation that
  %DIF >  is interpreted as the identity semantic perturbation, and likewise given two
  %DIF >  syntactic perturbations $\delta_1$ and $\delta_2$ we need a syntactic
  %DIF >  pertubration $\delta_{1,2}$ that is interpreted as the composition of the
  %DIF >  semantic perturbations $i_A(\delta_1)$ and $i_A(\delta_2)$.
  %DIF >  These properties are used to define the identity and composition of value and computation relations
}\end{remark}

%DIF >  It doesn't work: not all types have such an action, e.g., the function type, because of the definition of the perturbations
%DIF >  The fact that we cannot define the morphisms this way is unfortunate, since doing so would simplify
%DIF >  some aspects of the model construction.
\begin{remark}
  \DIFadd{It might seem strange that we equip objects with a monoid of
  perturbations but we do not require morphisms to respect this
  structure in any way. A natural attempt would be defining a morphism
  between value objects $(A, M_A, i_A)$ and $(A', M_{A'}, i_{A'})$ to
  be a predomain morphism $f \colon A \to A'$ along with a monoid
  homomorphism $h \colon M_A \to M_{A'}$ such that $f$ and $h$ satisfy
  a compatibility condition with respect to the interpretations $i_A$
  and $i_{A'}$. However, this condition simply fails to hold for all
  terms, in particular the property is preserved by $\lambda$
  abstraction.
  %DIF > % with this definition we cannot interpret the
  %DIF > % function type, since the morphisms defined in this way do not
  %DIF > % themselves form a value type.
  %DIF >  This is due to the way that the syntactic perturbations are defined for the function type.
  %DIF > 
  %DIF > % The fact that we cannot define the morphisms in this way is unfortunate, since
  %DIF > % doing so would simplify some aspects of the model construction.
  }

  %DIF >  However, with this definition of morphism, the category of value
  %DIF >  types is no longer Cartesian closed and is not a model of call-by-push-value.
  %DIF >  This is unfortunate, as it would simplify some aspects of the model construction.
\end{remark} 

\DIFaddend % every error domain $B$ with a monoid $M_B$ and interpretation homomorphism $i_B
% : M_B \to \morbisimid{B}$. This grouping of data deserves its own definition: We
% call the triple $(A, M_A, iA)$ a \textbf{value object} as it will be the
% denotation of value types in our final model. Likewise, the triple $(B, M_B,
% i_B)$ will be called a \textbf{computation object}.\max{make this a proper \textit{Definition} and define morphisms as well}

Having defined syntactic perturbations and value and computation objects, we can
now make formal the final aspect of the model: the notion of
\DIFdelbegin \DIFdel{quasi-represntability }\DIFdelend \DIFaddbegin \DIFadd{quasi-representability }\DIFaddend of relations.
%
\begin{definition}[quasi-left-representable relations]\label{def:quasi-left-representable}
Let $(A, M_A, i_A)$ and $(A', M_{A'}, i_{A'})$ be value objects, and let $c : A
\rel A'$ be a predomain relation. We say that $c$ is
\emph{quasi-left-representable} by a predomain morphism $e : A \to A'$ if there
are perturbations $\delle_c \in M_A$ and $\delre_c \in M_{A'}$ such that the
following two squares exist: (1) $\upl$: $e \ltsq{c}{\ltdyn_{A'}} i_{A'}(\delre_c)$, and
(2) $\upr$: $i_A(\delle_c) \ltsq{\ltdyn_{A}}{c} e$.
\end{definition}
%
Observe that this weakens the previous notion of representability, since under
that definition the perturbations were required to be the identity.
%
We make the analogous definition of quasi-left-representability for error domain
relations $d$.

Likewise, we define quasi-right-representability as follows:
%
\DIFdelbegin %DIFDELCMD < \begin{definition}[quasi-right-represntable relations]%%%
\DIFdelend \DIFaddbegin \begin{definition}[quasi-right-representable relations]\DIFaddend \label{def:quasi-right-representable}
Let $(B, M_B, i_B)$ and $(B', M_{B'}, i_{B'})$ be computation objects, and let
$d : B \rel B'$ be an error domain relation. We say that $d$ is
\emph{quasi-right-representable} by an error domain morphism $p : B' \multimap B$ if
there are perturbations $\dellp_d \in M_B$ and $\delrp_d \in M_{B'}$ such that
the following two squares exist: 
(1) $\dnl$: $p \ltsq{\ltdyn_{B'}}{d} i_{B'}(\delrp_d)$ and
(2) $\dnr$: $i_B(\dellp_d) \ltsq{d}{\ltdyn_{B}} p$.
Quasi-right-representability for predomain relations is defined similarly.
\end{definition}

For ordinary representability, the relation is \emph{\DIFdelbegin \DIFdel{uniqueley
}\DIFdelend \DIFaddbegin \DIFadd{uniquely
}\DIFaddend determined} by any morphism that left- or right-represents it: if two
relations are left-represented by the same morphism, then they are
equal. But for quasi-representability, this is no longer the case due
to the presence of perturbations. Instead we have the following weaker property:
\begin{definition}[quasi-equivalence]\label{def:quasi-equivalent}
  Let $A$ and $A'$ be value objects, and let $c, c' : A \rel A'$ be relations
  between the underlying predomains. We say that $c$ and $c'$ are
  \textbf{quasi-equivalent}, written $c \qordeq c'$, if there exist
  syntactic perturbations $\delta^l_1, \delta^l_2 \in M_A$ and $\delta^r_1,
  \delta^r_2 \in M_{A'}$ such that there is a square $i_A(\delta^l_1)
  \ltsq{c}{c'} i_{A'}(\delta^r_1)$ and a square $i_A(\delta^l_2) \ltsq{c'}{c}
  i_{A'}(\delta^r_2)$. Given computation objects $B$ and $B$, we make the
  analogous definition for relations $d, d' : B \rel B'$ between the underlying
  error domains.
\end{definition}

% \max{idea: move this after the quasi-representability part and say that this structure is what's needed to make them compose}
Given two quasi-representable relations $c$ and $c'$, it is not the case in
general that $cc'$ is quasi-representable. We need an additional condition that
specifies how relations interact with perturbations. In particular, we need to
be able to ``push'' and ``pull'' perturbations along relations $c$ and $d$.
%
The intuition for this requirement comes from the construction of the square
corresponding to the original \DIFdelbegin \DIFdel{UpL }\DIFdelend \DIFaddbegin \DIFadd{$\upl$ }\DIFaddend rule from the square for the simplified
version as was shown in Section \ref{sec:towards-relational-model}. Recall that
we horizontally composed the simplified \DIFdelbegin \DIFdel{UpL }\DIFdelend \DIFaddbegin \DIFadd{$\upl$ }\DIFaddend square for $c$ with the identity
square for $c'$ on the right. Now that the square for \DIFdelbegin \DIFdel{UpL }\DIFdelend \DIFaddbegin \DIFadd{$\upl$ }\DIFaddend involves a
perturbation on the right instead of the identity, this construction will not
work unless we can ``push'' the perturbation on $A_2$ to one on $A_3$. That is,
given a relation $c : A \rel A'$ and a syntactic perturbation $m_A \in M_A$, we
need to be able to turn it into a syntactic perturbation $m_{A'} \in M_{A'}$
such that the resulting semantic perturbations obtained by applying the
respective homomorphisms $i_A$ and $i_{A'}$ form a square. Dually, we will need
to ``pull'' a perturbation from right to left. This notion is made formal by the
following definition:
%
\begin{definition}
    Let $(A, M_A, i_A)$ and $(A', M_{A'}, i_{A'})$ be value objects,
    and let $c : A \rel A'$ be a predomain relation. A
    \textbf{push-pull structure} for $c$ (with respect to $M_A$ and
    $M_{A'}$) consists of monoid homomorphisms $\push : M_A \to M_A'$
    and $\pull : M_A' \to M_A$ such that for all $m_A \in M_A$ there
    is a square $i_A(m_A) \ltsq{c}{c} i_{A'}(\push\, m_A)$ and for all
    $m_{A'} \in M_{A'}$ there is a square $i_A(\pull\, m_{A'})
    \ltsq{c}{c} i_{A'}(m_{A'})$. A push-pull structure for error
    domain relations is defined similarly.
\end{definition}

We can now give the final definition of value and computation relations for our
model, which serve as the denotation of our type precision derivations:
%
\begin{definition}[value relations]
A \textbf{value relation} between value objects $(A, M_A, i_A)$ and $(A',
M_{A'}, i_{A'})$ consists of a predomain relation $c : A \rel A'$ that has a
push-pull structure, is quasi-left-representable by a morphism $e_c : A \to A'$,
and is such that $\li c$ is quasi-right-representable by an error domain morphism
$p_c : \li A' \to \li A$.

A \textbf{computation relation} between computation objects $(B, M_B, i_B)$ and
$(B', M_{B'}, i_{B'})$ consists of an error domain relation $d$ that has a
push-pull structure, is quasi-right-representable by a morphism $p_d : B' \to
B$, and is such that $Ud$ is quasi-left-representable by a morphism $e_d : UB
\to UB'$
\end{definition}

It is essential in the definition of a value relation that we have not
just that $c$ is quasi-left-representable but also that $\li c$ is
quasi-right-representable so that we can define the action of the
functor $\li$ on relations, taking a value relation to a computation
relation. Likewise, we require $Ud$ to be quasi-left-representable in
the definition of computation relations so that we can define the
action of $U$ on relations.

% Because we have augmented the definitions of object and relation for our model,
% we need to specify the actions of the functors $F$, $U$, $\times$, and $\arr$ on
% these objects and relations. 
% For instance, the functor $F$ now acts on \emph{value objects} $(A, M_A, i_A)$,
% which means we need to specify the monoid $M_{FA}$ and interpretation $i_{FA}$
% corresponding to the error domain $FA$.

To model the term precision in a way that is oblivious to the stepping
behavior of terms, we need to weaken the notion of square to allow for the
morphisms on either side to be ``out of sync'' rather than in lock step.
\DIFaddbegin \DIFadd{In the below, $f' \bisimsq{A_i}{A_o} f$ is the natural extension of
bisimilarity to morphisms: given bisimilar inputs, the functions have
bisimilar outputs.
}\DIFaddend \begin{definition}
  Let $c_i : A_i \rel A_i'$, $c_o : A_o \rel A_o'$ $f : A_i \to A_o$
  and $g : A_i' \to A_o'$.  An extensional square $f
  \ltsqbisim{c_i}{c_o} g$ holds if there \emph{exist} $f'
  \bisimsq{A_i}{A_o} f$ and $g' \bisimsq{A_i'}{A_o'} g$ such that $f' \ltsq{c_i}{c_o} g'$.
\end{definition}
\DIFdelbegin \DIFdel{Here $f' \bisimsq{A_i}{A_o} f$ is the natural extension of
bisimilarity to morphisms: given bisimilar inputs the functions have
bisimilar outputs.
}\DIFdelend 

\subsection{Constructing the \DIFdelbegin \DIFdel{denotational semantics}\DIFdelend \DIFaddbegin \DIFadd{Denotational Semantics}\DIFaddend }\DIFaddbegin \label{sec:relational-model-construction}
\DIFaddend 

With these concepts now defined, we can give an overview of our
denotational model:
\DIFdelbegin %DIFDELCMD < \begin{theorem}
%DIFDELCMD <   %%%
\DIFdelend \begin{itemize}
  \item For every type $A$ we define a value type $\sem{A}$, and define $\sem{x_1:A_1,\ldots} = \sem{A_1}\times\cdots$
  \item For every term $\Gamma \vdash M : A$, we define a value morphism $\sem{M} : \sem{\Gamma} \to U\li\sem{A}$.
  \item If $M = M'$ are equal according to our equational theory, then $\sem{M} = \sem{M'}$
  \item For every type precision derivation $c : A \ltdyn A'$ we
    define a value relation $\sem{c} : \sem{A} \rel \sem{A'}$.
  \item For every type precision equivalence $c \equiv c'$ we have a
    \DIFdelbegin \DIFdel{quasi-order-equivalence }\DIFdelend \DIFaddbegin \DIFadd{quasi-equivalence }\DIFaddend of relations $\sem{c} \bisim \sem{c'}$.
  \item For every term precision $\Delta \vdash M \ltdyn M' : c$ we
    have an \emph{extensional} square
    $\sem{M}\ltsqbisim{\sem{\Delta}}{U\li\sem{c}} \sem{M'}$
\end{itemize}
\DIFdelbegin %DIFDELCMD < \end{theorem}
%DIFDELCMD < %%%
\DIFdelend 

\DIFdelbegin \DIFdel{For the }\DIFdelend \DIFaddbegin \DIFadd{The }\DIFaddend term semantics $\sem{M}$ \DIFdelbegin \DIFdel{, nothing really differs from }\DIFdelend \DIFaddbegin \DIFadd{is nearly identical to }\DIFaddend the term semantics in
Section~\ref{sec:concrete-term-model}. \DIFdelbegin \DIFdel{The only }\DIFdelend \DIFaddbegin \DIFadd{One }\DIFaddend difference is that the \DIFdelbegin \DIFdel{definitions of $\sem{\upc c}$ and $\sem{\dnc c}$ can be
defined as the morphisms that left- and right-quasi-represent $\sem
c$, but upon inspection this definition agrees with the
one given
before. Similarly, the }\DIFdelend \DIFaddbegin \DIFadd{interpretation
of the function type $A_i \ra A_o$ in the term semantics of Section
\ref{sec:concrete-term-model} involves }\emph{\DIFadd{arbitrary functions}}\DIFadd{, while the
interpretation of the function type in the term semantics defined here involves
only those functions that are morphisms of predomains. Additionally, the
semantics of casts is slightly different. In the semantics of Section
\ref{sec:concrete-term-model}, the projection corresponding to the type
precision $cc'$ is simply the composition of the projections. In the new
semantics, the corresponding projection involves a perturbation. This is a
consequence of the construction of the composition of value/computation
relations (the relevant projection is defined in the proof of Lemma
\ref{lem:representation-comp-F-U} in the Appendix). However, the two projections
are extensionally equivalent since they are the same up to a perturbation.
}

 %DIF >  nothing really differs from the term semantics in Section~\ref{sec:concrete-term-model}. The only difference is
%DIF >  One difference is that the definitions of $\sem{\upc c}$ and $\sem{\dnc c}$ can be
%DIF >  defined as the morphisms that left- and right-quasi-represent $\sem
%DIF >  c$, but upon inspection this definition agrees with the one given
%DIF >  before. 

\DIFadd{The }\DIFaddend validity of the equational theory follows easily: $\beta\eta$ holds in our
model. \DIFdelbegin \DIFdel{We next }\DIFdelend \DIFaddbegin \DIFadd{In the remainder of this section, we }\DIFaddend elaborate on each of the remaining
components \DIFdelbegin \DIFdel{.
}\DIFdelend \DIFaddbegin \DIFadd{of the model.
%DIF >  Interpretations of functions are different: one is arbitrary fucntions, the other is monotone functions that preserve bisimilarity
%DIF >  Interpretation of the projection is different since we add a delay
}\DIFaddend 

\subsubsection{Interpreting types}\label{sec:perturbation-constructions}
% \max{this section needs prose}

% \eric{May be able to move some of this to the appendix}
For our value type semantics, we already have an interpretation as predomains\DIFdelbegin \DIFdel{,
}\DIFdelend \DIFaddbegin \DIFadd{;
}\DIFaddend we describe how we enhance each of the constructions we use \DIFaddbegin \DIFadd{(}\DIFaddend $U,\li,\to,\times,D$\DIFdelbegin \DIFdel{with a notion }\DIFdelend \DIFaddbegin \DIFadd{)
with a monoid }\DIFaddend of syntactic perturbation \DIFaddbegin \DIFadd{and a means of interpreting them as semantic
perturbations}\DIFaddend .

First consider the free error domain $\li-$. Given a value object $(A, M_A,
i_A)$ we define the syntactic perturbations $M_{\li A}$ to be $\mathbb{N} \oplus
M_A$, where $\oplus$ denotes the free product of monoids (the coproduct in the
category of monoids). The intuition behind needing to include $\mathbb{N}$ comes
from the example of the downcast for $\iarr$. Recall that we adjusted the
squares corresponding to the \DIFdelbegin \DIFdel{DnL and DnR }\DIFdelend \DIFaddbegin \DIFadd{$\dnl$ and $\dnr$ }\DIFaddend rules by adding $(\delta \circ
\eta)^\dagger$ on the side opposite the downcast. Thus, we need a syntactic
perturbation in $M_{\li A}$ that will be interpreted as $(\delta \circ
\eta)^\dagger$. We accomplish this by taking the coproduct with $\mathbb{N}$.
\DIFaddbegin \DIFadd{Since $\mathbb{N}$ is the free monoid on one generator, defining a homomorphism
from $\mathbb{N}$ to $\morbisimid{\li A}$ is equivalent to a choice of an
element in $\morbisimid{\li A}$.
}\DIFaddend %
The interpretation of the perturbations as endomorphisms is defined using the
universal property of the coproduct of monoids\DIFdelbegin \DIFdel{(see appendix for details). }\DIFdelend \DIFaddbegin \DIFadd{. Specifically, to define the
homomorphism $i_{\li A}$, by the universal property of the coproduct of monoids
it suffices to define two homomorphisms, one from $\mathbb{N} \to
\morbisimid{\li A}$ and one from $M_A \to \morbisimid{\li A}$. The first
homomorphism sends the generator $1 \in \mathbb{N}$ to the error domain
endomorpism $(\delta \circ \eta)^\dagger$, and the second sends $m_A \in M_A$ to
$\li(i_A(m_A))$. We observe in both cases that the resulting endomorphisms are
bisimilar to the identity, since in the first case, $-^\dagger$ preserves
bisimilarity, and in the second case, the action of $\li$ on morphisms preserves
bisimilarity.
}\DIFaddend 

Now we \DIFdelbegin \DIFdel{discuss }\DIFdelend \DIFaddbegin \DIFadd{define }\DIFaddend the action of $U$ on objects. Given a computation object $(B, M_B,
i_B)$, we define $M_{UB}$ to be $\mathbb{N} \oplus M_B$. The reason for
requiring $\mathbb{N}$ is related to the Kleisli \DIFaddbegin \DIFadd{action of the }\DIFaddend arrow functor: to
establish quasi-representability of $U(c \arr d)$ given quasi-representability
of $c$ and $d$, we will need to turn a perturbation on $\li A$ into a
perturbation on $U(A \arr B)$. Since the perturbations for $\li A$ involve
$\mathbb{N}$, so must the perturbations for $UB$. The interpretation $i_{UB}$ of
the perturbations on $UB$ works in \DIFdelbegin \DIFdel{the same manner as }\DIFdelend \DIFaddbegin \DIFadd{an analogous manner to }\DIFaddend that of $\li A$\DIFaddbegin \DIFadd{: we send
the generator $1 \in \mathbb{N}$ to the delay morphism $\delta_B \colon UB \to
UB$, and we send a perturbation $m_B \in M_B$ to $U(i_B(m_B))$}\DIFaddend .

The action of $\times$ on objects is as follows. Given $(A_1, M_{A_1}, i_{A_1})$
and $(A_2, M_{A_2}, i_{A_2})$, we define the monoid $M_{A_1 \times A_2} =
M_{A_1} \oplus M_{A_2}$. The interpretation \DIFaddbegin \DIFadd{$i_{A_1 \times A_2}$ }\DIFaddend is defined
using the universal property of the coproduct \DIFdelbegin \DIFdel{(see appendix for the details)}\DIFdelend \DIFaddbegin \DIFadd{of monoids. That is, we define
homomorphisms $M_{A_1} \to \morbisimid{A_1 \times A_2}$ and $M_{A_2} \to
\morbisimid{A_1 \times A_2}$. The former takes $m_1 \in M_{A_1}$ to the morphism
$(i_{A_1}(m_1), \id)$, while the latter takes $m_2 \in M_{A_2}$ to the morphism
$(\id, i_{A_2}(m_2))$}\DIFaddend .
% Could state this using the functorial action of \times

The action of $\arr$ on objects is as follows. Given $(A, M_A, iA)$ and $(B,
M_B, i_B)$ we define the monoid $M_{A \arr B} = M_A^{op} \oplus M_B$\DIFdelbegin \DIFdel{and the interpretation
is defined analogously to that of $\times$}\DIFdelend \DIFaddbegin \DIFadd{. (Here,
$M^{op}$ denotes the monoid with the same elements but with the order of
multiplication reversed, i.e., $x \cdot_{M^{op}} y := y \cdot_M x$.) To define
the interpretation $i_{A \arr B}$, it suffices to define a homomorphism
$M_A^{op} \to \morbisimid{A \arr B}$ and $M_B \to \morbisimid{A \arr B}$. The
former takes an element $m_A$ and a predomain morphism $f : A \to UB$ and
returns the composition $f \circ i_A(m_A)$, while the latter is defined
similarly by post-composition with $i_B$}\DIFaddend .
% Could state this using the functorial action of \arr

Finally, the most interesting type is of course the dynamic type. 
\DIFdelbegin \DIFdel{Now we define the perturbations for $D$. Recall that to each predomain $A$ we
associate a monoid $M_A$ of syntactic perturbations and a homomorphism into the
monoid of endomorphisms bisimilar to the identity, and likewise for error
domains. }\DIFdelend It may seem as though we need to define the perturbations for $D$ by
guarded recursion, but in fact we define them via least-fixpoint
in the category of monoids:
%
\( M_D \cong (M_{D \times D}) \oplus M_{U(D \to \li D)}. \)
%
\DIFaddbegin \begin{remark}
  \DIFadd{Here it is crucial to define the perturbations of the dynamic type
  }\emph{\DIFadd{inductively}} \DIFadd{rather than by guarded recursion. If we were to
  define the perturbations by guarded recursion, then it would be
  impossible to define the pull homomorphism for the $\iarr$ relation.
}\end{remark}
%DIF > 
\DIFaddend We now explain how to interpret these perturbations as endomorphisms.
\DIFdelbegin %DIFDELCMD < 

%DIFDELCMD < %%%
\DIFdelend We define $i_D : M_D \to \{ f : D \to D \mid f \bisim \id \}$ by induction as follows:
%
\begin{align*}
  i_D(m_\times) &= \lambda d. \text{case $d$ of }
    \{ \itimes (d_1, d_2) \To \itimes (i_{D \times D}(m_\times)(d_1, d_2)) \alt x \To x \} \\
  i_D(m_\to) &= \lambda d. \text{case $d$ of }
    \{ \iarr (\tilde{f}) \To \iarr (\lambda t. i_{U(D \to \li D)}(m_\to)(\tilde{f}_t)) \alt x \To x \}
\end{align*}
%
In the case of a perturbation on $D \times D$, we use the interpretation for the
perturbations of $D \times D$ as defined earlier in this section, which in turn
will use $i_D$ inductively. Similarly, for a perturbation on $U(D \to \li D)$ we
use the interpretation for perturbations on $U(D \to \li D)$.
One can verify that this defines a homomorphism from $M_D \to \{ f : D \to D : f
\bisim \id \}$.

% \begin{align*}
%  i_D(m_{\text{times}}, m_{\text{fun}}) &= \lambda d.\text{case $d$ of}  \\
%  &\alt \tnat(m) \mapsto \tnat(m) \\
%     &\alt \ttimes(d_1, d_2) \mapsto {\ttimes(i_{D \times D}(p_\text{times})(d_1, d_2))} \\
%     &\alt \tfun(\tilde{f}) \mapsto {\tfun(\lambda t. i_{U(D \to \li D)}(\tilde{f}_t))}
% \end{align*}

\subsubsection{Interpreting type precision}

Next we need to confirm that our interpretation of type precision
derivations as predomain relations indeed extends to value relations,
i.e., that the relations are quasi-representable and that they satisfy
the push-pull property for perturbations. We include most of these technical
verifications in the appendix (\DIFdelbegin \DIFdel{Definition
\ref{def:value-computation-rel-comp} ,
}\DIFdelend \DIFaddbegin \DIFadd{Definitions
\ref{def:value-computation-rel-comp} and
}\DIFaddend \ref{def:functorial-actions-on-relations}).

We go into more detail on the three relations involving the dynamic
type $\inat$, $\itimes$, and $\iarr$. As an illustrative case, we
establish the push-pull property for the relation $\iarr$. We define
$\pull_{\iarr} : M_D \to M_{U(D \arr \li D)}$ by giving a homomorphism
from $M_{D \times D} \to M_{U(D \arr \li D)}$ and from $M_{U(D \to \li
  D)} \to M_{U(D \arr \li D)}$. The former is the trivial homomorphism
sending everything to the identity element, while the latter is the
identity homomorphism.
%
Conversely, we define $\push_{\iarr} : M_{U(D \arr \li D)} \to M_D$ as the
inclusion into the coproduct.
%
Showing that the relevant squares exist involving push and pull is
straightforward.

It is easy to see that the relations $\inat$, $\itimes$, and $\iarr$
are quasi-left-representable, where all perturbations are taken to be the
identity elements of the respective monoids.
%
For quasi-right-representability, the only nontrivial case is $\li(\iarr)$.
Recall the \DIFdelbegin \DIFdel{defintion }\DIFdelend \DIFaddbegin \DIFadd{definition }\DIFaddend of the downcast for $\iarr$ given in Section
\ref{sec:term-interpretation}. We know that in order for the \DIFdelbegin \DIFdel{DnR and DnL }\DIFdelend \DIFaddbegin \DIFadd{$\dnr$ and $\dnl$ }\DIFaddend squares
for $\li \iarr$ to exist, we must insert a delay on the side opposite the
downcast. In terms of syntactic perturbations, this means that we take $\dellp$
and $\delrp$ in the definition of quasi-right-representability to both be
$\inl\, 1$, where $\inl$ is the left injection into the coproduct of monoids. We
recall that the definition of syntactic perturbations for $\li$ involves a
coproduct with $\mathbb{N}$, and that the interpretation homomorphism maps the
generator $1$ to the endomorphism of error domains $(\delta \circ
\eta)^\dagger$.
%
With this choice for the perturbations, it is straightforward to show that the
\DIFdelbegin \DIFdel{DnL and DnR }\DIFdelend \DIFaddbegin \DIFadd{$\dnl$ and $\dnr$ }\DIFaddend squares exist using the definition of the downcast and the
definition of the relation $\iarr$.

\DIFdelbegin \DIFdel{The proofs that the precision equivalence imply }\DIFdelend \DIFaddbegin \DIFadd{Lastly, the interpretation of the type precision equivalence rules $c \equiv c'$
as }\DIFaddend quasi-equivalence of \DIFdelbegin \DIFdel{relations are in the appendix (Lemma \ref{lem:quasi-order-equiv-functors} )}\DIFdelend \DIFaddbegin \DIFadd{the denoted relations $\sem{c} \bisim \sem{c'}$ is
justified by Lemma \ref{lem:quasi-order-equiv-functors} in the Appendix}\DIFaddend .

\subsubsection{Interpreting Term Precision: Extensional Squares}

Finally, we establish the core result of the semantics: interpreting
term precision as extensional squares. Fortunately\DIFaddbegin \DIFadd{, }\DIFaddend the hard work is
now complete\DIFaddbegin \DIFadd{, }\DIFaddend and the term precision rules follow from our
compositional constructions.
%
\DIFdelbegin \DIFdel{The }\DIFdelend \DIFaddbegin \DIFadd{In particular, the }\DIFaddend congruence rules of term precision are modeled easily: congruence
squares \DIFdelbegin \DIFdel{hold }\DIFdelend \DIFaddbegin \DIFadd{exist }\DIFaddend both for bisimilarity \DIFdelbegin \DIFdel{as well as }\DIFdelend \DIFaddbegin \DIFadd{and for the }\DIFaddend error ordering, and so
\DIFdelbegin \DIFdel{by combination }\DIFdelend we also have these for extensional squares.
%
The \DIFdelbegin \DIFdel{type precision equivalence rule }\DIFdelend \DIFaddbegin \textsc{\DIFadd{ErrBot}} \DIFadd{rule holds because $\mho$ is the least element in the semantics of the . 
%DIF > 
Next, the }\textsc{\DIFadd{EquivTyPrec}} \DIFadd{rule }\DIFaddend follows because $c \equiv c'$
implies quasi-equivalence \DIFdelbegin \DIFdel{, and that quasi-equivalence }\DIFdelend \DIFaddbegin \DIFadd{of the denoted relations, and this in turn }\DIFaddend implies that
$\id \ltsqbisim{\sem{c}}{\sem{c'}} \id$.
\DIFdelbegin %DIFDELCMD < 

%DIFDELCMD < %%%
\DIFdel{Most crucially, we }\DIFdelend %DIF > 
\DIFaddbegin \DIFadd{Lastly, it remains to }\DIFaddend verify the existence of the squares for the cast rules. For \DIFdelbegin \DIFdel{UpL}\DIFdelend \DIFaddbegin \DIFadd{$\upl$}\DIFaddend , the
relevant extensional square is obtained as follows:
%
% https://q.uiver.app/#q=WzAsNixbMCwwLCJBXzEiXSxbMSwwLCJBXzIiXSxbMiwwLCJBXzMiXSxbMCwxLCJBXzIiXSxbMSwxLCJBXzIiXSxbMiwxLCJBXzMiXSxbMCwzLCJlX2MiLDIseyJjdXJ2ZSI6MX1dLFsxLDQsIlxcZGVscmVfYyIsMl0sWzIsNSwiXFxwdXNoX3tjJ30oXFxkZWxyZV9jKSIsMix7ImN1cnZlIjoxfV0sWzAsMSwiYyIsMCx7InN0eWxlIjp7ImJvZHkiOnsibmFtZSI6ImJhcnJlZCJ9LCJoZWFkIjp7Im5hbWUiOiJub25lIn19fV0sWzEsMiwiYyciLDAseyJzdHlsZSI6eyJib2R5Ijp7Im5hbWUiOiJiYXJyZWQifSwiaGVhZCI6eyJuYW1lIjoibm9uZSJ9fX1dLFszLDQsInIoQV8yKSIsMix7InN0eWxlIjp7ImJvZHkiOnsibmFtZSI6ImJhcnJlZCJ9LCJoZWFkIjp7Im5hbWUiOiJub25lIn19fV0sWzQsNSwiYyciLDIseyJzdHlsZSI6eyJib2R5Ijp7Im5hbWUiOiJiYXJyZWQifSwiaGVhZCI6eyJuYW1lIjoibm9uZSJ9fX1dLFsyLDUsIlxcaWQiLDAseyJjdXJ2ZSI6LTF9XSxbMCwzLCJlX2MiLDAseyJjdXJ2ZSI6LTF9XSxbOCwxMywiXFxiaXNpbSIsMSx7InNob3J0ZW4iOnsic291cmNlIjoyMCwidGFyZ2V0IjoyMH0sInN0eWxlIjp7ImJvZHkiOnsibmFtZSI6Im5vbmUifSwiaGVhZCI6eyJuYW1lIjoibm9uZSJ9fX1dLFs2LDE0LCJcXGJpc2ltIiwxLHsic2hvcnRlbiI6eyJzb3VyY2UiOjIwLCJ0YXJnZXQiOjIwfSwic3R5bGUiOnsiYm9keSI6eyJuYW1lIjoibm9uZSJ9LCJoZWFkIjp7Im5hbWUiOiJub25lIn19fV1d
\[\begin{tikzcd}[ampersand replacement=\&,column sep=large]
	{A_1} \& {A_2} \& {A_3} \\
	{A_2} \& {A_2} \& {A_3}
	\arrow["c", "\shortmid"{marking}, no head, from=1-1, to=1-2]
	\arrow[""{name=0, anchor=center, inner sep=0}, "{e_c}"', curve={height=6pt}, from=1-1, to=2-1]
	\arrow[""{name=1, anchor=center, inner sep=0}, "{e_c}", curve={height=-6pt}, from=1-1, to=2-1]
	\arrow["{c'}", "\shortmid"{marking}, no head, from=1-2, to=1-3]
	\arrow["{\delre_c}"', from=1-2, to=2-2]
	\arrow[""{name=2, anchor=center, inner sep=0}, "{\push_{c'}(\delre_c)}"', curve={height=6pt}, from=1-3, to=2-3]
	\arrow[""{name=3, anchor=center, inner sep=0}, "\id", curve={height=-6pt}, from=1-3, to=2-3]
	\arrow["{r(A_2)}"', "\shortmid"{marking}, no head, from=2-1, to=2-2]
	\arrow["{c'}"', "\shortmid"{marking}, no head, from=2-2, to=2-3]
	\arrow["\bisim"{description}, draw=none, from=0, to=1]
	\arrow["\bisim"{description}, draw=none, from=2, to=3]
\end{tikzcd}\]
%
% Note that the horizontal composition of the squares is defined, as they are
% intensional squares.
We have used the push-pull structure on $c'$ to turn the perturbation on $A_2$
into a perturbation on $A_3$. Lastly we note that the bottom edge of the square
is equal to $c'$ by the fact that $c'$ is downward closed. The squares for \DIFdelbegin \DIFdel{UpR,
DnL, and DnR }\DIFdelend \DIFaddbegin \DIFadd{$\upr$,
$\dnl$, and $\dnr$ }\DIFaddend are obtained analogously.

% Now that we have defined an intensional model with an interpretation for the
% dynamic type, we can apply the abstract constructions introduced in Section
% \ref{sec:extensional-model-construction}. Doing so, we obtain an extensional
% model of gradual typing, where the squares are given by the ``bisimilarity
% closure'' of the intensional error ordering.

\subsection{Relational Adequacy}\label{sec:adequacy}

Now that we have established our denotational model, we can, as in
Section~\ref{sec:big-step-term-semantics}, extract a big step semantics. Our final
goal then is to prove that the extensional squares in our model imply
the graduality property for this final big step semantics.
\begin{comment}
First we establish some notation. Fix a morphism $f : 1 \to \li \mathbb{N} \cong
\li \mathbb{N}$. We write $f \da n$ to mean that there exists $m$ such that $f =
\delta^m(\eta n)$ and $f \da \mho$ to mean that there exists $m$ such that $f =
\delta^m(\mho)$.

Recall that $\ltbisim$ denotes the relation on value morphisms defined as the
closure of the intensional error-ordering on morphisms under weak bisimilarity
on bboth sides. Specialized to the setting of morphisms from $1 \to U \li
\mathbb{N}$, we have $f \ltbisim g$ iff there exist $f'$ and $g'$ with
%
\[ f \bisim_{\li \mathbb{N}} f' \le_{\li \mathbb{N}} g' \bisim_{\li \mathbb{N}} g. \]\max{$\le$? I thought it was $\ltdyn$}
%
We first note that in this ordering, the semantics of error is not equivalent to
the semantics of the diverging term. The main result we would like to show is as
follows: If $f \ltbisim g : \li \mathbb{N}$, then:
\begin{itemize}
  \item If $f \da n$ then $g \da n$.
  \item If $g \da \mho$ then $f \da \mho$.
  \item If $g \da n$ then $f \da n$ or $f \da \mho$.
\end{itemize}
%
However, this is actually not provable! Intuitively, the issue is
that to prove this result we need to the above result is a ``global''\max{what does ``global'' mean?} result, and it is not possible to prove such results
in the guarded setting\max{isn't everything we do in a ``guarded setting''?}. In particular, if we tried to prove the above result in
the guarded setting we would run into a problem where we would have a natural
number ``stuck'' under a $\later$, with no way to get out the underlying number.
%
Thus, to prove our adequacy result, we need to leave the guarded
setting\max{``pass to leave the guarded setting''?}. We have seen in Section \ref{sec:big-step-term-semantics} that clock
quantification makes this possible.
\end{comment}
The result we seek to show is the following:
% If $f \bisim^{gl} f' \ltls^{gl} g' \bisim^{gl} g$, then
\begin{theorem}[adequacy]
  If $\cdot \vdash M \ltdyn N : \nat$, then:
  \begin{itemize}
    \item If $M \Downarrow n$ then $N \Downarrow n$.
    \item If $N \Downarrow \mho$ then $M \Downarrow \mho$.
    \item If $N \Downarrow n$ then $M \Downarrow n$ or $M \Downarrow \mho$.
  \end{itemize}
\end{theorem}
%
% Recall as discussed in that section that the definitions of our model are all
% parameterized by a clock $k$. 
%
It therefore remains to show that the denotation of term precision in our model
$f \ltsqbisim{}{} g$, where $f, g : \li \mathbb{N}$, implies a relationship
between the partial elements $\mathbb{N} + {\mho}$ corresponding to $f$ and $g$. Recall
that to define the big-step semantics we used the equivalence of the
globalization of the free error domain and Capretta's coinductive delay monad:
$\li^{gl} \mathbb{N} \cong \delay(\mathbb{N} + {\mho})$. We now extend this
correspondence to a relationship between a global version of our denotation of
term precision and the corresponding relation on the delay monad.

% We saw in Section \ref{sec:big-step-term-semantics} that there is an isomorphism
% between the globalization of the free error domain and Capretta's coinductive
% delay monad: $\li^{gl} \mathbb{N} \cong \delay(\mathbb{N} + {\mho})$. We now
% seek to extend this to a relatinoship between a global version of our semantic
% notion of term precision and the corresponding relation for the delay monad.

To that end, we define\DIFaddbegin \DIFadd{, using clock quantification, }\DIFaddend a version of our denotation of term precision
that relates terms of type $\li^{gl} \mathbb{N}$. More concretely, consider the
predomain $\mathbb{N}$ with ordering and bisimilarity both the equality
relation. We define a global version of the lock-step error ordering and the
weak bisimilarity relation on elements of $\li^{gl} \mathbb{N}$; the former is
defined by
%
\( x \ltls^{gl} y := \forall k. x[k] \ltls y[k], \)
%
and the latter is defined by
%
\( x \bisim^{gl} y := \forall k. x[k] \bisim y[k]. \)
%
\DIFaddbegin \DIFadd{In these two definitions, we apply the clock $k$ to the elements $x$ and $y$ of
type $\li^{gl} \mathbb{N} := \forall (k : \Clock). \li^k \mathbb{N}$.
}\DIFaddend 

Correspondingly, we define a version of the lock-step error ordering and weak
bisimilarity relation on $\delay(\mathbb{N} + {\mho})$ 
(see Section \ref{sec:relations-on-delay-monad} in the appendix for the definitions).
% \max{do these intermediate details matter for the main body of the paper?}
%
%
By adapting Theorem 4.3 of \cite{kristensen-mogelberg-vezzosi2022} to the
setting of inductively-defined relations, we can show that both the global
lock-step error ordering and the global weak bisimilarity admit coinductive
definitions. In particular, modulo the above isomorphism between $\li^{gl} X$
and $\delay(\mathbb{N} + {\mho})$, the global version of the lock-step error
ordering is equivalent to the lock-step error ordering on $\delay(\mathbb{N} +
{\mho})$, and the global version of weak bisimilarity is equivalent to weak
bisimilarity on $\delay(\mathbb{N} + {\mho})$.
%
This implies that the global denotation of the term precision semantics for
$\li^{gl} \mathbb{N}$ agrees with the corresponding relation $\ltsqbisim{}{}^\text{Delay}$ for $\delay(\mathbb{N} + {\mho})$, i.e.,
the closure of its lock-step ordering under weak bisimilarity on both sides.
%
Then to obtain the desired result, it suffices to show being related in
$\ltsqbisim{}{}^\text{Delay}$ implies that the corresponding partial elements of
$\mathbb{N} + {\mho}$ are related in the expected manner. That is, if $d
\ltsqbisim{}{}^\text{Delay} d'$, then $p(d) \ltdyn^{Partial} p(d')$ where $p$ is
a partial function from delays into $\mathbb{N} + {\mho}$ defined by existential
quantification over the number of steps $i$ in which $d$ terminates, and the
relation $\ltdyn^\text{Partial}$ is defined such that $\mho$ is the least
element, $\bot \ltdyn^\text{Partial} \bot$, and $n \ltdyn^\text{Partial} n$.
This implication follows from the definition of $\ltsqbisim{}{}^\text{Delay}$.

% Then our result follows by proving the corresponding result for
% $\delay(\mathbb{N} + {\mho})$ which in turn follows from the definitions of the
% relations.

% \max{this doesn't actually relate to our notion of big step semantics}

% We have been writing the type as $\li X$, but it is perhaps more accurate to write it as $\li^k X$ to
% emphasize that the construction is parameterized by a clock $k$.

% Need : nat is clock irrelevant, as well as the inputs and outputs of effects
% Axioms about forcing clock
% Adapt prior argument to get that the defining of the global bisim
% and global lock-step error ordering are coinductive


\section{Discussion}\label{sec:discussion}

%DIF <  \eric{I moved this from the intro. We may need to revise it.}
\DIFdelbegin %DIFDELCMD < 

%DIFDELCMD < %%%
\DIFdelend \subsection{Related Work}

We give an overview of current approaches to proving graduality and discuss
their limitations where appropriate.

%%\subsubsection{From Static to Gradual}

The methods of Abstracting Gradual Typing \cite{garcia-clark-tanter2016} and the
formal tools of the Gradualizer \cite{cimini-siek2016} provide a way to derive
from a statically-typed language a language that satisfies graduality by
construction. The main downside to these approaches lies in their inflexibility:
since the process in entirely mechanical, the language designer must adhere to
the predefined framework.  Many gradually typed languages do not fit into either
framework, e.g., Typed Racket \cite{tobin-hochstadt06, tobin-hochstadt08}\DIFaddbegin \DIFadd{, }\DIFaddend and
the semantics produced is not always the desired one.
%
Furthermore, while these frameworks do prove graduality of the resulting
languages, they do not show the correctness of the equational theory, which is
equally important to sound gradual typing.


%DIF < \subsubsection{Double-Categorical Semantics}
%DIF >  \subsubsection{Graduality via Embedding-Projection Pairs}

%DIF <  New and Licata \cite{new-licata18} developed an axiomatic account of
%DIF <  the graduality relation on a call-by-name cast calculus terms and
%DIF <  showed that the graduality proof could be modeled using semantics in
%DIF <  certain kinds of \emph{double categories}, categories internal to the
%DIF <  category of categories. A double category extends a category with a
%DIF <  second notion of morphism, often a notion of ``relation'' to be paired
%DIF <  with the notion of functional morphism, as well as a notion of
%DIF <  functional morphisms preserving relations. In gradual typing the
%DIF <  notion of relation models type precision and the squares model the
%DIF <  term precision relation. This approach was influenced by the semantics
%DIF <  of parametricity using reflexive graph categories
%DIF <  \cite{ma-reynolds,dunphythesis,reynoldsprogramme}: reflexive graph
%DIF <  categories are essentially double categories without a notion of
%DIF <  relational composition. In addition to capturing the notions of type
%DIF <  and term precision, the double categorical approach allows for a
%DIF <  \emph{universal property} for casts: upcasts are the \emph{universal}
%DIF <  way to turn a relation arrow into a function in a forward direction
%DIF <  and downcasts are the dual universal arrow.  
\DIFdelbegin %DIFDELCMD < 

%DIFDELCMD < %%%
\DIFdel{New and Licata \mbox{%DIFAUXCMD
\cite{new-licata18} }\hskip0pt%DIFAUXCMD
developed an axiomatic account of the
graduality relation on a call-by-name cast calculus and showed that the
graduality proof could be modeled using semantics in certain kinds of double
categories.
%DIF < 
Later, New, Licata and Ahmed \mbox{%DIFAUXCMD
\cite{new-licata-ahmed2019} }\hskip0pt%DIFAUXCMD
extended this axiomatic
treatment from call-by-name to call-by-value as well by giving an axiomatic
theory of type and term precision in call-by-push-value. The double-categorical
framework serves as a foundation on which we base our model construction in this
paper. However, because of the lack of transitivity of the relation modeling
term precision, our model is not a double category but instead a reflexive graph
category, that is, a double category without a notion of relational composition.
}%DIFDELCMD < 

%DIFDELCMD < %%%
%DIF < % With this notion of abstract categorical model in hand, denotational
%DIF < % semantics is then the work of constructing concrete models that
%DIF < % exhibit the categorical construction. New and Licata
%DIF < % \cite{new-licata18} present such a model using categories of
%DIF < % $\omega$-CPOs, and this model was extended by Lennon-Bertrand,
%DIF < % Maillard, Tabareau and Tanter to prove graduality of a gradual
%DIF < % dependently typed calculus $\textrm{CastCIC}^{\mathcal G}$. This
%DIF < % domain-theoretic approach meets our criteria of being a semantic
%DIF < % framework for proving graduality, but suffers from the limitations of
%DIF < % classical domain theory: the inability to model viciously
%DIF < % self-referential structures such as higher-order extensible state and
%DIF < % similar features such as runtime-extensible dynamic types. Since these
%DIF < % features are quite common in dynamically typed languages, we seek a
%DIF < % new denotational framework that can model these type system features.
%DIFDELCMD < 

%DIFDELCMD < %%%
%DIF < \subsubsection{Embedding-Projection Pairs}
%DIFDELCMD < 

%DIFDELCMD < %%%
\DIFdelend A line of work by New, Licata and Ahmed
\DIFaddbegin \DIFadd{\mbox{%DIFAUXCMD
\cite{new-ahmed2018,new-licata18,new-licata-ahmed2019} }\hskip0pt%DIFAUXCMD
develops an axiomatic
account of gradual typing and }\DIFaddend proves graduality by interpreting type precision
$c : A \ltdyn A'$ as an \emph{embedding-projection pair}, that is\DIFdelbegin \DIFdel{a pure }\DIFdelend \DIFaddbegin \DIFadd{, a pure
embedding }\DIFaddend function $e : A \to A'$ and a possibly erroring projection $p : \li A'
\multimap \li A$ such that $p \circ e = \id$ and $e \circ p \ltdyn \id$\DIFdelbegin \DIFdel{\mbox{%DIFAUXCMD
\cite{new-ahmed2018,new-licata18,new-licata-ahmed2019}}\hskip0pt%DIFAUXCMD
}\DIFdelend . At
first glance, our model of precision looks different: it is a \emph{relation}
$\sem{c} : A \rel A'$ with a quasi-representability structure that gives us
something \DIFdelbegin \DIFdel{like }\DIFdelend \DIFaddbegin \DIFadd{close to }\DIFaddend $e$ and $p$ in an embedding-projection pair. The relationship
between these models is that \emph{if} the relation were \emph{truly}
representable we would have that $e$ and $p$ form a \emph{Galois connection} $e
\circ p \ltdyn \id$ and $\id \ltdyn p \circ e$. However, we have dropped the
stronger property of retraction from our analysis in this work. With true
representability, the relation $c$ is \emph{uniquely determined} by the
embedding $e$, but since we only have quasi-representability, we need to keep
the relation around explicitly. So the extra complexity of managing explicit
relations is a cost of the intensional reasoning that guarded type theory
introduces. 
%DIF > 
\DIFaddbegin \DIFadd{The line of work by New, Licata, and Ahmed involved both denotational methods
and operational logical relations approaches.
%DIF >  The New-Licata-Ahmed axiomatic approach to gradual typing can be
%DIF >  applied to both operational and denotational semantics.
}\DIFaddend 


%DIF < \subsubsection{Step-Indexed Logical Relations Models}
\DIFdelbegin \DIFdel{The standard alternative to domain theory that scales to essentially
arbitrary }\DIFdelend %DIF >  Later, New, Licata and Ahmed \cite{new-licata-ahmed2019} extended this axiomatic
%DIF >  treatment from call-by-name to call-by-value as well by giving an axiomatic
%DIF >  theory of type and term precision in call-by-push-value. 
\DIFaddbegin 

%DIF >  \subsubsection{Denotational and Operational Approaches}


%DIF > \subsubsection{Double-Categorical Semantics}

\DIFadd{On the denotational side, New and Licata \mbox{%DIFAUXCMD
\cite{new-licata18} }\hskip0pt%DIFAUXCMD
showed that the
graduality proof could be modeled using semantics in certain kinds of double
categories, which are categories internal to the category of categories. A
double category extends a category with a second notion of morphism, often a
notion of ``relation'' to be paired with the notion of functional morphism, as
well as a notion of functional morphisms preserving relations. In gradual
typing, the notion of relation models type precision and the squares model the
term precision relation.
%DIF > 
%DIF >  This approach was influenced by the semantics of parametricity using reflexive
%DIF >  graph categories \cite{ma-reynolds,dunphythesis,reynoldsprogramme}: reflexive
%DIF >  graph categories are essentially double categories without a notion of
%DIF >  relational composition. In addition to capturing the notions of type and term
%DIF >  precision, the double categorical approach allows for a \emph{universal
%DIF >  property} for casts: upcasts are the \emph{universal} way to turn a relation
%DIF >  arrow into a function in a forward direction and downcasts are the dual
%DIF >  universal arrow.
%DIF > 
The double-categorical framework serves as a foundation on which we base our
model construction in this paper. However, our model is not actually a double
category: the lack of transitivity of bisimilarity means that we cannot compose
extensional squares horizontally.
}

\DIFadd{With the notion of abstract categorical model for gradual typing in hand,
denotational semantics is then the work of constructing concrete models that
exhibit the categorical construction. New and Licata \mbox{%DIFAUXCMD
\cite{new-licata18} }\hskip0pt%DIFAUXCMD
present
such a model using categories of $\omega$-CPOs, and this model was extended by
Lennon-Bertrand, Maillard, Tabareau and Tanter \mbox{%DIFAUXCMD
\cite{gradualizing-cic} }\hskip0pt%DIFAUXCMD
to prove
graduality of a gradual dependently typed calculus $\textrm{CastCIC}^{\mathcal
G}$. This domain-theoretic approach meets our criterion of being a semantic
framework for proving graduality, but suffers from the limitations of classical
domain theory: the inability to model viciously }\DIFaddend self-referential \DIFdelbegin \DIFdel{definitions is }\emph{\DIFdel{step-indexing}} %DIFAUXCMD
\DIFdel{or its
synthetic form of }\emph{\DIFdel{guarded recursion}}%DIFAUXCMD
\DIFdel{.
A }\DIFdelend \DIFaddbegin \DIFadd{structures such
as higher-order extensible state and similar features such as runtime-extensible
dynamic types.
}

%DIF >  Since these features are quite common in dynamically typed languages, we have
%DIF >  developed in this work a new denotational framework that can be extended to
%DIF >  model these type system features.

%DIF > \subsubsection{Step-Indexed Logical Relations Models}
%DIF >  The standard alternative to domain theory that scales to essentially
%DIF >  arbitrary self-referential definitions is \emph{step-indexing} or its
%DIF >  synthetic form of \emph{guarded recursion}. 

\DIFadd{On the operational side, a }\DIFaddend series of works
\cite{new-ahmed2018, new-licata-ahmed2019, new-jamner-ahmed19}
developed step-indexed logical relations models of gradually typed languages\DIFdelbegin \DIFdel{based on operational semantics}\DIFdelend .
Unlike classical domain theory, such step-indexed techniques \DIFdelbegin \DIFdel{are capable of modeling
}\DIFdelend \DIFaddbegin \DIFadd{can scale to
essentially arbitrary self-referential definitions, which means they can model
}\DIFaddend higher-order store and runtime-extensible dynamic types
\cite{appelmcallester01,ahmed06,neis09,new-jamner-ahmed19}. However, \DIFaddbegin \DIFadd{as is
common with operational approaches, }\DIFaddend their proof developments are highly
repetitive and technical, with each development formulating a logical relation
from first-principles and proving many of the same tedious lemmas without
reusable mathematical abstractions.
%
%DIF > 
% Additionally, the logical relations constructed to prove graduality in prior
% work \cite{new-ahmed2018,new-licata-ahmed2019,new-giovannini-licata-2022} suffer
% from technical complications of requiring two separate expression relations, one
% that counts steps on the left and the other on the right, and there is no
% analogue of this in our approach. However, using two expression relations allows
% some but not all transitive reasoning of term precision to be recovered. In the
% future we aim to explore if this approach is feasible in guarded semantics.
%
This is addressed somewhat by Siek and Chen \cite{siek-chen2021}, who give a
proof in Agda of graduality for an operational semantics \DIFdelbegin \DIFdel{who use }\DIFdelend \DIFaddbegin \DIFadd{using }\DIFaddend a \emph{guarded
logic of propositions} shallowly embedded in Agda. The guarded logic simplifies
the treatment of step-indexed logical relations, but the approach is still
fundamentally operational, and so the main lemmas of the work are still tied to
the particular operational syntactic calculus being modeled.
\DIFdelbegin \DIFdel{A further advantage of the denotational approach is that it
easily validates equational reasoning, not just graduality, and it is
completely independent of any particular syntax of gradual typing. 
}\DIFdelend %DIF >  A further advantage of the denotational approach is that it easily validates
%DIF >  equational reasoning, not just graduality, and it is completely independent of
%DIF >  any particular syntax of gradual typing.

%DIF < \subsubsection{Gradual Dependent Types}
%DIF <  Discuss Joey Eremondi's thesis on gradual dependent types
\DIFaddbegin \DIFadd{Our work combines the benefits of both the denotational and step-indexed
approaches. We take as a starting point the denotational approach of New and
Licata, but we work with guarded type theory (a synthetic form of step-indexing)
rather than classical domain theory. As we have seen, in the guarded setting the compositional
theory of New-Licata breaks down. 
%DIF >  when accounting for steps, the casts are no
%DIF >  longer embedding-projection pairs, and we cannot carry out compositional
%DIF >  reasoning on terms when those terms differ by an arbitrary unknown number of
%DIF >  steps. 
Our work provides a novel refinement to their theory that allows for some
compositional reasoning to be recovered even in the guarded setting.
}

%DIF >  \subsubsection{Gradual Dependent Types}
\DIFaddend Eremondi \cite{Eremondi_2023} uses guarded type theory to
define a syntactic model for a gradually-typed source
language with dependent types. By working in guarded type theory, they are
able to give an exact semantics to non-terminating computations in a language
which is logically consistent, allowing for metatheoretic results about the
source language to be proven.
%
Similarly to our approach, they define a guarded lift monad to model potentially-
nonterminating computations and use guarded recursion to model the dynamic type.
However, they do not give a denotational semantics to term precision and it is unclear
how to prove graduality in this setting.
\DIFdelbegin \DIFdel{The }\DIFdelend \DIFaddbegin \DIFadd{Their }\DIFaddend work includes a formalization of the syntactic model in Guarded Cubical Agda.

%DIF <  Discuss Jeremy Siek's work on graduality in Agda
\DIFaddbegin \subsection{\DIFadd{Comparison of Denotational and Operational Approaches}}
%DIF >  \eric{Some of this could be folded into the previous section, but this was one of the
%DIF >  big-picture questions from one of the reviewers.}
\DIFaddend 

\DIFaddbegin \DIFadd{Denotational methods have several benefits compared to operational approaches,
and this carries over to the setting of gradual typing. First, denotational
methods are compositional and reusable, and allow for the use of existing
mathematical constructs and theorems, e.g., partial orderings, monads, etc,
while operational methods tend to require a significant amount of boilerplate
work to be done from scratch in each new development.
%DIF > 
As a specific example of the compositional nature of our approach, the treatment
of the cast rules in our work is more compositional than in previous work using
operational semantics. Recall from Section \ref{sec:towards-relational-model}
that the cast rules needed for the proof of graduality build in composition of
type precision derivations. Rather than proving these from scratch in the model,
we are able to take as primitive simpler versions of the cast rules whose
validity in the model is easier to establish. Then from these simpler rules, we
derive the original ones using compositional reasoning.
}

\DIFadd{A further advantage of the denotational approach is that it is completely
independent of any particular syntax of gradual typing. This allows for reuse
across multiple languages and makes it particularly straightforward to
accommodate additions to a language: adding support for a new type amounts to
defining a new object and extending the dynamic type accordingly. It also may be
possible to model alternative cast semantics by changing the definition of the
dynamic type and some type casts.
%DIF > 
In contrast, operational methods are not as readily extensible, generally
requiring adding cases to the logical relations and the inductive proofs. In a
mechanized metatheory, this manifests as a ``copy-paste'' reusability rather than
the true code reuse one obtains with a denotational semantics.
%DIF > 
An additional benefit to the denotational approach is that it is trivial to
establish the validity of the $\beta$ and $\eta$ principles, because they are
equalities in the semantics, whereas they require tedious manual proof in the
operational approach.
}

\DIFadd{Finally, while we advocate for a denotational approach, the overall structure of
our semantics could be applied to an operational logical relations proof. For
example, there could be a step-indexed logical relation for strong error
ordering and one for weak bisimilarity, corresponding to the fact that objects
in our denotational model have an error ordering and a bisimilarity relation.
Then the precision rules for casts could be similarly validated using
perturbations as we did in this work.
}


\DIFaddend \subsection{Mechanization}\DIFaddbegin \label{sec:mechanization}
\DIFaddend % Discuss Guarded Cubical Agda and mechanization efforts
In parallel with developing the theory discussed in this paper, we have
developed a partial formalization of our results in Guarded Cubical Agda
\cite{veltri-vezzosi2020}.
%
We formalized the major components of the definition of the concrete model
described in Section \ref{sec:concrete-relational-model}: predomains, error
domains, morphisms, relations, squares, using the guarded
features to define the free error domain, the lock-step error ordering, and weak
bisimilarity. We also formalized the no-go theorem from Section
\DIFdelbegin \DIFdel{\ref{sec:concrete-term-model}}\DIFdelend \DIFaddbegin \DIFadd{\ref{sec:towards-relational-model}}\DIFaddend .

Building on these definitions, we formalized the notion of semantic
perturbations and push-pull structures as well as quasi-representable relations,
culminating in the definition of value and computation types and relations as
introduced in Section \ref{sec:concrete-relational-model}.
%
% culminating in the definition of value and computation types as predomains
% (resp. error domains) equipped with a monoid of syntactic perturbations with an
% interpretation homomorphism into the semantic perturbations. 
%
We constructed the predomain for the dynamic type via mixed induction and
guarded recursion, and defined its monoid of perturbations. We also defined the
relations $\inat$, $\itimes$, and $\iarr$, and proved that they are
quasi-representable.
%DIF < 
\DIFdelbegin \DIFdel{We have not yet formalized the constructions involving value and computation
relations: showing that these relations compose, and defining the actions of the
functors $\li$, $U$, $\times$, and $\arr$ on the relations. We have also not yet
formalized the syntax-to-semantics translation.
}\DIFdelend 

%DIF >  We have not yet formalized the constructions involving value and computation
%DIF >  relations: showing that these relations compose, and defining the actions of the
%DIF >  functors $\li$, $U$, $\times$, and $\arr$ on the relations. We have also not yet
%DIF >  formalized the syntax-to-semantics translation.
\DIFaddbegin 

\DIFaddend % \max{need to say what's not formalized}

% We plan to formalize the construction of perturbations and 
% quasi-representable relations, but we have yet to decide
% whether to follow the approach we take in this work and define
% the abstract notion of intensional model and formalize the constructions in that setting,
% and then apply those abstract constructions to the concrete model.
% Alternatively, it may be better from a mechanization standpoint
% to carry out those abstract constructions explicitly in the concrete model,
% i.e., our representation of objects in the concrete model of predomains
% would include a field for the perturbations and our notion of relations
% would include fields for the push-pull property and quasi-representability.
% We leave this investigation to future work.

Lastly, we have formalized the big-step term semantics discussed in Section
\ref{sec:big-step-term-semantics} and the adequacy of the relational model
discussed in Section \ref{sec:adequacy}. This required us to add axioms about
clock quantification as well as axioms asserting the \emph{clock-irrelevance} of
booleans and natural numbers since as of this writing these axioms are not
built-in to Guarded Cubical Agda. These axioms are discussed in prior work on
guarded type theory \cite{atkey-mcbride2013, kristensen-mogelberg-vezzosi2022}.
The mechanization of adequacy also required us to formalize some essential
lemmas involving clocks and clock-irrelevance; we are considering later
refactoring these as part of a ``standard library'' for Guarded Cubical
Agda.
% \max{say these axioms are taken from prior work and cite that}

\DIFaddbegin \DIFadd{The remaining formalization work is the following:
}\begin{enumerate}
    \item \DIFadd{Showing that the functors $\times$ and $\arr$ preserve
    quasi-representability, which requires tedious reasoning about the Kleisli actions of
    $\times$ and $\arr$. These results are proved in the Appendix.
    }\item \DIFadd{Verifying the rules in the model corresponding to the equations for
    type precision derivations (see the bottom of Figure \ref{fig:gtlc-syntax}).
    These are also included in the Appendix.
    }\item \DIFadd{Formalizing the syntax-to-semantics translation.
}\end{enumerate} 

\DIFaddend \begin{comment}
% prove graduality in the syntax of 
% GTLC, which involves the construction of the abstract model described in 
% \ref{sec:concrete-model} and the extensional model with external dynamic 
% type. We also plan to formalize the adequacy result in \ref{sec:appendix-adequacy}.

% step-2 
Then we plan to construct the step-2 intensional model. Besides all the 
data in step-1, we need to include perturbations, functors $\times$, $\arr$, $U$, and $F$ that preserve 
perturbations and push/pull properties for all morphisms on value and 
computation types. Notice that for any object $A$ which has value type, 
we will take not only the monoid of perturbations $P^V_A$ and the monoid 
homomorphism $\ptbv_A : \pv_A \to \vf(A,A)$ on itself, but also $P^C_{F A}
$ and $\ptbe_{F A} : \pe_{F A} \to \ef(F A,F A)$ on $F A$, which have 
computation types. Similarly, for any computation object $B$, we will 
construct the perturbations on $U B$ besides the monoid $P^C_B$ and 
monoid homomorphism $\ptbe_B : \pe_B \to \ef(B,B)$. Also, for functors 
that preserves perturbations, we need to include the ones in the context 
of Kleisli category. For this part, we need to define the perturbation on 
not only the objects itself, but also the global lift and delay of objects, 
which requires us to provide each piece of supporting constructor. This step 
and futher steps towards to the model construction are still 
work-in-progress, but once it's finished, we will provide a complete 
framework which takes formalization on an explicit type and obtains an 
extensional model.

% step-3
In the step-3 intensional model, we will enhance it with 
quasi-representability. For any value relation $c : A \rel A'$, we need 
to show that there exists a left-representation structure for $c$ and a 
right-representation structure for $F\ c$. Correspondingly, for any 
computation relation $d : B \rel B'$, we will show there exists a 
right-representation structure for $d$ and a left-representation 
structure for $U\ d$. As we define the quasi-representability for value 
and computation relation, we will construct the quasi-representability on 
the function and product of the relation, which makes it necessary to 
have the dual version of quasi-representability.

% step-4 construct a concrete dynamic type and apply it to the abstract model
After defining the abstract model and its interface, we will model GTLC 
by providing explicit construction triples of dynamic type at each step, 
which includes defining Dyn as a predomain, its pure and Kleisli 
perturbation monoids, push/pull property for pure and Kleisli 
perturbation, as well as quasi-representability. The 
quasi-representability involves explicit rules which show that Nat is 
more precise than Dyn (Inj-Nat) and Dyn $\to$ Dyn is more precise than 
Dyn (Inj-Arr). Currently, we have formalized the concrete construction of 
Dyn in Cubical Agda and it was more challenging than expected because we 
define Dyn using the technique of guarded recursion and fixed point, which 
means that every time we analyze the case inside of Dyn, we need to unfold 
it and add corresponding proof. 

% adequacy
Besides the abstract model and its concrete construction on dynamic type, 
we will also formalize the adequacy result in \ref{sec:appendix-adequacy}, 
which involves clock quantification of the lift monad, the weak bisim 
relation, and the lock-step error ordering. In order to prove adequacy, 
we will first prove that the global lift of X is isomorphic to Delay(1 + X)
whether X is clock-irrelevant or not. Then, we aim to prove the equivalence 
between the global lock-step error ordering and the error ordering observed 
in Delay(1 + X) and equivalence between the global weak bisimilarity 
relation and the weak bisimilarity relation on Delay(1 + X). We have 
finished some prerequisite proofs on clock quantification and postulated 
some theorems on clock globalization.
\end{comment}



% \subsection{Synthetic Ordering}
% \max{cut this subsection if we need space}
% A key to managing the complexity of our concrete construction is in
% using a \emph{synthetic} approach to step-indexing rather than working
% analytically with presheaves. This has helped immensely in our ongoing
% mechanization in cubical Agda as it sidesteps the need to formalize
% these constructions internally. 
% %
% However, there are other aspects of the model, the bisimilarity and
% the monotonicity, which are treated analytically and are similarly
% tedious.
% %
% It may be possible to utilize further synthetic techniques to reduce
% this burden as well, and have all type intrinsically carry a notion of
% bisimilarity and ordering relation, and all constructions to
% automatically preserve them.
% %
% A synthetic approach to ordering is common in (non-guarded) synthetic
% domain theory and has also been used for synthetic reasoning for cost
% models \cite{fiore_1997,GrodinNSH24}.

\subsection{Future Work}

% \max{Expand on this, e.g. what some of the challenges might be and what would be reusable}
% In the future, we plan to apply our approach to give a denotational
% semantics for languages that feature higher-order state or
% runtime-extensible dynamic typing
% \cite{DBLP:journals/corr/abs-2210-02169} as well as richer type
% disciplines such as gradual dependent types and effect systems.

%DIF <  \max{Expand on this, e.g. what some of the challenges might be and what would be reusable}
\DIFdelbegin \DIFdel{In the future, we plan }\DIFdelend \DIFaddbegin \DIFadd{Our immediate next step is }\DIFaddend to apply our approach to give a denotational
semantics \DIFdelbegin \DIFdel{for languages that feature }\DIFdelend \DIFaddbegin \DIFadd{to gradually-typed languages with advanced type systems including
}\DIFaddend higher-order state or runtime-extensible dynamic typing
\cite{DBLP:journals/corr/abs-2210-02169}\DIFaddbegin \DIFadd{, }\DIFaddend as well as richer type disciplines
such as gradual dependent types and effect systems. This will involve adapting
prior operational models based on step-indexing (e.g.,
\cite{new-giovannini-licata-2022}) to the denotational setting and using guarded
type theory to obtain solutions to guarded domain equations as we did in this
work for the denotation of the dynamic type. The generality of our techniques
should allow for many of the constructions to be reused across different
languages. For instance, the free error domain construction can be easily
extended \DIFdelbegin \DIFdel{with additional cases }\DIFdelend to model effects besides error and stepping
\DIFdelbegin \DIFdel{.
Additionally, we }\DIFdelend \DIFaddbegin \DIFadd{\mbox{%DIFAUXCMD
\cite{guarded-interaction-trees, probabilistic-fpc-guarded}}\hskip0pt%DIFAUXCMD
.
We also }\DIFaddend aim to complete our Agda formalization and evolve it into a
reusable framework for mechanized denotational semantics of gradually-typed
languages.

\DIFaddbegin \DIFadd{This work has focused on the ``Natural'' semantics of casts , which
validate the full call-by-value $\eta$ laws for types, but other cast
semantics have been proposed such as eager
(\mbox{%DIFAUXCMD
\cite{herman-tomb-flanagan-2010}}\hskip0pt%DIFAUXCMD
) and transient (\mbox{%DIFAUXCMD
\cite{transient}}\hskip0pt%DIFAUXCMD
)
cast semantics which trade off certain $\eta$ equalities for other
benefits such as early error detection or reduced runtime
overhead\mbox{%DIFAUXCMD
\cite{deepandshallowtypes,new-licata-ahmed2019}}\hskip0pt%DIFAUXCMD
. It would be a
good test of the generality of our semantic framework to see if it can
model these alternative semantics by adjusting the semantics of types
and casts, and providing proofs of weakened $\eta$ principles.
}

\DIFadd{Additionally, while our focus in this work has been on reasoning up to
weak bisimilarity, the presence of explicit step counting in the model
could be viewed as a form of }\emph{\DIFadd{cost semantics}}\DIFadd{, where a runtime
cost is incurred from inspecting the dynamic type. This could possibly
be used to verify that cast optimizations such as space efficient
implementations \mbox{%DIFAUXCMD
\cite{herman-tomb-flanagan-2010} }\hskip0pt%DIFAUXCMD
are not only
extensionally correct, but have a lower abstract cost.
}

%DIF >  \max{Expand on this, e.g. what some of the challenges might be and what would be reusable}
%DIF >  In the future, we plan to apply our approach to give a denotational semantics
%DIF >  to gradually-typed languages that feature higher-order state or runtime-extensible dynamic
%DIF >  typing \cite{DBLP:journals/corr/abs-2210-02169} as well as richer type
%DIF >  disciplines such as gradual dependent types and effect systems. This will
%DIF >  involve adapting prior operational models based on step-indexing (e.g.,
%DIF >  \cite{new-giovannini-licata-2022}) to the denotational setting and using guarded
%DIF >  type theory to obtain solutions to guarded domain equations as we did in this
%DIF >  work for the denotation of the dynamic type. The generality of our techniques
%DIF >  should allow for many of the constructions to be reused across different
%DIF >  languages. For instance, the free error domain construction can be easily
%DIF >  extended with additional cases to model effects besides error and stepping.
%DIF >  Additionally, we aim to complete our Agda formalization and evolve it into a
%DIF >  reusable framework for mechanized denotational semantics of gradually-typed
%DIF >  languages.

\DIFaddend %% to gradually-typed
%% languages with algebraic effects, building on prior work on gradual typing for effect handlers
%% \cite{greff}. In particular, that work proves graduality via a complicated step-indexed logical relation,
%% and we hope to prove their results by building a denotational model for GrEff.
%% This would serve as a step towards applying our techniques to prove graduality for languages
%% with other advanced features.

%% The extensional model we construct differs from the usual notion of extensional
%% model considered in prior work on gradual typing in that it lacks horizontal composition of squares.
%% We would like to clarify the relationship between our notion of model and prior extensional models,
%% with the aim of determining whether our approach could allow for the construction of such a model.


\bibliographystyle{ACM-Reference-Format}
\DIFdelbegin %DIFDELCMD < \begin{DIFnomarkup}
%DIFDELCMD < %%% -*-BibTeX-*-
%DIFDELCMD < %%% Do NOT edit. File created by BibTeX with style
%DIFDELCMD < %%% ACM-Reference-Format-Journals [18-Jan-2012].
%DIFDELCMD < 

%DIFDELCMD < \begin{thebibliography}{32}
%DIFDELCMD < 

%DIFDELCMD < %%% ====================================================================
%DIFDELCMD < %%% NOTE TO THE USER: you can override these defaults by providing
%DIFDELCMD < %%% customized versions of any of these macros before the \bibliography
%DIFDELCMD < %%% command.  Each of them MUST provide its own final punctuation,
%DIFDELCMD < %%% except for \shownote{}, \showDOI{}, and \showURL{}.  The latter two
%DIFDELCMD < %%% do not use final punctuation, in order to avoid confusing it with
%DIFDELCMD < %%% the Web address.
%DIFDELCMD < %%%
%DIFDELCMD < %%% To suppress output of a particular field, define its macro to expand
%DIFDELCMD < %%% to an empty string, or better, \unskip, like this:
%DIFDELCMD < %%%
%DIFDELCMD < %%% \newcommand{\showDOI}[1]{\unskip}   % LaTeX syntax
%DIFDELCMD < %%%
%DIFDELCMD < %%% \def \showDOI #1{\unskip}           % plain TeX syntax
%DIFDELCMD < %%%
%DIFDELCMD < %%% ====================================================================
%DIFDELCMD < 

%DIFDELCMD < \ifx \showCODEN    \undefined \def \showCODEN     #1{\unskip}     \fi
%DIFDELCMD < \ifx \showDOI      \undefined \def \showDOI       #1{#1}\fi
%DIFDELCMD < \ifx \showISBNx    \undefined \def \showISBNx     #1{\unskip}     \fi
%DIFDELCMD < \ifx \showISBNxiii \undefined \def \showISBNxiii  #1{\unskip}     \fi
%DIFDELCMD < \ifx \showISSN     \undefined \def \showISSN      #1{\unskip}     \fi
%DIFDELCMD < \ifx \showLCCN     \undefined \def \showLCCN      #1{\unskip}     \fi
%DIFDELCMD < \ifx \shownote     \undefined \def \shownote      #1{#1}          \fi
%DIFDELCMD < \ifx \showarticletitle \undefined \def \showarticletitle #1{#1}   \fi
%DIFDELCMD < \ifx \showURL      \undefined \def \showURL       {\relax}        \fi
%DIFDELCMD < % The following commands are used for tagged output and should be
%DIFDELCMD < % invisible to TeX
%DIFDELCMD < \providecommand\bibfield[2]{#2}
%DIFDELCMD < \providecommand\bibinfo[2]{#2}
%DIFDELCMD < \providecommand\natexlab[1]{#1}
%DIFDELCMD < \providecommand\showeprint[2][]{arXiv:#2}
%DIFDELCMD < 

%DIFDELCMD < \bibitem[\protect\citeauthoryear{Ahmed}{Ahmed}{2006}]%
%DIFDELCMD <         {ahmed06}
%DIFDELCMD < \bibfield{author}{\bibinfo{person}{Amal~J. Ahmed}.}
%DIFDELCMD <   \bibinfo{year}{2006}\natexlab{}.
%DIFDELCMD < \newblock \showarticletitle{Step-Indexed Syntactic Logical Relations for
%DIFDELCMD <   Recursive and Quantified Types}. In \bibinfo{booktitle}{\emph{Programming
%DIFDELCMD <   Languages and Systems, 15th European Symposium on Programming, {ESOP} 2006,
%DIFDELCMD <   Held as Part of the Joint European Conferences on Theory and Practice of
%DIFDELCMD <   Software, {ETAPS} 2006, Vienna, Austria, March 27-28, 2006, Proceedings}}
%DIFDELCMD <   \emph{(\bibinfo{series}{Lecture Notes in Computer Science},
%DIFDELCMD <   Vol.~\bibinfo{volume}{3924})}, \bibfield{editor}{\bibinfo{person}{Peter
%DIFDELCMD <   Sestoft}} (Ed.). \bibinfo{publisher}{Springer}, \bibinfo{pages}{69--83}.
%DIFDELCMD < \newblock
%DIFDELCMD < \urldef\tempurl%
%DIFDELCMD < \url{https://doi.org/10.1007/11693024\_6}
%DIFDELCMD < \showDOI{\tempurl}
%DIFDELCMD < 

%DIFDELCMD < \bibitem[\protect\citeauthoryear{Appel and McAllester}{Appel and
%DIFDELCMD <   McAllester}{2001}]%
%DIFDELCMD <         {appelmcallester01}
%DIFDELCMD < \bibfield{author}{\bibinfo{person}{Andrew~W. Appel} {and}
%DIFDELCMD <   \bibinfo{person}{David McAllester}.} \bibinfo{year}{2001}\natexlab{}.
%DIFDELCMD < \newblock \showarticletitle{An indexed model of recursive types for
%DIFDELCMD <   foundational proof-carrying code}.
%DIFDELCMD < \newblock \bibinfo{journal}{\emph{ACM Trans. Program. Lang. Syst.}}
%DIFDELCMD <   \bibinfo{volume}{23}, \bibinfo{number}{5} (\bibinfo{date}{sep}
%DIFDELCMD <   \bibinfo{year}{2001}), \bibinfo{pages}{657–683}.
%DIFDELCMD < \newblock
%DIFDELCMD < \showISSN{0164-0925}
%DIFDELCMD < \urldef\tempurl%
%DIFDELCMD < \url{https://doi.org/10.1145/504709.504712}
%DIFDELCMD < \showDOI{\tempurl}
%DIFDELCMD < 

%DIFDELCMD < \bibitem[\protect\citeauthoryear{Atkey and McBride}{Atkey and McBride}{2013}]%
%DIFDELCMD <         {atkey-mcbride2013}
%DIFDELCMD < \bibfield{author}{\bibinfo{person}{Robert Atkey} {and} \bibinfo{person}{Conor
%DIFDELCMD <   McBride}.} \bibinfo{year}{2013}\natexlab{}.
%DIFDELCMD < \newblock \showarticletitle{Productive Coprogramming with Guarded Recursion}.
%DIFDELCMD <   In \bibinfo{booktitle}{\emph{Proceedings of the 18th ACM SIGPLAN
%DIFDELCMD <   International Conference on Functional Programming}} (Boston, Massachusetts,
%DIFDELCMD <   USA) \emph{(\bibinfo{series}{ICFP '13})}. \bibinfo{publisher}{Association for
%DIFDELCMD <   Computing Machinery}, \bibinfo{address}{New York, NY, USA},
%DIFDELCMD <   \bibinfo{pages}{197--208}.
%DIFDELCMD < \newblock
%DIFDELCMD < \showISBNx{9781450323260}
%DIFDELCMD < \urldef\tempurl%
%DIFDELCMD < \url{https://doi.org/10.1145/2500365.2500597}
%DIFDELCMD < \showDOI{\tempurl}
%DIFDELCMD < 

%DIFDELCMD < \bibitem[\protect\citeauthoryear{Bahr, Grathwohl, and Møgelberg}{Bahr
%DIFDELCMD <   et~al\mbox{.}}{2017}]%
%DIFDELCMD <         {bahr-grathwohl-bugge-mogelberg2017}
%DIFDELCMD < \bibfield{author}{\bibinfo{person}{Patrick Bahr}, \bibinfo{person}{Hans~Bugge
%DIFDELCMD <   Grathwohl}, {and} \bibinfo{person}{Rasmus~Ejlers Møgelberg}.}
%DIFDELCMD <   \bibinfo{year}{2017}\natexlab{}.
%DIFDELCMD < \newblock \showarticletitle{The clocks are ticking: No more delays!}. In
%DIFDELCMD <   \bibinfo{booktitle}{\emph{2017 32nd Annual ACM/IEEE Symposium on Logic in
%DIFDELCMD <   Computer Science (LICS)}}. \bibinfo{pages}{1--12}.
%DIFDELCMD < \newblock
%DIFDELCMD < \urldef\tempurl%
%DIFDELCMD < \url{https://doi.org/10.1109/LICS.2017.8005097}
%DIFDELCMD < \showDOI{\tempurl}
%DIFDELCMD < 

%DIFDELCMD < \bibitem[\protect\citeauthoryear{Baunsgaard~Kristensen, Mogelberg, and
%DIFDELCMD <   Vezzosi}{Baunsgaard~Kristensen et~al\mbox{.}}{2022}]%
%DIFDELCMD <         {kristensen-mogelberg-vezzosi2022}
%DIFDELCMD < \bibfield{author}{\bibinfo{person}{Magnus Baunsgaard~Kristensen},
%DIFDELCMD <   \bibinfo{person}{Rasmus~Ejlers Mogelberg}, {and} \bibinfo{person}{Andrea
%DIFDELCMD <   Vezzosi}.} \bibinfo{year}{2022}\natexlab{}.
%DIFDELCMD < \newblock \showarticletitle{Greatest HITs: Higher Inductive Types in
%DIFDELCMD <   Coinductive Definitions via Induction under Clocks}. In
%DIFDELCMD <   \bibinfo{booktitle}{\emph{Proceedings of the 37th Annual ACM/IEEE Symposium
%DIFDELCMD <   on Logic in Computer Science}} (Haifa, Israel) \emph{(\bibinfo{series}{LICS
%DIFDELCMD <   '22})}. \bibinfo{publisher}{Association for Computing Machinery},
%DIFDELCMD <   \bibinfo{address}{New York, NY, USA}, Article \bibinfo{articleno}{42},
%DIFDELCMD <   \bibinfo{numpages}{13}~pages.
%DIFDELCMD < \newblock
%DIFDELCMD < \showISBNx{9781450393515}
%DIFDELCMD < \urldef\tempurl%
%DIFDELCMD < \url{https://doi.org/10.1145/3531130.3533359}
%DIFDELCMD < \showDOI{\tempurl}
%DIFDELCMD < 

%DIFDELCMD < \bibitem[\protect\citeauthoryear{Birkedal, Mogelberg, Schwinghammer, and
%DIFDELCMD <   Stovring}{Birkedal et~al\mbox{.}}{2011}]%
%DIFDELCMD <         {birkedal-mogelberg-schwinghammer-stovring2011}
%DIFDELCMD < \bibfield{author}{\bibinfo{person}{Lars Birkedal},
%DIFDELCMD <   \bibinfo{person}{Rasmus~Ejlers Mogelberg}, \bibinfo{person}{Jan
%DIFDELCMD <   Schwinghammer}, {and} \bibinfo{person}{Kristian Stovring}.}
%DIFDELCMD <   \bibinfo{year}{2011}\natexlab{}.
%DIFDELCMD < \newblock \showarticletitle{First Steps in Synthetic Guarded Domain Theory:
%DIFDELCMD <   Step-Indexing in the Topos of Trees}. In \bibinfo{booktitle}{\emph{2011 IEEE
%DIFDELCMD <   26th Annual Symposium on Logic in Computer Science}}.
%DIFDELCMD <   \bibinfo{pages}{55--64}.
%DIFDELCMD < \newblock
%DIFDELCMD < \urldef\tempurl%
%DIFDELCMD < \url{https://doi.org/10.1109/LICS.2011.16}
%DIFDELCMD < \showDOI{\tempurl}
%DIFDELCMD < 

%DIFDELCMD < \bibitem[\protect\citeauthoryear{Bizjak, Birkedal, and Miculan}{Bizjak
%DIFDELCMD <   et~al\mbox{.}}{2014}]%
%DIFDELCMD <         {10.1007/978-3-319-08918-8_8}
%DIFDELCMD < \bibfield{author}{\bibinfo{person}{Ale{\v{s}} Bizjak}, \bibinfo{person}{Lars
%DIFDELCMD <   Birkedal}, {and} \bibinfo{person}{Marino Miculan}.}
%DIFDELCMD <   \bibinfo{year}{2014}\natexlab{}.
%DIFDELCMD < \newblock \showarticletitle{A Model of Countable Nondeterminism in Guarded Type
%DIFDELCMD <   Theory}. In \bibinfo{booktitle}{\emph{Rewriting and Typed Lambda Calculi}},
%DIFDELCMD <   \bibfield{editor}{\bibinfo{person}{Gilles Dowek}} (Ed.).
%DIFDELCMD <   \bibinfo{publisher}{Springer International Publishing},
%DIFDELCMD <   \bibinfo{address}{Cham}, \bibinfo{pages}{108--123}.
%DIFDELCMD < \newblock
%DIFDELCMD < \showISBNx{978-3-319-08918-8}
%DIFDELCMD < 

%DIFDELCMD < \bibitem[\protect\citeauthoryear{Capretta}{Capretta}{2005}]%
%DIFDELCMD <         {lmcs:2265}
%DIFDELCMD < \bibfield{author}{\bibinfo{person}{Venanzio Capretta}.}
%DIFDELCMD <   \bibinfo{year}{2005}\natexlab{}.
%DIFDELCMD < \newblock \showarticletitle{{General Recursion via Coinductive Types}}.
%DIFDELCMD < \newblock \bibinfo{journal}{\emph{{Logical Methods in Computer Science}}}
%DIFDELCMD <   \bibinfo{volume}{{Volume 1, Issue 2}} (\bibinfo{date}{July}
%DIFDELCMD <   \bibinfo{year}{2005}).
%DIFDELCMD < \newblock
%DIFDELCMD < \urldef\tempurl%
%DIFDELCMD < \url{https://doi.org/10.2168/LMCS-1(2:1)2005}
%DIFDELCMD < \showDOI{\tempurl}
%DIFDELCMD < 

%DIFDELCMD < \bibitem[\protect\citeauthoryear{Cimini and Siek}{Cimini and Siek}{2016}]%
%DIFDELCMD <         {cimini-siek2016}
%DIFDELCMD < \bibfield{author}{\bibinfo{person}{Matteo Cimini} {and}
%DIFDELCMD <   \bibinfo{person}{Jeremy~G. Siek}.} \bibinfo{year}{2016}\natexlab{}.
%DIFDELCMD < \newblock \showarticletitle{The Gradualizer: A Methodology and Algorithm for
%DIFDELCMD <   Generating Gradual Type Systems}. In \bibinfo{booktitle}{\emph{Proceedings of
%DIFDELCMD <   the 43rd Annual ACM SIGPLAN-SIGACT Symposium on Principles of Programming
%DIFDELCMD <   Languages}} (St. Petersburg, FL, USA) \emph{(\bibinfo{series}{POPL '16})}.
%DIFDELCMD <   \bibinfo{publisher}{Association for Computing Machinery},
%DIFDELCMD <   \bibinfo{address}{New York, NY, USA}, \bibinfo{pages}{443--455}.
%DIFDELCMD < \newblock
%DIFDELCMD < \showISBNx{9781450335492}
%DIFDELCMD < \urldef\tempurl%
%DIFDELCMD < \url{https://doi.org/10.1145/2837614.2837632}
%DIFDELCMD < \showDOI{\tempurl}
%DIFDELCMD < 

%DIFDELCMD < \bibitem[\protect\citeauthoryear{Cohen, Coquand, Huber, and M\"{o}rtberg}{Cohen
%DIFDELCMD <   et~al\mbox{.}}{2017}]%
%DIFDELCMD <         {CohenCoquandHuberMortberg2017}
%DIFDELCMD < \bibfield{author}{\bibinfo{person}{Cyril Cohen}, \bibinfo{person}{Thierry
%DIFDELCMD <   Coquand}, \bibinfo{person}{Simon Huber}, {and} \bibinfo{person}{Anders
%DIFDELCMD <   M\"{o}rtberg}.} \bibinfo{year}{2017}\natexlab{}.
%DIFDELCMD < \newblock \showarticletitle{Cubical type theory: A constructive interpretation
%DIFDELCMD <   of the univalence axiom}.
%DIFDELCMD < \newblock \bibinfo{journal}{\emph{IFCoLog J. Log. Appl.}}
%DIFDELCMD <   (\bibinfo{year}{2017}).
%DIFDELCMD < \newblock
%DIFDELCMD < 

%DIFDELCMD < \bibitem[\protect\citeauthoryear{Eremondi}{Eremondi}{2023}]%
%DIFDELCMD <         {Eremondi_2023}
%DIFDELCMD < \bibfield{author}{\bibinfo{person}{Joseph~S. Eremondi}.}
%DIFDELCMD <   \bibinfo{year}{2023}\natexlab{}.
%DIFDELCMD < \newblock \emph{\bibinfo{title}{On the design of a gradual dependently typed
%DIFDELCMD <   language for programming}}.
%DIFDELCMD < \newblock \bibinfo{thesistype}{Ph.D. Dissertation}. \bibinfo{school}{University
%DIFDELCMD <   of British Columbia}.
%DIFDELCMD < \newblock
%DIFDELCMD < \urldef\tempurl%
%DIFDELCMD < \url{https://doi.org/10.14288/1.0428823}
%DIFDELCMD < \showDOI{\tempurl}
%DIFDELCMD < 

%DIFDELCMD < \bibitem[\protect\citeauthoryear{Garcia, Clark, and Tanter}{Garcia
%DIFDELCMD <   et~al\mbox{.}}{2016}]%
%DIFDELCMD <         {garcia-clark-tanter2016}
%DIFDELCMD < \bibfield{author}{\bibinfo{person}{Ronald Garcia}, \bibinfo{person}{Alison~M.
%DIFDELCMD <   Clark}, {and} \bibinfo{person}{\'{E}ric Tanter}.}
%DIFDELCMD <   \bibinfo{year}{2016}\natexlab{}.
%DIFDELCMD < \newblock \showarticletitle{Abstracting Gradual Typing}. In
%DIFDELCMD <   \bibinfo{booktitle}{\emph{Proceedings of the 43rd Annual ACM SIGPLAN-SIGACT
%DIFDELCMD <   Symposium on Principles of Programming Languages}} (St. Petersburg, FL, USA)
%DIFDELCMD <   \emph{(\bibinfo{series}{POPL '16})}. \bibinfo{publisher}{Association for
%DIFDELCMD <   Computing Machinery}, \bibinfo{address}{New York, NY, USA},
%DIFDELCMD <   \bibinfo{pages}{429--442}.
%DIFDELCMD < \newblock
%DIFDELCMD < \showISBNx{9781450335492}
%DIFDELCMD < \urldef\tempurl%
%DIFDELCMD < \url{https://doi.org/10.1145/2837614.2837670}
%DIFDELCMD < \showDOI{\tempurl}
%DIFDELCMD < 

%DIFDELCMD < \bibitem[\protect\citeauthoryear{Lennon-Bertrand, Maillard, Tabareau, and
%DIFDELCMD <   Tanter}{Lennon-Bertrand et~al\mbox{.}}{2022}]%
%DIFDELCMD <         {10.1145/3495528}
%DIFDELCMD < \bibfield{author}{\bibinfo{person}{Meven Lennon-Bertrand},
%DIFDELCMD <   \bibinfo{person}{Kenji Maillard}, \bibinfo{person}{Nicolas Tabareau}, {and}
%DIFDELCMD <   \bibinfo{person}{\'{E}ric Tanter}.} \bibinfo{year}{2022}\natexlab{}.
%DIFDELCMD < \newblock \showarticletitle{Gradualizing the Calculus of Inductive
%DIFDELCMD <   Constructions}.
%DIFDELCMD < \newblock \bibinfo{journal}{\emph{ACM Trans. Program. Lang. Syst.}}
%DIFDELCMD <   \bibinfo{volume}{44}, \bibinfo{number}{2}, Article \bibinfo{articleno}{7}
%DIFDELCMD <   (\bibinfo{date}{apr} \bibinfo{year}{2022}), \bibinfo{numpages}{82}~pages.
%DIFDELCMD < \newblock
%DIFDELCMD < \showISSN{0164-0925}
%DIFDELCMD < \urldef\tempurl%
%DIFDELCMD < \url{https://doi.org/10.1145/3495528}
%DIFDELCMD < \showDOI{\tempurl}
%DIFDELCMD < 

%DIFDELCMD < \bibitem[\protect\citeauthoryear{Levy}{Levy}{1999}]%
%DIFDELCMD <         {levy99}
%DIFDELCMD < \bibfield{author}{\bibinfo{person}{Paul~Blain Levy}.}
%DIFDELCMD <   \bibinfo{year}{1999}\natexlab{}.
%DIFDELCMD < \newblock \showarticletitle{Call-by-Push-Value: {A} Subsuming Paradigm}. In
%DIFDELCMD <   \bibinfo{booktitle}{\emph{Typed Lambda Calculi and Applications, 4th
%DIFDELCMD <   International Conference, TLCA'99, L'Aquila, Italy, April 7-9, 1999,
%DIFDELCMD <   Proceedings}} \emph{(\bibinfo{series}{Lecture Notes in Computer Science},
%DIFDELCMD <   Vol.~\bibinfo{volume}{1581})}, \bibfield{editor}{\bibinfo{person}{Jean{-}Yves
%DIFDELCMD <   Girard}} (Ed.). \bibinfo{publisher}{Springer}, \bibinfo{pages}{228--242}.
%DIFDELCMD < \newblock
%DIFDELCMD < \urldef\tempurl%
%DIFDELCMD < \url{https://doi.org/10.1007/3-540-48959-2\_17}
%DIFDELCMD < \showDOI{\tempurl}
%DIFDELCMD < 

%DIFDELCMD < \bibitem[\protect\citeauthoryear{M\o{}gelberg and Paviotti}{M\o{}gelberg and
%DIFDELCMD <   Paviotti}{2016}]%
%DIFDELCMD <         {mogelberg-paviotti2016}
%DIFDELCMD < \bibfield{author}{\bibinfo{person}{Rasmus~Ejlers M\o{}gelberg} {and}
%DIFDELCMD <   \bibinfo{person}{Marco Paviotti}.} \bibinfo{year}{2016}\natexlab{}.
%DIFDELCMD < \newblock \showarticletitle{Denotational Semantics of Recursive Types in
%DIFDELCMD <   Synthetic Guarded Domain Theory}. In \bibinfo{booktitle}{\emph{Proceedings of
%DIFDELCMD <   the 31st Annual ACM/IEEE Symposium on Logic in Computer Science}} (New York,
%DIFDELCMD <   NY, USA) \emph{(\bibinfo{series}{LICS '16})}. \bibinfo{publisher}{Association
%DIFDELCMD <   for Computing Machinery}, \bibinfo{address}{New York, NY, USA},
%DIFDELCMD <   \bibinfo{pages}{317--326}.
%DIFDELCMD < \newblock
%DIFDELCMD < \showISBNx{9781450343916}
%DIFDELCMD < \urldef\tempurl%
%DIFDELCMD < \url{https://doi.org/10.1145/2933575.2934516}
%DIFDELCMD < \showDOI{\tempurl}
%DIFDELCMD < 

%DIFDELCMD < \bibitem[\protect\citeauthoryear{M\o{}gelberg and Veltri}{M\o{}gelberg and
%DIFDELCMD <   Veltri}{2019}]%
%DIFDELCMD <         {mogelberg-veltri2019}
%DIFDELCMD < \bibfield{author}{\bibinfo{person}{Rasmus~Ejlers M\o{}gelberg} {and}
%DIFDELCMD <   \bibinfo{person}{Niccol\`{o} Veltri}.} \bibinfo{year}{2019}\natexlab{}.
%DIFDELCMD < \newblock \showarticletitle{Bisimulation as Path Type for Guarded Recursive
%DIFDELCMD <   Types}.
%DIFDELCMD < \newblock \bibinfo{journal}{\emph{Proc. ACM Program. Lang.}}
%DIFDELCMD <   \bibinfo{volume}{3}, \bibinfo{number}{POPL}, Article \bibinfo{articleno}{4}
%DIFDELCMD <   (\bibinfo{date}{jan} \bibinfo{year}{2019}), \bibinfo{numpages}{29}~pages.
%DIFDELCMD < \newblock
%DIFDELCMD < \urldef\tempurl%
%DIFDELCMD < \url{https://doi.org/10.1145/3290317}
%DIFDELCMD < \showDOI{\tempurl}
%DIFDELCMD < 

%DIFDELCMD < \bibitem[\protect\citeauthoryear{Moggi}{Moggi}{1991}]%
%DIFDELCMD <         {MOGGI199155}
%DIFDELCMD < \bibfield{author}{\bibinfo{person}{Eugenio Moggi}.}
%DIFDELCMD <   \bibinfo{year}{1991}\natexlab{}.
%DIFDELCMD < \newblock \showarticletitle{Notions of computation and monads}.
%DIFDELCMD < \newblock \bibinfo{journal}{\emph{Information and Computation}}
%DIFDELCMD <   \bibinfo{volume}{93}, \bibinfo{number}{1} (\bibinfo{year}{1991}),
%DIFDELCMD <   \bibinfo{pages}{55--92}.
%DIFDELCMD < \newblock
%DIFDELCMD < \showISSN{0890-5401}
%DIFDELCMD < \urldef\tempurl%
%DIFDELCMD < \url{https://doi.org/10.1016/0890-5401(91)90052-4}
%DIFDELCMD < \showDOI{\tempurl}
%DIFDELCMD < \newblock
%DIFDELCMD < \shownote{Selections from 1989 IEEE Symposium on Logic in Computer Science.}
%DIFDELCMD < 

%DIFDELCMD < \bibitem[\protect\citeauthoryear{Nakano}{Nakano}{2000}]%
%DIFDELCMD <         {Nakano2000}
%DIFDELCMD < \bibfield{author}{\bibinfo{person}{H. Nakano}.}
%DIFDELCMD <   \bibinfo{year}{2000}\natexlab{}.
%DIFDELCMD < \newblock \showarticletitle{A modality for recursion}. In
%DIFDELCMD <   \bibinfo{booktitle}{\emph{Proceedings Fifteenth Annual IEEE Symposium on
%DIFDELCMD <   Logic in Computer Science (Cat. No.99CB36332)}}. \bibinfo{pages}{255--266}.
%DIFDELCMD < \newblock
%DIFDELCMD < \urldef\tempurl%
%DIFDELCMD < \url{https://doi.org/10.1109/LICS.2000.855774}
%DIFDELCMD < \showDOI{\tempurl}
%DIFDELCMD < 

%DIFDELCMD < \bibitem[\protect\citeauthoryear{Neis, Dreyer, and Rossberg}{Neis
%DIFDELCMD <   et~al\mbox{.}}{2009}]%
%DIFDELCMD <         {neis09}
%DIFDELCMD < \bibfield{author}{\bibinfo{person}{Georg Neis}, \bibinfo{person}{Derek Dreyer},
%DIFDELCMD <   {and} \bibinfo{person}{Andreas Rossberg}.} \bibinfo{year}{2009}\natexlab{}.
%DIFDELCMD < \newblock \showarticletitle{Non-Parametric Parametricity}. In
%DIFDELCMD <   \bibinfo{booktitle}{\emph{{I}nternational {C}onference on {F}unctional
%DIFDELCMD <   {P}rogramming ({ICFP})}}. \bibinfo{pages}{135--148}.
%DIFDELCMD < \newblock
%DIFDELCMD < \urldef\tempurl%
%DIFDELCMD < \url{https://doi.org/10.1145/1596550.1596572}
%DIFDELCMD < \showDOI{\tempurl}
%DIFDELCMD < 

%DIFDELCMD < \bibitem[\protect\citeauthoryear{New and Ahmed}{New and Ahmed}{2018}]%
%DIFDELCMD <         {new-ahmed2018}
%DIFDELCMD < \bibfield{author}{\bibinfo{person}{Max~S. New} {and} \bibinfo{person}{Amal
%DIFDELCMD <   Ahmed}.} \bibinfo{year}{2018}\natexlab{}.
%DIFDELCMD < \newblock \showarticletitle{Graduality from Embedding-Projection Pairs}.
%DIFDELCMD < \newblock \bibinfo{journal}{\emph{Proc. ACM Program. Lang.}}
%DIFDELCMD <   \bibinfo{volume}{2}, \bibinfo{number}{ICFP}, Article \bibinfo{articleno}{73}
%DIFDELCMD <   (\bibinfo{date}{jul} \bibinfo{year}{2018}), \bibinfo{numpages}{30}~pages.
%DIFDELCMD < \newblock
%DIFDELCMD < \urldef\tempurl%
%DIFDELCMD < \url{https://doi.org/10.1145/3236768}
%DIFDELCMD < \showDOI{\tempurl}
%DIFDELCMD < 

%DIFDELCMD < \bibitem[\protect\citeauthoryear{New, Giovannini, and Licata}{New
%DIFDELCMD <   et~al\mbox{.}}{2023}]%
%DIFDELCMD <         {new-giovannini-licata-2022}
%DIFDELCMD < \bibfield{author}{\bibinfo{person}{Max~S. New}, \bibinfo{person}{Eric
%DIFDELCMD <   Giovannini}, {and} \bibinfo{person}{Daniel~R. Licata}.}
%DIFDELCMD <   \bibinfo{year}{2023}\natexlab{}.
%DIFDELCMD < \newblock \showarticletitle{Gradual Typing for Effect Handlers}.
%DIFDELCMD < \newblock \bibinfo{journal}{\emph{Proc. {ACM} Program. Lang.}}
%DIFDELCMD <   \bibinfo{volume}{7}, \bibinfo{number}{{OOPSLA2}} (\bibinfo{year}{2023}),
%DIFDELCMD <   \bibinfo{pages}{1758--1786}.
%DIFDELCMD < \newblock
%DIFDELCMD < \urldef\tempurl%
%DIFDELCMD < \url{https://doi.org/10.1145/3622860}
%DIFDELCMD < \showDOI{\tempurl}
%DIFDELCMD < 

%DIFDELCMD < \bibitem[\protect\citeauthoryear{New, Jamner, and Ahmed}{New
%DIFDELCMD <   et~al\mbox{.}}{2019a}]%
%DIFDELCMD <         {new-jamner-ahmed19}
%DIFDELCMD < \bibfield{author}{\bibinfo{person}{Max~S. New}, \bibinfo{person}{Dustin
%DIFDELCMD <   Jamner}, {and} \bibinfo{person}{Amal Ahmed}.}
%DIFDELCMD <   \bibinfo{year}{2019}\natexlab{a}.
%DIFDELCMD < \newblock \showarticletitle{Graduality and Parametricity: Together Again for
%DIFDELCMD <   the First Time}.
%DIFDELCMD < \newblock \bibinfo{journal}{\emph{Proc. ACM Program. Lang.}}
%DIFDELCMD <   \bibinfo{volume}{4}, \bibinfo{number}{POPL}, Article \bibinfo{articleno}{46}
%DIFDELCMD <   (\bibinfo{date}{dec} \bibinfo{year}{2019}), \bibinfo{numpages}{32}~pages.
%DIFDELCMD < \newblock
%DIFDELCMD < \urldef\tempurl%
%DIFDELCMD < \url{https://doi.org/10.1145/3371114}
%DIFDELCMD < \showDOI{\tempurl}
%DIFDELCMD < 

%DIFDELCMD < \bibitem[\protect\citeauthoryear{New and Licata}{New and Licata}{2018}]%
%DIFDELCMD <         {new-licata18}
%DIFDELCMD < \bibfield{author}{\bibinfo{person}{Max~S. New} {and} \bibinfo{person}{Daniel~R.
%DIFDELCMD <   Licata}.} \bibinfo{year}{2018}\natexlab{}.
%DIFDELCMD < \newblock \showarticletitle{Call-by-name Gradual Type Theory}. In
%DIFDELCMD <   \bibinfo{booktitle}{\emph{Formal Structures for Computation and Deduction,
%DIFDELCMD <   Oxford England}}.
%DIFDELCMD < \newblock
%DIFDELCMD < \urldef\tempurl%
%DIFDELCMD < \url{https://doi.org/10.4230/LIPIcs.FSCD.2018.24}
%DIFDELCMD < \showDOI{\tempurl}
%DIFDELCMD < 

%DIFDELCMD < \bibitem[\protect\citeauthoryear{New, Licata, and Ahmed}{New
%DIFDELCMD <   et~al\mbox{.}}{2019b}]%
%DIFDELCMD <         {new-licata-ahmed2019}
%DIFDELCMD < \bibfield{author}{\bibinfo{person}{Max~S. New}, \bibinfo{person}{Daniel~R.
%DIFDELCMD <   Licata}, {and} \bibinfo{person}{Amal Ahmed}.}
%DIFDELCMD <   \bibinfo{year}{2019}\natexlab{b}.
%DIFDELCMD < \newblock \showarticletitle{Gradual Type Theory}.
%DIFDELCMD < \newblock \bibinfo{journal}{\emph{Proc. ACM Program. Lang.}}
%DIFDELCMD <   \bibinfo{volume}{3}, \bibinfo{number}{POPL}, Article \bibinfo{articleno}{15}
%DIFDELCMD <   (\bibinfo{date}{jan} \bibinfo{year}{2019}), \bibinfo{numpages}{31}~pages.
%DIFDELCMD < \newblock
%DIFDELCMD < \urldef\tempurl%
%DIFDELCMD < \url{https://doi.org/10.1145/3290328}
%DIFDELCMD < \showDOI{\tempurl}
%DIFDELCMD < 

%DIFDELCMD < \bibitem[\protect\citeauthoryear{Shulman}{Shulman}{2007}]%
%DIFDELCMD <         {shulman2007}
%DIFDELCMD < \bibfield{author}{\bibinfo{person}{Michael Shulman}.}
%DIFDELCMD <   \bibinfo{year}{2007}\natexlab{}.
%DIFDELCMD < \newblock \showarticletitle{Framed bicategories and monoidal fibrations}.
%DIFDELCMD < \newblock \bibinfo{journal}{\emph{Theory and Applications of Categories}}
%DIFDELCMD <   \bibinfo{volume}{20} (\bibinfo{date}{06} \bibinfo{year}{2007}).
%DIFDELCMD < \newblock
%DIFDELCMD < 

%DIFDELCMD < \bibitem[\protect\citeauthoryear{Siek and Chen}{Siek and Chen}{2021}]%
%DIFDELCMD <         {siek-chen2021}
%DIFDELCMD < \bibfield{author}{\bibinfo{person}{Jeremy~G. Siek} {and}
%DIFDELCMD <   \bibinfo{person}{Tianyu Chen}.} \bibinfo{year}{2021}\natexlab{}.
%DIFDELCMD < \newblock \showarticletitle{Parameterized cast calculi and reusable meta-theory
%DIFDELCMD <   for gradually typed lambda calculi}.
%DIFDELCMD < \newblock \bibinfo{journal}{\emph{J. Funct. Program.}}  \bibinfo{volume}{31}
%DIFDELCMD <   (\bibinfo{year}{2021}), \bibinfo{pages}{e30}.
%DIFDELCMD < \newblock
%DIFDELCMD < \urldef\tempurl%
%DIFDELCMD < \url{https://doi.org/10.1017/S0956796821000241}
%DIFDELCMD < \showDOI{\tempurl}
%DIFDELCMD < 

%DIFDELCMD < \bibitem[\protect\citeauthoryear{Siek and Taha}{Siek and Taha}{2006}]%
%DIFDELCMD <         {siek-taha06}
%DIFDELCMD < \bibfield{author}{\bibinfo{person}{Jeremy~G. Siek} {and} \bibinfo{person}{Walid
%DIFDELCMD <   Taha}.} \bibinfo{year}{2006}\natexlab{}.
%DIFDELCMD < \newblock \showarticletitle{Gradual Typing for Functional Languages}. In
%DIFDELCMD <   \bibinfo{booktitle}{\emph{Scheme and Functional Programming Workshop
%DIFDELCMD <   (Scheme)}}. \bibinfo{pages}{81--92}.
%DIFDELCMD < \newblock
%DIFDELCMD < 

%DIFDELCMD < \bibitem[\protect\citeauthoryear{Siek, Vitousek, Cimini, and Boyland}{Siek
%DIFDELCMD <   et~al\mbox{.}}{2015}]%
%DIFDELCMD <         {siek_et_al:LIPIcs:2015:5031}
%DIFDELCMD < \bibfield{author}{\bibinfo{person}{Jeremy~G. Siek}, \bibinfo{person}{Michael~M.
%DIFDELCMD <   Vitousek}, \bibinfo{person}{Matteo Cimini}, {and} \bibinfo{person}{John~Tang
%DIFDELCMD <   Boyland}.} \bibinfo{year}{2015}\natexlab{}.
%DIFDELCMD < \newblock \showarticletitle{{Refined Criteria for Gradual Typing}}. In
%DIFDELCMD <   \bibinfo{booktitle}{\emph{1st Summit on Advances in Programming Languages
%DIFDELCMD <   (SNAPL 2015)}} \emph{(\bibinfo{series}{Leibniz International Proceedings in
%DIFDELCMD <   Informatics (LIPIcs)}, Vol.~\bibinfo{volume}{32})},
%DIFDELCMD <   \bibfield{editor}{\bibinfo{person}{Thomas Ball}, \bibinfo{person}{Rastislav
%DIFDELCMD <   Bodik}, \bibinfo{person}{Shriram Krishnamurthi}, \bibinfo{person}{Benjamin~S.
%DIFDELCMD <   Lerner}, {and} \bibinfo{person}{Greg Morrisett}} (Eds.).
%DIFDELCMD <   \bibinfo{publisher}{Schloss Dagstuhl--Leibniz-Zentrum fuer Informatik},
%DIFDELCMD <   \bibinfo{address}{Dagstuhl, Germany}, \bibinfo{pages}{274--293}.
%DIFDELCMD < \newblock
%DIFDELCMD < \showISBNx{978-3-939897-80-4}
%DIFDELCMD < \showISSN{1868-8969}
%DIFDELCMD < \urldef\tempurl%
%DIFDELCMD < \url{https://doi.org/10.4230/LIPIcs.SNAPL.2015.274}
%DIFDELCMD < \showDOI{\tempurl}
%DIFDELCMD < 

%DIFDELCMD < \bibitem[\protect\citeauthoryear{Sterling, Gratzer, and Birkedal}{Sterling
%DIFDELCMD <   et~al\mbox{.}}{2022}]%
%DIFDELCMD <         {DBLP:journals/corr/abs-2210-02169}
%DIFDELCMD < \bibfield{author}{\bibinfo{person}{Jonathan Sterling}, \bibinfo{person}{Daniel
%DIFDELCMD <   Gratzer}, {and} \bibinfo{person}{Lars Birkedal}.}
%DIFDELCMD <   \bibinfo{year}{2022}\natexlab{}.
%DIFDELCMD < \newblock \showarticletitle{Denotational semantics of general store and
%DIFDELCMD <   polymorphism}.
%DIFDELCMD < \newblock \bibinfo{journal}{\emph{CoRR}}  \bibinfo{volume}{abs/2210.02169}
%DIFDELCMD <   (\bibinfo{year}{2022}).
%DIFDELCMD < \newblock
%DIFDELCMD < \urldef\tempurl%
%DIFDELCMD < \url{https://doi.org/10.48550/ARXIV.2210.02169}
%DIFDELCMD < \showDOI{\tempurl}
%DIFDELCMD < \showeprint[arXiv]{2210.02169}
%DIFDELCMD < 

%DIFDELCMD < \bibitem[\protect\citeauthoryear{Tobin-Hochstadt and Felleisen}{Tobin-Hochstadt
%DIFDELCMD <   and Felleisen}{2006}]%
%DIFDELCMD <         {tobin-hochstadt06}
%DIFDELCMD < \bibfield{author}{\bibinfo{person}{Sam Tobin-Hochstadt} {and}
%DIFDELCMD <   \bibinfo{person}{Matthias Felleisen}.} \bibinfo{year}{2006}\natexlab{}.
%DIFDELCMD < \newblock \showarticletitle{Interlanguage Migration: From Scripts to Programs}.
%DIFDELCMD <   In \bibinfo{booktitle}{\emph{Dynamic Languages Symposium (DLS)}}.
%DIFDELCMD <   \bibinfo{pages}{964--974}.
%DIFDELCMD < \newblock
%DIFDELCMD < 

%DIFDELCMD < \bibitem[\protect\citeauthoryear{Tobin-Hochstadt and Felleisen}{Tobin-Hochstadt
%DIFDELCMD <   and Felleisen}{2008}]%
%DIFDELCMD <         {tobin-hochstadt08}
%DIFDELCMD < \bibfield{author}{\bibinfo{person}{Sam Tobin-Hochstadt} {and}
%DIFDELCMD <   \bibinfo{person}{Matthias Felleisen}.} \bibinfo{year}{2008}\natexlab{}.
%DIFDELCMD < \newblock \showarticletitle{The Design and Implementation of Typed Scheme}. In
%DIFDELCMD <   \bibinfo{booktitle}{\emph{{ACM} {S}ymposium on {P}rinciples of {P}rogramming
%DIFDELCMD <   {L}anguages ({POPL}), San Francisco, California}}.
%DIFDELCMD < \newblock
%DIFDELCMD < 

%DIFDELCMD < \bibitem[\protect\citeauthoryear{Veltri and Vezzosi}{Veltri and
%DIFDELCMD <   Vezzosi}{2020}]%
%DIFDELCMD <         {veltri-vezzosi2020}
%DIFDELCMD < \bibfield{author}{\bibinfo{person}{Niccol\`{o} Veltri} {and}
%DIFDELCMD <   \bibinfo{person}{Andrea Vezzosi}.} \bibinfo{year}{2020}\natexlab{}.
%DIFDELCMD < \newblock \showarticletitle{Formalizing $\pi$-Calculus in Guarded Cubical
%DIFDELCMD <   Agda}. In \bibinfo{booktitle}{\emph{Proceedings of the 9th ACM SIGPLAN
%DIFDELCMD <   International Conference on Certified Programs and Proofs}} (New Orleans, LA,
%DIFDELCMD <   USA) \emph{(\bibinfo{series}{CPP 2020})}. \bibinfo{publisher}{Association for
%DIFDELCMD <   Computing Machinery}, \bibinfo{address}{New York, NY, USA},
%DIFDELCMD <   \bibinfo{pages}{270--283}.
%DIFDELCMD < \newblock
%DIFDELCMD < \showISBNx{9781450370974}
%DIFDELCMD < \urldef\tempurl%
%DIFDELCMD < \url{https://doi.org/10.1145/3372885.3373814}
%DIFDELCMD < \showDOI{\tempurl}
%DIFDELCMD < 

%DIFDELCMD < \end{thebibliography}
%DIFDELCMD < 

%DIFDELCMD < \end{DIFnomarkup}
%DIFDELCMD < %%%
\DIFdelend \DIFaddbegin \begin{DIFnomarkup}
%%% -*-BibTeX-*-
%%% Do NOT edit. File created by BibTeX with style
%%% ACM-Reference-Format-Journals [18-Jan-2012].

\begin{thebibliography}{40}

%%% ====================================================================
%%% NOTE TO THE USER: you can override these defaults by providing
%%% customized versions of any of these macros before the \bibliography
%%% command.  Each of them MUST provide its own final punctuation,
%%% except for \shownote{}, \showDOI{}, and \showURL{}.  The latter two
%%% do not use final punctuation, in order to avoid confusing it with
%%% the Web address.
%%%
%%% To suppress output of a particular field, define its macro to expand
%%% to an empty string, or better, \unskip, like this:
%%%
%%% \newcommand{\showDOI}[1]{\unskip}   % LaTeX syntax
%%%
%%% \def \showDOI #1{\unskip}           % plain TeX syntax
%%%
%%% ====================================================================

\ifx \showCODEN    \undefined \def \showCODEN     #1{\unskip}     \fi
\ifx \showDOI      \undefined \def \showDOI       #1{#1}\fi
\ifx \showISBNx    \undefined \def \showISBNx     #1{\unskip}     \fi
\ifx \showISBNxiii \undefined \def \showISBNxiii  #1{\unskip}     \fi
\ifx \showISSN     \undefined \def \showISSN      #1{\unskip}     \fi
\ifx \showLCCN     \undefined \def \showLCCN      #1{\unskip}     \fi
\ifx \shownote     \undefined \def \shownote      #1{#1}          \fi
\ifx \showarticletitle \undefined \def \showarticletitle #1{#1}   \fi
\ifx \showURL      \undefined \def \showURL       {\relax}        \fi
% The following commands are used for tagged output and should be
% invisible to TeX
\providecommand\bibfield[2]{#2}
\providecommand\bibinfo[2]{#2}
\providecommand\natexlab[1]{#1}
\providecommand\showeprint[2][]{arXiv:#2}

\bibitem[Ahmed(2006)]%
        {ahmed06}
\bibfield{author}{\bibinfo{person}{Amal~J. Ahmed}.} \bibinfo{year}{2006}\natexlab{}.
\newblock \showarticletitle{Step-Indexed Syntactic Logical Relations for Recursive and Quantified Types}. In \bibinfo{booktitle}{\emph{Programming Languages and Systems, 15th European Symposium on Programming, {ESOP} 2006, Held as Part of the Joint European Conferences on Theory and Practice of Software, {ETAPS} 2006, Vienna, Austria, March 27-28, 2006, Proceedings}} \emph{(\bibinfo{series}{Lecture Notes in Computer Science}, Vol.~\bibinfo{volume}{3924})}, \bibfield{editor}{\bibinfo{person}{Peter Sestoft}} (Ed.). \bibinfo{publisher}{Springer}, \bibinfo{pages}{69--83}.
\newblock
\urldef\tempurl%
\url{https://doi.org/10.1007/11693024\_6}
\showDOI{\tempurl}


\bibitem[Appel and McAllester(2001)]%
        {appelmcallester01}
\bibfield{author}{\bibinfo{person}{Andrew~W. Appel} {and} \bibinfo{person}{David McAllester}.} \bibinfo{year}{2001}\natexlab{}.
\newblock \showarticletitle{An indexed model of recursive types for foundational proof-carrying code}.
\newblock \bibinfo{journal}{\emph{ACM Trans. Program. Lang. Syst.}} \bibinfo{volume}{23}, \bibinfo{number}{5} (\bibinfo{date}{sep} \bibinfo{year}{2001}), \bibinfo{pages}{657–683}.
\newblock
\showISSN{0164-0925}
\urldef\tempurl%
\url{https://doi.org/10.1145/504709.504712}
\showDOI{\tempurl}


\bibitem[Atkey and McBride(2013)]%
        {atkey-mcbride2013}
\bibfield{author}{\bibinfo{person}{Robert Atkey} {and} \bibinfo{person}{Conor McBride}.} \bibinfo{year}{2013}\natexlab{}.
\newblock \showarticletitle{Productive Coprogramming with Guarded Recursion}. In \bibinfo{booktitle}{\emph{Proceedings of the 18th ACM SIGPLAN International Conference on Functional Programming}} (Boston, Massachusetts, USA) \emph{(\bibinfo{series}{ICFP '13})}. \bibinfo{publisher}{Association for Computing Machinery}, \bibinfo{address}{New York, NY, USA}, \bibinfo{pages}{197--208}.
\newblock
\showISBNx{9781450323260}
\urldef\tempurl%
\url{https://doi.org/10.1145/2500365.2500597}
\showDOI{\tempurl}


\bibitem[Bahr et~al\mbox{.}(2017)]%
        {bahr-grathwohl-bugge-mogelberg2017}
\bibfield{author}{\bibinfo{person}{Patrick Bahr}, \bibinfo{person}{Hans~Bugge Grathwohl}, {and} \bibinfo{person}{Rasmus~Ejlers Møgelberg}.} \bibinfo{year}{2017}\natexlab{}.
\newblock \showarticletitle{The clocks are ticking: No more delays!}. In \bibinfo{booktitle}{\emph{2017 32nd Annual ACM/IEEE Symposium on Logic in Computer Science (LICS)}}. \bibinfo{pages}{1--12}.
\newblock
\urldef\tempurl%
\url{https://doi.org/10.1109/LICS.2017.8005097}
\showDOI{\tempurl}


\bibitem[Baunsgaard~Kristensen et~al\mbox{.}(2022)]%
        {kristensen-mogelberg-vezzosi2022}
\bibfield{author}{\bibinfo{person}{Magnus Baunsgaard~Kristensen}, \bibinfo{person}{Rasmus~Ejlers Mogelberg}, {and} \bibinfo{person}{Andrea Vezzosi}.} \bibinfo{year}{2022}\natexlab{}.
\newblock \showarticletitle{Greatest HITs: Higher Inductive Types in Coinductive Definitions via Induction under Clocks}. In \bibinfo{booktitle}{\emph{Proceedings of the 37th Annual ACM/IEEE Symposium on Logic in Computer Science}} (Haifa, Israel) \emph{(\bibinfo{series}{LICS '22})}. \bibinfo{publisher}{Association for Computing Machinery}, \bibinfo{address}{New York, NY, USA}, Article \bibinfo{articleno}{42}, \bibinfo{numpages}{13}~pages.
\newblock
\showISBNx{9781450393515}
\urldef\tempurl%
\url{https://doi.org/10.1145/3531130.3533359}
\showDOI{\tempurl}


\bibitem[Birkedal et~al\mbox{.}(2011)]%
        {birkedal-mogelberg-schwinghammer-stovring2011}
\bibfield{author}{\bibinfo{person}{Lars Birkedal}, \bibinfo{person}{Rasmus~Ejlers Mogelberg}, \bibinfo{person}{Jan Schwinghammer}, {and} \bibinfo{person}{Kristian Stovring}.} \bibinfo{year}{2011}\natexlab{}.
\newblock \showarticletitle{First Steps in Synthetic Guarded Domain Theory: Step-Indexing in the Topos of Trees}. In \bibinfo{booktitle}{\emph{2011 IEEE 26th Annual Symposium on Logic in Computer Science}}. \bibinfo{pages}{55--64}.
\newblock
\urldef\tempurl%
\url{https://doi.org/10.1109/LICS.2011.16}
\showDOI{\tempurl}


\bibitem[Birkedal et~al\mbox{.}(2009)]%
        {Birkedal-Stovring-Thamsborg-2009}
\bibfield{author}{\bibinfo{person}{Lars Birkedal}, \bibinfo{person}{Kristian St{\o}vring}, {and} \bibinfo{person}{Jacob Thamsborg}.} \bibinfo{year}{2009}\natexlab{}.
\newblock \showarticletitle{Realizability Semantics of Parametric Polymorphism, General References, and Recursive Types}. In \bibinfo{booktitle}{\emph{Foundations of Software Science and Computational Structures}}, \bibfield{editor}{\bibinfo{person}{Luca de~Alfaro}} (Ed.). \bibinfo{publisher}{Springer Berlin Heidelberg}, \bibinfo{address}{Berlin, Heidelberg}, \bibinfo{pages}{456--470}.
\newblock
\showISBNx{978-3-642-00596-1}


\bibitem[Capretta(2005)]%
        {lmcs:2265}
\bibfield{author}{\bibinfo{person}{Venanzio Capretta}.} \bibinfo{year}{2005}\natexlab{}.
\newblock \showarticletitle{{General Recursion via Coinductive Types}}.
\newblock \bibinfo{journal}{\emph{{Logical Methods in Computer Science}}}  \bibinfo{volume}{{Volume 1, Issue 2}} (\bibinfo{date}{July} \bibinfo{year}{2005}).
\newblock
\urldef\tempurl%
\url{https://doi.org/10.2168/LMCS-1(2:1)2005}
\showDOI{\tempurl}


\bibitem[Cimini and Siek(2016)]%
        {cimini-siek2016}
\bibfield{author}{\bibinfo{person}{Matteo Cimini} {and} \bibinfo{person}{Jeremy~G. Siek}.} \bibinfo{year}{2016}\natexlab{}.
\newblock \showarticletitle{The Gradualizer: A Methodology and Algorithm for Generating Gradual Type Systems}. In \bibinfo{booktitle}{\emph{Proceedings of the 43rd Annual ACM SIGPLAN-SIGACT Symposium on Principles of Programming Languages}} (St. Petersburg, FL, USA) \emph{(\bibinfo{series}{POPL '16})}. \bibinfo{publisher}{Association for Computing Machinery}, \bibinfo{address}{New York, NY, USA}, \bibinfo{pages}{443--455}.
\newblock
\showISBNx{9781450335492}
\urldef\tempurl%
\url{https://doi.org/10.1145/2837614.2837632}
\showDOI{\tempurl}


\bibitem[Cohen et~al\mbox{.}(2017)]%
        {CohenCoquandHuberMortberg2017}
\bibfield{author}{\bibinfo{person}{Cyril Cohen}, \bibinfo{person}{Thierry Coquand}, \bibinfo{person}{Simon Huber}, {and} \bibinfo{person}{Anders M\"{o}rtberg}.} \bibinfo{year}{2017}\natexlab{}.
\newblock \showarticletitle{Cubical type theory: A constructive interpretation of the univalence axiom}.
\newblock \bibinfo{journal}{\emph{IFCoLog J. Log. Appl.}} (\bibinfo{year}{2017}).
\newblock


\bibitem[Dunphy(2002)]%
        {dunphythesis}
\bibfield{author}{\bibinfo{person}{Brian~Patrick Dunphy}.} \bibinfo{year}{2002}\natexlab{}.
\newblock \emph{\bibinfo{title}{Parametricity As a Notion of Uniformity in Reflexive Graphs}}.
\newblock \bibinfo{thesistype}{Ph.\,D. Dissertation}. \bibinfo{school}{University of Illinois at Urbana-Champaign}, \bibinfo{address}{Champaign, IL, USA}.
\newblock Advisor(s) Reddy, Uday.
\newblock


\bibitem[Eremondi(2023)]%
        {Eremondi_2023}
\bibfield{author}{\bibinfo{person}{Joseph~S. Eremondi}.} \bibinfo{year}{2023}\natexlab{}.
\newblock \emph{\bibinfo{title}{On the design of a gradual dependently typed language for programming}}.
\newblock \bibinfo{thesistype}{Ph.\,D. Dissertation}. \bibinfo{school}{University of British Columbia}.
\newblock
\urldef\tempurl%
\url{https://doi.org/10.14288/1.0428823}
\showDOI{\tempurl}


\bibitem[Frumin et~al\mbox{.}(2024)]%
        {guarded-interaction-trees}
\bibfield{author}{\bibinfo{person}{Dan Frumin}, \bibinfo{person}{Amin Timany}, {and} \bibinfo{person}{Lars Birkedal}.} \bibinfo{year}{2024}\natexlab{}.
\newblock \showarticletitle{Modular Denotational Semantics for Effects with Guarded Interaction Trees}.
\newblock \bibinfo{journal}{\emph{Proc. ACM Program. Lang.}} \bibinfo{volume}{8}, \bibinfo{number}{POPL}, Article \bibinfo{articleno}{12} (\bibinfo{date}{Jan.} \bibinfo{year}{2024}), \bibinfo{numpages}{30}~pages.
\newblock
\urldef\tempurl%
\url{https://doi.org/10.1145/3632854}
\showDOI{\tempurl}


\bibitem[Garcia et~al\mbox{.}(2016)]%
        {garcia-clark-tanter2016}
\bibfield{author}{\bibinfo{person}{Ronald Garcia}, \bibinfo{person}{Alison~M. Clark}, {and} \bibinfo{person}{\'{E}ric Tanter}.} \bibinfo{year}{2016}\natexlab{}.
\newblock \showarticletitle{Abstracting Gradual Typing}. In \bibinfo{booktitle}{\emph{Proceedings of the 43rd Annual ACM SIGPLAN-SIGACT Symposium on Principles of Programming Languages}} (St. Petersburg, FL, USA) \emph{(\bibinfo{series}{POPL '16})}. \bibinfo{publisher}{Association for Computing Machinery}, \bibinfo{address}{New York, NY, USA}, \bibinfo{pages}{429--442}.
\newblock
\showISBNx{9781450335492}
\urldef\tempurl%
\url{https://doi.org/10.1145/2837614.2837670}
\showDOI{\tempurl}


\bibitem[Greenman(2022)]%
        {deepandshallowtypes}
\bibfield{author}{\bibinfo{person}{Ben Greenman}.} \bibinfo{year}{2022}\natexlab{}.
\newblock \showarticletitle{Deep and shallow types for gradual languages}. In \bibinfo{booktitle}{\emph{Proceedings of the 43rd ACM SIGPLAN International Conference on Programming Language Design and Implementation}} (San Diego, CA, USA) \emph{(\bibinfo{series}{PLDI 2022})}. \bibinfo{publisher}{Association for Computing Machinery}, \bibinfo{address}{New York, NY, USA}, \bibinfo{pages}{580–593}.
\newblock
\showISBNx{9781450392655}
\urldef\tempurl%
\url{https://doi.org/10.1145/3519939.3523430}
\showURL{%
\tempurl}


\bibitem[Herman et~al\mbox{.}(2010)]%
        {herman-tomb-flanagan-2010}
\bibfield{author}{\bibinfo{person}{David Herman}, \bibinfo{person}{Aaron Tomb}, {and} \bibinfo{person}{Cormac Flanagan}.} \bibinfo{year}{2010}\natexlab{}.
\newblock \showarticletitle{Space-efficient gradual typing}.
\newblock \bibinfo{journal}{\emph{Higher Order Symbol. Comput.}} \bibinfo{volume}{23}, \bibinfo{number}{2} (\bibinfo{date}{June} \bibinfo{year}{2010}), \bibinfo{pages}{167–189}.
\newblock
\showISSN{1388-3690}
\urldef\tempurl%
\url{https://doi.org/10.1007/s10990-011-9066-z}
\showDOI{\tempurl}


\bibitem[Lennon-Bertrand et~al\mbox{.}(2022)]%
        {gradualizing-cic}
\bibfield{author}{\bibinfo{person}{Meven Lennon-Bertrand}, \bibinfo{person}{Kenji Maillard}, \bibinfo{person}{Nicolas Tabareau}, {and} \bibinfo{person}{\'{E}ric Tanter}.} \bibinfo{year}{2022}\natexlab{}.
\newblock \showarticletitle{Gradualizing the Calculus of Inductive Constructions}.
\newblock \bibinfo{journal}{\emph{ACM Trans. Program. Lang. Syst.}} \bibinfo{volume}{44}, \bibinfo{number}{2}, Article \bibinfo{articleno}{7} (\bibinfo{date}{apr} \bibinfo{year}{2022}), \bibinfo{numpages}{82}~pages.
\newblock
\showISSN{0164-0925}
\urldef\tempurl%
\url{https://doi.org/10.1145/3495528}
\showDOI{\tempurl}


\bibitem[Levy(1999)]%
        {levy99}
\bibfield{author}{\bibinfo{person}{Paul~Blain Levy}.} \bibinfo{year}{1999}\natexlab{}.
\newblock \showarticletitle{Call-by-Push-Value: {A} Subsuming Paradigm}. In \bibinfo{booktitle}{\emph{Typed Lambda Calculi and Applications, 4th International Conference, TLCA'99, L'Aquila, Italy, April 7-9, 1999, Proceedings}} \emph{(\bibinfo{series}{Lecture Notes in Computer Science}, Vol.~\bibinfo{volume}{1581})}, \bibfield{editor}{\bibinfo{person}{Jean{-}Yves Girard}} (Ed.). \bibinfo{publisher}{Springer}, \bibinfo{pages}{228--242}.
\newblock
\urldef\tempurl%
\url{https://doi.org/10.1007/3-540-48959-2\_17}
\showDOI{\tempurl}


\bibitem[Levy(2003)]%
        {levystacks}
\bibfield{author}{\bibinfo{person}{Paul~Blain Levy}.} \bibinfo{year}{2003}\natexlab{}.
\newblock \showarticletitle{Adjunction Models For Call-By-Push-Value With Stacks}.
\newblock \bibinfo{journal}{\emph{Electronic Notes in Theoretical Computer Science}}  \bibinfo{volume}{69} (\bibinfo{year}{2003}), \bibinfo{pages}{248--271}.
\newblock
\showISSN{1571-0661}
\urldef\tempurl%
\url{https://doi.org/10.1016/S1571-0661(04)80568-1}
\showDOI{\tempurl}
\newblock
\shownote{CTCS'02, Category Theory and Computer Science}.


\bibitem[M\o{}gelberg and Paviotti(2016)]%
        {mogelberg-paviotti2016}
\bibfield{author}{\bibinfo{person}{Rasmus~Ejlers M\o{}gelberg} {and} \bibinfo{person}{Marco Paviotti}.} \bibinfo{year}{2016}\natexlab{}.
\newblock \showarticletitle{Denotational Semantics of Recursive Types in Synthetic Guarded Domain Theory}. In \bibinfo{booktitle}{\emph{Proceedings of the 31st Annual ACM/IEEE Symposium on Logic in Computer Science}} (New York, NY, USA) \emph{(\bibinfo{series}{LICS '16})}. \bibinfo{publisher}{Association for Computing Machinery}, \bibinfo{address}{New York, NY, USA}, \bibinfo{pages}{317--326}.
\newblock
\showISBNx{9781450343916}
\urldef\tempurl%
\url{https://doi.org/10.1145/2933575.2934516}
\showDOI{\tempurl}


\bibitem[M\o{}gelberg and Veltri(2019)]%
        {mogelberg-veltri2019}
\bibfield{author}{\bibinfo{person}{Rasmus~Ejlers M\o{}gelberg} {and} \bibinfo{person}{Niccol\`{o} Veltri}.} \bibinfo{year}{2019}\natexlab{}.
\newblock \showarticletitle{Bisimulation as Path Type for Guarded Recursive Types}.
\newblock \bibinfo{journal}{\emph{Proc. ACM Program. Lang.}} \bibinfo{volume}{3}, \bibinfo{number}{POPL}, Article \bibinfo{articleno}{4} (\bibinfo{date}{jan} \bibinfo{year}{2019}), \bibinfo{numpages}{29}~pages.
\newblock
\urldef\tempurl%
\url{https://doi.org/10.1145/3290317}
\showDOI{\tempurl}


\bibitem[Moggi(1989)]%
        {39155}
\bibfield{author}{\bibinfo{person}{E. Moggi}.} \bibinfo{year}{1989}\natexlab{}.
\newblock \showarticletitle{Computational lambda-calculus and monads}. In \bibinfo{booktitle}{\emph{[1989] Proceedings. Fourth Annual Symposium on Logic in Computer Science}}. \bibinfo{pages}{14--23}.
\newblock
\urldef\tempurl%
\url{https://doi.org/10.1109/LICS.1989.39155}
\showDOI{\tempurl}


\bibitem[Moggi(1991)]%
        {MOGGI199155}
\bibfield{author}{\bibinfo{person}{Eugenio Moggi}.} \bibinfo{year}{1991}\natexlab{}.
\newblock \showarticletitle{Notions of computation and monads}.
\newblock \bibinfo{journal}{\emph{Information and Computation}} \bibinfo{volume}{93}, \bibinfo{number}{1} (\bibinfo{year}{1991}), \bibinfo{pages}{55--92}.
\newblock
\showISSN{0890-5401}
\urldef\tempurl%
\url{https://doi.org/10.1016/0890-5401(91)90052-4}
\showDOI{\tempurl}
\newblock
\shownote{Selections from 1989 IEEE Symposium on Logic in Computer Science}.


\bibitem[Nakano(2000)]%
        {Nakano2000}
\bibfield{author}{\bibinfo{person}{H. Nakano}.} \bibinfo{year}{2000}\natexlab{}.
\newblock \showarticletitle{A modality for recursion}. In \bibinfo{booktitle}{\emph{Proceedings Fifteenth Annual IEEE Symposium on Logic in Computer Science (Cat. No.99CB36332)}}. \bibinfo{pages}{255--266}.
\newblock
\urldef\tempurl%
\url{https://doi.org/10.1109/LICS.2000.855774}
\showDOI{\tempurl}


\bibitem[Neis et~al\mbox{.}(2009)]%
        {neis09}
\bibfield{author}{\bibinfo{person}{Georg Neis}, \bibinfo{person}{Derek Dreyer}, {and} \bibinfo{person}{Andreas Rossberg}.} \bibinfo{year}{2009}\natexlab{}.
\newblock \showarticletitle{Non-Parametric Parametricity}. In \bibinfo{booktitle}{\emph{{I}nternational {C}onference on {F}unctional {P}rogramming ({ICFP})}}. \bibinfo{pages}{135--148}.
\newblock
\urldef\tempurl%
\url{https://doi.org/10.1145/1596550.1596572}
\showDOI{\tempurl}


\bibitem[New and Ahmed(2018)]%
        {new-ahmed2018}
\bibfield{author}{\bibinfo{person}{Max~S. New} {and} \bibinfo{person}{Amal Ahmed}.} \bibinfo{year}{2018}\natexlab{}.
\newblock \showarticletitle{Graduality from Embedding-Projection Pairs}.
\newblock \bibinfo{journal}{\emph{Proc. ACM Program. Lang.}} \bibinfo{volume}{2}, \bibinfo{number}{ICFP}, Article \bibinfo{articleno}{73} (\bibinfo{date}{jul} \bibinfo{year}{2018}), \bibinfo{numpages}{30}~pages.
\newblock
\urldef\tempurl%
\url{https://doi.org/10.1145/3236768}
\showDOI{\tempurl}


\bibitem[New et~al\mbox{.}(2023)]%
        {new-giovannini-licata-2022}
\bibfield{author}{\bibinfo{person}{Max~S. New}, \bibinfo{person}{Eric Giovannini}, {and} \bibinfo{person}{Daniel~R. Licata}.} \bibinfo{year}{2023}\natexlab{}.
\newblock \showarticletitle{Gradual Typing for Effect Handlers}.
\newblock \bibinfo{journal}{\emph{Proc. {ACM} Program. Lang.}} \bibinfo{volume}{7}, \bibinfo{number}{{OOPSLA2}} (\bibinfo{year}{2023}), \bibinfo{pages}{1758--1786}.
\newblock
\urldef\tempurl%
\url{https://doi.org/10.1145/3622860}
\showDOI{\tempurl}


\bibitem[New et~al\mbox{.}(2019a)]%
        {new-jamner-ahmed19}
\bibfield{author}{\bibinfo{person}{Max~S. New}, \bibinfo{person}{Dustin Jamner}, {and} \bibinfo{person}{Amal Ahmed}.} \bibinfo{year}{2019}\natexlab{a}.
\newblock \showarticletitle{Graduality and Parametricity: Together Again for the First Time}.
\newblock \bibinfo{journal}{\emph{Proc. ACM Program. Lang.}} \bibinfo{volume}{4}, \bibinfo{number}{POPL}, Article \bibinfo{articleno}{46} (\bibinfo{date}{dec} \bibinfo{year}{2019}), \bibinfo{numpages}{32}~pages.
\newblock
\urldef\tempurl%
\url{https://doi.org/10.1145/3371114}
\showDOI{\tempurl}


\bibitem[New and Licata(2018)]%
        {new-licata18}
\bibfield{author}{\bibinfo{person}{Max~S. New} {and} \bibinfo{person}{Daniel~R. Licata}.} \bibinfo{year}{2018}\natexlab{}.
\newblock \showarticletitle{Call-by-name Gradual Type Theory}. In \bibinfo{booktitle}{\emph{Formal Structures for Computation and Deduction, Oxford England}}.
\newblock
\urldef\tempurl%
\url{https://doi.org/10.4230/LIPIcs.FSCD.2018.24}
\showDOI{\tempurl}


\bibitem[New et~al\mbox{.}(2019b)]%
        {new-licata-ahmed2019}
\bibfield{author}{\bibinfo{person}{Max~S. New}, \bibinfo{person}{Daniel~R. Licata}, {and} \bibinfo{person}{Amal Ahmed}.} \bibinfo{year}{2019}\natexlab{b}.
\newblock \showarticletitle{Gradual Type Theory}.
\newblock \bibinfo{journal}{\emph{Proc. ACM Program. Lang.}} \bibinfo{volume}{3}, \bibinfo{number}{POPL}, Article \bibinfo{articleno}{15} (\bibinfo{date}{jan} \bibinfo{year}{2019}), \bibinfo{numpages}{31}~pages.
\newblock
\urldef\tempurl%
\url{https://doi.org/10.1145/3290328}
\showDOI{\tempurl}


\bibitem[Shulman(2007)]%
        {shulman2007}
\bibfield{author}{\bibinfo{person}{Michael Shulman}.} \bibinfo{year}{2007}\natexlab{}.
\newblock \showarticletitle{Framed bicategories and monoidal fibrations}.
\newblock \bibinfo{journal}{\emph{Theory and Applications of Categories}}  \bibinfo{volume}{20} (\bibinfo{date}{06} \bibinfo{year}{2007}).
\newblock


\bibitem[Siek and Chen(2021)]%
        {siek-chen2021}
\bibfield{author}{\bibinfo{person}{Jeremy~G. Siek} {and} \bibinfo{person}{Tianyu Chen}.} \bibinfo{year}{2021}\natexlab{}.
\newblock \showarticletitle{Parameterized cast calculi and reusable meta-theory for gradually typed lambda calculi}.
\newblock \bibinfo{journal}{\emph{J. Funct. Program.}}  \bibinfo{volume}{31} (\bibinfo{year}{2021}), \bibinfo{pages}{e30}.
\newblock
\urldef\tempurl%
\url{https://doi.org/10.1017/S0956796821000241}
\showDOI{\tempurl}


\bibitem[Siek and Taha(2006)]%
        {siek-taha06}
\bibfield{author}{\bibinfo{person}{Jeremy~G. Siek} {and} \bibinfo{person}{Walid Taha}.} \bibinfo{year}{2006}\natexlab{}.
\newblock \showarticletitle{Gradual Typing for Functional Languages}. In \bibinfo{booktitle}{\emph{Scheme and Functional Programming Workshop (Scheme)}}. \bibinfo{pages}{81--92}.
\newblock


\bibitem[Siek et~al\mbox{.}(2015)]%
        {siek_et_al:LIPIcs:2015:5031}
\bibfield{author}{\bibinfo{person}{Jeremy~G. Siek}, \bibinfo{person}{Michael~M. Vitousek}, \bibinfo{person}{Matteo Cimini}, {and} \bibinfo{person}{John~Tang Boyland}.} \bibinfo{year}{2015}\natexlab{}.
\newblock \showarticletitle{{Refined Criteria for Gradual Typing}}. In \bibinfo{booktitle}{\emph{1st Summit on Advances in Programming Languages (SNAPL 2015)}} \emph{(\bibinfo{series}{Leibniz International Proceedings in Informatics (LIPIcs)}, Vol.~\bibinfo{volume}{32})}, \bibfield{editor}{\bibinfo{person}{Thomas Ball}, \bibinfo{person}{Rastislav Bodik}, \bibinfo{person}{Shriram Krishnamurthi}, \bibinfo{person}{Benjamin~S. Lerner}, {and} \bibinfo{person}{Greg Morrisett}} (Eds.). \bibinfo{publisher}{Schloss Dagstuhl--Leibniz-Zentrum fuer Informatik}, \bibinfo{address}{Dagstuhl, Germany}, \bibinfo{pages}{274--293}.
\newblock
\showISBNx{978-3-939897-80-4}
\showISSN{1868-8969}
\urldef\tempurl%
\url{https://doi.org/10.4230/LIPIcs.SNAPL.2015.274}
\showDOI{\tempurl}


\bibitem[Stassen et~al\mbox{.}(2024)]%
        {probabilistic-fpc-guarded}
\bibfield{author}{\bibinfo{person}{Philipp Jan~Andries Stassen}, \bibinfo{person}{Rasmus~Ejlers Møgelberg}, \bibinfo{person}{Maaike Zwart}, \bibinfo{person}{Alejandro Aguirre}, {and} \bibinfo{person}{Lars Birkedal}.} \bibinfo{year}{2024}\natexlab{}.
\newblock \bibinfo{title}{Modelling Probabilistic FPC in Guarded Type Theory}.
\newblock
\newblock
\showeprint[arxiv]{2408.04455}~[cs.PL]
\urldef\tempurl%
\url{https://arxiv.org/abs/2408.04455}
\showURL{%
\tempurl}


\bibitem[Sterling et~al\mbox{.}(2022)]%
        {DBLP:journals/corr/abs-2210-02169}
\bibfield{author}{\bibinfo{person}{Jonathan Sterling}, \bibinfo{person}{Daniel Gratzer}, {and} \bibinfo{person}{Lars Birkedal}.} \bibinfo{year}{2022}\natexlab{}.
\newblock \showarticletitle{Denotational semantics of general store and polymorphism}.
\newblock \bibinfo{journal}{\emph{CoRR}}  \bibinfo{volume}{abs/2210.02169} (\bibinfo{year}{2022}).
\newblock
\urldef\tempurl%
\url{https://doi.org/10.48550/ARXIV.2210.02169}
\showDOI{\tempurl}
\showeprint[arXiv]{2210.02169}


\bibitem[Tobin-Hochstadt and Felleisen(2006)]%
        {tobin-hochstadt06}
\bibfield{author}{\bibinfo{person}{Sam Tobin-Hochstadt} {and} \bibinfo{person}{Matthias Felleisen}.} \bibinfo{year}{2006}\natexlab{}.
\newblock \showarticletitle{Interlanguage Migration: From Scripts to Programs}. In \bibinfo{booktitle}{\emph{Dynamic Languages Symposium (DLS)}}. \bibinfo{pages}{964--974}.
\newblock


\bibitem[Tobin-Hochstadt and Felleisen(2008)]%
        {tobin-hochstadt08}
\bibfield{author}{\bibinfo{person}{Sam Tobin-Hochstadt} {and} \bibinfo{person}{Matthias Felleisen}.} \bibinfo{year}{2008}\natexlab{}.
\newblock \showarticletitle{The Design and Implementation of Typed Scheme}. In \bibinfo{booktitle}{\emph{{ACM} {S}ymposium on {P}rinciples of {P}rogramming {L}anguages ({POPL}), San Francisco, California}}.
\newblock


\bibitem[Veltri and Vezzosi(2020)]%
        {veltri-vezzosi2020}
\bibfield{author}{\bibinfo{person}{Niccol\`{o} Veltri} {and} \bibinfo{person}{Andrea Vezzosi}.} \bibinfo{year}{2020}\natexlab{}.
\newblock \showarticletitle{Formalizing $\pi$-Calculus in Guarded Cubical Agda}. In \bibinfo{booktitle}{\emph{Proceedings of the 9th ACM SIGPLAN International Conference on Certified Programs and Proofs}} (New Orleans, LA, USA) \emph{(\bibinfo{series}{CPP 2020})}. \bibinfo{publisher}{Association for Computing Machinery}, \bibinfo{address}{New York, NY, USA}, \bibinfo{pages}{270--283}.
\newblock
\showISBNx{9781450370974}
\urldef\tempurl%
\url{https://doi.org/10.1145/3372885.3373814}
\showDOI{\tempurl}


\bibitem[Vitousek et~al\mbox{.}(2017)]%
        {transient}
\bibfield{author}{\bibinfo{person}{Michael~M. Vitousek}, \bibinfo{person}{Cameron Swords}, {and} \bibinfo{person}{Jeremy~G. Siek}.} \bibinfo{year}{2017}\natexlab{}.
\newblock \showarticletitle{Big types in little runtime: open-world soundness and collaborative blame for gradual type systems}. In \bibinfo{booktitle}{\emph{Proceedings of the 44th ACM SIGPLAN Symposium on Principles of Programming Languages}} (Paris, France) \emph{(\bibinfo{series}{POPL '17})}. \bibinfo{publisher}{Association for Computing Machinery}, \bibinfo{address}{New York, NY, USA}, \bibinfo{pages}{762–774}.
\newblock
\showISBNx{9781450346603}
\urldef\tempurl%
\url{https://doi.org/10.1145/3009837.3009849}
\showDOI{\tempurl}


\end{thebibliography}

\end{DIFnomarkup}

\newpage
\DIFaddend \appendix
\section{Gradual Typing Syntax}\DIFaddbegin \label{sec:appendix-gtlc-syntax}
\DIFaddend 

\newcommand{\dynof}[1]{\textrm{dyn}(#1)}
\begin{theorem}
  For every $A$, there is a derivation $\dynof A : A \ltdyn D$
\end{theorem}
\begin{proof}
  By induction on $A$:
  \begin{enumerate}
  \item $A = \dyn$, then $r(\dyn) : \dyn \ltdyn \dyn$
  \item $A = \nat$, then $\inat : \nat \ltdyn \dyn$
  \item $A = A_1 \ra A_2$ then $(\dynof {A_1} \ra \dynof {A_2})\iarr : A_1 \ra A_2 \ltdyn \dyn$ 
  \item $A = A_1 \times A_2$ then $(\dynof {A_1} \times \dynof{A_2})\itimes : A_1 \times A_2 \ltdyn \dyn$ 
  \end{enumerate}
\end{proof}

\begin{theorem}
  \label{thm:thin}
  For any two derivations $c,c' : A \ltdyn A'$ of the same precision
  $c \equiv c'$
\end{theorem}
\begin{proof}
  \begin{enumerate}
  \item We show this by showing that derivations have a canonical
    form.

    The following presentation of precision derivations has unique derivations
    \begin{mathpar}
      \inferrule{}{\textrm{refl}(D) : D \ltdyn D}\and
      \inferrule{}{\textsf{Inj}_{\text{nat}} : \nat \ltdyn D}\and
      \inferrule{}{\textrm{refl}(\nat) : \nat \ltdyn \nat}\and
      \inferrule{c : A_i \ra A_o \ltdyn D\ra D}{c(\textsf{Inj}_{\text{arr}}) : A_i \ra A_o \ltdyn \nat}\and
      \inferrule{c : A_i \ltdyn A_i' \and d : A_o \ltdyn A_o'}{c \ra d : A_i \ra A_o \ltdyn A_i'\ra A_o'}
      %% \inferrule{A_1 \times A_2 \ltdyn D\times D}{A_1 \times A_2 \ltdyn \nat}
      %% \inferrule{A_1 \ltdyn A_1' \and A_2 \ltdyn A_2'}{A_1 \times A_2 \ltdyn A_1'\times A_2'}
    \end{mathpar}
    Since it satisfies reflexivity, cut-elimination and congruence, it
    is a model of the original theory. Since it is a sub-theory of the
    original theory, it is equivalent.
  \end{enumerate}
\end{proof}

As mentioned in Section~\ref{sec:GTLC}, the cast calculus
separates gradual type casts into downcasts and upcasts, whereas many
gradual calculi are built out of a single notion of cast
\[ \inferrule{\Gamma \vdash M : A}{\Gamma \vdash (M :: A') : A' }\]
The term precision rules for this style of calculus are as follows\cite{siek_et_al:LIPIcs:2015:5031}:
\begin{mathpar}
  \inferrule*[Right=CastRight]
  {\Delta \vdash M \ltdyn N_1 : c_1 \and
    c_{1} : A \ltdyn A_1\and
    c_2 : A \ltdyn A_2
  }
  {\Delta \vdash M \ltdyn (N :: A_2) : c_2}

  \inferrule*[Right=CastLeft]
  {\Delta \vdash M_1 \ltdyn N : c_1 \and
    c_1 : A_1 \ltdyn A\and
    c_2 : A_2 \ltdyn A
  }
  {\Delta \vdash (M_1 :: A_2) \ltdyn N : c_2}
\end{mathpar}
New-Ahmed and New-Licata showed that such if the cast $M_1 :: A_2$ for $M_1 : A_1$ is interpreted as
\[ \dnc {(\dynof {A_2})}{\upc {(\dynof{A_1})} M} \]
then the CastRight and CastLeft rules are derivable from the UpL/UpR/DnL/DnR rules
and the \emph{retraction} property of casts: that a downcast after an upcast
is equal to the identity.
%
In the intensional setting, this retraction property needs to be
weakened: a downcast after an upcast is merely \emph{weakly bisimilar}
to the identity.
%
But this weak bisimilarity is no longer strong enough to show the validity
of CastRight and CastLeft, due to the lack of transitivity of weak
bisimilarity.

Fortunately there is an alternative argument that makes CastRight and
CastLeft valid if we change slightly the interpretation of casts. To
show this alternative translation we need to first introduce the
\emph{least upper bound} of types:
\begin{lemma}
  The preordered set of types with type precision as ordering has all
  binary least upper bounds defined as follows:
  \begin{align*}
    (A_1 \times A_2) \sqcup (A_1' \times A_2') &= (A_1 \sqcup A_1') \times (A_2 \sqcup A_2')\\
    (A_1 \ra A_2) \sqcup (A_1' \ra A_2') &= (A_1 \ra A_1') \times (A_2 \ra A_2')\\
    \nat \sqcup \nat &= \nat\\
    A \sqcup A' &= \dyn \text{ otherwise }
  \end{align*}
\end{lemma}

Then for $M_1 : A_1$ we translate $M_1 :: A_2$ to
\[ \dnc {c^\sqcup_2}\upc{c^\sqcup_1} M \]
where $c^\sqcup_1 : A_1 \ltdyn A_1\sqcup A_2$ and $c^\sqcup_2 : A_2 \ltdyn
A_1\sqcup A_2$. In the presence of the strong retraction property, New
and Ahmed showed that this is equal to the prior interpretation. With
a weak retraction property, it is still weakly bisimilar to the prior
interpretation, so for the extensional properties we care about
ultimately this semantics is equivalent. However our compositional
arguments do care about this difference, and so we use instead this
translation that avoids the need for any retraction:
\begin{theorem}
  If $M_1 :: A_2$ is interpreted as $\dnc {c^\sqcup_2}\upc{c^\sqcup_1} M$,
  then CastRight and CastLeft are admissible.
\end{theorem}
\begin{proof}
  We show CastLeft, as this is the one that required retraction in
  prior work. CastRight follows similarly.

  Assume $M_1 \ltdyn N : c_1$, we seek to prove that
  \[ \dnc {c^\sqcup_2}\upc{c^\sqcup_1} M \ltdyn N : c_2 \]
  %
  Since $A \sqcup A_2$ is the
  least upper bound, we have there must be a derivation $c^{\sqcup} :
  A_1\sqcup A_2 \ltdyn A$.
  %
  Further by Theorem~\ref{thm:thin}, we have
  $c^\sqcup_2c^{\sqcup} \equiv c_2$, so it suffices to show
  \[M_1 \ltdyn \dnc {c^\sqcup_2}\upc{c^\sqcup} N : c^\sqcup_2c^{\sqcup} \]
  %
  Then applying DnL and UpL it is sufficient to show
  \[ M_1 \ltdyn N : c^\sqcup_1c^\sqcup \]
  But again by Theorem~\ref{thm:thin} we have $c^\sqcup_1c^\sqcup
  \equiv c_1$ so this follows by our assumption that $M_1 \ltdyn N
  :c_1$.
\end{proof}

\section{Kleisli Actions}\label{sec:kleisli-actions}

We now define the Kleisli actions of the type constructors. These are the following
\begin{enumerate}
\item $- \tok B : \op\PreDom_k \to \ErrDom_k$
\item $A \tok - : \ErrDom_k \to \ErrDom_k$
\item $- \timesk A_2 : \PreDom_k \to \PreDom_k$
\item $A_1 \timesk - : \PreDom_k \to \PreDom_k$
\end{enumerate}
While these can be defined in any CBPV model, for simplicity we just
give their definition explicitly for predomains and error domains:
\begin{definition}{~}
  Given $\phi : \li A_i \multimap \li A_o$, we define $\phi \tok B : U(A_o \to B) \to U(A_i \to B)$ as
  \[ (\phi \tok B)(f)(x) = f^\dagger(\phi(\eta(x))) \]

  This is functorial in that $\id \tok B = \id$ and $(\phi \circ
  \phi') \tok B = (\phi' \tok B) \circ (\phi \tok B)$.

  Further, this preserves squares in that if $\phi \ltsq{\li c_i}{\li c_o} \phi'$ then $\phi \tok B \ltsq{U(c_o \to d)}{U(c_i \to d)} \phi' \tok B'$.
\end{definition}
\begin{proof}
  {~}
  \begin{itemize}
  \item For identity
    \[ (\id \tok B)(f)(x) = f^\dagger(\eta(x)) = f(x) \]

  \item 
    For composition, expanding definitions it suffices to show
    \[ f^\dagger(\phi(\phi'(\eta(x)))) = ((\phi \tok B)(f))^\dagger(\phi'(\eta(x))) \]
    Which follows if we can show
    \[ ((\phi \tok B)(f))^\dagger = f^\dagger \circ \phi \]
    By the freeness of $\li -$ it suffices to show it for inputs of the form $\eta y$:
    \[ ((\phi \tok B)(f))^\dagger(\eta y) = ((\phi \tok B)(f))(y) = f^\dagger(\phi(\eta y))\]
  \item For squares, we assume $f \binrel{U(c_o\to d)} f'$ and $x
    \binrel {c_i} x'$ and we need to show $f^\dagger(\phi(x)) \binrel
    d (f')^\dagger(\phi'(x))$. This follows easily since $f^\dagger
    \ltsq{\li c_o}{d} f^{\prime\dagger}$.
  \end{itemize}
\end{proof}

\begin{definition}
  Given $f : UB \to UB'$, define $A \tok f : U(A \to B) \to U(A \to B')$ as
    \[ (A \tok f)(g)(x) = f(g(x)) \]

  This is functorial in that $A \tok \id = \id$ and $A \tok (f \circ
  f') = (A \tok f) \circ (A \tok f')$.

  Further this preserves squares in that if $f \ltsq{Ud_i}{Ud_o} f'$
  then $A \tok f \ltsq{U(c \to d_i)}{U(c \to d_o)} A' \tok f'$
\end{definition}
\begin{proof}
  \begin{itemize}
  \item For identity
    \[ (A \tok \id)(g)(x) = g(x) \]
  \item 
    for composition
    \[  (A \tok (f \circ f'))(g)(x) = f(f'(g(x))) = (A\tok f)((A \tok f')(g))(x) \]
  \item 
   for squares, given $g \binrel{U(c \to d_i)} g'$ and $x \binrel{c}
   x'$ we have $f(g(x)) \binrel{d_o} f'(g'(x'))$.
  \end{itemize}
\end{proof}

\begin{definition}[Kleisli Actions of $\times$]
  Given $\phi_1 : \li A_1 \multimap \li A_1'$ we define $\phi_1 \timesk A_2 : \li (A_1 \times A_2) \multimap (A_1' \times A_2)$ as the unique extension of the following to a homomorphism:
  \[ (\phi_1 \timesk A_2)(\eta (x_1,x_2)) = (\lambda x_1'. \eta(x_1',x_2))^\dagger(\phi_1(\eta(x_1))) \]
  and similarly define
  \[ (A_1 \timesk \phi_2)(\eta (x_1,x_2)) = (\lambda x_2'. \eta(x_1,x_2'))^\dagger(\phi_2(\eta(x_2)))\]
\end{definition}
\begin{proof}
  We show cases only for $-\timesk A_2$ as the other is symmetric.
  \begin{itemize}
  \item For identity it is sufficient to consider inputs of the form $\eta(x_1,x_2)$, then indeed
    \[ (\id \timesk A_2)(\eta(x_1,x_2)) = (\lambda x_1'. \eta(x_1',x_2))^\dagger(\eta(x_1)) = \eta(x_1,x_2) \]
  \item For composition, again for pure inputs this reduces to
    \[ (\lambda x_1'. \eta(x_1',x_2))^\dagger(\phi_1'(\phi_1(\eta(x_1))))
    = (\phi_1'\timesk A_2)((\lambda x_1'. \eta(x_1',x_2))^\dagger(\phi_1(\eta(x_1)))) \]
    So it suffices to show
    \[ (\lambda x_1'. \eta(x_1',x_2))^\dagger \circ \phi_1' = (\phi_1'\timesk A_2) \circ (\lambda x_1'. \eta(x_1',x_2))^\dagger \]
    Since both sides are homomorphisms out of the free error domain by freeness it suffices to show they are equal for inputs of the form $\eta(x_1')$:
    \[  (\lambda x_1'. \eta(x_1',x_2))^\dagger(\phi_1'(\eta(x_1')))
    = (\phi_1'\timesk A_2)(\eta(x_1',x_2))
    = (\phi_1'\timesk A_2)((\lambda x_1'. \eta(x_1',x_2))^\dagger(\eta(x_1')))
    \]

  \item Finally, for squares, we need to show if $\phi_l \ltsq{\li
    c_1}{\li c_1'} \phi_r$ then $\phi_l \timesk A_{2l} \ltsq{\li (c_1
    \times c_2)}{\li (c_1' \times c_2)} \phi_r \timesk A_{2r}$.

    By the universal property of $\li-$ on relations, it suffices to
    assume $x_{1l} \binrel{c_1} x_{1r}$ and $x_{2l} \binrel{c_2} x_{2r}$ and prove
    \[ (\phi_l \timesk A_{2l})(\eta(x_{1l},x_{2l})) \binrel{(\li (c_1' \times c_2))} (\phi_r \timesk A_{2r})(\eta(x_{1r},x_{2r})) \]
    expanding definitions this is
    \[ (\lambda x_{1l}'. \eta(x_{1l}',x_{2l}))^\dagger\phi_l(\eta(x_{1l}))
    \binrel{(\li (c_1' \times c_2))}
    (\lambda x_{1r}'. \eta(x_{1r}',x_{2r}))^\dagger\phi_r(\eta(x_{1r}))
    \]
    By our assumptions we have $\phi_l(\eta(x_{1l}) \binrel {\li c_1'} \phi_r(\eta(x_{1r})$ so since $-^\dagger$ preserves squares it suffices to show
    \[ (\lambda x_{1l}'. \eta(x_{1l}',x_{2l})) \ltsq{\li c_1}{\li (c_1' \times c_2)} (\lambda x_{1r}'. \eta(x_{1r}',x_{2r})) \]
    To show this assume $x_{1l}' \binrel{c_1'} x_{1r}'$. Then we have
    \[ \eta(x_{1l}',x_{2l}) \binrel{\li (c_1' \times c_2)} \eta(x_{1r'},x_{2l})  \]
    holds.
  \end{itemize}
\end{proof}

\section{Omitted Proofs Section \ref{sec:towards-relational-model}}

\begin{theorem}
  Let $R$ be a binary relation on the free error domain $U(\li A)$. If
  $R$ satisfies the following properties
  \begin{enumerate}
  \item Transitivity
  \item $\theta$-congruence: If $\later_t (\tilde{x}_t \binrel{R} \tilde{y}_t)$, then $\theta(\tilde{x}) \binrel{R} \theta(\tilde{y})$.
  \item Right step-insensitivity: If $x \binrel{R} y$ then $x \binrel{R} \delta y$.
  \end{enumerate}
  Then for any $l : U(\li A)$, we have $l \binrel{R} \Omega$. If $R$ is left
  step-insensitive instead then $\Omega \binrel{R} x$.
\end{theorem}
\begin{proof}
  By \lob-induction: we assume that $\laterhs (l \binrel{R} \Omega)$, and we
  show $(l \binrel{R} \Omega)$. We have $l \binrel{R} \delta l$ by right
  step-insensitivity.
  %
  By transitivity, it will therefore suffice to show that $\delta l \binrel{R}
  \Omega$.
  %
  Then since $\Omega = \delta \Omega$, it suffices to show $\delta l \binrel{R}
  \delta\Omega$.
  %
  By $\theta$-congruence, it suffices to show that
  $\later_t [(\nxt\, l)_t \binrel{R} (\nxt\, \Omega)_t]$, i.e., $\later (l \binrel{R} \Omega)$,
  which is our induction hypothesis.

  The proof when $R$ is left step-insensitive is symmetric.
\end{proof}

\section{Details of the Relational Model Construction}

In Section \ref{sec:concrete-relational-model} we described our model of gradual
typing. We omitted several technical proofs involving the well-definedness of
composition of relations and the actions of functors. We provide those proofs in
this section.

\subsection{Free Composition of Error Domain Relations}
\begin{definition}\label{def:free-comp-ed-rel}
  Let $d : B_1 \rel B_2$ and $d' : B_2 \rel B_3$ be error domain relations. We
  define the composition $dd'$ inductively by the following rules:
  \begin{mathpar}
      \inferrule*[right = Comp]
      {b_1 \mathbin{d} b_2 \and b_2 \mathbin{d'} b_3}
      {b_1 \mathbin{d \relcomp d'} b_3}

      \inferrule*[right = DnClosed]
      {b_1' \le_{B_1} b_1 \and b_1 \mathbin{d \relcomp d'} b_3}
      {b_1' \mathbin{d \relcomp d'} b_3}

      \inferrule*[right = UpClosed]
      {b_1 \mathbin{d \relcomp d'} b_3 \and b_3 \le_{B_3} b_3'}
      {b_1 \mathbin{d \relcomp d'} b_3'}

      \inferrule*[right = PresErr]
      { }
      {\mho_{B_1} \mathbin{d \relcomp d'} b_3}

      \inferrule*[right = PresTheta]
      {\later_t( \tilde{b_1} \mathbin{d \relcomp d'} \tilde{b_3} ) }
      {\theta_{B_1}(\tilde{b_1}) \mathbin{d \relcomp d'} \theta_{B_3}(\tilde{b_3}) }
  \end{mathpar}

  This has the universal property that $\phi \ltsq{d \relcomp d'}{d''}
  \phi'$ if and only if for any $b \binrel{d} b'$ and $b' \binrel{d'}
  b''$ that $\phi(b) \binrel{d''} \phi'(b'')$.
\end{definition}

\subsection{The Dynamic Type}
The ordering $\le_D$ on the predomain $D$ is defined by
\begin{align*}
  \inat(n) \le \inat(n') 
      &\iff n = n' \\
  \itimes (d_1, d_2) \le \itimes (d_1', d_2')
      &\iff d_1 \le_D d_2 \text{ and } d_1' \le_D d_2'\\
  \iarr(\tilde{f}) \le \iarr(\tilde{f'}) 
      &\iff \later_t(\tilde{f}_t \le \tilde{f'}_t),
\end{align*}

and the bisimilarity relation is analogous.


\subsection{Functorial Actions on Syntactic Perturbations}

\DIFaddbegin \DIFadd{We provide details on the definition of the homomorphisms specifying the
interpretation of perturbations corresponding for the functors $\li$, $U$,
$\times$, and $\arr$.
}

\DIFaddend \begin{itemize}
  %DIF <  \li
  \item \DIFdelbegin \DIFdel{Then to interpret the perturbations as endomorphisms}\DIFdelend \DIFaddbegin \DIFadd{To define the homomorphism $i_{\li A}$}\DIFaddend , by the universal property
  of the coproduct of monoids it suffices to define two homomorphisms, one from
  $\mathbb{N} \to \morbisimid{\li A}$ and one from $M_A \to \morbisimid{\li A}$.
  The first homomorphism sends the generator $1 \in \mathbb{N}$ to the error
  domain endomorpism $(\delta \circ \eta)^\dagger$, and the second sends $m_A \in
  M_A$ to $\li(i_A(m_A))$ and we observe in both cases that the resulting
  endomorphisms are bisimilar to the identity since the action of $\li$ on
  morphisms preserves bisimilarity.

  \item To \DIFdelbegin \DIFdel{construct }\DIFdelend \DIFaddbegin \DIFadd{define }\DIFaddend the homomorphism $i_{A_1 \times A_2}$, we
  appeal to the universal property and give homomorphisms $M_{A_1} \to
  \morbisimid{A_1 \times A_2}$ and $M_{A_2} \to \morbisimid{A_1 \times A_2}$. The
  former applies $i_{A_1}$ in the first component and the identity in the second
  component, while the latter does the reverse with $i_{A_2}$.

  \item To \DIFdelbegin \DIFdel{construct the interpretation as endomorphisms }\DIFdelend \DIFaddbegin \DIFadd{define the homomorphism }\DIFaddend $i_{A \arr B}$ it suffices to define a
  homomorphism $M_A^{op} \to \morbisimid{A \arr B}$ and $M_B \to \morbisimid{A
  \arr B}$. The former takes an element $m_A$ and a predomain morphism $f : A \to
  UB$ and returns the composition $f \circ i_A(m_A)$, while the latter is defined
  similarly by post-composition with $i_B$.

\end{itemize}

\subsection{Lemmas about Perturbations}

The goal of this section is to prove the following lemmas:

%%%%%%%%%%%%%%%
% Composition %
%%%%%%%%%%%%%%%
\begin{lemma}\label{lem:push-pull-comp}
  Let $(A_1, M_{A_1}, i_{A_1})$, $(A_2, M_{A_2}, i_{A_2})$, and $(A_3, M_{A_3},
  i_{A_3})$ be value objects, and let $c : A_1 \rel A_2$ and $c' : A_2 \rel A_3$
  be predomain relations.

  Given a push-pull structure $\Pi_c$ for $c$ and $\Pi_{c'}$ for $c'$, we can
  define a push-pull structure $\Pi_{c \comp c'}$ for $c \comp c'$.

  Likewise, we can define a push-pull structure for the composition of error
  domain relations.
\end{lemma}
\begin{proof}
  We define $\piv_{c \comp c'}$ as follows: We first define
  %
  $\push_{c \comp c'} = \push_{c'} \circ \push_{c}$ and 
  $\pull_{c \comp c'} = \pull_{c} \circ \pull_{c'}$.
  %
  We then observe that the required squares exist for both push and pull. In
  particular, for push we have that $i_{A_1}(m_{A_1}) \ltdyn_c^c
  i_{A_2}(\push_c(m_{A_1}))$ using the push property for $c$, and then using the
  push property for $c'$ we have $i_{A_2}(\push_c(m_{A_1})) \ltdyn_{c'}^{c'}
  i_{A_3}(\push_{c'}(\push_c(m_{A_1})))$. We can then compose these squares
  horizontally to obtain the desired square. The pull property follows
  similarly.

  The push-pull structure for the composition of computation relations is
  defined analogously.
\end{proof}


%%%%%%%%%%%
% U and F %
%%%%%%%%%%%
\begin{lemma}\label{lem:push-pull-U-F} 

  Let $A$ and $A'$ be value objects and let $c : A \rel A'$ be a relation on the
  underlying predomains. Given a push-pull structure $\Pi_c$ for $c$ we can
  define a push-pull structure $\Pi_{\li c}$ for $\li c$ with respect to $\li A$
  and $\li A'$.

  Likewise, let $B$ and $B'$ be computation objects and let $d : B \rel B'$ be
  an error domain relation. Given a push-pull structure for $d$ with respect to
  $B$ and $B'$, we can define a push-pull structure for $Ud$ with respect to
  $UB$ and $UB'$.
\end{lemma}
\begin{proof}

  We define $\push_{\li c} : \mathbb{N} \oplus M_{A} \to \mathbb{N} \oplus
  M_{A'}$ by the universal property of the coproduct of monoids. In particular,
  it suffices to define a homomorphism of monoids $\mathbb{N} \to \mathbb{N}
  \oplus M_{A'}$ and $M_{A} \to \mathbb{N} \oplus M_{A'}$.

  The former is simply $\inl$ and the latter is $\inr \circ \push_c$.

  Then to establish the push property, it suffices by the universal property of
  the coproduct, and the fact that $\mathbb{N}$ is the free monoid on one
  generator, to show that the following two squares exist:

% https://q.uiver.app/#q=WzAsNCxbMCwwLCJcXGxpIEEiXSxbMSwwLCJcXGxpIEEnIl0sWzAsMSwiXFxsaSBBIl0sWzEsMSwiXFxsaSBBJyJdLFswLDEsIlxcbGkgYyIsMCx7InN0eWxlIjp7ImJvZHkiOnsibmFtZSI6ImJhcnJlZCJ9LCJoZWFkIjp7Im5hbWUiOiJub25lIn19fV0sWzIsMywiXFxsaSBjIiwyLHsic3R5bGUiOnsiYm9keSI6eyJuYW1lIjoiYmFycmVkIn0sImhlYWQiOnsibmFtZSI6Im5vbmUifX19XSxbMCwyLCJcXGRlbHRhXipfQSIsMl0sWzEsMywiXFxkZWx0YV4qX3tBJ30iXV0=
\[\begin{tikzcd}[ampersand replacement=\&]
	{\li A} \& {\li A'} \\
	{\li A} \& {\li A'}
	\arrow["{\li c}", "\shortmid"{marking}, no head, from=1-1, to=1-2]
	\arrow["{\delta^*_A}"', from=1-1, to=2-1]
	\arrow["{\delta^*_{A'}}", from=1-2, to=2-2]
	\arrow["{\li c}"', "\shortmid"{marking}, no head, from=2-1, to=2-2]
\end{tikzcd}\]

% https://q.uiver.app/#q=WzAsNCxbMCwwLCJcXGxpIEEiXSxbMCwxLCJcXGxpIEEiXSxbMSwwLCJcXGxpIEEnIl0sWzEsMSwiXFxsaSBBJyJdLFswLDIsIlxcbGkgYyIsMCx7InN0eWxlIjp7ImJvZHkiOnsibmFtZSI6ImJhcnJlZCJ9LCJoZWFkIjp7Im5hbWUiOiJub25lIn19fV0sWzEsMywiXFxsaSBjIiwyLHsic3R5bGUiOnsiYm9keSI6eyJuYW1lIjoiYmFycmVkIn0sImhlYWQiOnsibmFtZSI6Im5vbmUifX19XSxbMCwxLCJcXGxpKGlfQShtX0EpKSIsMl0sWzIsMywiXFxsaSAoaV97QSd9KFxccHVzaF9jKG1fQSkpKSJdXQ==
\[\begin{tikzcd}[ampersand replacement=\&]
	{\li A} \& {\li A'} \\
	{\li A} \& {\li A'}
	\arrow["{\li c}", "\shortmid"{marking}, no head, from=1-1, to=1-2]
	\arrow["{\li(i_A(m_A))}"', from=1-1, to=2-1]
	\arrow["{\li (i_{A'}(\push_c(m_A)))}", from=1-2, to=2-2]
	\arrow["{\li c}"', "\shortmid"{marking}, no head, from=2-1, to=2-2]
\end{tikzcd}\]

where $\delta* = (\delta \circ \eta)^\dagger$

The first square exists by the fact that the monadic extension operation
$-^\dagger$ is monotone in its argument, and the fact that there is a
square $\delta_A \circ \eta_A \ltsq{}{} \delta_{A'} \circ \eta_{A'}$.

The second square exists by the push property for $c$ and the fact that $\li$
acts on squares.

The proof for the pull property is dual, and the construction of a push-pull
structure for $Ud$ is analogous.
\end{proof}


%%%%%%%%%%%%
% Products %
%%%%%%%%%%%%
\begin{lemma}\label{lem:push-pull-times} Let $A_1$, $A_1'$, $A_2$, and $A_2'$ be
  value objects. Let $c_1 : A_1 \rel A_1'$ and $c_2 : A_2 \rel A_2'$ be
  relations on the underlying predomains. Given push-pull structures $\Pi_{c_1}$
  for $c_1$ and $\Pi_{c_2}$ for $c_2$ we can define a push-pull structure
  $\Pi_{c_1 \times c_2}$ for $c_1 \times c_2$.
\end{lemma}
\begin{proof}
Recall that the monoid of syntactic perturbations for $A_1 \times A_2$ is
$M_{A_1} \oplus M_{A_2}$ and the monoid for $A_1' \times A_2'$ is $M_{A_1'}
\oplus M_{A_2'}$. Thus to define the push homomorphism for $c_1 \times c_2$, it
suffices by the universal property of the coproduct of monoids to define a
homomorphism $M_{A_1} \to M_{A_1'} \oplus M_{A_2'}$ and $M_{A_2} \to M_{A_1'}
\oplus M_{A_2'}$. The former is $\inl \circ \push_{c_1}$, and the latter is
$\inr \circ \push_{c_2}$. The pull homomorphism is defined similarly.

For the push property, we need to show that there is a square 
$i_{A_1 \times A_2}(m_1) \ltsq{}{} i_{A_1' \times A_2'}(\push_{c_1 \times c_2}(m_1))$.
%
By the universal property of the coproduct, it suffices to consider two cases: a
perturbation $m_1 \in M_{A_1}$ and a perturbation $m_2 \in M_{A_2}$. In the
first case, using the definition of the interpretation homomorphism for the
product, the square becomes $(i_{A_1}(m_1) \times \id_{A_2}) \ltsq{}{}
(i_{A_1'}(\push_{c_1}(m_1)) \times \id_{A_2'})$.
%
Then by the action of $\times$ on squares, it suffices to show that there are
squares $i_{A_1}(m_1) \ltsq{}{} i_{A_1'}(\push_{c_1}(m_1))$ and $\id_{A_2}
\ltsq{}{} \id_{A_2'}$. The former is the push property for $c_1$, and the latter
is an identity square.

The second case, i.e., $m_2 \in M_{A_2}$, is symmetric, and the proof of the
pull property is dual.
\end{proof}

%%%%%%%%%
% Arrow %
%%%%%%%%%
\begin{lemma}\label{lem:push-pull-arrow} Let $A$ and $A'$ be value objects and
  $B$ and $B'$ be computation objects. Let $c : A \rel A'$ be a relation on the
  underlying predomains and $d : B \rel B'$ a relation on the underlying error
  domains. Given push-pull structures $\Pi_c$ for $c$ and $\Pi_d$ for $d$, we
  can define a push-pull structure $\Pi_{c \arr d}$ for $c \arr d$.
\end{lemma}
\begin{proof}
  Similar to the product case above.
\end{proof}


We next define the actions of the Kleisli arrow and product functors on
syntactic perturbations.

% ``Kleisli arrow'' action on syntactic perturbations that takes a
% perturbation in $M_{UB}$ to a perturbation in $M_{U(A \to B)}$, as follows:

\begin{definition}\label{def:kleisli-arrow-perturbations} Let $A$ be a value
  object and $B$ be a computation object. We define a homomorphism of monoids
  $id \tok - \colon M_{UB} \to M_{U(A \arr B)}$ via the universal property of the
  coproduct of monoids by sending $\mathbb{N}$ to the first injection in $M_{U(A
  \arr B)} = \mathbb{N} \oplus M_A^{op} \oplus M_B$, and sending $M_B$ to the
  third injection.

  Likewise, we define $- \tok \id \colon M_{\li A}^{op} \to M_{U(A \arr B)}$ as
  follows. First note that $M_{\li A}^{op} = \mathbb{N} \oplus M_A^{op}$. Then
  we construct the homomorphism via the universal property, sending $\mathbb{N}$
  to the first injection $M_A^{op}$ to the second injection.
\end{definition}

We can similarly construct a Kleisli product operation on perturbations:

\begin{definition}\label{def:kleisli-product-perturbations}
  Let $A_1$ and $A_2$ be value objects. We define a homomorphism of monoids 
  $- \timesk \id \colon M_{\li A_1} \to M_{\li (A_1 \times A_2)}$ using the universal property,
  sending $\mathbb{N}$ to the first injection and $M_{A_1}$ to the second injection.

  Likewise, we define $\id \timesk - \colon M_{\li A_2} \to M_{\li (A_1 \times A_2)}$
  by sending $\mathbb{N}$ to the first injection and $M_{A_2}$ to the third injection.

  %  Given a perturbation $(n, a_1, a_2) \in $ we define $(n, a_1, a_2)
  % \timesk \id = (n, a_1, \id_{P_{A_2}})$, and likewise we define $\id \timesk (n,
  % a_1, a_2) = (n, \id_{P_{A_1}}, a_2)$.
\end{definition}

It is then easy to verify that the Kleisli arrow action on perturbation is well-defined 
in that for any $m \in M_{UB}$, we have
%
\[ \ptb_{U(A \arr B)}(\id \tok m) = \id \tok i_{UB}(m), \]
%
where the LHS is the Kleisli action on perturbation and the RHS is the Kleisli
action on morphisms. 
%
The proof uses the fact that the morphism $\delta$ commutes with all error
domain morphisms, which is a consequence of the definition of error domain
morphism.
%
The other verification involving the Kleisli arrow action is similar, as are the
two Kleisli product actions.

% %
% \[ (\delta_{U(A \arr B)}^*)^n \circ U(\id_{A} \arr \i_B(b)) = 
%     \id \tok (\delta_{UB}^*)^n \circ U(\ptb_B(b)) \]
% %
% This can be checked by unfolding the definition of $\tok$; the proof uses the
% fact that $\delta_{UB}^*$ commutes with computation morphisms.
% %

% Given a perturbation $q = (n, a) \in P_{FA}$ we define $(n, a) \tok \id = (n, a,
% \id_{P_B}) \in P_{U(A \to B)}$. We need to show that
% %
% \[ \ptb_{U(A \arr B)}((n, a) \tok \id) = \ptb_{FA}(n, a) \tok \id. \]
% %
% By similar reasoning to the above, this holds because $\delta_{FA}^*$ commutes
% with computation morphisms.



% We have $\ptb_{U(A \arr B)}(\id \tok (n, b)) = \ptb_{U(A \arr B)}(n, \id, b)$, 
% which is equal to $\delta_{U(A \arr B)}^n (\id \to \ptb_B(b))$


% %%%%%%%%%%%%%%%%%%%%%%
% % Model Construction %
% %%%%%%%%%%%%%%%%%%%%%%
% We now proceed with the construction of the model:

%     Define a step-2 model $\mathcal M'$ as follows:
%     \begin{itemize}
%       \item Value objects are pairs consisting of:
%       \begin{itemize}
%         \item A value object $A$ in $\vf$.
%         \item A monoid $P_A$ of perturbations
%         and a monoid homomorphism $\ptb_A : P_A \to \{ f \in \vf(A, A) \mid f \bisim \id_A \}$.

%       \end{itemize}  

%       \item Computation objects are pairs consisting of:
%       \begin{itemize}
%         \item A computation object $B$ in $\ef$.
%         \item A monoid $P_B$ of perturbations and a monoid homomorphism
%         $\ptb_B : P_B \to \{ \phi \in \ef(B, B) \mid \phi \bisim \id_B \}$.
%       \end{itemize}

%       \item Morphisms are given by morphisms of the underlying objects in $\vf$ and $\ef$, respectively.
%       %, i.e.,
%       % \[ \vf'((A, P_A, \ptb_A, P^K_A, \ptbk_A), (A', P_{A'}, \ptb_{A'}, P^K_{A'}, \ptbk_{A'})) = \vf(A, A') \]
%       %
%       % and likewise for computations.

%       \item Given objects $(A, P_A, \ptb_A)$ and $(A', P_{A'}, \ptb_{A'})$ a value relation
%       is a pair consisting of
%       \begin{itemize}
%         \item A value relation $c \in \vsq$.
%         \item A push-pull structure $\piv_c$ for $c$ with respect to $P_A$ and $P_{A'}$.
%         %\item A push-pull structure $\pie_{Fc}$ for $Fc$ with respect to $P^K_A$ and $P^K_{A'}$.
%       \end{itemize}

%       \item Similarly, a computation relation between $(B, P_B, \ptb_B)$ and $(B', P_{B'}, \ptb_{B'})$
%       consists of
%       \begin{itemize}
%         \item A computation relation $d \in \esq$.
%         \item A push-pull structure $\pie_d$ for $d$ with respect to $P_B$ and $P_{B'}$.
%         %\item A push-pull structure $\piv_{Ud}$ for $Ud$ with respect to $P^K_B$ and $P^K_{B'}$.
%       \end{itemize}

%       \item The squares are the same as the squares of the original model $\mathcal M$
%       %morphisms of $\vsq'$ and $\esq'$ are given by the morphisms of $\vsq$ and $\esq$.

%       \item We define composition of relations $(c, \piv_c)$ and $(c', \piv_{c'})$
%       as $(c \comp c', \piv_{c \comp c'})$ where $\piv_{c \comp c'}$ is as in Lemma
%       \ref{lem:push-pull-comp}, and likewise for computation relations.

%     \end{itemize}

%       % Functors \times, +, F, U, arrow
%       Now we define the actions of the functors:
%       \begin{itemize}

%       \item We define $\times$ on objects by

%       \[ (A_1, P_{A_1}, \ptb_{A_1}) \times (A_2, P_{A_2}, \ptb_{A_2}) =
%         (A_1 \times A_2, P_{A_1} \times P_{A_2}, \ptb_{A_1 \times A_2}) \]

%       where $\ptb_{A_1 \times A_2}(p_1, p_2) = \ptb_{A_1}(p_1) \times \ptb_{A_2}(p_2)$.

%       Using Lemma \ref{lem:push-pull-times}, we define $\times$ on relations by

%       \[ (c_1, \piv_{c_1}), (c_2 \piv_{c_2}) = 
%         (c_1 \times c_2, \piv_{c_1 \times c_2}). \]

%       \item We define $F$ on objects by 

%       \[ F(A, P_A, \ptb_A) = (FA, \mathbb{N} \times P_A, \ptb_{FA}), \]

%       where we define $\ptb_{FA}(n, a) = (\delta_{FA}^*)^n \circ F(\ptb_A(a))$.

%       Using Lemma \ref{lem:push-pull-U-F}, we define $F$ on relations by

%       \[ F(c, \piv_c) = (Fc, \pie_{Fc}). \]

%       \item We define $U$ on objects by

%       \[ U(B, P_B, \ptb_B) = (UB, \mathbb{N} \times P_B, \ptb_{UB}), \]

%       where we define $\ptb_{UB}(n, b) = (\delta_{UB}^*)^n \circ U(\ptb_B(b))$.

%       Using Lemma \ref{lem:push-pull-U-F}, we define $U$ on relations by

%       \[ U(d, \pie_d) = (Ud, \piv_{Ud}).  \]

%       We define $\arr$ on objects by

%       \[ (A, P_A, \ptb_A) \arr (B, P_B, \ptb_B) = 
%         (A \arr B, P_A^{op} \times P_B, \ptb_{A \arr B}), \]

%       where we define $\ptb_{A \arr B}(a, b) = \ptb_A(a) \arr \ptb_B(b)$.

%       Using Lemma \ref{lem:push-pull-arrow}, we define $\arr$ on relations by

%       \[ (c, \piv_c) \arr (d, \pie_d) =
%         (c \arr d, \pie_{c \arr d}). \]

%     \end{itemize}


    


%%%%%%%%%%%%%%%%%%%%%%%%%%%%%%%%%%%%%%%%%%%%%%%%%%%%%%%%%%%%%%%%%%%%%%%%%%%%%%

\subsection{Lemmas involving Quasi-Representable Relations}

In this section we prove lemmas about quasi-representability needed for showing
that our notions of value and computation relation are compositional, and for
defining the action of the functors $U$, $\li$, $\arr$, and $\times$ on value
and computation relations.

Recall the notion of quasi-equivalence of relations as defined in Definition
\ref{def:quasi-equivalent}. The following lemma will be useful in showing
that two relations are quasi-equivalent.

%%%%%%%%%%%%%%%%%%%%%%%%%%%%%%%%%%%%%%%%%%%%%%%%%%%%%%%%%%%%%%
% quasi-representable by same morphism --> quasi-equivalent
%%%%%%%%%%%%%%%%%%%%%%%%%%%%%%%%%%%%%%%%%%%%%%%%%%%%%%%%%%%%%%
\begin{lemma}\label{lem:left-rep-by-same-morphism} Let $A$ and $A'$ be value
  objects, and let $c, c' : A \rel A'$ be relations between the underlying
  predomains. If $c$ and $c'$ are both quasi-left-representable by the same
  predomain morphism $f : A \to A'$, then $c \qordeq c'$. If $c$ and $c'$ are
  both quasi-right-representable by the same predomain morphism $g : A' \to A$
  then $c \qordeq c'$.

  Dually, let $B$ and $B'$ be computation objects, and let $d, d' : B \rel B'$
  be relations between the underlying error domains. If $d$ and $d'$ are both
  quasi-right-representable by the same error domain morphism $\phi : B'
  \multimap B$, then $d \qordeq d'$. If $d$ and $d'$ are both
  quasi-left-representable by the same $\phi : B \multimap B'$
\end{lemma}
\begin{proof}
  We show the result for $c$ and $c'$ that are quasi-left-representable by the same $f$;
  the other proofs are analogous. 

  By $\upr$ for $c'$, there exists a perturbation $\delle_{c'}$ and a square
  $\upr_{c'} : \delle_{c'} \ltdyn_{c'}^{r(A)} f$.

  By $\upl$ for $c$, there exists a perturbation $\delre_c$ and a square 
  $\upl_c : f \ltdyn_{r(A')}^{c} \delre_c$.

  Composing these horizontally we get the following square:

  % https://q.uiver.app/#q=WzAsNixbMCwwLCJBIl0sWzEsMCwiQSJdLFsyLDAsIkEnIl0sWzAsMSwiQSJdLFsxLDEsIkEnIl0sWzIsMSwiQSciXSxbMCwxLCJyKEEpIiwwLHsic3R5bGUiOnsiYm9keSI6eyJuYW1lIjoiYmFycmVkIn0sImhlYWQiOnsibmFtZSI6Im5vbmUifX19XSxbMSwyLCJjIiwwLHsic3R5bGUiOnsiYm9keSI6eyJuYW1lIjoiYmFycmVkIn0sImhlYWQiOnsibmFtZSI6Im5vbmUifX19XSxbMyw0LCJjJyIsMix7InN0eWxlIjp7ImJvZHkiOnsibmFtZSI6ImJhcnJlZCJ9LCJoZWFkIjp7Im5hbWUiOiJub25lIn19fV0sWzQsNSwicihBJykiLDIseyJzdHlsZSI6eyJib2R5Ijp7Im5hbWUiOiJiYXJyZWQifSwiaGVhZCI6eyJuYW1lIjoibm9uZSJ9fX1dLFswLDMsImlfQShcXGRlbGxlX3tjJ30pIiwyXSxbMSw0LCJmIiwyXSxbMiw1LCJpX3tBJ30oXFxkZWxyZV9jKSJdXQ==
  \[\begin{tikzcd}[ampersand replacement=\&]
    A \& A \& {A'} \\
    A \& {A'} \& {A'}
    \arrow["{r(A)}", "\shortmid"{marking}, no head, from=1-1, to=1-2]
    \arrow["{i_A(\delle_{c'})}"', from=1-1, to=2-1]
    \arrow["c", "\shortmid"{marking}, no head, from=1-2, to=1-3]
    \arrow["f"', from=1-2, to=2-2]
    \arrow["{i_{A'}(\delre_c)}", from=1-3, to=2-3]
    \arrow["{c'}"', "\shortmid"{marking}, no head, from=2-1, to=2-2]
    \arrow["{r(A')}"', "\shortmid"{marking}, no head, from=2-2, to=2-3]
  \end{tikzcd}\]

  Then since $r(A) \comp c = c$ and $c' \comp r(A') = c'$, we obtain the desired square.
  %
  The other square (i.e., with $c'$ on top) is constructed in an analogous
  manner.

\end{proof}


Next we show that the functors $\li$ and $U$ preserve quasi-representability.

%%%%%%%%%%%%%
% U and F
%%%%%%%%%%%%%
\begin{lemma}\label{lem:representation-U-F}
  Let $A$ and $A'$ be value objects and let $c : A \rel A'$ be a predomain relation.
  If $c$ is quasi-left-representable, then $\li c$ is quasi-left-representable.
  Likewise, if $c$ is quasi-right-representable, then $\li c$ is quasi-right-representable.

  %Then we can define a left-representation structure $\rho^L_{UF(c)}$ for $UF(c)$.

  Similarly, let $B$ and $B'$ be computation objects and let $d : B \rel B'$ be an error domain relation.
  If $d$ is quasi-right-representable, then $Ud$ is quasi-right-representable.

  % Then we can define a right-representation structure $\rho^R_{FU(d)}$ for $FU(d)$.
\end{lemma}
\begin{proof}
  To show that $\li c$ is quasi-left-representable, we define
  \begin{itemize}
    \item $e_{\li c} = \li e_c$
    \item $\delre_{\li c} = \inr (\delre_c)$ 
    (recalling that the monoid $M_{\li A'} = \mathbb{N} \oplus M_{A'}$)
    \item $\delle_{\li c} = \inr (\delle_c)$
  \end{itemize}
  We obtain the two squares $\upl$ and $\upr$ using the definition of the
  interpretation of syntactic perturbations for $\li$ and the functorial action
  of $\li$ on the corresponding squares for $c$.

  The proof for $Ud$ is analogous.
\end{proof}



%%%%%%%%%%%%
% Identity %
%%%%%%%%%%%%

Next we show that the reflexive relation is quasi-representable.

\begin{lemma}\label{lem:reflexive-rel-quasi-rep}
  Let $A$ be a value object. Then $r(A)$ is quasi-left-representable and quasi-right-representable.
  Likewise, $r(B)$ is quasi-left- and quasi-right-representable for any computation object $B$.
\end{lemma}
\begin{proof}
  To show $r(A)$ is quasi-left-representable, we take the embedding to be the
  identity morphism and the perturbations $\delle = \delre = \id_{M_{A}}$, i.e.,
  the identtiy of the monoid. Then because $i_A$ is a homomorphism of monoids,
  it sends the identity element to the identity morphism, so the $\upl$ and
  $\upr$ squares are just identity squares.

  The same argument shows that $r(A)$ is quasi-right-representable and that
  $r(B)$ is quasi-left- and quasi-right-representable.
\end{proof}


%%%%%%%%%%%%%%%
% Composition %
%%%%%%%%%%%%%%%

We begin by showing that the composition of quasi-representable relations with
push-pull structures is quasi-representable.

\begin{lemma}\label{lem:representation-comp}
  Let $A$, $A'$, and $A''$ be value objects, and $B$, $B'$, and
  $B''$ be computation objects. Let $c : A \rel A'$ and $c' : A' \rel A''$ be
  predomain relations with push-pull structures, and let $d : B \rel B'$ and $d'
  : B' \rel B''$ be error domain relations with push-pull structures.
  \begin{enumerate}
    \item If $c$ and $c'$ are quasi-left- (resp. right)-representable, then $c
    \comp c'$ is quasi-left- (resp. right)-representable.

    \item If $d$ and $d'$ are quasi-left- (resp. right)-representable, then
    $d \comp d'$ is quasi-left (resp. right)-representable.
  \end{enumerate}
\end{lemma}
\begin{proof}
    % 1.
    \item To show $c \comp c'$ is quasi-left-representable we define
    \begin{itemize}
      \item $e_{c \comp c'} = e_{c'} \circ e_c$
      \item $\delre_{c \comp c'} = \delre_{c'} \cdot \push_{c'}(\delre_c)$ where
      $\cdot$ denotes multiplication in the monoid $M_A$
      \item $\delle_{c \comp c'} = \pull_c(\delle_{c'}) \cdot \delle_c$
      \item $\upl$ is the following square: % https://q.uiver.app/#q=WzAsMTAsWzAsMCwiQSJdLFswLDEsIkEnIl0sWzAsMiwiQSciXSxbMiwwLCJBJyciXSxbMiwxLCJBJyciXSxbMiwyLCJBJyciXSxbMCwzLCJBJyciXSxbMiwzLCJBJyciXSxbMSwxLCJBJyJdLFsxLDAsIkEnIl0sWzAsMSwiZV9jIiwyXSxbMSwyLCJcXGlkIiwyXSxbMiw2LCJlX3tjJ30iLDJdLFs0LDUsIlxcaWQiXSxbMSw4LCJyKEEnKSIsMCx7InN0eWxlIjp7ImJvZHkiOnsibmFtZSI6ImJhcnJlZCJ9LCJoZWFkIjp7Im5hbWUiOiJub25lIn19fV0sWzgsNCwiYyciLDAseyJzdHlsZSI6eyJib2R5Ijp7Im5hbWUiOiJiYXJyZWQifSwiaGVhZCI6eyJuYW1lIjoibm9uZSJ9fX1dLFszLDQsIlxccHVzaF97Yyd9KFxcZGVscmVfYykiXSxbMCw5LCJjIiwwLHsic3R5bGUiOnsiYm9keSI6eyJuYW1lIjoiYmFycmVkIn0sImhlYWQiOnsibmFtZSI6Im5vbmUifX19XSxbOSwzLCJjJyIsMCx7InN0eWxlIjp7ImJvZHkiOnsibmFtZSI6ImJhcnJlZCJ9LCJoZWFkIjp7Im5hbWUiOiJub25lIn19fV0sWzksOCwiXFxkZWxyZV9jIl0sWzUsNywiXFxkZWxyZV97Yyd9ICJdLFsyLDUsImMnIiwwLHsic3R5bGUiOnsiYm9keSI6eyJuYW1lIjoiYmFycmVkIn0sImhlYWQiOnsibmFtZSI6Im5vbmUifX19XSxbNiw3LCJyKEEnJykiLDIseyJzdHlsZSI6eyJib2R5Ijp7Im5hbWUiOiJiYXJyZWQifSwiaGVhZCI6eyJuYW1lIjoibm9uZSJ9fX1dLFsxMCwxOSwiXFx1cGxfYyIsMSx7InNob3J0ZW4iOnsic291cmNlIjoyMCwidGFyZ2V0IjoyMH0sInN0eWxlIjp7ImJvZHkiOnsibmFtZSI6Im5vbmUifSwiaGVhZCI6eyJuYW1lIjoibm9uZSJ9fX1dLFsxMSwxMywiKCoqKSIsMSx7InNob3J0ZW4iOnsic291cmNlIjoyMCwidGFyZ2V0IjoyMH0sInN0eWxlIjp7ImJvZHkiOnsibmFtZSI6Im5vbmUifSwiaGVhZCI6eyJuYW1lIjoibm9uZSJ9fX1dLFsxOSwxNiwiKCopIiwxLHsic2hvcnRlbiI6eyJzb3VyY2UiOjIwLCJ0YXJnZXQiOjIwfSwic3R5bGUiOnsiYm9keSI6eyJuYW1lIjoibm9uZSJ9LCJoZWFkIjp7Im5hbWUiOiJub25lIn19fV0sWzEyLDIwLCJcXHVwbF97Yyd9IiwxLHsic2hvcnRlbiI6eyJzb3VyY2UiOjIwLCJ0YXJnZXQiOjIwfSwic3R5bGUiOnsiYm9keSI6eyJuYW1lIjoibm9uZSJ9LCJoZWFkIjp7Im5hbWUiOiJub25lIn19fV1d
\[\begin{tikzcd}[ampersand replacement=\&,column sep=3.15em]
	A \& {A'} \& {A''} \\
	{A'} \& {A'} \& {A''} \\
	{A'} \&\& {A''} \\
	{A''} \&\& {A''}
	\arrow[""{name=0, anchor=center, inner sep=0}, "{e_c}"', from=1-1, to=2-1]
	\arrow[""{name=1, anchor=center, inner sep=0}, "\id"', from=2-1, to=3-1]
	\arrow[""{name=2, anchor=center, inner sep=0}, "{e_{c'}}"', from=3-1, to=4-1]
	\arrow[""{name=3, anchor=center, inner sep=0}, "\id", from=2-3, to=3-3]
	\arrow["{r(A')}", "\shortmid"{marking}, no head, from=2-1, to=2-2]
	\arrow["{c'}", "\shortmid"{marking}, no head, from=2-2, to=2-3]
	\arrow[""{name=4, anchor=center, inner sep=0}, "{\push_{c'}(\delre_c)}", from=1-3, to=2-3]
	\arrow["c", "\shortmid"{marking}, no head, from=1-1, to=1-2]
	\arrow["{c'}", "\shortmid"{marking}, no head, from=1-2, to=1-3]
	\arrow[""{name=5, anchor=center, inner sep=0}, "{\delre_c}", from=1-2, to=2-2]
	\arrow[""{name=6, anchor=center, inner sep=0}, "{\delre_{c'} }", from=3-3, to=4-3]
	\arrow["{c'}", "\shortmid"{marking}, no head, from=3-1, to=3-3]
	\arrow["{r(A'')}"', "\shortmid"{marking}, no head, from=4-1, to=4-3]
	\arrow["{\upl_c}"{description}, draw=none, from=0, to=5]
	\arrow["{(**)}"{description}, draw=none, from=1, to=3]
	\arrow["{(*)}"{description}, draw=none, from=5, to=4]
	\arrow["{\upl_{c'}}"{description}, draw=none, from=2, to=6]
\end{tikzcd}\]

      The square $(*)$ exists by the push-pull property for $c'$, and the square
      $(**)$ exists by the downward closure of $c'$.

      \item $\upr$ is the following square: % https://q.uiver.app/#q=WzAsMTAsWzAsMCwiQSJdLFswLDEsIkEiXSxbMCwyLCJBIl0sWzIsMCwiQSJdLFsyLDEsIkEnIl0sWzIsMiwiQSciXSxbMCwzLCJBIl0sWzIsMywiQScnIl0sWzEsMiwiQSciXSxbMSwzLCJBJyJdLFswLDEsIlxcZGVsbGVfYyIsMl0sWzEsMiwiXFxpZCIsMl0sWzIsNiwiXFxwdWxsX2MoXFxkZWxsZV97Yyd9KSIsMl0sWzQsNSwiXFxpZCJdLFszLDQsImVfYyJdLFs1LDcsImVfe2MnfSJdLFswLDMsInIoQSkiLDAseyJzdHlsZSI6eyJib2R5Ijp7Im5hbWUiOiJiYXJyZWQifSwiaGVhZCI6eyJuYW1lIjoibm9uZSJ9fX1dLFsxLDQsImMiLDAseyJzdHlsZSI6eyJib2R5Ijp7Im5hbWUiOiJiYXJyZWQifSwiaGVhZCI6eyJuYW1lIjoibm9uZSJ9fX1dLFsyLDgsImMiLDAseyJzdHlsZSI6eyJib2R5Ijp7Im5hbWUiOiJiYXJyZWQifSwiaGVhZCI6eyJuYW1lIjoibm9uZSJ9fX1dLFs4LDUsInIoQScpIiwwLHsic3R5bGUiOnsiYm9keSI6eyJuYW1lIjoiYmFycmVkIn0sImhlYWQiOnsibmFtZSI6Im5vbmUifX19XSxbNiw5LCJjIiwyLHsic3R5bGUiOnsiYm9keSI6eyJuYW1lIjoiYmFycmVkIn0sImhlYWQiOnsibmFtZSI6Im5vbmUifX19XSxbOSw3LCJjJyIsMix7InN0eWxlIjp7ImJvZHkiOnsibmFtZSI6ImJhcnJlZCJ9LCJoZWFkIjp7Im5hbWUiOiJub25lIn19fV0sWzgsOSwiXFxkZWxsZV97Yyd9IiwyXSxbMTEsMTMsIigqKSIsMSx7InNob3J0ZW4iOnsic291cmNlIjoyMCwidGFyZ2V0IjoyMH0sInN0eWxlIjp7ImJvZHkiOnsibmFtZSI6Im5vbmUifSwiaGVhZCI6eyJuYW1lIjoibm9uZSJ9fX1dLFsxMCwxNCwiXFx1cHJfYyIsMSx7InNob3J0ZW4iOnsic291cmNlIjoyMCwidGFyZ2V0IjoyMH0sInN0eWxlIjp7ImJvZHkiOnsibmFtZSI6Im5vbmUifSwiaGVhZCI6eyJuYW1lIjoibm9uZSJ9fX1dLFsyMiwxNSwiIiwxLHsic2hvcnRlbiI6eyJzb3VyY2UiOjIwLCJ0YXJnZXQiOjIwfSwic3R5bGUiOnsiYm9keSI6eyJuYW1lIjoibm9uZSJ9LCJoZWFkIjp7Im5hbWUiOiJub25lIn19fV0sWzEyLDIyLCIoKiopIiwxLHsic2hvcnRlbiI6eyJzb3VyY2UiOjIwLCJ0YXJnZXQiOjIwfSwic3R5bGUiOnsiYm9keSI6eyJuYW1lIjoibm9uZSJ9LCJoZWFkIjp7Im5hbWUiOiJub25lIn19fV0sWzIyLDE1LCJcXHVwcl97Yyd9IiwxLHsic2hvcnRlbiI6eyJzb3VyY2UiOjIwLCJ0YXJnZXQiOjIwfSwic3R5bGUiOnsiYm9keSI6eyJuYW1lIjoibm9uZSJ9LCJoZWFkIjp7Im5hbWUiOiJub25lIn19fV1d
\[\begin{tikzcd}[ampersand replacement=\&,column sep=3.15em]
	A \&\& A \\
	A \&\& {A'} \\
	A \& {A'} \& {A'} \\
	A \& {A'} \& {A''}
	\arrow[""{name=0, anchor=center, inner sep=0}, "{\delle_c}"', from=1-1, to=2-1]
	\arrow[""{name=1, anchor=center, inner sep=0}, "\id"', from=2-1, to=3-1]
	\arrow[""{name=2, anchor=center, inner sep=0}, "{\pull_c(\delle_{c'})}"', from=3-1, to=4-1]
	\arrow[""{name=3, anchor=center, inner sep=0}, "\id", from=2-3, to=3-3]
	\arrow[""{name=4, anchor=center, inner sep=0}, "{e_c}", from=1-3, to=2-3]
	\arrow[""{name=5, anchor=center, inner sep=0}, "{e_{c'}}", from=3-3, to=4-3]
	\arrow["{r(A)}", "\shortmid"{marking}, no head, from=1-1, to=1-3]
	\arrow["c", "\shortmid"{marking}, no head, from=2-1, to=2-3]
	\arrow["c", "\shortmid"{marking}, no head, from=3-1, to=3-2]
	\arrow["{r(A')}", "\shortmid"{marking}, no head, from=3-2, to=3-3]
	\arrow["c"', "\shortmid"{marking}, no head, from=4-1, to=4-2]
	\arrow["{c'}"', "\shortmid"{marking}, no head, from=4-2, to=4-3]
	\arrow[""{name=6, anchor=center, inner sep=0}, "{\delle_{c'}}"', from=3-2, to=4-2]
	\arrow["{(*)}"{description}, draw=none, from=1, to=3]
	\arrow["{\upr_c}"{description}, draw=none, from=0, to=4]
	\arrow[draw=none, from=6, to=5]
	\arrow["{(**)}"{description}, draw=none, from=2, to=6]
	\arrow["{\upr_{c'}}"{description}, draw=none, from=6, to=5]
\end{tikzcd}\]
    \end{itemize}

    To show $c \comp c'$ is quasi-right-representable we define
    \begin{itemize}
      \item $p_{c \comp c'} = p_c \circ p_{c'}$
      \item $\dellp_{c \comp c'} = \dellp_c \cdot \pull_c(\dellp_{c'})$
      \item $\delrp_{c \comp c'} = \push_{c'}(\delrp_c) \cdot \delrp_{c'}$
      \item $\dnr$ is the following square: % https://q.uiver.app/#q=WzAsMTAsWzAsMCwiQSJdLFsyLDAsIkEnJyJdLFsyLDEsIkEnIl0sWzIsMiwiQSciXSxbMiwzLCJBIl0sWzAsMSwiQSJdLFswLDIsIkEiXSxbMCwzLCJBIl0sWzEsMCwiQSciXSxbMSwxLCJBJyJdLFswLDgsImMiLDAseyJzdHlsZSI6eyJib2R5Ijp7Im5hbWUiOiJiYXJyZWQifSwiaGVhZCI6eyJuYW1lIjoibm9uZSJ9fX1dLFs4LDEsImMnIiwwLHsic3R5bGUiOnsiYm9keSI6eyJuYW1lIjoiYmFycmVkIn0sImhlYWQiOnsibmFtZSI6Im5vbmUifX19XSxbNiwzLCJjIiwwLHsic3R5bGUiOnsiYm9keSI6eyJuYW1lIjoiYmFycmVkIn0sImhlYWQiOnsibmFtZSI6Im5vbmUifX19XSxbNyw0LCJyKEEpIiwyLHsic3R5bGUiOnsiYm9keSI6eyJuYW1lIjoiYmFycmVkIn0sImhlYWQiOnsibmFtZSI6Im5vbmUifX19XSxbMSwyLCJwX3tjJ30iXSxbMiwzLCJcXGlkIl0sWzMsNCwicF9jIl0sWzAsNSwiaV9BKFxccHVsbF9jKFxcZGVsbHBfe2MnfSkpIiwyXSxbNSw2LCJcXGlkIiwyXSxbNiw3LCJpX0EoXFxkZWxscF9kKSIsMl0sWzUsOSwiYyIsMCx7InN0eWxlIjp7ImJvZHkiOnsibmFtZSI6ImJhcnJlZCJ9LCJoZWFkIjp7Im5hbWUiOiJub25lIn19fV0sWzksMiwicihBJykiLDAseyJzdHlsZSI6eyJib2R5Ijp7Im5hbWUiOiJiYXJyZWQifSwiaGVhZCI6eyJuYW1lIjoibm9uZSJ9fX1dLFs4LDksIlxcZGVsbHBfe2MnfSIsMl0sWzE3LDIyLCIoKikiLDEseyJzaG9ydGVuIjp7InNvdXJjZSI6MjAsInRhcmdldCI6MjB9LCJzdHlsZSI6eyJib2R5Ijp7Im5hbWUiOiJub25lIn0sImhlYWQiOnsibmFtZSI6Im5vbmUifX19XSxbMjIsMTQsIlxcZG5yX3tjJ30iLDEseyJzaG9ydGVuIjp7InNvdXJjZSI6MjAsInRhcmdldCI6MjB9LCJzdHlsZSI6eyJib2R5Ijp7Im5hbWUiOiJub25lIn0sImhlYWQiOnsibmFtZSI6Im5vbmUifX19XSxbMTksMTYsIlxcZG5yX2MiLDEseyJzaG9ydGVuIjp7InNvdXJjZSI6MjAsInRhcmdldCI6MjB9LCJzdHlsZSI6eyJib2R5Ijp7Im5hbWUiOiJub25lIn0sImhlYWQiOnsibmFtZSI6Im5vbmUifX19XV0=
\[\begin{tikzcd}[ampersand replacement=\&]
	A \& {A'} \& {A''} \\
	A \& {A'} \& {A'} \\
	A \&\& {A'} \\
	A \&\& A
	\arrow["c", "\shortmid"{marking}, no head, from=1-1, to=1-2]
	\arrow[""{name=0, anchor=center, inner sep=0}, "{i_A(\pull_c(\dellp_{c'}))}"', from=1-1, to=2-1]
	\arrow["{c'}", "\shortmid"{marking}, no head, from=1-2, to=1-3]
	\arrow[""{name=1, anchor=center, inner sep=0}, "{\dellp_{c'}}"', from=1-2, to=2-2]
	\arrow[""{name=2, anchor=center, inner sep=0}, "{p_{c'}}", from=1-3, to=2-3]
	\arrow["c", "\shortmid"{marking}, no head, from=2-1, to=2-2]
	\arrow["\id"', from=2-1, to=3-1]
	\arrow["{r(A')}", "\shortmid"{marking}, no head, from=2-2, to=2-3]
	\arrow["\id", from=2-3, to=3-3]
	\arrow["c", "\shortmid"{marking}, no head, from=3-1, to=3-3]
	\arrow[""{name=3, anchor=center, inner sep=0}, "{i_A(\dellp_d)}"', from=3-1, to=4-1]
	\arrow[""{name=4, anchor=center, inner sep=0}, "{p_c}", from=3-3, to=4-3]
	\arrow["{r(A)}"', "\shortmid"{marking}, no head, from=4-1, to=4-3]
	\arrow["{(*)}"{description}, draw=none, from=0, to=1]
	\arrow["{\dnr_{c'}}"{description}, draw=none, from=1, to=2]
	\arrow["{\dnr_c}"{description}, draw=none, from=3, to=4]
\end{tikzcd}\]

      Here the square $(*)$ exists by the push-pull property for $c$.

      \item $\dnl$ is the following square: % https://q.uiver.app/#q=WzAsMTAsWzAsMCwiQScnIl0sWzIsMCwiQScnIl0sWzAsMSwiQSciXSxbMCwyLCJBJyJdLFswLDMsIkEiXSxbMiwxLCJBJyciXSxbMiwyLCJBJyciXSxbMiwzLCJBJyciXSxbMSwzLCJBJyJdLFsxLDIsIkEnIl0sWzAsMiwicF97Yyd9IiwyXSxbMiwzLCJcXGlkIiwyXSxbMyw0LCJwX2MiLDJdLFsxLDUsImlfe0EnJ30oXFxkZWxycF97Yyd9KSJdLFs1LDYsIlxcaWQiXSxbNiw3LCJpX3tBJyd9KFxccHVzaF97Yyd9KFxcZGVscnBfYykpIl0sWzAsMSwicihBJycpIiwwLHsic3R5bGUiOnsiYm9keSI6eyJuYW1lIjoiYmFycmVkIn0sImhlYWQiOnsibmFtZSI6Im5vbmUifX19XSxbMiw1LCJjJyIsMCx7InN0eWxlIjp7ImJvZHkiOnsibmFtZSI6ImJhcnJlZCJ9LCJoZWFkIjp7Im5hbWUiOiJub25lIn19fV0sWzMsOSwicihBJykiLDAseyJzdHlsZSI6eyJib2R5Ijp7Im5hbWUiOiJiYXJyZWQifSwiaGVhZCI6eyJuYW1lIjoibm9uZSJ9fX1dLFs5LDYsImMnIiwwLHsic3R5bGUiOnsiYm9keSI6eyJuYW1lIjoiYmFycmVkIn0sImhlYWQiOnsibmFtZSI6Im5vbmUifX19XSxbNCw4LCJjIiwyLHsic3R5bGUiOnsiYm9keSI6eyJuYW1lIjoiYmFycmVkIn0sImhlYWQiOnsibmFtZSI6Im5vbmUifX19XSxbOCw3LCJjJyIsMix7InN0eWxlIjp7ImJvZHkiOnsibmFtZSI6ImJhcnJlZCJ9LCJoZWFkIjp7Im5hbWUiOiJub25lIn19fV0sWzksOCwiXFxkZWxycF9jIl0sWzEwLDEzLCJcXGRubF97Yyd9IiwxLHsic2hvcnRlbiI6eyJzb3VyY2UiOjIwLCJ0YXJnZXQiOjIwfSwic3R5bGUiOnsiYm9keSI6eyJuYW1lIjoibm9uZSJ9LCJoZWFkIjp7Im5hbWUiOiJub25lIn19fV0sWzIyLDE1LCIoKikiLDEseyJzaG9ydGVuIjp7InNvdXJjZSI6MjAsInRhcmdldCI6MjB9LCJzdHlsZSI6eyJib2R5Ijp7Im5hbWUiOiJub25lIn0sImhlYWQiOnsibmFtZSI6Im5vbmUifX19XSxbMTIsMjIsIlxcZG5sX2MiLDEseyJzaG9ydGVuIjp7InNvdXJjZSI6MjAsInRhcmdldCI6MjB9LCJzdHlsZSI6eyJib2R5Ijp7Im5hbWUiOiJub25lIn0sImhlYWQiOnsibmFtZSI6Im5vbmUifX19XV0=
\[\begin{tikzcd}[ampersand replacement=\&]
	{A''} \&\& {A''} \\
	{A'} \&\& {A''} \\
	{A'} \& {A'} \& {A''} \\
	A \& {A'} \& {A''}
	\arrow["{r(A'')}", "\shortmid"{marking}, no head, from=1-1, to=1-3]
	\arrow[""{name=0, anchor=center, inner sep=0}, "{p_{c'}}"', from=1-1, to=2-1]
	\arrow[""{name=1, anchor=center, inner sep=0}, "{i_{A''}(\delrp_{c'})}", from=1-3, to=2-3]
	\arrow["{c'}", "\shortmid"{marking}, no head, from=2-1, to=2-3]
	\arrow["\id"', from=2-1, to=3-1]
	\arrow["\id", from=2-3, to=3-3]
	\arrow["{r(A')}", "\shortmid"{marking}, no head, from=3-1, to=3-2]
	\arrow[""{name=2, anchor=center, inner sep=0}, "{p_c}"', from=3-1, to=4-1]
	\arrow["{c'}", "\shortmid"{marking}, no head, from=3-2, to=3-3]
	\arrow[""{name=3, anchor=center, inner sep=0}, "{\delrp_c}", from=3-2, to=4-2]
	\arrow[""{name=4, anchor=center, inner sep=0}, "{i_{A''}(\push_{c'}(\delrp_c))}", from=3-3, to=4-3]
	\arrow["c"', "\shortmid"{marking}, no head, from=4-1, to=4-2]
	\arrow["{c'}"', "\shortmid"{marking}, no head, from=4-2, to=4-3]
	\arrow["{\dnl_{c'}}"{description}, draw=none, from=0, to=1]
	\arrow["{\dnl_c}"{description}, draw=none, from=2, to=3]
	\arrow["{(*)}"{description}, draw=none, from=3, to=4]
\end{tikzcd}\]
    \end{itemize}

    % 2.
    \item Same as (1), replacing predomains/morphisms/relations/squares with error domains.
\end{proof}


The next two lemmas concern quasi-equivalence and the functors $\li$ and
$U$.

\begin{lemma}\label{lem:Fcc-equiv-FcFc'}

  Let $A$, $A'$, and $A''$ be value objects and let $c : A \rel A'$ and $c' : A'
  \rel A''$ be predomain relations. Then we have $\li(c \comp c') \bisim \li c
  \comp \li c'$.
\end{lemma}
\begin{proof}
  First, we claim that $\li(c \comp c')$ and $\li c \comp \li c'$ are both
  quasi-left-represented by $\li e_{c'} \circ \li e_c$. Indeed, we have by part
  (1) of Lemma \ref{lem:representation-comp} that $e_c' \circ e_c$
  quasi-left-represents $c \comp c'$, and then by Lemma
  \ref{lem:representation-U-F} we have $\li(e_{c'} \circ e_c) = \li e_{c'} \circ \li e_c$
  quasi-left-represents $\li (c \comp c')$.
  %
  On the other hand, we also know that $\li e_c$ quasi-left-represents $\li c$ and
  $\li e_{c'}$ quasi-left-represents $\li c'$ again by Lemma
  \ref{lem:representation-U-F}. Then by part (2) of Lemma
  \ref{lem:representation-comp}, their composition quasi-left-represents $\li c
  \comp \li c'$.

  The result now follows by Lemma \ref{lem:left-rep-by-same-morphism}.
\end{proof}

\begin{lemma}\label{lem:Udd-equiv-UdUd'}
  Let $B$, $B'$, and $B''$ be value objects and let $d : B \rel B'$ and $d' : B'
  \rel B''$ be error domain relations. Then we have $U(d \comp d') \bisim Ud \comp Ud'$.
\end{lemma}
\begin{proof}
  Analogous to the proof of the previous lemma.
\end{proof}

Now we combine the above lemmas to show a result about the functors $\li$ and
$U$ applied to a composition of relations:

\begin{lemma}\label{lem:representation-comp-F-U}
  Let $A$, $A'$, and $A''$ be value objects, and $B$, $B'$, and
  $B''$ be computation objects. Let $c : A \rel A'$ and $c' : A' \rel A''$ be
  predomain relations with push-pull structures, and let $d : B \rel B'$ and $d'
  : B' \rel B''$ be error domain relations with push-pull structures.
  \begin{enumerate}
    \item If $\li c$ and $\li c'$ are quasi-right-representable, then $\li (c
    \comp c')$ is quasi-right-representable.

    \item If $Ud$ and $Ud'$ are quasi-left-representable, then $U(d \comp d')$
    is quasi-left-representable.

    % \item Given left-representation structures $\rho^L_c$ for $c$ and $\rho^L_{c'}$ for $c'$,
    % we can define a left-representation structure for the composition $c \comp c'$
    % \item Given right-representation structures $\rho^R_d$ for $d$ and $\rho^R_{d'}$ for $d'$,
    % we can define a right-representation structure for the composition $d \comp d'$
    % \item Given right-representation structures $\rho^R_{Fc}$ for $Fc$ and $\rho^R_{Fc'}$ for $Fc'$,
    % we can define a right-representation structure $\rho^R_{F(c \comp c')}$ for $F(c \comp c')$.
    % \item Given left-representation structures $\rho^L_{Ud}$ for $Ud$ and $\rho^L_{Ud'}$ for $Ud'$,
    % we can define a left-representation structure $\rho^L_{U(d \comp d')}$ for $U(d \comp d')$.
  \end{enumerate}

\end{lemma}
\begin{proof}
  \begin{enumerate}
    % 1.
    \item We show that $\li (c \comp c')$ is quasi-right-representable.

    First, by Lemma \ref{lem:Fcc-equiv-FcFc'} there is a square $\alpha$ of the form

    % https://q.uiver.app/#q=WzAsNSxbMCwwLCJcXGxpIEEiXSxbMiwwLCJcXGxpIEEnJyJdLFswLDEsIlxcbGkgQSJdLFsxLDEsIlxcbGkgQSciXSxbMiwxLCJcXGxpIEEnJyJdLFswLDEsIlxcbGkgKGMgXFxjb21wIGMnKSIsMCx7InN0eWxlIjp7ImJvZHkiOnsibmFtZSI6ImJhcnJlZCJ9LCJoZWFkIjp7Im5hbWUiOiJub25lIn19fV0sWzIsMywiXFxsaSBjIiwwLHsic3R5bGUiOnsiYm9keSI6eyJuYW1lIjoiYmFycmVkIn0sImhlYWQiOnsibmFtZSI6Im5vbmUifX19XSxbMyw0LCJcXGxpIGMnIiwwLHsic3R5bGUiOnsiYm9keSI6eyJuYW1lIjoiYmFycmVkIn0sImhlYWQiOnsibmFtZSI6Im5vbmUifX19XSxbMCwyLCJpX3tcXGxpIEF9KFxcZGVsdGFebCkiLDJdLFsxLDQsImlfe1xcbGkgQScnfShcXGRlbHRhXnIpIl0sWzgsOSwiXFxhbHBoYSIsMSx7InNob3J0ZW4iOnsic291cmNlIjoyMCwidGFyZ2V0IjoyMH0sInN0eWxlIjp7ImJvZHkiOnsibmFtZSI6Im5vbmUifSwiaGVhZCI6eyJuYW1lIjoibm9uZSJ9fX1dXQ==
    \[\begin{tikzcd}[ampersand replacement=\&]
      {\li A} \&\& {\li A''} \\
      {\li A} \& {\li A'} \& {\li A''}
      \arrow["{\li (c \comp c')}", "\shortmid"{marking}, no head, from=1-1, to=1-3]
      \arrow[""{name=0, anchor=center, inner sep=0}, "{i_{\li A}(\delta^l)}"', from=1-1, to=2-1]
      \arrow[""{name=1, anchor=center, inner sep=0}, "{i_{\li A''}(\delta^r)}", from=1-3, to=2-3]
      \arrow["{\li c}", "\shortmid"{marking}, no head, from=2-1, to=2-2]
      \arrow["{\li c'}", "\shortmid"{marking}, no head, from=2-2, to=2-3]
      \arrow["\alpha"{description}, draw=none, from=0, to=1]
    \end{tikzcd}\]
    where $\delta^l$ and $\delta^r$ are syntactic perturbations.

    We define the projection $p_{\li (c \comp c')}$ to be $p_{\li c \comp \li c'} \circ i_{A''}(\delta^r)$.
    We define $\dellp_{\li (c \comp c')}$ to be $\dellp_{\li c \comp \li c'} \cdot \delta^l$.

    Then we can build the $\dnr$ square by pasting the square $\alpha$ on top of
    the $\dnr$ square for the composition $\li c \comp \li c'$, as shown below:

    % https://q.uiver.app/#q=WzAsNyxbMCwwLCJcXGxpIEEiXSxbMiwwLCJcXGxpIEEnJyJdLFswLDEsIlxcbGkgQSJdLFsxLDEsIlxcbGkgQSciXSxbMiwxLCJcXGxpIEEnJyJdLFswLDIsIlxcbGkgQSJdLFsyLDIsIlxcbGkgQSJdLFswLDEsIlxcbGkgKGMgXFxjb21wIGMnKSIsMCx7InN0eWxlIjp7ImJvZHkiOnsibmFtZSI6ImJhcnJlZCJ9LCJoZWFkIjp7Im5hbWUiOiJub25lIn19fV0sWzIsMywiXFxsaSBjIiwwLHsic3R5bGUiOnsiYm9keSI6eyJuYW1lIjoiYmFycmVkIn0sImhlYWQiOnsibmFtZSI6Im5vbmUifX19XSxbMyw0LCJcXGxpIGMnIiwwLHsic3R5bGUiOnsiYm9keSI6eyJuYW1lIjoiYmFycmVkIn0sImhlYWQiOnsibmFtZSI6Im5vbmUifX19XSxbNSw2LCJyKFxcbGkgQSkiLDIseyJzdHlsZSI6eyJib2R5Ijp7Im5hbWUiOiJiYXJyZWQifSwiaGVhZCI6eyJuYW1lIjoibm9uZSJ9fX1dLFswLDIsImlfe1xcbGkgQX0oXFxkZWx0YV5sKSIsMl0sWzIsNSwiaV9BKFxcZGVsbHBfe1xcbGkgYyBcXGNvbXAgXFxsaSBjJ30pIiwyXSxbMSw0LCJpX3tcXGxpIEEnJ30oXFxkZWx0YV5yKSJdLFs0LDYsInBfe1xcbGkgYyBcXGNvbXAgXFxsaSBjJ30iXSxbMTIsMTQsIlxcZG5yX3tcXGxpIGMgXFxjb21wIFxcbGkgYyd9IiwxLHsic2hvcnRlbiI6eyJzb3VyY2UiOjIwLCJ0YXJnZXQiOjIwfSwic3R5bGUiOnsiYm9keSI6eyJuYW1lIjoibm9uZSJ9LCJoZWFkIjp7Im5hbWUiOiJub25lIn19fV0sWzExLDEzLCJcXGFscGhhIiwzLHsic2hvcnRlbiI6eyJzb3VyY2UiOjIwLCJ0YXJnZXQiOjIwfSwic3R5bGUiOnsiYm9keSI6eyJuYW1lIjoibm9uZSJ9LCJoZWFkIjp7Im5hbWUiOiJub25lIn19fV1d
    \[\begin{tikzcd}[ampersand replacement=\&]
      {\li A} \&\& {\li A''} \\
      {\li A} \& {\li A'} \& {\li A''} \\
      {\li A} \&\& {\li A}
      \arrow["{\li (c \comp c')}", "\shortmid"{marking}, no head, from=1-1, to=1-3]
      \arrow[""{name=0, anchor=center, inner sep=0}, "{i_{\li A}(\delta^l)}"', from=1-1, to=2-1]
      \arrow[""{name=1, anchor=center, inner sep=0}, "{i_{\li A''}(\delta^r)}", from=1-3, to=2-3]
      \arrow["{\li c}", "\shortmid"{marking}, no head, from=2-1, to=2-2]
      \arrow[""{name=2, anchor=center, inner sep=0}, "{i_A(\dellp_{\li c \comp \li c'})}"', from=2-1, to=3-1]
      \arrow["{\li c'}", "\shortmid"{marking}, no head, from=2-2, to=2-3]
      \arrow[""{name=3, anchor=center, inner sep=0}, "{p_{\li c \comp \li c'}}", from=2-3, to=3-3]
      \arrow["{r(\li A)}"', "\shortmid"{marking}, no head, from=3-1, to=3-3]
      \arrow["\alpha"{marking, allow upside down}, draw=none, from=0, to=1]
      \arrow["{\dnr_{\li c \comp \li c'}}"{description}, draw=none, from=2, to=3]
    \end{tikzcd}\]

    We define $\delrp_{\li (c \comp c')}$ to be $\delrp_{\li c \comp \li c'} \cdot \delta^r$.
    For $\dnl$, we paste the identity square $\delta^r \ltdyn \delta^r$ on top of
    the $\dnl$ square for the composition $\li c \comp \li c'$, and below that we paste
    the square $\id \ltdyn_{\li (c \comp c')}^{\li c \comp \li c'} \id$ which we get from
    the fact that $\li$ is lax.

    % https://q.uiver.app/#q=WzAsOSxbMCwwLCJcXGxpIEEnJyJdLFsyLDAsIlxcbGkgQScnIl0sWzAsMSwiXFxsaSBBJyciXSxbMiwxLCJcXGxpIEEnJyJdLFswLDIsIlxcbGkgQSJdLFsyLDIsIlxcbGkgQScnIl0sWzAsMywiXFxsaSBBIl0sWzIsMywiXFxsaSBBJyciXSxbMSwyLCJcXGxpIEEnIl0sWzAsMSwicihcXGxpIEEnJykiLDAseyJzdHlsZSI6eyJib2R5Ijp7Im5hbWUiOiJiYXJyZWQifSwiaGVhZCI6eyJuYW1lIjoibm9uZSJ9fX1dLFsyLDMsInIoXFxsaSBBJycpIiwwLHsic3R5bGUiOnsiYm9keSI6eyJuYW1lIjoiYmFycmVkIn0sImhlYWQiOnsibmFtZSI6Im5vbmUifX19XSxbNiw3LCJcXGxpIChjIFxcY29tcCBjJykiLDIseyJzdHlsZSI6eyJib2R5Ijp7Im5hbWUiOiJiYXJyZWQifSwiaGVhZCI6eyJuYW1lIjoibm9uZSJ9fX1dLFswLDIsImlfe0EnJ30oXFxkZWx0YV5yKSIsMl0sWzIsNCwicF97XFxsaSBjIFxcY29tcCBcXGxpIGMnfSIsMl0sWzQsNiwiXFxpZCIsMl0sWzEsMywiaV97QScnfShcXGRlbHRhXnIpIl0sWzMsNSwiaV97QScnfShcXGRlbHJwX3tcXGxpIGMgXFxjb21wIFxcbGkgYyd9KSJdLFs1LDcsIlxcaWQiXSxbNCw4LCJcXGxpIGMiLDAseyJzdHlsZSI6eyJib2R5Ijp7Im5hbWUiOiJiYXJyZWQifSwiaGVhZCI6eyJuYW1lIjoibm9uZSJ9fX1dLFs4LDUsIlxcbGkgYyciLDAseyJzdHlsZSI6eyJib2R5Ijp7Im5hbWUiOiJiYXJyZWQifSwiaGVhZCI6eyJuYW1lIjoibm9uZSJ9fX1dLFsxMiwxNSwiXFxpZCIsMyx7InNob3J0ZW4iOnsic291cmNlIjoyMCwidGFyZ2V0IjoyMH0sInN0eWxlIjp7ImJvZHkiOnsibmFtZSI6Im5vbmUifSwiaGVhZCI6eyJuYW1lIjoibm9uZSJ9fX1dLFsxMywxNiwiXFxkbmxfe1xcbGkgYyBcXGNvbXAgXFxsaSBjJ30iLDEseyJzaG9ydGVuIjp7InNvdXJjZSI6MjAsInRhcmdldCI6MjB9LCJzdHlsZSI6eyJib2R5Ijp7Im5hbWUiOiJub25lIn0sImhlYWQiOnsibmFtZSI6Im5vbmUifX19XV0=
    \[\begin{tikzcd}[ampersand replacement=\&]
      {\li A''} \&\& {\li A''} \\
      {\li A''} \&\& {\li A''} \\
      {\li A} \& {\li A'} \& {\li A''} \\
      {\li A} \&\& {\li A''}
      \arrow["{r(\li A'')}", "\shortmid"{marking}, no head, from=1-1, to=1-3]
      \arrow[""{name=0, anchor=center, inner sep=0}, "{i_{A''}(\delta^r)}"', from=1-1, to=2-1]
      \arrow[""{name=1, anchor=center, inner sep=0}, "{i_{A''}(\delta^r)}", from=1-3, to=2-3]
      \arrow["{r(\li A'')}", "\shortmid"{marking}, no head, from=2-1, to=2-3]
      \arrow[""{name=2, anchor=center, inner sep=0}, "{p_{\li c \comp \li c'}}"', from=2-1, to=3-1]
      \arrow[""{name=3, anchor=center, inner sep=0}, "{i_{A''}(\delrp_{\li c \comp \li c'})}", from=2-3, to=3-3]
      \arrow["{\li c}", "\shortmid"{marking}, no head, from=3-1, to=3-2]
      \arrow["\id"', from=3-1, to=4-1]
      \arrow["{\li c'}", "\shortmid"{marking}, no head, from=3-2, to=3-3]
      \arrow["\id", from=3-3, to=4-3]
      \arrow["{\li (c \comp c')}"', "\shortmid"{marking}, no head, from=4-1, to=4-3]
      \arrow["\id"{marking, allow upside down}, draw=none, from=0, to=1]
      \arrow["{\dnl_{\li c \comp \li c'}}"{description}, draw=none, from=2, to=3]
    \end{tikzcd}\]


    % 2.
    \item The proof that $U(d \comp d')$ is quasi-left-representable is analogous.

  \end{enumerate}
\end{proof}

% Want to show: U(d \comp d') is weakly equivalent to U(d) \comp U(d').
% This holds because both are quasi-representable by the same projection
% \begin{lemma}
%   Let $\mathcal M$ be a step-3 intensional model, and let
%   $d : B \rel B'$ and $d' : B' \rel B''$.
%   Let $\rho_\text{comp}$ be an arbitrary right-representation for $d \comp d'$.
%   Then the projection $p_{d \circ d'}$ is weakly equivalent to
% \end{lemma}




%%%%%%%%%%%%
% Products %
%%%%%%%%%%%%
\begin{lemma}\label{lem:representation-product} 
  Let $c_1 : A_1 \rel A_1'$ and $c_2 : A_2 \rel A_2'$. 
  \begin{enumerate}
    \item If $c_1$ and $c_2$ are
    quasi-left-representable, then $c_1 \times c_2$ is quasi-left-representable.
    \item If $\li c_1$ and $\li c_2$ are quasi-right-representable, then so
    is $\li (c_1 \times c_2)$.
  \end{enumerate}

  % Let $\rho^L_{c_1}$ be a left-representation structure for $c_1$, and
  % let $\rho^L_{c_2}$ be a left-representation structure for $c_2$.
  % Then we can define a left-representation structure for $c_1 \times c_2$.

  % Likewise, let $\rho^R_{Fc_1}$ and $\rho^R_{Fc_2}$ be right-representation
  % structures for $Fc_1$ and $Fc_2$ respectively.
  % Then we can define a right-representation structure for $F(c_1 \times c_2)$.
\end{lemma}
\begin{proof}
  \begin{enumerate}
    \item To show $c_1 \times c_2$ is quasi-left-representable, we define
    \begin{itemize}
      \item $e_{c_1 \times c_2} = e_{c_1} \times e_{c_2}$
      \item $\delre_{c_1 \times c_2} = \inl (\delre_{c_1}) \cdot \inr
      (\delre_{c_2})$ (recall from Section \ref{sec:perturbation-constructions}
      that $M_{A_1' \times A_2'}$ is the coproduct $M_{A_1'} \oplus M_{A_2'}$)
      \item $\delle_{c_1 \times c_2}$ is defined similarly
      \item The $\upl$ square is constructed using the functorial action of
      $\times$ and the corresponding squares for $c_1$ and $c_2$:
      % https://q.uiver.app/#q=WzAsNixbMCwwLCJBXzEgXFx0aW1lcyBBXzIiXSxbMCwxLCJBXzEnIFxcdGltZXMgQV8yIl0sWzAsMiwiQV8xJyBcXHRpbWVzIEFfMiciXSxbMSwwLCJBXzEnIFxcdGltZXMgQV8yJyJdLFsxLDEsIkFfMScgXFx0aW1lcyBBXzInIl0sWzEsMiwiQV8xJyBcXHRpbWVzIEFfMiciXSxbMCwzLCJjXzEgXFx0aW1lcyBjXzIiLDAseyJzdHlsZSI6eyJib2R5Ijp7Im5hbWUiOiJiYXJyZWQifSwiaGVhZCI6eyJuYW1lIjoibm9uZSJ9fX1dLFsxLDQsInIoQV8xJykgXFx0aW1lcyBjXzIiLDAseyJzdHlsZSI6eyJib2R5Ijp7Im5hbWUiOiJiYXJyZWQifSwiaGVhZCI6eyJuYW1lIjoibm9uZSJ9fX1dLFsyLDUsInIoQV8xJykgXFx0aW1lcyByKEFfMicpIiwyLHsic3R5bGUiOnsiYm9keSI6eyJuYW1lIjoiYmFycmVkIn0sImhlYWQiOnsibmFtZSI6Im5vbmUifX19XSxbMCwxLCJlX3tjXzF9IFxcdGltZXMgXFxpZCIsMl0sWzEsMiwiXFxpZCBcXHRpbWVzIGVfe2NfMn0iLDJdLFszLDQsImlfe0FfMSd9KFxcZGVscmVfe2NfMX0pIFxcdGltZXMgXFxpZCJdLFs0LDUsIlxcaWQgXFx0aW1lcyBpX3tBXzInfShcXGRlbHJlX3tjXzJ9KSJdLFs5LDExLCJcXHVwbF97Y18xfSBcXHRpbWVzIFxcaWQiLDEseyJzaG9ydGVuIjp7InNvdXJjZSI6MjAsInRhcmdldCI6MjB9LCJzdHlsZSI6eyJib2R5Ijp7Im5hbWUiOiJub25lIn0sImhlYWQiOnsibmFtZSI6Im5vbmUifX19XSxbMTAsMTIsIlxcaWQgXFx0aW1lcyBcXHVwbF97Y18yfSIsMSx7InNob3J0ZW4iOnsic291cmNlIjoyMCwidGFyZ2V0IjoyMH0sInN0eWxlIjp7ImJvZHkiOnsibmFtZSI6Im5vbmUifSwiaGVhZCI6eyJuYW1lIjoibm9uZSJ9fX1dXQ==
        \[\begin{tikzcd}[ampersand replacement=\&]
          {A_1 \times A_2} \& {A_1' \times A_2'} \\
          {A_1' \times A_2} \& {A_1' \times A_2'} \\
          {A_1' \times A_2'} \& {A_1' \times A_2'}
          \arrow["{c_1 \times c_2}", "\shortmid"{marking}, no head, from=1-1, to=1-2]
          \arrow[""{name=0, anchor=center, inner sep=0}, "{e_{c_1} \times \id}"', from=1-1, to=2-1]
          \arrow[""{name=1, anchor=center, inner sep=0}, "{i_{A_1'}(\delre_{c_1}) \times \id}", from=1-2, to=2-2]
          \arrow["{r(A_1') \times c_2}", "\shortmid"{marking}, no head, from=2-1, to=2-2]
          \arrow[""{name=2, anchor=center, inner sep=0}, "{\id \times e_{c_2}}"', from=2-1, to=3-1]
          \arrow[""{name=3, anchor=center, inner sep=0}, "{\id \times i_{A_2'}(\delre_{c_2})}", from=2-2, to=3-2]
          \arrow["{r(A_1') \times r(A_2')}"', "\shortmid"{marking}, no head, from=3-1, to=3-2]
          \arrow["{\upl_{c_1} \times \id}"{description}, draw=none, from=0, to=1]
          \arrow["{\id \times \upl_{c_2}}"{description}, draw=none, from=2, to=3]
        \end{tikzcd}\]
      \item The $\upr$ square is defined similarly
    \end{itemize}

    \item To show that $\li(c_1 \times c_2)$ is quasi-right-representable, we define
    \begin{itemize}
      \item $p_{\li (c_1 \times c_2)} = (p_{\li c_1} \timesk A_2) \circ (A_1' \timesk p_{\li c_2})$
      \item $\dellp_{\li (c_1 \times c_2)} = (\dellp_{\li c_1} \timesk \id) \cdot (\id \timesk \dellp_{\li c_2})$ 
      using the Kleisli product actions on syntactic perturbations in Definition \ref{def:kleisli-product-perturbations}
      \item $\delrp_{\li (c_1 \times c_2)} = (\delrp_{\li c_1} \timesk \id) \cdot (\id \timesk \delrp_{\li c_2})$
      \item The squares are obtained via the functorial action of $\timesk$ on the squares for $\li c_1$ and $\li c_2$.
    \end{itemize}
  \end{enumerate}
\end{proof}


%%%%%%%%%
% Arrow %
%%%%%%%%%
\begin{lemma}\label{lem:representation-arrow}
  Let $c : A \rel A'$ and $d : B \rel B'$.
  %
  \begin{enumerate}
    \item If $c$ is quasi-left-representable and $d$ is quasi-right-representable, then
    $c \arr d$ is quasi-right-representable.

    \item If $\li c$ is quasi-right-representable and $Ud$ is
    quasi-left-representable, then $U(c \arr d)$ is quasi-left-representable.
  \end{enumerate}

  % Let $\rho^L_c$ be a left-representation structure for $c$,
  % and let $\rho^R_d$ be a right-representation structure for $d$.
  % Then we can define a right-representation structure for $c \arr d$.

  % Likewise, let $\rho^R_{Fc}$ be a right-representation structure for $Fc$,
  % and let $\rho^L_{Ud}$ be a left-representation structure for $Ud$.
  % Then we can define a left-representation structure for $U(c \arr d)$.
\end{lemma}
\begin{proof}
  \begin{enumerate}
    \item 
      To show $c \arr d$ is quasi-right-representable, we define:
      \begin{itemize}
        \item $p_{c \arr d} = e_c \arr p_d$
        \item $\dellp_{c \arr d} = \inl(\delle_c) \cdot \inr(\dellp_d)$
        \item $\delrp_{c \arr d} = \inl(\delre_c) \cdot \inr(\delrp_d)$
        \item The squares $\dnr$ and $\dnl$ are obtained via the functorial action of $\arr$ on squares.

      \end{itemize}

    % TODO check this
    \item
      We show $U(c \arr d)$ is quasi-left-representable as follows:
      \begin{itemize}
        \item $e_{U(c \arr d)} = (p_{\li c} \tok B') \circ (A \tok e_{Ud})$
        \item $\delre_{U(c \arr d)} = (\delrp_{\li c} \tok B') \cdot (A' \tok
              \delre_{Ud})$ using the Kleisli arrow actions on syntactic
              perturbations in Definition \ref{def:kleisli-arrow-perturbations}.
        \item $\delle_{U(c \arr d)} = (\dellp_{\li c} \tok B) \cdot (A \tok \delle_{Ud})$
        \item The squares $\upl$ and $\upr$ are obtained via the functorial action of $\tok$ on squares.
        For instance, $\upl$ is given by the following square:

        % https://q.uiver.app/#q=WzAsNixbMCwwLCJVKEEgXFx0byBCKSJdLFsyLDAsIlUoQScgXFx0byBCJykiXSxbMiwxLCJVKEEnIFxcdG8gQicpIl0sWzIsMiwiVShBJyBcXHRvIEInKSJdLFswLDEsIlUoQSBcXHRvIEInKSJdLFswLDIsIlUoQScgXFx0byBCJykiXSxbMCwxLCJVKGMgXFx0byBkKSJdLFs0LDIsIlUoYyBcXHRvIHIoQicpKSJdLFs1LDMsIlUocihBJykgXFx0byByKEInKSkiXSxbNCw1LCJwX3tGY30gXFx0b2sgQiciLDJdLFsxLDIsIkEnIFxcdG9rIFxcZGVscmVfe1VkfSJdLFsyLDMsIlxcZGVscnBfe0ZjfSBcXHRvayBCJyJdLFswLDQsIkEgXFx0b2sgZV97VWR9IiwyXSxbMTIsMTAsIlxcaWRfe0ZjfSBcXHRvayBcXHVwbF97VWR9IiwxLHsic2hvcnRlbiI6eyJzb3VyY2UiOjIwLCJ0YXJnZXQiOjIwfSwic3R5bGUiOnsiYm9keSI6eyJuYW1lIjoibm9uZSJ9LCJoZWFkIjp7Im5hbWUiOiJub25lIn19fV0sWzksMTEsIlxcZG5sX3tGY30gXFx0b2sgXFxpZF97cihCJyl9IiwxLHsic2hvcnRlbiI6eyJzb3VyY2UiOjIwLCJ0YXJnZXQiOjIwfSwic3R5bGUiOnsiYm9keSI6eyJuYW1lIjoibm9uZSJ9LCJoZWFkIjp7Im5hbWUiOiJub25lIn19fV1d
        \[\begin{tikzcd}[ampersand replacement=\&,row sep=large]
          {U(A \to B)} \&\& {U(A' \to B')} \\
          {U(A \to B')} \&\& {U(A' \to B')} \\
          {U(A' \to B')} \&\& {U(A' \to B')}
          \arrow["{U(c \to d)}", from=1-1, to=1-3]
          \arrow["{U(c \to r(B'))}", from=2-1, to=2-3]
          \arrow["{U(r(A') \to r(B'))}", from=3-1, to=3-3]
          \arrow[""{name=0, anchor=center, inner sep=0}, "{p_{Fc} \tok B'}"', from=2-1, to=3-1]
          \arrow[""{name=1, anchor=center, inner sep=0}, "{A' \tok \delre_{Ud}}", from=1-3, to=2-3]
          \arrow[""{name=2, anchor=center, inner sep=0}, "{\delrp_{Fc} \tok B'}", from=2-3, to=3-3]
          \arrow[""{name=3, anchor=center, inner sep=0}, "{A \tok e_{Ud}}"', from=1-1, to=2-1]
          \arrow["{\id_{Fc} \tok \upl_{Ud}}"{description}, draw=none, from=3, to=1]
          \arrow["{\dnl_{Fc} \tok \id_{r(B')}}"{description}, draw=none, from=0, to=2]
        \end{tikzcd}\]
      \end{itemize}

    The construction of $\upr$ is similar.
\end{enumerate}
\end{proof}


\subsection{Model Construction: Composition and Functorial Actions on Relations}

We now show that our notions of value and computation relations compose.
\begin{definition}[composition of value (resp. computation) relations]\label{def:value-computation-rel-comp}
  Let $A$, $A'$ and $A''$ be value objects and let $c$ and $c'$ be \emph{value
  relations} between $A$ and $A'$ and $A'$ and $A''$ respectively. Then we
  define their composition $cc'$ as follows:

  \begin{itemize}
    \item The predomain relation is the composition of the underlying predomain
    relations $cc'$
    \item The push-pull structure is given by Lemma \ref{lem:push-pull-comp}.
    \item Quasi-left-representability of $cc'$ follows from Lemma
    \ref{lem:representation-comp} and the fact that $c$ and $c'$ are
    quasi-left-representable
    \item Quasi-right-representability of $\li(cc')$ holds by Lemma
    \ref{lem:representation-comp-F-U} and the quasi-right-representability of
    $\li c$ and $\li c'$
  \end{itemize}

  Likewise, let $B$, $B'$ and $B''$ be computation objects and let $d$ and $d'$ be
  \emph{computation relations} between $B$ and $B'$ and $B'$ and $B''$
  respectively. We define their composition $dd'$ as follows:

  \begin{itemize}
    \item The error domain relation is the composition of the underlying error
    domain relations $dd'$
    \item The push-pull structure is given by Lemma \ref{lem:push-pull-comp}.
    \item Quasi-right-representability of $dd'$ follows from Lemma
    \ref{lem:representation-comp} and the fact that $d$ and $d'$ are
    quasi-right-representable
    \item Quasi-left-representability of $U(dd')$ holds by Lemma
    \ref{lem:representation-comp-F-U} and the quasi-left-representability of
    $Ud$ and $Ud'$
  \end{itemize}
\end{definition}


We can define the functorial actions of $\li$, $U$, $\times$, and $\arr$ on
value and computation relations as follows:
\begin{definition}[functorial actions on relations]\label{def:functorial-actions-on-relations}

  \begin{enumerate}

    \item Let $A$ and $A'$ be value objects and $c : A \rel A'$ a value relation.
          Then we define the computation relation $\li c$ as follows:
      \begin{itemize}
        \item The error domain relation is given by the action of $\li$ on the
        underlying predomain relation of $c$
        \item The push-pull structure is given by Lemma \ref{lem:push-pull-U-F}
        \item Quasi-right-representability of $\li c$ holds because it is part
        of the definition of the value relation $c$
        \item Quasi-left-representability of $U \li c$ follows from Lemma
        \ref{lem:representation-U-F} and the quasi-left-representability of $c$
      \end{itemize}

    \item Let $B$ and $B'$ be computation objects and $d : B \rel B'$ a computation relation.
          We define $Ud$ as follows:
          \begin{itemize}
            \item The predomain relation is given by the action of $U$ on the
            error domain relation of $d$
            \item The push-pull structure is given by Lemma
            \ref{lem:push-pull-U-F}
            \item Quasi-left-representability of $Ud$ holds because it is part
            of the definition of the computation relation $d$
            \item Quasi-right-representability of $\li(Ud)$ follows from Lemma
            \ref{lem:representation-U-F} and the quasi-right-representability of
            $d$
          \end{itemize}

    \item Let $c_1 : A_1 \rel A_1'$ and $c_2 : A_2 \rel A_2'$ be value relations.
          We define $c_1 \times c_2$ as follows:
          \begin{itemize}
            \item The relation is given by the action of $\times$ on the
            predomain relations of $c_1$ and $c_2$
            \item The push-pull structure is given by Lemma
            \ref{lem:push-pull-times}
            \item Quasi-left-representability of $c_1 \times c_2$ and
            quasi-right-representabiity of $\li(c_1 \times c_2)$ are given by
            Lemma \ref{lem:representation-product}
          \end{itemize}

    \item Let $c : A \rel A'$ be a value relation and $d : B \rel B'$ be a computation relation.
          We define $c \arr d$ as follows:
          \begin{itemize}
            \item The error domain relation is given by the action of $\arr$ on
            the underlying relations of $c$ and $d$
            \item The push-pull structure is given by Lemma
            \ref{lem:push-pull-arrow}
            \item Quasi-right-representability of $c \arr d$, and
                  quasi-left-representability of $U(c \arr d)$ follow from Lemma
                  \ref{lem:representation-arrow}
          \end{itemize}

  \end{enumerate}
\end{definition}

Lastly, we prove quasi-equivalence of relations involving composition and the
functors $U$, $\li$, $\to$ and $\times$.

\begin{lemma}\label{lem:quasi-order-equiv-functors}
  The following hold:
  \begin{itemize}
    \item $U(d \comp d') \qordeq U(d) \comp U(d')$
    \item $\li(c \comp c') \qordeq \li (c) \comp \li (c')$
    \item $(c \comp c') \to (d \comp d') \qordeq (c \to d) \comp (c' \to d')$
    \item $(c_1 \comp c_1') \times (c_2 \comp c_2') \qordeq (c_1 \times c_2) \comp (c_1'\times c_2')$
  \end{itemize}
\end{lemma}
\begin{proof}
  (1) and (2) were already shown in Lemmas \ref{lem:Fcc-equiv-FcFc'} and
  \ref{lem:Udd-equiv-UdUd'}. (3) follows from observing that both relations are
  right-represented by the morphism $e_c'e_c \arr p_{d}p_{d'}$ and (4) from the
  fact that both are left-represented by the morphism $e_{c_1'} e_{c_1} \times
  e_{c_2'} e_{c_2}$
\end{proof}

% We now establish the quasi-order-equivalence for the functors.
% We already showed that $U(d \comp d') \qordeq U(d)U(d')$ and $F(c \comp c'') \qordeq F(c)F(c')$
% in the proof of Lemma \ref{lem:representation-comp}.
% The other two properties
% $(cc') \to (dd') \qordeq (c \to d)(c' \to d')$ and
% $(c_1c_1') \times (c_2c_2') \qordeq (c_1 \times c_2)(c_1'\times c_2')$
% are proved similarly, noting in both cases that the both relations are
% quasi-represented by the same morphism.

\subsection{Definitions of Error Ordering and Weak Bisimilarity for the Delay Monad}
\label{sec:relations-on-delay-monad}

We define a notion of lock-step error ordering and weak bisimilarity relation
for the coinductive delay monad $\delay(\mathbb{N} + {\mho})$:

The lock-step error ordering is defined coinductively by the following rules:

%
\begin{mathpar}
  \inferrule*[]
  { }
  {\tnow (\inr\, 1) \ledelay d}

  \inferrule*[]
  {x_1 \le_X x_2}
  {\tnow (\inl\, x_1) \ledelay \tnow (\inl\, x_2)}

  \inferrule*[]
  {d_1 \ledelay d_2}
  {\tlater\, d_1 \ledelay \tlater\, d_2}
\end{mathpar}
%
And we similarly define by coinduction a ``weak bisimilarity'' relation on
$\delay(\mathbb{N} + {\mho})$. This uses a relation $d \Da x_?$ between
$\delay(\mathbb{N} + {\mho})$ and $\mathbb{N} + {\mho}$ that is defined as $d
\Da n_? := \Sigma_{i \in \mathbb{N}} d = \tlater^i(\tnow\, n_?)$. Then weak
bisimilarity for the delay monad is defined coinductively by the rules
%
\begin{mathpar}
  \inferrule*[]
  {n_? \bisim_{\mathbb{N} + {\mho}} m_?}
  {\tnow\, n_? \bisimdelay \tnow\, m_? }

  \inferrule*[leftskip=1.5em]
  {d_1 \Da n_? \and n_? \bisim_{\mathbb{N} + {\mho}} m_?}
  {\tlater\, d_1 \bisimdelay \tnow\, m_? }

  \inferrule*[leftskip=1.5em]
  {d_2 \Da m_? \and n_? \bisim_{\mathbb{N} + {\mho}} m_?}
  {\tnow\, n_? \bisimdelay \tlater\, d_2}

  \inferrule*[leftskip=1.5em]
  {d_1 \bisimdelay d_2}
  {\tlater\, d_1 \bisimdelay \tlater\, d_2 }
  %
  % \inferrule*[]
  % {d_1 \Da x_? \and d_2 \Da y_? \and x_? \bisim_{X + 1} y_?}
  % {d_1 \bisimdelay d_2}
  %
  % \inferrule*[]
  % {d_1 \bisimdelay d_2}
  % {\tlater d_1 \bisimdelay \tlater d_2 }
\end{mathpar}



% %%%%%%%%%%%%%%%%%%%%%%
% % Model Construction %
% %%%%%%%%%%%%%%%%%%%%%%
% Now we can give the proof of the main lemma:
%   % Write 
%   % %
%   % \[ \mathcal M = (\vf, \vsq, \ef, \esq, \Ff, \Fsq, \Uf, \Usq, \arrf, \arrsq). \] 
%   % %

%   We define a step-3 model $\mathcal M'$ as follows:
%   \begin{itemize}
%     \item The objects of $\mathcal M'$ are defined to be the same as the objects of $\mathcal M$.
%     \item The value and computation morphisms in $\mathcal M'$ are the same as those of $\mathcal M$.
%     \item A value relation is defined to be a tuple $(c, \rho^L_c, \rho^R_{Fc})$ with
%     \begin{itemize}
%       \item $c$ a value relation in $\mathcal M$,
%       \item $\rho^L_c$ a left-representation structure for $c$, and
%       \item $\rho^R_{Fc}$ a right-representation structure for $Fc$
%     \end{itemize}
%     \item Likewise, a computation relation is defined to be a tuple $(d, \rho^R_d, \rho^L_{Ud})$ with
%     \begin{itemize}
%       \item $d$ a computation relation in $\mathcal M$,
%       \item $\rho^R_d$ a right-representation structure for $d$, and
%       \item $\rho^L_{Ud}$ a left-representation structure for $Ud$.
%     \end{itemize}
%     \item Morphisms of value relations (i.e., the value squares) are defined by simply
%     ignoring the representation structures. That is, a morphism of value relations
%     $\alpha \in \vsq'((c, \rho^L_c, \rho^R_{Fc}), (c' \rho^L_{c'}, \rho^R_{Fc'}))$ is simply a morphism of value
%     relations in $\vsq(c, c')$. Likewise for computations.
%   \end{itemize}

% We define horizontal composition of relations and squares as follows:
% Let $c : A \rel A'$ and $c' : A' \rel A''$. We define
% %
% \begin{align*} 
%   ((c, \rho^L_c, &\rho^R_{Fc}) \comp (c', \rho^L_{c'}, \rho^R_{Fc'})) =
%   (c \comp c', \rho^L_{c \comp c'}, \rho^R_{F(c \comp c')}), 
% \end{align*}
% %
% and
% %
% \begin{align*}
%   ((d, \rho^R_d, &\rho^L_{Ud}) \comp (d', \rho^R_{d'}, \rho^L_{Ud'})) =
%   (d \comp d', \rho^R_{d \comp d'}, \rho^L_{U(d \comp d')}),
% \end{align*}
% %
% where $\rho^L_{c \comp c'}, \rho^R_{F(c \comp c')}, \rho^R_{d \comp d'},$
% and $\rho^L_{U(d \comp d')}$ are as defined in Lemma \ref{lem:representation-comp}.

% Now we define the functors $F$, $U$, $\times$, and $\arr$.
% On objects, the behavior is the same as the respective functors in $\mathcal M$.
% For relations, we define

% \[ F(c, \rho^L_c, \rho^R_{Fc}) = 
%   (F c, \rho^R_{Fc}, \rho^L_{UF(c)}), \]
% and
% \[ U(d, \rho^R_d, \rho^L_{Ud}) = 
%   (U d, \rho^L_{Ud}, \rho^R_{FU(d)}), \]

% where $\rho^L_{UF(c)}$ and $\rho^R_{FU(d)}$ are as defined in the proof of Lemma
% \ref{lem:representation-UF-FU}.

% We define \[ (c_1, \rho^L_{c_1}, \rho^R_{Fc_1}) \times (c_2, \rho^L_{c_2}, \rho^R_{Fc_2}) = 
%              (c_1 \times c_2, \rho^L_{c_1 \times c_2}, \rho^R_{F(c_1 \times c_2)}), \]

% where $\rho^L_{c_1 \times c_2}$ and $\rho^R_{F(c_1 \times c_2)}$ are as defined in the
% proof of Lemma \ref{lem:representation-product}.

% Lastly, we define

% \[ (c, \rho^L_c, \rho^R_{Fc}) \arr (d, \rho^R_d, \rho^L_{Ud}) =
%     (c \arr d, \rho^R_{c \arr d}, \rho^L_{U(c \arr d)}), \]

% where $\rho^R_{c \arr d}$ and $\rho^L_{U(c \arr d)}$ are as defined in the proof
% of Lemma \ref{lem:representation-arrow}.

% % We define $(c, \rho^L_c) \arr (d, \rho^R_d) = (c \arr d, \rho^R_{c \arr d})$.




%%%%%%%%%%%%%%%%%%%%%%%%%%%%%%%%%%%%%%%%%%%%%%%%%%%%%%%%%%%%%%%%%%%%%%%%%%%%%%



\end{document}
